\section{Dataset}
\label{sec:data}
\ihead{Dataset}
The dataset combines electricity demand in the 50Hertz control area with power generation forecast by source, market prices, holiday indicators, and weather data from multiple meteorological stations.

Since the influence of exogenous variables on electricity demand is unknown a priori, all available features are retained during data preparation. This includes generation variables and a broad set of meteorological observations. Retaining the full feature set allows for exploratory analysis and model-based feature selection in later stages.

During data inspection, duplicate and missing hourly timestamps were observed due to \ac{DST} transitions, during which one hour is repeated in autumn and omitted in spring. To ensure a unique and continuous time index required for forecasting, all timestamps are converted from local time to \ac{UTC}. Furthermore, hourly resampling is performed using left-closed, right-labeled intervals to prevent data leakage, ensuring that each aggregated value depends exclusively on past observations and is therefore suitable for forecasting applications.

\paragraph*{Energy Demand}
Electricity demand and generation data for the 50Hertz control area from 2015 to 2024 are provided by the German Federal Network Agency (Bundesnetzagentur) via the SMARD API~\cite{bundesnetzagentur_2025_smard}. The primary forecasting target is the grid load, representing total electricity demand in the region.

To gain an initial understanding of the data, the time series are visualized. The annual profiles exhibit a high degree of similarity, indicating pronounced periodic behavior. To analyze characteristic intra-annual variations, the demand is averaged across all years and shown in Figure~\ref{fig:grid_load_year}. The resulting profile reveals a strong and stable seasonal pattern, with higher demand in winter and lower demand in summer, driven by increased heating, lighting, and indoor activity during colder months. This structure remains consistent throughout the observation period, suggesting no major structural changes in demand behavior between 2015 and 2024.

In addition to the annual cycle, a clear reduction in demand is observed around the Christmas and New Year period. This temporary decrease aligns with reduced industrial and commercial activity during public holidays and underlines the importance of calendar-related features in electricity demand forecasting. Furthermore, in Figure~\ref{fig:average_load_by_day_and_season} it is clear that during the weekend the demand decreases for the same reason. The projects objective is a prediction of a complete day. Hence, the average grid load over a day is useful to see beforehand. Figure~\ref{fig:average_load_by_hour_and_season} illustrates a clear demand pattern, with low consumption at night and two peaks around 9am and 5pm. The morning peak results from the concurrent start of industrial and residential activity, while the evening peak reflects continued industrial demand combined with increased household consumption. As activity declines overnight, electricity demand reaches its daily minimum. This pattern is consistent across all seasons.
\begin{figure}[ht]
    \centering
    \begin{subfigure}[a]{0.49\linewidth}
        \centering
        \begin{tikzpicture}[>=Stealth]
    \begin{scope}[yshift=-1cm]
    % Plot 1
    \begin{axis}[
        small axis,
        width=\linewidth,
        height=\imageheight,
        grid=both,
        date coordinates in=x,
        xlabel={Year},
        ylabel={Grid Load in MWh},
        ylabel style={yshift=-8pt},
        ymin=6000, ymax=18000,
        xmin=2015-01-01, xmax=2024-12-31,
        xtick={ 2015-01-01, 2016-01-01, 2017-01-01, 2018-01-01, 2019-01-01, 2020-01-01, 2021-01-01, 2022-01-01, 2023-01-01, 2024-01-01},
        xticklabels={$2015$, $2016$, $2017$, $2018$, $2019$, $2020$, $2021$, $2022$, $2023$, $2024$},
        xticklabel style={rotate=45}
    ]
    
        % OG Plot
        \addplot+[color=tabblue, no markers] coordinates {
  (2015-01-01 00:00:00+00:00,8981.2500000000)
  (2015-01-01 10:00:00+00:00,11049.7500000000)
  (2015-01-01 20:00:00+00:00,10182.5000000000)
  (2015-01-02 06:00:00+00:00,10610.7500000000)
  (2015-01-02 16:00:00+00:00,13134.0000000000)
  (2015-01-03 02:00:00+00:00,9371.7500000000)
  (2015-01-03 12:00:00+00:00,12160.2500000000)
  (2015-01-03 22:00:00+00:00,10164.7500000000)
  (2015-01-04 08:00:00+00:00,10771.7500000000)
  (2015-01-04 18:00:00+00:00,12307.2500000000)
  (2015-01-05 04:00:00+00:00,10715.2500000000)
  (2015-01-05 14:00:00+00:00,14403.5000000000)
  (2015-01-06 00:00:00+00:00,10538.7500000000)
  (2015-01-06 10:00:00+00:00,14944.7500000000)
  (2015-01-06 20:00:00+00:00,12816.2500000000)
  (2015-01-07 06:00:00+00:00,14313.2500000000)
  (2015-01-07 16:00:00+00:00,15303.7500000000)
  (2015-01-08 02:00:00+00:00,10778.0000000000)
  (2015-01-08 12:00:00+00:00,14286.5000000000)
  (2015-01-08 22:00:00+00:00,11373.2500000000)
  (2015-01-09 08:00:00+00:00,14664.7500000000)
  (2015-01-09 18:00:00+00:00,14394.0000000000)
  (2015-01-10 04:00:00+00:00,9750.7500000000)
  (2015-01-10 14:00:00+00:00,11965.0000000000)
  (2015-01-11 00:00:00+00:00,9666.0000000000)
  (2015-01-11 10:00:00+00:00,12958.5000000000)
  (2015-01-11 20:00:00+00:00,11386.0000000000)
  (2015-01-12 06:00:00+00:00,14118.0000000000)
  (2015-01-12 16:00:00+00:00,15080.7500000000)
  (2015-01-13 02:00:00+00:00,10351.7500000000)
  (2015-01-13 12:00:00+00:00,14968.2500000000)
  (2015-01-13 22:00:00+00:00,10903.0000000000)
  (2015-01-14 08:00:00+00:00,14602.0000000000)
  (2015-01-14 18:00:00+00:00,14608.5000000000)
  (2015-01-15 04:00:00+00:00,11578.7500000000)
  (2015-01-15 14:00:00+00:00,14187.7500000000)
  (2015-01-16 00:00:00+00:00,10313.2500000000)
  (2015-01-16 10:00:00+00:00,15275.0000000000)
  (2015-01-16 20:00:00+00:00,12621.5000000000)
  (2015-01-17 06:00:00+00:00,11131.7500000000)
  (2015-01-17 16:00:00+00:00,13233.7500000000)
  (2015-01-18 02:00:00+00:00,9842.5000000000)
  (2015-01-18 12:00:00+00:00,12614.5000000000)
  (2015-01-18 22:00:00+00:00,11078.5000000000)
  (2015-01-19 08:00:00+00:00,15000.7500000000)
  (2015-01-19 18:00:00+00:00,14851.7500000000)
  (2015-01-20 04:00:00+00:00,11950.2500000000)
  (2015-01-20 14:00:00+00:00,14679.2500000000)
  (2015-01-21 00:00:00+00:00,11105.7500000000)
  (2015-01-21 10:00:00+00:00,15488.0000000000)
  (2015-01-21 20:00:00+00:00,13349.0000000000)
  (2015-01-22 06:00:00+00:00,14747.7500000000)
  (2015-01-22 16:00:00+00:00,15381.7500000000)
  (2015-01-23 02:00:00+00:00,11368.2500000000)
  (2015-01-23 12:00:00+00:00,14881.5000000000)
  (2015-01-23 22:00:00+00:00,11966.5000000000)
  (2015-01-24 08:00:00+00:00,13216.2500000000)
  (2015-01-24 18:00:00+00:00,13364.2500000000)
  (2015-01-25 04:00:00+00:00,9961.5000000000)
  (2015-01-25 14:00:00+00:00,11839.7500000000)
  (2015-01-26 00:00:00+00:00,10488.0000000000)
  (2015-01-26 10:00:00+00:00,15153.7500000000)
  (2015-01-26 20:00:00+00:00,13123.5000000000)
  (2015-01-27 06:00:00+00:00,14554.0000000000)
  (2015-01-27 16:00:00+00:00,14844.2500000000)
  (2015-01-28 02:00:00+00:00,10716.2500000000)
  (2015-01-28 12:00:00+00:00,15068.5000000000)
  (2015-01-28 22:00:00+00:00,11857.2500000000)
  (2015-01-29 08:00:00+00:00,15257.0000000000)
  (2015-01-29 18:00:00+00:00,14721.2500000000)
  (2015-01-30 04:00:00+00:00,11933.0000000000)
  (2015-01-30 14:00:00+00:00,14052.7500000000)
  (2015-01-31 00:00:00+00:00,10564.5000000000)
  (2015-01-31 10:00:00+00:00,13443.2500000000)
  (2015-01-31 20:00:00+00:00,11301.0000000000)
  (2015-02-01 06:00:00+00:00,9920.5000000000)
  (2015-02-01 16:00:00+00:00,12314.0000000000)
  (2015-02-02 02:00:00+00:00,11117.5000000000)
  (2015-02-02 12:00:00+00:00,14889.0000000000)
  (2015-02-02 22:00:00+00:00,12486.2500000000)
  (2015-02-03 08:00:00+00:00,15251.7500000000)
  (2015-02-03 18:00:00+00:00,15395.0000000000)
  (2015-02-04 04:00:00+00:00,12656.5000000000)
  (2015-02-04 14:00:00+00:00,14806.2500000000)
  (2015-02-05 00:00:00+00:00,11955.2500000000)
  (2015-02-05 10:00:00+00:00,15015.2500000000)
  (2015-02-05 20:00:00+00:00,13574.5000000000)
  (2015-02-06 06:00:00+00:00,14782.0000000000)
  (2015-02-06 16:00:00+00:00,14955.0000000000)
  (2015-02-07 02:00:00+00:00,11254.0000000000)
  (2015-02-07 12:00:00+00:00,12829.0000000000)
  (2015-02-07 22:00:00+00:00,11143.5000000000)
  (2015-02-08 08:00:00+00:00,11741.7500000000)
  (2015-02-08 18:00:00+00:00,13001.5000000000)
  (2015-02-09 04:00:00+00:00,11814.2500000000)
  (2015-02-09 14:00:00+00:00,14243.5000000000)
  (2015-02-10 00:00:00+00:00,11427.2500000000)
  (2015-02-10 10:00:00+00:00,15548.7500000000)
  (2015-02-10 20:00:00+00:00,13497.2500000000)
  (2015-02-11 06:00:00+00:00,14560.0000000000)
  (2015-02-11 16:00:00+00:00,15020.7500000000)
  (2015-02-12 02:00:00+00:00,11717.0000000000)
  (2015-02-12 12:00:00+00:00,14843.5000000000)
  (2015-02-12 22:00:00+00:00,12465.0000000000)
  (2015-02-13 08:00:00+00:00,15121.2500000000)
  (2015-02-13 18:00:00+00:00,15169.7500000000)
  (2015-02-14 04:00:00+00:00,11277.7500000000)
  (2015-02-14 14:00:00+00:00,12164.0000000000)
  (2015-02-15 00:00:00+00:00,10677.2500000000)
  (2015-02-15 10:00:00+00:00,12658.5000000000)
  (2015-02-15 20:00:00+00:00,11703.5000000000)
  (2015-02-16 06:00:00+00:00,14281.7500000000)
  (2015-02-16 16:00:00+00:00,14443.0000000000)
  (2015-02-17 02:00:00+00:00,11579.5000000000)
  (2015-02-17 12:00:00+00:00,14817.5000000000)
  (2015-02-17 22:00:00+00:00,12231.0000000000)
  (2015-02-18 08:00:00+00:00,15213.5000000000)
  (2015-02-18 18:00:00+00:00,15124.5000000000)
  (2015-02-19 04:00:00+00:00,12274.5000000000)
  (2015-02-19 14:00:00+00:00,13299.5000000000)
  (2015-02-20 00:00:00+00:00,11280.7500000000)
  (2015-02-20 10:00:00+00:00,14601.2500000000)
  (2015-02-20 20:00:00+00:00,12905.7500000000)
  (2015-02-21 06:00:00+00:00,11350.2500000000)
  (2015-02-21 16:00:00+00:00,12092.5000000000)
  (2015-02-22 02:00:00+00:00,10272.5000000000)
  (2015-02-22 12:00:00+00:00,11866.7500000000)
  (2015-02-22 22:00:00+00:00,11446.0000000000)
  (2015-02-23 08:00:00+00:00,14556.2500000000)
  (2015-02-23 18:00:00+00:00,15303.5000000000)
  (2015-02-24 04:00:00+00:00,12268.2500000000)
  (2015-02-24 14:00:00+00:00,13999.7500000000)
  (2015-02-25 00:00:00+00:00,11355.2500000000)
  (2015-02-25 10:00:00+00:00,15104.7500000000)
  (2015-02-25 20:00:00+00:00,13639.5000000000)
  (2015-02-26 06:00:00+00:00,14701.0000000000)
  (2015-02-26 16:00:00+00:00,14214.5000000000)
  (2015-02-27 02:00:00+00:00,11790.2500000000)
  (2015-02-27 12:00:00+00:00,14323.7500000000)
  (2015-02-27 22:00:00+00:00,12051.0000000000)
  (2015-02-28 08:00:00+00:00,12926.2500000000)
  (2015-02-28 18:00:00+00:00,13504.2500000000)
  (2015-03-01 04:00:00+00:00,9624.2500000000)
  (2015-03-01 14:00:00+00:00,10763.5000000000)
  (2015-03-02 00:00:00+00:00,9494.5000000000)
  (2015-03-02 10:00:00+00:00,14204.7500000000)
  (2015-03-02 20:00:00+00:00,12764.5000000000)
  (2015-03-03 06:00:00+00:00,13775.2500000000)
  (2015-03-03 16:00:00+00:00,13828.2500000000)
  (2015-03-04 02:00:00+00:00,10750.5000000000)
  (2015-03-04 12:00:00+00:00,14120.2500000000)
  (2015-03-04 22:00:00+00:00,11431.7500000000)
  (2015-03-05 08:00:00+00:00,14347.0000000000)
  (2015-03-05 18:00:00+00:00,14845.5000000000)
  (2015-03-06 04:00:00+00:00,11905.0000000000)
  (2015-03-06 14:00:00+00:00,13834.2500000000)
  (2015-03-07 00:00:00+00:00,10359.5000000000)
  (2015-03-07 10:00:00+00:00,12680.5000000000)
  (2015-03-07 20:00:00+00:00,10794.2500000000)
  (2015-03-08 06:00:00+00:00,9398.5000000000)
  (2015-03-08 16:00:00+00:00,10490.0000000000)
  (2015-03-09 02:00:00+00:00,9729.7500000000)
  (2015-03-09 12:00:00+00:00,13556.2500000000)
  (2015-03-09 22:00:00+00:00,11223.2500000000)
  (2015-03-10 08:00:00+00:00,13898.2500000000)
  (2015-03-10 18:00:00+00:00,14196.0000000000)
  (2015-03-11 04:00:00+00:00,11447.0000000000)
  (2015-03-11 14:00:00+00:00,13614.0000000000)
  (2015-03-12 00:00:00+00:00,10509.7500000000)
  (2015-03-12 10:00:00+00:00,14561.2500000000)
  (2015-03-12 20:00:00+00:00,12513.7500000000)
  (2015-03-13 06:00:00+00:00,13569.2500000000)
  (2015-03-13 16:00:00+00:00,13468.5000000000)
  (2015-03-14 02:00:00+00:00,9790.5000000000)
  (2015-03-14 12:00:00+00:00,12059.2500000000)
  (2015-03-14 22:00:00+00:00,9964.0000000000)
  (2015-03-15 08:00:00+00:00,10890.0000000000)
  (2015-03-15 18:00:00+00:00,12192.2500000000)
  (2015-03-16 04:00:00+00:00,11093.2500000000)
  (2015-03-16 14:00:00+00:00,13178.0000000000)
  (2015-03-17 00:00:00+00:00,9938.0000000000)
  (2015-03-17 10:00:00+00:00,13625.7500000000)
  (2015-03-17 20:00:00+00:00,12289.5000000000)
  (2015-03-18 06:00:00+00:00,12973.0000000000)
  (2015-03-18 16:00:00+00:00,13208.5000000000)
  (2015-03-19 02:00:00+00:00,10368.7500000000)
  (2015-03-19 12:00:00+00:00,13573.7500000000)
  (2015-03-19 22:00:00+00:00,11017.2500000000)
  (2015-03-20 08:00:00+00:00,13784.0000000000)
  (2015-03-20 18:00:00+00:00,14098.5000000000)
  (2015-03-21 04:00:00+00:00,10258.7500000000)
  (2015-03-21 14:00:00+00:00,11815.2500000000)
  (2015-03-22 00:00:00+00:00,9351.5000000000)
  (2015-03-22 10:00:00+00:00,12051.0000000000)
  (2015-03-22 20:00:00+00:00,10988.5000000000)
  (2015-03-23 06:00:00+00:00,13307.2500000000)
  (2015-03-23 16:00:00+00:00,12932.5000000000)
  (2015-03-24 02:00:00+00:00,10343.2500000000)
  (2015-03-24 12:00:00+00:00,13642.2500000000)
  (2015-03-24 22:00:00+00:00,10994.2500000000)
  (2015-03-25 08:00:00+00:00,13340.7500000000)
  (2015-03-25 18:00:00+00:00,13959.7500000000)
  (2015-03-26 04:00:00+00:00,10807.5000000000)
  (2015-03-26 14:00:00+00:00,13174.0000000000)
  (2015-03-27 00:00:00+00:00,10295.7500000000)
  (2015-03-27 10:00:00+00:00,13866.5000000000)
  (2015-03-27 20:00:00+00:00,12161.0000000000)
  (2015-03-28 06:00:00+00:00,10269.2500000000)
  (2015-03-28 16:00:00+00:00,11246.0000000000)
  (2015-03-29 02:00:00+00:00,8870.0000000000)
  (2015-03-29 12:00:00+00:00,10763.2500000000)
  (2015-03-29 22:00:00+00:00,9759.7500000000)
  (2015-03-30 08:00:00+00:00,13947.7500000000)
  (2015-03-30 18:00:00+00:00,13503.0000000000)
  (2015-03-31 04:00:00+00:00,12652.7500000000)
  (2015-03-31 14:00:00+00:00,12972.0000000000)
  (2015-04-01 00:00:00+00:00,10087.2500000000)
  (2015-04-01 10:00:00+00:00,14079.0000000000)
  (2015-04-01 20:00:00+00:00,12141.0000000000)
  (2015-04-02 06:00:00+00:00,13985.5000000000)
  (2015-04-02 16:00:00+00:00,12964.0000000000)
  (2015-04-03 02:00:00+00:00,9575.7500000000)
  (2015-04-03 12:00:00+00:00,10403.7500000000)
  (2015-04-03 22:00:00+00:00,9740.7500000000)
  (2015-04-04 08:00:00+00:00,11945.0000000000)
  (2015-04-04 18:00:00+00:00,11525.0000000000)
  (2015-04-05 04:00:00+00:00,8805.0000000000)
  (2015-04-05 14:00:00+00:00,8928.2500000000)
  (2015-04-06 00:00:00+00:00,8577.0000000000)
  (2015-04-06 10:00:00+00:00,10515.2500000000)
  (2015-04-06 20:00:00+00:00,10741.2500000000)
  (2015-04-07 06:00:00+00:00,13760.2500000000)
  (2015-04-07 16:00:00+00:00,13122.5000000000)
  (2015-04-08 02:00:00+00:00,10302.0000000000)
  (2015-04-08 12:00:00+00:00,13453.5000000000)
  (2015-04-08 22:00:00+00:00,10149.5000000000)
  (2015-04-09 08:00:00+00:00,13532.0000000000)
  (2015-04-09 18:00:00+00:00,12791.0000000000)
  (2015-04-10 04:00:00+00:00,12069.5000000000)
  (2015-04-10 14:00:00+00:00,12500.7500000000)
  (2015-04-11 00:00:00+00:00,9329.2500000000)
  (2015-04-11 10:00:00+00:00,12111.0000000000)
  (2015-04-11 20:00:00+00:00,10093.7500000000)
  (2015-04-12 06:00:00+00:00,9700.7500000000)
  (2015-04-12 16:00:00+00:00,10588.2500000000)
  (2015-04-13 02:00:00+00:00,9691.0000000000)
  (2015-04-13 12:00:00+00:00,13224.0000000000)
  (2015-04-13 22:00:00+00:00,10553.0000000000)
  (2015-04-14 08:00:00+00:00,13752.0000000000)
  (2015-04-14 18:00:00+00:00,13259.2500000000)
  (2015-04-15 04:00:00+00:00,12381.5000000000)
  (2015-04-15 14:00:00+00:00,12690.2500000000)
  (2015-04-16 00:00:00+00:00,9607.7500000000)
  (2015-04-16 10:00:00+00:00,13586.5000000000)
  (2015-04-16 20:00:00+00:00,11669.0000000000)
  (2015-04-17 06:00:00+00:00,13573.2500000000)
  (2015-04-17 16:00:00+00:00,12780.5000000000)
  (2015-04-18 02:00:00+00:00,9433.7500000000)
  (2015-04-18 12:00:00+00:00,10729.2500000000)
  (2015-04-18 22:00:00+00:00,8859.0000000000)
  (2015-04-19 08:00:00+00:00,10667.2500000000)
  (2015-04-19 18:00:00+00:00,10595.5000000000)
  (2015-04-20 04:00:00+00:00,11681.5000000000)
  (2015-04-20 14:00:00+00:00,12480.0000000000)
  (2015-04-21 00:00:00+00:00,9065.7500000000)
  (2015-04-21 10:00:00+00:00,13287.5000000000)
  (2015-04-21 20:00:00+00:00,11396.7500000000)
  (2015-04-22 06:00:00+00:00,13210.7500000000)
  (2015-04-22 16:00:00+00:00,12518.7500000000)
  (2015-04-23 02:00:00+00:00,10038.5000000000)
  (2015-04-23 12:00:00+00:00,12684.2500000000)
  (2015-04-23 22:00:00+00:00,9769.0000000000)
  (2015-04-24 08:00:00+00:00,13235.5000000000)
  (2015-04-24 18:00:00+00:00,11959.5000000000)
  (2015-04-25 04:00:00+00:00,9114.7500000000)
  (2015-04-25 14:00:00+00:00,10655.5000000000)
  (2015-04-26 00:00:00+00:00,8032.5000000000)
  (2015-04-26 10:00:00+00:00,10685.2500000000)
  (2015-04-26 20:00:00+00:00,10470.5000000000)
  (2015-04-27 06:00:00+00:00,13651.7500000000)
  (2015-04-27 16:00:00+00:00,13430.0000000000)
  (2015-04-28 02:00:00+00:00,10313.5000000000)
  (2015-04-28 12:00:00+00:00,13746.7500000000)
  (2015-04-28 22:00:00+00:00,10527.0000000000)
  (2015-04-29 08:00:00+00:00,13706.0000000000)
  (2015-04-29 18:00:00+00:00,12843.7500000000)
  (2015-04-30 04:00:00+00:00,12077.0000000000)
  (2015-04-30 14:00:00+00:00,12743.5000000000)
  (2015-05-01 00:00:00+00:00,8308.5000000000)
  (2015-05-01 10:00:00+00:00,10310.5000000000)
  (2015-05-01 20:00:00+00:00,9391.2500000000)
  (2015-05-02 06:00:00+00:00,10245.2500000000)
  (2015-05-02 16:00:00+00:00,10380.2500000000)
  (2015-05-03 02:00:00+00:00,8233.5000000000)
  (2015-05-03 12:00:00+00:00,9969.0000000000)
  (2015-05-03 22:00:00+00:00,8763.7500000000)
  (2015-05-04 08:00:00+00:00,13628.7500000000)
  (2015-05-04 18:00:00+00:00,12349.7500000000)
  (2015-05-05 04:00:00+00:00,11802.2500000000)
  (2015-05-05 14:00:00+00:00,12988.7500000000)
  (2015-05-06 00:00:00+00:00,9092.7500000000)
  (2015-05-06 10:00:00+00:00,13569.7500000000)
  (2015-05-06 20:00:00+00:00,11655.7500000000)
  (2015-05-07 06:00:00+00:00,13343.7500000000)
  (2015-05-07 16:00:00+00:00,12844.2500000000)
  (2015-05-08 02:00:00+00:00,9657.7500000000)
  (2015-05-08 12:00:00+00:00,12363.0000000000)
  (2015-05-08 22:00:00+00:00,9461.0000000000)
  (2015-05-09 08:00:00+00:00,11453.0000000000)
  (2015-05-09 18:00:00+00:00,10174.5000000000)
  (2015-05-10 04:00:00+00:00,8147.7500000000)
  (2015-05-10 14:00:00+00:00,9502.0000000000)
  (2015-05-11 00:00:00+00:00,8776.7500000000)
  (2015-05-11 10:00:00+00:00,13193.7500000000)
  (2015-05-11 20:00:00+00:00,11483.5000000000)
  (2015-05-12 06:00:00+00:00,13055.0000000000)
  (2015-05-12 16:00:00+00:00,13061.2500000000)
  (2015-05-13 02:00:00+00:00,9418.5000000000)
  (2015-05-13 12:00:00+00:00,12478.5000000000)
  (2015-05-13 22:00:00+00:00,9117.2500000000)
  (2015-05-14 08:00:00+00:00,10325.5000000000)
  (2015-05-14 18:00:00+00:00,9923.2500000000)
  (2015-05-15 04:00:00+00:00,9850.7500000000)
  (2015-05-15 14:00:00+00:00,11183.5000000000)
  (2015-05-16 00:00:00+00:00,7917.2500000000)
  (2015-05-16 10:00:00+00:00,11054.0000000000)
  (2015-05-16 20:00:00+00:00,9658.5000000000)
  (2015-05-17 06:00:00+00:00,9330.2500000000)
  (2015-05-17 16:00:00+00:00,10181.0000000000)
  (2015-05-18 02:00:00+00:00,8975.5000000000)
  (2015-05-18 12:00:00+00:00,13147.5000000000)
  (2015-05-18 22:00:00+00:00,9680.5000000000)
  (2015-05-19 08:00:00+00:00,13631.0000000000)
  (2015-05-19 18:00:00+00:00,12301.2500000000)
  (2015-05-20 04:00:00+00:00,11692.7500000000)
  (2015-05-20 14:00:00+00:00,12719.0000000000)
  (2015-05-21 00:00:00+00:00,9475.7500000000)
  (2015-05-21 10:00:00+00:00,13450.2500000000)
  (2015-05-21 20:00:00+00:00,11430.0000000000)
  (2015-05-22 06:00:00+00:00,13343.7500000000)
  (2015-05-22 16:00:00+00:00,12211.7500000000)
  (2015-05-23 02:00:00+00:00,8770.7500000000)
  (2015-05-23 12:00:00+00:00,10430.7500000000)
  (2015-05-23 22:00:00+00:00,8368.2500000000)
  (2015-05-24 08:00:00+00:00,9660.7500000000)
  (2015-05-24 18:00:00+00:00,9316.2500000000)
  (2015-05-25 04:00:00+00:00,7563.2500000000)
  (2015-05-25 14:00:00+00:00,9623.5000000000)
  (2015-05-26 00:00:00+00:00,8432.2500000000)
  (2015-05-26 10:00:00+00:00,13559.0000000000)
  (2015-05-26 20:00:00+00:00,11355.2500000000)
  (2015-05-27 06:00:00+00:00,13286.7500000000)
  (2015-05-27 16:00:00+00:00,13038.0000000000)
  (2015-05-28 02:00:00+00:00,9584.5000000000)
  (2015-05-28 12:00:00+00:00,13083.7500000000)
  (2015-05-28 22:00:00+00:00,9601.7500000000)
  (2015-05-29 08:00:00+00:00,13209.5000000000)
  (2015-05-29 18:00:00+00:00,11759.2500000000)
  (2015-05-30 04:00:00+00:00,9090.2500000000)
  (2015-05-30 14:00:00+00:00,10674.0000000000)
  (2015-05-31 00:00:00+00:00,7956.0000000000)
  (2015-05-31 10:00:00+00:00,10601.5000000000)
  (2015-05-31 20:00:00+00:00,9941.2500000000)
  (2015-06-01 06:00:00+00:00,12999.2500000000)
  (2015-06-01 16:00:00+00:00,12645.7500000000)
  (2015-06-02 02:00:00+00:00,9323.5000000000)
  (2015-06-02 12:00:00+00:00,13429.0000000000)
  (2015-06-02 22:00:00+00:00,9784.2500000000)
  (2015-06-03 08:00:00+00:00,13295.0000000000)
  (2015-06-03 18:00:00+00:00,12305.7500000000)
  (2015-06-04 04:00:00+00:00,11501.2500000000)
  (2015-06-04 14:00:00+00:00,12370.0000000000)
  (2015-06-05 00:00:00+00:00,9194.2500000000)
  (2015-06-05 10:00:00+00:00,13348.2500000000)
  (2015-06-05 20:00:00+00:00,11323.0000000000)
  (2015-06-06 06:00:00+00:00,10837.5000000000)
  (2015-06-06 16:00:00+00:00,11078.7500000000)
  (2015-06-07 02:00:00+00:00,7741.5000000000)
  (2015-06-07 12:00:00+00:00,9800.7500000000)
  (2015-06-07 22:00:00+00:00,9377.0000000000)
  (2015-06-08 08:00:00+00:00,13507.7500000000)
  (2015-06-08 18:00:00+00:00,12444.0000000000)
  (2015-06-09 04:00:00+00:00,11660.0000000000)
  (2015-06-09 14:00:00+00:00,12808.5000000000)
  (2015-06-10 00:00:00+00:00,9491.2500000000)
  (2015-06-10 10:00:00+00:00,13553.0000000000)
  (2015-06-10 20:00:00+00:00,11387.5000000000)
  (2015-06-11 06:00:00+00:00,12968.2500000000)
  (2015-06-11 16:00:00+00:00,12737.7500000000)
  (2015-06-12 02:00:00+00:00,9562.0000000000)
  (2015-06-12 12:00:00+00:00,13112.0000000000)
  (2015-06-12 22:00:00+00:00,9556.0000000000)
  (2015-06-13 08:00:00+00:00,12263.7500000000)
  (2015-06-13 18:00:00+00:00,10731.5000000000)
  (2015-06-14 04:00:00+00:00,7888.5000000000)
  (2015-06-14 14:00:00+00:00,9500.0000000000)
  (2015-06-15 00:00:00+00:00,8444.0000000000)
  (2015-06-15 10:00:00+00:00,13474.5000000000)
  (2015-06-15 20:00:00+00:00,11293.2500000000)
  (2015-06-16 06:00:00+00:00,13268.2500000000)
  (2015-06-16 16:00:00+00:00,12311.0000000000)
  (2015-06-17 02:00:00+00:00,9212.0000000000)
  (2015-06-17 12:00:00+00:00,12695.2500000000)
  (2015-06-17 22:00:00+00:00,9540.5000000000)
  (2015-06-18 08:00:00+00:00,13326.7500000000)
  (2015-06-18 18:00:00+00:00,12162.7500000000)
  (2015-06-19 04:00:00+00:00,11411.0000000000)
  (2015-06-19 14:00:00+00:00,12835.7500000000)
  (2015-06-20 00:00:00+00:00,9051.0000000000)
  (2015-06-20 10:00:00+00:00,12053.5000000000)
  (2015-06-20 20:00:00+00:00,10016.0000000000)
  (2015-06-21 06:00:00+00:00,9455.5000000000)
  (2015-06-21 16:00:00+00:00,10529.5000000000)
  (2015-06-22 02:00:00+00:00,9259.7500000000)
  (2015-06-22 12:00:00+00:00,13222.0000000000)
  (2015-06-22 22:00:00+00:00,9699.2500000000)
  (2015-06-23 08:00:00+00:00,13824.5000000000)
  (2015-06-23 18:00:00+00:00,12420.2500000000)
  (2015-06-24 04:00:00+00:00,11913.5000000000)
  (2015-06-24 14:00:00+00:00,12968.0000000000)
  (2015-06-25 00:00:00+00:00,9440.2500000000)
  (2015-06-25 10:00:00+00:00,13472.0000000000)
  (2015-06-25 20:00:00+00:00,11310.7500000000)
  (2015-06-26 06:00:00+00:00,13324.7500000000)
  (2015-06-26 16:00:00+00:00,12589.7500000000)
  (2015-06-27 02:00:00+00:00,8823.5000000000)
  (2015-06-27 12:00:00+00:00,11206.0000000000)
  (2015-06-27 22:00:00+00:00,8606.0000000000)
  (2015-06-28 08:00:00+00:00,10121.7500000000)
  (2015-06-28 18:00:00+00:00,10261.0000000000)
  (2015-06-29 04:00:00+00:00,11306.0000000000)
  (2015-06-29 14:00:00+00:00,12942.2500000000)
  (2015-06-30 00:00:00+00:00,9283.5000000000)
  (2015-06-30 10:00:00+00:00,13567.7500000000)
  (2015-06-30 20:00:00+00:00,11299.5000000000)
  (2015-07-01 06:00:00+00:00,13058.5000000000)
  (2015-07-01 16:00:00+00:00,13059.0000000000)
  (2015-07-02 02:00:00+00:00,9892.0000000000)
  (2015-07-02 12:00:00+00:00,13510.5000000000)
  (2015-07-02 22:00:00+00:00,10403.2500000000)
  (2015-07-03 08:00:00+00:00,13416.5000000000)
  (2015-07-03 18:00:00+00:00,11962.0000000000)
  (2015-07-04 04:00:00+00:00,9250.2500000000)
  (2015-07-04 14:00:00+00:00,10862.5000000000)
  (2015-07-05 00:00:00+00:00,8024.7500000000)
  (2015-07-05 10:00:00+00:00,10833.7500000000)
  (2015-07-05 20:00:00+00:00,10364.2500000000)
  (2015-07-06 06:00:00+00:00,12988.5000000000)
  (2015-07-06 16:00:00+00:00,12668.7500000000)
  (2015-07-07 02:00:00+00:00,9646.7500000000)
  (2015-07-07 12:00:00+00:00,13317.2500000000)
  (2015-07-07 22:00:00+00:00,9888.0000000000)
  (2015-07-08 08:00:00+00:00,13431.0000000000)
  (2015-07-08 18:00:00+00:00,11928.2500000000)
  (2015-07-09 04:00:00+00:00,11079.7500000000)
  (2015-07-09 14:00:00+00:00,12516.2500000000)
  (2015-07-10 00:00:00+00:00,9596.2500000000)
  (2015-07-10 10:00:00+00:00,13440.5000000000)
  (2015-07-10 20:00:00+00:00,10702.5000000000)
  (2015-07-11 06:00:00+00:00,10679.7500000000)
  (2015-07-11 16:00:00+00:00,10645.0000000000)
  (2015-07-12 02:00:00+00:00,7967.7500000000)
  (2015-07-12 12:00:00+00:00,9934.0000000000)
  (2015-07-12 22:00:00+00:00,9143.5000000000)
  (2015-07-13 08:00:00+00:00,13600.2500000000)
  (2015-07-13 18:00:00+00:00,12265.7500000000)
  (2015-07-14 04:00:00+00:00,11379.5000000000)
  (2015-07-14 14:00:00+00:00,12974.7500000000)
  (2015-07-15 00:00:00+00:00,9367.5000000000)
  (2015-07-15 10:00:00+00:00,13516.7500000000)
  (2015-07-15 20:00:00+00:00,11299.0000000000)
  (2015-07-16 06:00:00+00:00,12955.7500000000)
  (2015-07-16 16:00:00+00:00,12208.0000000000)
  (2015-07-17 02:00:00+00:00,9465.2500000000)
  (2015-07-17 12:00:00+00:00,12918.7500000000)
  (2015-07-17 22:00:00+00:00,9303.7500000000)
  (2015-07-18 08:00:00+00:00,11880.7500000000)
  (2015-07-18 18:00:00+00:00,10162.5000000000)
  (2015-07-19 04:00:00+00:00,8097.7500000000)
  (2015-07-19 14:00:00+00:00,10008.0000000000)
  (2015-07-20 00:00:00+00:00,8368.0000000000)
  (2015-07-20 10:00:00+00:00,12936.5000000000)
  (2015-07-20 20:00:00+00:00,11030.7500000000)
  (2015-07-21 06:00:00+00:00,13051.5000000000)
  (2015-07-21 16:00:00+00:00,12760.2500000000)
  (2015-07-22 02:00:00+00:00,9584.2500000000)
  (2015-07-22 12:00:00+00:00,13016.2500000000)
  (2015-07-22 22:00:00+00:00,9984.2500000000)
  (2015-07-23 08:00:00+00:00,13195.7500000000)
  (2015-07-23 18:00:00+00:00,11879.2500000000)
  (2015-07-24 04:00:00+00:00,11128.7500000000)
  (2015-07-24 14:00:00+00:00,12200.0000000000)
  (2015-07-25 00:00:00+00:00,8596.2500000000)
  (2015-07-25 10:00:00+00:00,11434.5000000000)
  (2015-07-25 20:00:00+00:00,9485.2500000000)
  (2015-07-26 06:00:00+00:00,9031.7500000000)
  (2015-07-26 16:00:00+00:00,9545.7500000000)
  (2015-07-27 02:00:00+00:00,8150.0000000000)
  (2015-07-27 12:00:00+00:00,12547.2500000000)
  (2015-07-27 22:00:00+00:00,9083.7500000000)
  (2015-07-28 08:00:00+00:00,13039.2500000000)
  (2015-07-28 18:00:00+00:00,11464.5000000000)
  (2015-07-29 04:00:00+00:00,10526.7500000000)
  (2015-07-29 14:00:00+00:00,11957.2500000000)
  (2015-07-30 00:00:00+00:00,8898.7500000000)
  (2015-07-30 10:00:00+00:00,12996.7500000000)
  (2015-07-30 20:00:00+00:00,10550.5000000000)
  (2015-07-31 06:00:00+00:00,12600.2500000000)
  (2015-07-31 16:00:00+00:00,11989.7500000000)
  (2015-08-01 02:00:00+00:00,8044.7500000000)
  (2015-08-01 12:00:00+00:00,10437.5000000000)
  (2015-08-01 22:00:00+00:00,8392.7500000000)
  (2015-08-02 08:00:00+00:00,9860.7500000000)
  (2015-08-02 18:00:00+00:00,9914.7500000000)
  (2015-08-03 04:00:00+00:00,10450.5000000000)
  (2015-08-03 14:00:00+00:00,12681.2500000000)
  (2015-08-04 00:00:00+00:00,9145.7500000000)
  (2015-08-04 10:00:00+00:00,13472.7500000000)
  (2015-08-04 20:00:00+00:00,10969.5000000000)
  (2015-08-05 06:00:00+00:00,12799.5000000000)
  (2015-08-05 16:00:00+00:00,12518.2500000000)
  (2015-08-06 02:00:00+00:00,9424.2500000000)
  (2015-08-06 12:00:00+00:00,12914.5000000000)
  (2015-08-06 22:00:00+00:00,9959.0000000000)
  (2015-08-07 08:00:00+00:00,13600.7500000000)
  (2015-08-07 18:00:00+00:00,11375.2500000000)
  (2015-08-08 04:00:00+00:00,9116.7500000000)
  (2015-08-08 14:00:00+00:00,10417.5000000000)
  (2015-08-09 00:00:00+00:00,8161.7500000000)
  (2015-08-09 10:00:00+00:00,10566.7500000000)
  (2015-08-09 20:00:00+00:00,10156.7500000000)
  (2015-08-10 06:00:00+00:00,12569.0000000000)
  (2015-08-10 16:00:00+00:00,12684.2500000000)
  (2015-08-11 02:00:00+00:00,9551.0000000000)
  (2015-08-11 12:00:00+00:00,12902.5000000000)
  (2015-08-11 22:00:00+00:00,9854.7500000000)
  (2015-08-12 08:00:00+00:00,13291.7500000000)
  (2015-08-12 18:00:00+00:00,12298.2500000000)
  (2015-08-13 04:00:00+00:00,11278.0000000000)
  (2015-08-13 14:00:00+00:00,12463.2500000000)
  (2015-08-14 00:00:00+00:00,9027.5000000000)
  (2015-08-14 10:00:00+00:00,13279.2500000000)
  (2015-08-14 20:00:00+00:00,10659.7500000000)
  (2015-08-15 06:00:00+00:00,10495.0000000000)
  (2015-08-15 16:00:00+00:00,10591.0000000000)
  (2015-08-16 02:00:00+00:00,7870.7500000000)
  (2015-08-16 12:00:00+00:00,9816.7500000000)
  (2015-08-16 22:00:00+00:00,9121.2500000000)
  (2015-08-17 08:00:00+00:00,13601.7500000000)
  (2015-08-17 18:00:00+00:00,12710.0000000000)
  (2015-08-18 04:00:00+00:00,11515.0000000000)
  (2015-08-18 14:00:00+00:00,12802.2500000000)
  (2015-08-19 00:00:00+00:00,9335.2500000000)
  (2015-08-19 10:00:00+00:00,13597.5000000000)
  (2015-08-19 20:00:00+00:00,11064.0000000000)
  (2015-08-20 06:00:00+00:00,12969.5000000000)
  (2015-08-20 16:00:00+00:00,12619.2500000000)
  (2015-08-21 02:00:00+00:00,9393.2500000000)
  (2015-08-21 12:00:00+00:00,12815.7500000000)
  (2015-08-21 22:00:00+00:00,9370.5000000000)
  (2015-08-22 08:00:00+00:00,11658.2500000000)
  (2015-08-22 18:00:00+00:00,10444.2500000000)
  (2015-08-23 04:00:00+00:00,7533.0000000000)
  (2015-08-23 14:00:00+00:00,9493.7500000000)
  (2015-08-24 00:00:00+00:00,8495.2500000000)
  (2015-08-24 10:00:00+00:00,13652.7500000000)
  (2015-08-24 20:00:00+00:00,11088.2500000000)
  (2015-08-25 06:00:00+00:00,12824.5000000000)
  (2015-08-25 16:00:00+00:00,12757.7500000000)
  (2015-08-26 02:00:00+00:00,9372.0000000000)
  (2015-08-26 12:00:00+00:00,13099.0000000000)
  (2015-08-26 22:00:00+00:00,9824.0000000000)
  (2015-08-27 08:00:00+00:00,13577.0000000000)
  (2015-08-27 18:00:00+00:00,12979.5000000000)
  (2015-08-28 04:00:00+00:00,11719.5000000000)
  (2015-08-28 14:00:00+00:00,12566.0000000000)
  (2015-08-29 00:00:00+00:00,8845.0000000000)
  (2015-08-29 10:00:00+00:00,11572.0000000000)
  (2015-08-29 20:00:00+00:00,9692.0000000000)
  (2015-08-30 06:00:00+00:00,9271.7500000000)
  (2015-08-30 16:00:00+00:00,10149.2500000000)
  (2015-08-31 02:00:00+00:00,9166.5000000000)
  (2015-08-31 12:00:00+00:00,13197.2500000000)
  (2015-08-31 22:00:00+00:00,10118.7500000000)
  (2015-09-01 08:00:00+00:00,13495.0000000000)
  (2015-09-01 18:00:00+00:00,12727.0000000000)
  (2015-09-02 04:00:00+00:00,12039.7500000000)
  (2015-09-02 14:00:00+00:00,12357.7500000000)
  (2015-09-03 00:00:00+00:00,9165.5000000000)
  (2015-09-03 10:00:00+00:00,13345.0000000000)
  (2015-09-03 20:00:00+00:00,11134.0000000000)
  (2015-09-04 06:00:00+00:00,13235.7500000000)
  (2015-09-04 16:00:00+00:00,12328.2500000000)
  (2015-09-05 02:00:00+00:00,8546.2500000000)
  (2015-09-05 12:00:00+00:00,10825.7500000000)
  (2015-09-05 22:00:00+00:00,8593.5000000000)
  (2015-09-06 08:00:00+00:00,10749.0000000000)
  (2015-09-06 18:00:00+00:00,10843.5000000000)
  (2015-09-07 04:00:00+00:00,11257.7500000000)
  (2015-09-07 14:00:00+00:00,12318.2500000000)
  (2015-09-08 00:00:00+00:00,9233.5000000000)
  (2015-09-08 10:00:00+00:00,13496.0000000000)
  (2015-09-08 20:00:00+00:00,11133.7500000000)
  (2015-09-09 06:00:00+00:00,13108.7500000000)
  (2015-09-09 16:00:00+00:00,12606.5000000000)
  (2015-09-10 02:00:00+00:00,9662.5000000000)
  (2015-09-10 12:00:00+00:00,12793.2500000000)
  (2015-09-10 22:00:00+00:00,9919.0000000000)
  (2015-09-11 08:00:00+00:00,13299.7500000000)
  (2015-09-11 18:00:00+00:00,12708.0000000000)
  (2015-09-12 04:00:00+00:00,9454.2500000000)
  (2015-09-12 14:00:00+00:00,10561.2500000000)
  (2015-09-13 00:00:00+00:00,8027.7500000000)
  (2015-09-13 10:00:00+00:00,10731.0000000000)
  (2015-09-13 20:00:00+00:00,9919.5000000000)
  (2015-09-14 06:00:00+00:00,13152.0000000000)
  (2015-09-14 16:00:00+00:00,12791.7500000000)
  (2015-09-15 02:00:00+00:00,9781.0000000000)
  (2015-09-15 12:00:00+00:00,12834.5000000000)
  (2015-09-15 22:00:00+00:00,9992.2500000000)
  (2015-09-16 08:00:00+00:00,13351.5000000000)
  (2015-09-16 18:00:00+00:00,12827.0000000000)
  (2015-09-17 04:00:00+00:00,12062.2500000000)
  (2015-09-17 14:00:00+00:00,12527.0000000000)
  (2015-09-18 00:00:00+00:00,9028.2500000000)
  (2015-09-18 10:00:00+00:00,13026.7500000000)
  (2015-09-18 20:00:00+00:00,10625.2500000000)
  (2015-09-19 06:00:00+00:00,10760.7500000000)
  (2015-09-19 16:00:00+00:00,10905.0000000000)
  (2015-09-20 02:00:00+00:00,7913.5000000000)
  (2015-09-20 12:00:00+00:00,10245.2500000000)
  (2015-09-20 22:00:00+00:00,8856.2500000000)
  (2015-09-21 08:00:00+00:00,13213.7500000000)
  (2015-09-21 18:00:00+00:00,13107.0000000000)
  (2015-09-22 04:00:00+00:00,12369.5000000000)
  (2015-09-22 14:00:00+00:00,12847.0000000000)
  (2015-09-23 00:00:00+00:00,9411.0000000000)
  (2015-09-23 10:00:00+00:00,13512.2500000000)
  (2015-09-23 20:00:00+00:00,11354.2500000000)
  (2015-09-24 06:00:00+00:00,12936.7500000000)
  (2015-09-24 16:00:00+00:00,12396.0000000000)
  (2015-09-25 02:00:00+00:00,9787.5000000000)
  (2015-09-25 12:00:00+00:00,12865.7500000000)
  (2015-09-25 22:00:00+00:00,9472.7500000000)
  (2015-09-26 08:00:00+00:00,11974.5000000000)
  (2015-09-26 18:00:00+00:00,11347.0000000000)
  (2015-09-27 04:00:00+00:00,8603.5000000000)
  (2015-09-27 14:00:00+00:00,10073.2500000000)
  (2015-09-28 00:00:00+00:00,8715.0000000000)
  (2015-09-28 10:00:00+00:00,13394.0000000000)
  (2015-09-28 20:00:00+00:00,11304.2500000000)
  (2015-09-29 06:00:00+00:00,13588.0000000000)
  (2015-09-29 16:00:00+00:00,13095.2500000000)
  (2015-09-30 02:00:00+00:00,10166.5000000000)
  (2015-09-30 12:00:00+00:00,12860.5000000000)
  (2015-09-30 22:00:00+00:00,9936.2500000000)
  (2015-10-01 08:00:00+00:00,11158.7500000000)
  (2015-10-01 18:00:00+00:00,12966.2500000000)
  (2015-10-02 04:00:00+00:00,12380.5000000000)
  (2015-10-02 14:00:00+00:00,10775.5000000000)
  (2015-10-03 00:00:00+00:00,8792.5000000000)
  (2015-10-03 10:00:00+00:00,8298.5000000000)
  (2015-10-03 20:00:00+00:00,9631.0000000000)
  (2015-10-04 06:00:00+00:00,9173.5000000000)
  (2015-10-04 16:00:00+00:00,10768.5000000000)
  (2015-10-05 02:00:00+00:00,9169.2500000000)
  (2015-10-05 12:00:00+00:00,10764.5000000000)
  (2015-10-05 22:00:00+00:00,10092.2500000000)
  (2015-10-06 08:00:00+00:00,13385.7500000000)
  (2015-10-06 18:00:00+00:00,13113.7500000000)
  (2015-10-07 04:00:00+00:00,12626.5000000000)
  (2015-10-07 14:00:00+00:00,12920.7500000000)
  (2015-10-08 00:00:00+00:00,9428.2500000000)
  (2015-10-08 10:00:00+00:00,13379.7500000000)
  (2015-10-08 20:00:00+00:00,11526.2500000000)
  (2015-10-09 06:00:00+00:00,13712.7500000000)
  (2015-10-09 16:00:00+00:00,13109.0000000000)
  (2015-10-10 02:00:00+00:00,9355.5000000000)
  (2015-10-10 12:00:00+00:00,8811.7500000000)
  (2015-10-10 22:00:00+00:00,9022.7500000000)
  (2015-10-11 08:00:00+00:00,8823.2500000000)
  (2015-10-11 18:00:00+00:00,11267.2500000000)
  (2015-10-12 04:00:00+00:00,12419.5000000000)
  (2015-10-12 14:00:00+00:00,11721.2500000000)
  (2015-10-13 00:00:00+00:00,9768.5000000000)
  (2015-10-13 10:00:00+00:00,13726.7500000000)
  (2015-10-13 20:00:00+00:00,11855.0000000000)
  (2015-10-14 06:00:00+00:00,14598.2500000000)
  (2015-10-14 16:00:00+00:00,14458.7500000000)
  (2015-10-15 02:00:00+00:00,10624.0000000000)
  (2015-10-15 12:00:00+00:00,13996.2500000000)
  (2015-10-15 22:00:00+00:00,10481.7500000000)
  (2015-10-16 08:00:00+00:00,14254.0000000000)
  (2015-10-16 18:00:00+00:00,13235.0000000000)
  (2015-10-17 04:00:00+00:00,10091.2500000000)
  (2015-10-17 14:00:00+00:00,11381.2500000000)
  (2015-10-18 00:00:00+00:00,8924.7500000000)
  (2015-10-18 10:00:00+00:00,10728.2500000000)
  (2015-10-18 20:00:00+00:00,10683.7500000000)
  (2015-10-19 06:00:00+00:00,13887.0000000000)
  (2015-10-19 16:00:00+00:00,14050.5000000000)
  (2015-10-20 02:00:00+00:00,10391.5000000000)
  (2015-10-20 12:00:00+00:00,13468.7500000000)
  (2015-10-20 22:00:00+00:00,10361.0000000000)
  (2015-10-21 08:00:00+00:00,13814.0000000000)
  (2015-10-21 18:00:00+00:00,13481.5000000000)
  (2015-10-22 04:00:00+00:00,12593.2500000000)
  (2015-10-22 14:00:00+00:00,12935.0000000000)
  (2015-10-23 00:00:00+00:00,9826.2500000000)
  (2015-10-23 10:00:00+00:00,12041.2500000000)
  (2015-10-23 20:00:00+00:00,11421.7500000000)
  (2015-10-24 06:00:00+00:00,11265.5000000000)
  (2015-10-24 16:00:00+00:00,12523.5000000000)
  (2015-10-25 02:00:00+00:00,8871.5000000000)
  (2015-10-25 12:00:00+00:00,10253.7500000000)
  (2015-10-25 22:00:00+00:00,9988.7500000000)
  (2015-10-26 08:00:00+00:00,12961.5000000000)
  (2015-10-26 18:00:00+00:00,13973.7500000000)
  (2015-10-27 04:00:00+00:00,11434.2500000000)
  (2015-10-27 14:00:00+00:00,12683.7500000000)
  (2015-10-28 00:00:00+00:00,10316.2500000000)
  (2015-10-28 10:00:00+00:00,12420.5000000000)
  (2015-10-28 20:00:00+00:00,12750.2500000000)
  (2015-10-29 06:00:00+00:00,13399.2500000000)
  (2015-10-29 16:00:00+00:00,14366.5000000000)
  (2015-10-30 02:00:00+00:00,10204.5000000000)
  (2015-10-30 12:00:00+00:00,12919.0000000000)
  (2015-10-30 22:00:00+00:00,10461.0000000000)
  (2015-10-31 08:00:00+00:00,9980.7500000000)
  (2015-10-31 18:00:00+00:00,11519.5000000000)
  (2015-11-01 04:00:00+00:00,8676.7500000000)
  (2015-11-01 14:00:00+00:00,10338.0000000000)
  (2015-11-02 00:00:00+00:00,9420.0000000000)
  (2015-11-02 10:00:00+00:00,13886.0000000000)
  (2015-11-02 20:00:00+00:00,12589.0000000000)
  (2015-11-03 06:00:00+00:00,13758.2500000000)
  (2015-11-03 16:00:00+00:00,14562.7500000000)
  (2015-11-04 02:00:00+00:00,10566.2500000000)
  (2015-11-04 12:00:00+00:00,13825.7500000000)
  (2015-11-04 22:00:00+00:00,11022.2500000000)
  (2015-11-05 08:00:00+00:00,13426.7500000000)
  (2015-11-05 18:00:00+00:00,14268.5000000000)
  (2015-11-06 04:00:00+00:00,11181.2500000000)
  (2015-11-06 14:00:00+00:00,12902.0000000000)
  (2015-11-07 00:00:00+00:00,9187.0000000000)
  (2015-11-07 10:00:00+00:00,12545.2500000000)
  (2015-11-07 20:00:00+00:00,10590.0000000000)
  (2015-11-08 06:00:00+00:00,9257.2500000000)
  (2015-11-08 16:00:00+00:00,12052.7500000000)
  (2015-11-09 02:00:00+00:00,9074.5000000000)
  (2015-11-09 12:00:00+00:00,13635.5000000000)
  (2015-11-09 22:00:00+00:00,10472.2500000000)
  (2015-11-10 08:00:00+00:00,13926.0000000000)
  (2015-11-10 18:00:00+00:00,14132.0000000000)
  (2015-11-11 04:00:00+00:00,10975.2500000000)
  (2015-11-11 14:00:00+00:00,13514.5000000000)
  (2015-11-12 00:00:00+00:00,9747.7500000000)
  (2015-11-12 10:00:00+00:00,13845.0000000000)
  (2015-11-12 20:00:00+00:00,12280.7500000000)
  (2015-11-13 06:00:00+00:00,13666.0000000000)
  (2015-11-13 16:00:00+00:00,14165.0000000000)
  (2015-11-14 02:00:00+00:00,9574.7500000000)
  (2015-11-14 12:00:00+00:00,11789.0000000000)
  (2015-11-14 22:00:00+00:00,9552.5000000000)
  (2015-11-15 08:00:00+00:00,10870.5000000000)
  (2015-11-15 18:00:00+00:00,11774.7500000000)
  (2015-11-16 04:00:00+00:00,10538.5000000000)
  (2015-11-16 14:00:00+00:00,13439.7500000000)
  (2015-11-17 00:00:00+00:00,9608.7500000000)
  (2015-11-17 10:00:00+00:00,13990.5000000000)
  (2015-11-17 20:00:00+00:00,12247.0000000000)
  (2015-11-18 06:00:00+00:00,12933.0000000000)
  (2015-11-18 16:00:00+00:00,13695.0000000000)
  (2015-11-19 02:00:00+00:00,9923.0000000000)
  (2015-11-19 12:00:00+00:00,13699.7500000000)
  (2015-11-19 22:00:00+00:00,10699.0000000000)
  (2015-11-20 08:00:00+00:00,14013.0000000000)
  (2015-11-20 18:00:00+00:00,13977.5000000000)
  (2015-11-21 04:00:00+00:00,9714.0000000000)
  (2015-11-21 14:00:00+00:00,12018.2500000000)
  (2015-11-22 00:00:00+00:00,9697.7500000000)
  (2015-11-22 10:00:00+00:00,12349.0000000000)
  (2015-11-22 20:00:00+00:00,11170.2500000000)
  (2015-11-23 06:00:00+00:00,14037.7500000000)
  (2015-11-23 16:00:00+00:00,15296.5000000000)
  (2015-11-24 02:00:00+00:00,11177.2500000000)
  (2015-11-24 12:00:00+00:00,14177.0000000000)
  (2015-11-24 22:00:00+00:00,11456.5000000000)
  (2015-11-25 08:00:00+00:00,14961.5000000000)
  (2015-11-25 18:00:00+00:00,14982.0000000000)
  (2015-11-26 04:00:00+00:00,11821.5000000000)
  (2015-11-26 14:00:00+00:00,14436.2500000000)
  (2015-11-27 00:00:00+00:00,10896.7500000000)
  (2015-11-27 10:00:00+00:00,14434.7500000000)
  (2015-11-27 20:00:00+00:00,12514.2500000000)
  (2015-11-28 06:00:00+00:00,10965.2500000000)
  (2015-11-28 16:00:00+00:00,13224.5000000000)
  (2015-11-29 02:00:00+00:00,9526.7500000000)
  (2015-11-29 12:00:00+00:00,11695.0000000000)
  (2015-11-29 22:00:00+00:00,10410.7500000000)
  (2015-11-30 08:00:00+00:00,14503.7500000000)
  (2015-11-30 18:00:00+00:00,14688.7500000000)
  (2015-12-01 04:00:00+00:00,11642.2500000000)
  (2015-12-01 14:00:00+00:00,14347.7500000000)
  (2015-12-02 00:00:00+00:00,10265.7500000000)
  (2015-12-02 10:00:00+00:00,14620.2500000000)
  (2015-12-02 20:00:00+00:00,12895.5000000000)
  (2015-12-03 06:00:00+00:00,14061.5000000000)
  (2015-12-03 16:00:00+00:00,14954.7500000000)
  (2015-12-04 02:00:00+00:00,10398.2500000000)
  (2015-12-04 12:00:00+00:00,14136.5000000000)
  (2015-12-04 22:00:00+00:00,10827.5000000000)
  (2015-12-05 08:00:00+00:00,12184.7500000000)
  (2015-12-05 18:00:00+00:00,12794.0000000000)
  (2015-12-06 04:00:00+00:00,9219.7500000000)
  (2015-12-06 14:00:00+00:00,10986.2500000000)
  (2015-12-07 00:00:00+00:00,9188.2500000000)
  (2015-12-07 10:00:00+00:00,14536.7500000000)
  (2015-12-07 20:00:00+00:00,12983.7500000000)
  (2015-12-08 06:00:00+00:00,14017.2500000000)
  (2015-12-08 16:00:00+00:00,14895.2500000000)
  (2015-12-09 02:00:00+00:00,10452.7500000000)
  (2015-12-09 12:00:00+00:00,14342.5000000000)
  (2015-12-09 22:00:00+00:00,11158.2500000000)
  (2015-12-10 08:00:00+00:00,14179.7500000000)
  (2015-12-10 18:00:00+00:00,14679.7500000000)
  (2015-12-11 04:00:00+00:00,11640.2500000000)
  (2015-12-11 14:00:00+00:00,14362.0000000000)
  (2015-12-12 00:00:00+00:00,10070.0000000000)
  (2015-12-12 10:00:00+00:00,13063.0000000000)
  (2015-12-12 20:00:00+00:00,11166.5000000000)
  (2015-12-13 06:00:00+00:00,9566.7500000000)
  (2015-12-13 16:00:00+00:00,12687.5000000000)
  (2015-12-14 02:00:00+00:00,10354.0000000000)
  (2015-12-14 12:00:00+00:00,14532.5000000000)
  (2015-12-14 22:00:00+00:00,11340.2500000000)
  (2015-12-15 08:00:00+00:00,14606.0000000000)
  (2015-12-15 18:00:00+00:00,14763.2500000000)
  (2015-12-16 04:00:00+00:00,11488.7500000000)
  (2015-12-16 14:00:00+00:00,13775.2500000000)
  (2015-12-17 00:00:00+00:00,9527.2500000000)
  (2015-12-17 10:00:00+00:00,13919.7500000000)
  (2015-12-17 20:00:00+00:00,12136.7500000000)
  (2015-12-18 06:00:00+00:00,13315.5000000000)
  (2015-12-18 16:00:00+00:00,13836.0000000000)
  (2015-12-19 02:00:00+00:00,8801.7500000000)
  (2015-12-19 12:00:00+00:00,11223.5000000000)
  (2015-12-19 22:00:00+00:00,9343.2500000000)
  (2015-12-20 08:00:00+00:00,10333.2500000000)
  (2015-12-20 18:00:00+00:00,11501.5000000000)
  (2015-12-21 04:00:00+00:00,9700.7500000000)
  (2015-12-21 14:00:00+00:00,12767.7500000000)
  (2015-12-22 00:00:00+00:00,9225.2500000000)
  (2015-12-22 10:00:00+00:00,13252.7500000000)
  (2015-12-22 20:00:00+00:00,11486.7500000000)
  (2015-12-23 06:00:00+00:00,11695.7500000000)
  (2015-12-23 16:00:00+00:00,13264.2500000000)
  (2015-12-24 02:00:00+00:00,8146.2500000000)
  (2015-12-24 12:00:00+00:00,10292.7500000000)
  (2015-12-24 22:00:00+00:00,8687.7500000000)
  (2015-12-25 08:00:00+00:00,10000.7500000000)
  (2015-12-25 18:00:00+00:00,10131.2500000000)
  (2015-12-26 04:00:00+00:00,7759.5000000000)
  (2015-12-26 14:00:00+00:00,9122.5000000000)
  (2015-12-27 00:00:00+00:00,7480.2500000000)
  (2015-12-27 10:00:00+00:00,10341.7500000000)
  (2015-12-27 20:00:00+00:00,9713.2500000000)
  (2015-12-28 06:00:00+00:00,10342.7500000000)
  (2015-12-28 16:00:00+00:00,12618.5000000000)
  (2015-12-29 02:00:00+00:00,8290.5000000000)
  (2015-12-29 12:00:00+00:00,11213.7500000000)
  (2015-12-29 22:00:00+00:00,9209.0000000000)
  (2015-12-30 08:00:00+00:00,11111.7500000000)
  (2015-12-30 18:00:00+00:00,12375.0000000000)
  (2015-12-31 04:00:00+00:00,9041.2500000000)
  (2015-12-31 14:00:00+00:00,10429.7500000000)
  (2016-01-01 00:00:00+00:00,8698.2500000000)
  (2016-01-01 10:00:00+00:00,10019.5000000000)
  (2016-01-01 20:00:00+00:00,10134.7500000000)
  (2016-01-02 06:00:00+00:00,9388.5000000000)
  (2016-01-02 16:00:00+00:00,12752.2500000000)
  (2016-01-03 02:00:00+00:00,9493.7500000000)
  (2016-01-03 12:00:00+00:00,11394.2500000000)
  (2016-01-03 22:00:00+00:00,10648.5000000000)
  (2016-01-04 08:00:00+00:00,14148.0000000000)
  (2016-01-04 18:00:00+00:00,15062.2500000000)
  (2016-01-05 04:00:00+00:00,12136.7500000000)
  (2016-01-05 14:00:00+00:00,14743.0000000000)
  (2016-01-06 00:00:00+00:00,11588.7500000000)
  (2016-01-06 10:00:00+00:00,15176.2500000000)
  (2016-01-06 20:00:00+00:00,13518.7500000000)
  (2016-01-07 06:00:00+00:00,14334.2500000000)
  (2016-01-07 16:00:00+00:00,15273.2500000000)
  (2016-01-08 02:00:00+00:00,11079.0000000000)
  (2016-01-08 12:00:00+00:00,14172.2500000000)
  (2016-01-08 22:00:00+00:00,11281.7500000000)
  (2016-01-09 08:00:00+00:00,12166.0000000000)
  (2016-01-09 18:00:00+00:00,12910.5000000000)
  (2016-01-10 04:00:00+00:00,9516.2500000000)
  (2016-01-10 14:00:00+00:00,11162.5000000000)
  (2016-01-11 00:00:00+00:00,9752.7500000000)
  (2016-01-11 10:00:00+00:00,14940.0000000000)
  (2016-01-11 20:00:00+00:00,13159.2500000000)
  (2016-01-12 06:00:00+00:00,14470.7500000000)
  (2016-01-12 16:00:00+00:00,14942.0000000000)
  (2016-01-13 02:00:00+00:00,10704.0000000000)
  (2016-01-13 12:00:00+00:00,14780.5000000000)
  (2016-01-13 22:00:00+00:00,11685.0000000000)
  (2016-01-14 08:00:00+00:00,14401.2500000000)
  (2016-01-14 18:00:00+00:00,14844.2500000000)
  (2016-01-15 04:00:00+00:00,11684.0000000000)
  (2016-01-15 14:00:00+00:00,14059.5000000000)
  (2016-01-16 00:00:00+00:00,10951.2500000000)
  (2016-01-16 10:00:00+00:00,13590.7500000000)
  (2016-01-16 20:00:00+00:00,11797.0000000000)
  (2016-01-17 06:00:00+00:00,10153.5000000000)
  (2016-01-17 16:00:00+00:00,13127.5000000000)
  (2016-01-18 02:00:00+00:00,11084.5000000000)
  (2016-01-18 12:00:00+00:00,14784.0000000000)
  (2016-01-18 22:00:00+00:00,12968.5000000000)
  (2016-01-19 08:00:00+00:00,15143.5000000000)
  (2016-01-19 18:00:00+00:00,15551.0000000000)
  (2016-01-20 04:00:00+00:00,12575.7500000000)
  (2016-01-20 14:00:00+00:00,14864.7500000000)
  (2016-01-21 00:00:00+00:00,11638.2500000000)
  (2016-01-21 10:00:00+00:00,15331.2500000000)
  (2016-01-21 20:00:00+00:00,13892.7500000000)
  (2016-01-22 06:00:00+00:00,15469.5000000000)
  (2016-01-22 16:00:00+00:00,15904.7500000000)
  (2016-01-23 02:00:00+00:00,11223.7500000000)
  (2016-01-23 12:00:00+00:00,13304.7500000000)
  (2016-01-23 22:00:00+00:00,11107.7500000000)
  (2016-01-24 08:00:00+00:00,11792.7500000000)
  (2016-01-24 18:00:00+00:00,12823.5000000000)
  (2016-01-25 04:00:00+00:00,11294.2500000000)
  (2016-01-25 14:00:00+00:00,14290.0000000000)
  (2016-01-26 00:00:00+00:00,10449.7500000000)
  (2016-01-26 10:00:00+00:00,14318.5000000000)
  (2016-01-26 20:00:00+00:00,12726.7500000000)
  (2016-01-27 06:00:00+00:00,13986.2500000000)
  (2016-01-27 16:00:00+00:00,14717.7500000000)
  (2016-01-28 02:00:00+00:00,10275.5000000000)
  (2016-01-28 12:00:00+00:00,13814.7500000000)
  (2016-01-28 22:00:00+00:00,11485.7500000000)
  (2016-01-29 08:00:00+00:00,14439.7500000000)
  (2016-01-29 18:00:00+00:00,14267.0000000000)
  (2016-01-30 04:00:00+00:00,9968.0000000000)
  (2016-01-30 14:00:00+00:00,11935.2500000000)
  (2016-01-31 00:00:00+00:00,9604.2500000000)
  (2016-01-31 10:00:00+00:00,12299.7500000000)
  (2016-01-31 20:00:00+00:00,10852.0000000000)
  (2016-02-01 06:00:00+00:00,13139.5000000000)
  (2016-02-01 16:00:00+00:00,14622.7500000000)
  (2016-02-02 02:00:00+00:00,10011.2500000000)
  (2016-02-02 12:00:00+00:00,14012.0000000000)
  (2016-02-02 22:00:00+00:00,11244.5000000000)
  (2016-02-03 08:00:00+00:00,14117.0000000000)
  (2016-02-03 18:00:00+00:00,14450.2500000000)
  (2016-02-04 04:00:00+00:00,11488.7500000000)
  (2016-02-04 14:00:00+00:00,14016.5000000000)
  (2016-02-05 00:00:00+00:00,10539.5000000000)
  (2016-02-05 10:00:00+00:00,14838.5000000000)
  (2016-02-05 20:00:00+00:00,11951.7500000000)
  (2016-02-06 06:00:00+00:00,10325.0000000000)
  (2016-02-06 16:00:00+00:00,12108.2500000000)
  (2016-02-07 02:00:00+00:00,9052.2500000000)
  (2016-02-07 12:00:00+00:00,10695.0000000000)
  (2016-02-07 22:00:00+00:00,9948.2500000000)
  (2016-02-08 08:00:00+00:00,14136.2500000000)
  (2016-02-08 18:00:00+00:00,14374.5000000000)
  (2016-02-09 04:00:00+00:00,11549.0000000000)
  (2016-02-09 14:00:00+00:00,13710.7500000000)
  (2016-02-10 00:00:00+00:00,10576.5000000000)
  (2016-02-10 10:00:00+00:00,14360.0000000000)
  (2016-02-10 20:00:00+00:00,13003.0000000000)
  (2016-02-11 06:00:00+00:00,13882.5000000000)
  (2016-02-11 16:00:00+00:00,14075.7500000000)
  (2016-02-12 02:00:00+00:00,10884.0000000000)
  (2016-02-12 12:00:00+00:00,13990.7500000000)
  (2016-02-12 22:00:00+00:00,11491.7500000000)
  (2016-02-13 08:00:00+00:00,12492.2500000000)
  (2016-02-13 18:00:00+00:00,12903.7500000000)
  (2016-02-14 04:00:00+00:00,9271.0000000000)
  (2016-02-14 14:00:00+00:00,10778.7500000000)
  (2016-02-15 00:00:00+00:00,9497.2500000000)
  (2016-02-15 10:00:00+00:00,14908.0000000000)
  (2016-02-15 20:00:00+00:00,12975.5000000000)
  (2016-02-16 06:00:00+00:00,14205.7500000000)
  (2016-02-16 16:00:00+00:00,14551.2500000000)
  (2016-02-17 02:00:00+00:00,11403.0000000000)
  (2016-02-17 12:00:00+00:00,14640.0000000000)
  (2016-02-17 22:00:00+00:00,11945.2500000000)
  (2016-02-18 08:00:00+00:00,14621.0000000000)
  (2016-02-18 18:00:00+00:00,14906.2500000000)
  (2016-02-19 04:00:00+00:00,12031.0000000000)
  (2016-02-19 14:00:00+00:00,13885.7500000000)
  (2016-02-20 00:00:00+00:00,10236.7500000000)
  (2016-02-20 10:00:00+00:00,13548.2500000000)
  (2016-02-20 20:00:00+00:00,11299.0000000000)
  (2016-02-21 06:00:00+00:00,9586.7500000000)
  (2016-02-21 16:00:00+00:00,11731.5000000000)
  (2016-02-22 02:00:00+00:00,9806.5000000000)
  (2016-02-22 12:00:00+00:00,14221.0000000000)
  (2016-02-22 22:00:00+00:00,11576.2500000000)
  (2016-02-23 08:00:00+00:00,14238.0000000000)
  (2016-02-23 18:00:00+00:00,14591.7500000000)
  (2016-02-24 04:00:00+00:00,11886.2500000000)
  (2016-02-24 14:00:00+00:00,13754.7500000000)
  (2016-02-25 00:00:00+00:00,11245.5000000000)
  (2016-02-25 10:00:00+00:00,14445.5000000000)
  (2016-02-25 20:00:00+00:00,13449.7500000000)
  (2016-02-26 06:00:00+00:00,14105.7500000000)
  (2016-02-26 16:00:00+00:00,13955.2500000000)
  (2016-02-27 02:00:00+00:00,10581.0000000000)
  (2016-02-27 12:00:00+00:00,12093.0000000000)
  (2016-02-27 22:00:00+00:00,10571.0000000000)
  (2016-02-28 08:00:00+00:00,11087.5000000000)
  (2016-02-28 18:00:00+00:00,12373.2500000000)
  (2016-02-29 04:00:00+00:00,11289.5000000000)
  (2016-02-29 14:00:00+00:00,13786.0000000000)
  (2016-03-01 00:00:00+00:00,10756.7500000000)
  (2016-03-01 10:00:00+00:00,14722.7500000000)
  (2016-03-01 20:00:00+00:00,13251.5000000000)
  (2016-03-02 06:00:00+00:00,13900.5000000000)
  (2016-03-02 16:00:00+00:00,13999.7500000000)
  (2016-03-03 02:00:00+00:00,10648.0000000000)
  (2016-03-03 12:00:00+00:00,14010.2500000000)
  (2016-03-03 22:00:00+00:00,11398.5000000000)
  (2016-03-04 08:00:00+00:00,14116.2500000000)
  (2016-03-04 18:00:00+00:00,14526.5000000000)
  (2016-03-05 04:00:00+00:00,10149.0000000000)
  (2016-03-05 14:00:00+00:00,11584.2500000000)
  (2016-03-06 00:00:00+00:00,9277.0000000000)
  (2016-03-06 10:00:00+00:00,12374.0000000000)
  (2016-03-06 20:00:00+00:00,11089.7500000000)
  (2016-03-07 06:00:00+00:00,13562.0000000000)
  (2016-03-07 16:00:00+00:00,13483.5000000000)
  (2016-03-08 02:00:00+00:00,10755.5000000000)
  (2016-03-08 12:00:00+00:00,14306.5000000000)
  (2016-03-08 22:00:00+00:00,11763.2500000000)
  (2016-03-09 08:00:00+00:00,13848.5000000000)
  (2016-03-09 18:00:00+00:00,14469.5000000000)
  (2016-03-10 04:00:00+00:00,11457.0000000000)
  (2016-03-10 14:00:00+00:00,13989.7500000000)
  (2016-03-11 00:00:00+00:00,10661.5000000000)
  (2016-03-11 10:00:00+00:00,14710.5000000000)
  (2016-03-11 20:00:00+00:00,12814.7500000000)
  (2016-03-12 06:00:00+00:00,11200.2500000000)
  (2016-03-12 16:00:00+00:00,12217.0000000000)
  (2016-03-13 02:00:00+00:00,9597.2500000000)
  (2016-03-13 12:00:00+00:00,11531.2500000000)
  (2016-03-13 22:00:00+00:00,10718.7500000000)
  (2016-03-14 08:00:00+00:00,13829.7500000000)
  (2016-03-14 18:00:00+00:00,14356.5000000000)
  (2016-03-15 04:00:00+00:00,11501.5000000000)
  (2016-03-15 14:00:00+00:00,14033.7500000000)
  (2016-03-16 00:00:00+00:00,10584.5000000000)
  (2016-03-16 10:00:00+00:00,14517.7500000000)
  (2016-03-16 20:00:00+00:00,13245.0000000000)
  (2016-03-17 06:00:00+00:00,13754.2500000000)
  (2016-03-17 16:00:00+00:00,13199.2500000000)
  (2016-03-18 02:00:00+00:00,10859.0000000000)
  (2016-03-18 12:00:00+00:00,13857.0000000000)
  (2016-03-18 22:00:00+00:00,11304.7500000000)
  (2016-03-19 08:00:00+00:00,12402.0000000000)
  (2016-03-19 18:00:00+00:00,13174.2500000000)
  (2016-03-20 04:00:00+00:00,9464.0000000000)
  (2016-03-20 14:00:00+00:00,11000.2500000000)
  (2016-03-21 00:00:00+00:00,9722.2500000000)
  (2016-03-21 10:00:00+00:00,14249.0000000000)
  (2016-03-21 20:00:00+00:00,12736.2500000000)
  (2016-03-22 06:00:00+00:00,13550.5000000000)
  (2016-03-22 16:00:00+00:00,13844.5000000000)
  (2016-03-23 02:00:00+00:00,10431.0000000000)
  (2016-03-23 12:00:00+00:00,14092.5000000000)
  (2016-03-23 22:00:00+00:00,11438.2500000000)
  (2016-03-24 08:00:00+00:00,14157.2500000000)
  (2016-03-24 18:00:00+00:00,14039.5000000000)
  (2016-03-25 04:00:00+00:00,8885.5000000000)
  (2016-03-25 14:00:00+00:00,10631.0000000000)
  (2016-03-26 00:00:00+00:00,9048.5000000000)
  (2016-03-26 10:00:00+00:00,11701.2500000000)
  (2016-03-26 20:00:00+00:00,10508.7500000000)
  (2016-03-27 06:00:00+00:00,9332.0000000000)
  (2016-03-27 16:00:00+00:00,9630.7500000000)
  (2016-03-28 02:00:00+00:00,8271.0000000000)
  (2016-03-28 12:00:00+00:00,9832.2500000000)
  (2016-03-28 22:00:00+00:00,8890.0000000000)
  (2016-03-29 08:00:00+00:00,13728.0000000000)
  (2016-03-29 18:00:00+00:00,13141.5000000000)
  (2016-03-30 04:00:00+00:00,12057.7500000000)
  (2016-03-30 14:00:00+00:00,12879.5000000000)
  (2016-03-31 00:00:00+00:00,9901.5000000000)
  (2016-03-31 10:00:00+00:00,14024.2500000000)
  (2016-03-31 20:00:00+00:00,11959.0000000000)
  (2016-04-01 06:00:00+00:00,13851.2500000000)
  (2016-04-01 16:00:00+00:00,12753.7500000000)
  (2016-04-02 02:00:00+00:00,9585.7500000000)
  (2016-04-02 12:00:00+00:00,11017.0000000000)
  (2016-04-02 22:00:00+00:00,8957.2500000000)
  (2016-04-03 08:00:00+00:00,10855.7500000000)
  (2016-04-03 18:00:00+00:00,11128.7500000000)
  (2016-04-04 04:00:00+00:00,11809.2500000000)
  (2016-04-04 14:00:00+00:00,12423.7500000000)
  (2016-04-05 00:00:00+00:00,9239.7500000000)
  (2016-04-05 10:00:00+00:00,13324.5000000000)
  (2016-04-05 20:00:00+00:00,11411.2500000000)
  (2016-04-06 06:00:00+00:00,13489.2500000000)
  (2016-04-06 16:00:00+00:00,12696.2500000000)
  (2016-04-07 02:00:00+00:00,9886.5000000000)
  (2016-04-07 12:00:00+00:00,13287.7500000000)
  (2016-04-07 22:00:00+00:00,9995.7500000000)
  (2016-04-08 08:00:00+00:00,13513.2500000000)
  (2016-04-08 18:00:00+00:00,12471.2500000000)
  (2016-04-09 04:00:00+00:00,9697.0000000000)
  (2016-04-09 14:00:00+00:00,10486.5000000000)
  (2016-04-10 00:00:00+00:00,8231.2500000000)
  (2016-04-10 10:00:00+00:00,11058.7500000000)
  (2016-04-10 20:00:00+00:00,10273.5000000000)
  (2016-04-11 06:00:00+00:00,13109.5000000000)
  (2016-04-11 16:00:00+00:00,12730.0000000000)
  (2016-04-12 02:00:00+00:00,9640.0000000000)
  (2016-04-12 12:00:00+00:00,13468.0000000000)
  (2016-04-12 22:00:00+00:00,10107.2500000000)
  (2016-04-13 08:00:00+00:00,13545.5000000000)
  (2016-04-13 18:00:00+00:00,13200.7500000000)
  (2016-04-14 04:00:00+00:00,12288.7500000000)
  (2016-04-14 14:00:00+00:00,13147.0000000000)
  (2016-04-15 00:00:00+00:00,9675.5000000000)
  (2016-04-15 10:00:00+00:00,13841.2500000000)
  (2016-04-15 20:00:00+00:00,11179.5000000000)
  (2016-04-16 06:00:00+00:00,10954.2500000000)
  (2016-04-16 16:00:00+00:00,10893.5000000000)
  (2016-04-17 02:00:00+00:00,8408.0000000000)
  (2016-04-17 12:00:00+00:00,10246.7500000000)
  (2016-04-17 22:00:00+00:00,9365.2500000000)
  (2016-04-18 08:00:00+00:00,13360.5000000000)
  (2016-04-18 18:00:00+00:00,12822.0000000000)
  (2016-04-19 04:00:00+00:00,12091.5000000000)
  (2016-04-19 14:00:00+00:00,13100.0000000000)
  (2016-04-20 00:00:00+00:00,9597.7500000000)
  (2016-04-20 10:00:00+00:00,13553.2500000000)
  (2016-04-20 20:00:00+00:00,11854.7500000000)
  (2016-04-21 06:00:00+00:00,13603.2500000000)
  (2016-04-21 16:00:00+00:00,12825.2500000000)
  (2016-04-22 02:00:00+00:00,10417.5000000000)
  (2016-04-22 12:00:00+00:00,13103.5000000000)
  (2016-04-22 22:00:00+00:00,9859.5000000000)
  (2016-04-23 08:00:00+00:00,12026.2500000000)
  (2016-04-23 18:00:00+00:00,11137.7500000000)
  (2016-04-24 04:00:00+00:00,8530.2500000000)
  (2016-04-24 14:00:00+00:00,10150.5000000000)
  (2016-04-25 00:00:00+00:00,9170.7500000000)
  (2016-04-25 10:00:00+00:00,13931.5000000000)
  (2016-04-25 20:00:00+00:00,11875.5000000000)
  (2016-04-26 06:00:00+00:00,13836.2500000000)
  (2016-04-26 16:00:00+00:00,13212.5000000000)
  (2016-04-27 02:00:00+00:00,10760.2500000000)
  (2016-04-27 12:00:00+00:00,13613.2500000000)
  (2016-04-27 22:00:00+00:00,10743.2500000000)
  (2016-04-28 08:00:00+00:00,14099.5000000000)
  (2016-04-28 18:00:00+00:00,12948.2500000000)
  (2016-04-29 04:00:00+00:00,12505.7500000000)
  (2016-04-29 14:00:00+00:00,12821.0000000000)
  (2016-04-30 00:00:00+00:00,9237.0000000000)
  (2016-04-30 10:00:00+00:00,11789.7500000000)
  (2016-04-30 20:00:00+00:00,10366.2500000000)
  (2016-05-01 06:00:00+00:00,9910.7500000000)
  (2016-05-01 16:00:00+00:00,10578.2500000000)
  (2016-05-02 02:00:00+00:00,9745.0000000000)
  (2016-05-02 12:00:00+00:00,13595.2500000000)
  (2016-05-02 22:00:00+00:00,10979.0000000000)
  (2016-05-03 08:00:00+00:00,14337.5000000000)
  (2016-05-03 18:00:00+00:00,13539.2500000000)
  (2016-05-04 04:00:00+00:00,12998.7500000000)
  (2016-05-04 14:00:00+00:00,13515.7500000000)
  (2016-05-05 00:00:00+00:00,9811.0000000000)
  (2016-05-05 10:00:00+00:00,10658.7500000000)
  (2016-05-05 20:00:00+00:00,10720.7500000000)
  (2016-05-06 06:00:00+00:00,12334.5000000000)
  (2016-05-06 16:00:00+00:00,11800.5000000000)
  (2016-05-07 02:00:00+00:00,8923.0000000000)
  (2016-05-07 12:00:00+00:00,11164.7500000000)
  (2016-05-07 22:00:00+00:00,9088.5000000000)
  (2016-05-08 08:00:00+00:00,10763.0000000000)
  (2016-05-08 18:00:00+00:00,10687.2500000000)
  (2016-05-09 04:00:00+00:00,11853.5000000000)
  (2016-05-09 14:00:00+00:00,13062.5000000000)
  (2016-05-10 00:00:00+00:00,9791.5000000000)
  (2016-05-10 10:00:00+00:00,13787.5000000000)
  (2016-05-10 20:00:00+00:00,11945.0000000000)
  (2016-05-11 06:00:00+00:00,13845.2500000000)
  (2016-05-11 16:00:00+00:00,13360.2500000000)
  (2016-05-12 02:00:00+00:00,10083.7500000000)
  (2016-05-12 12:00:00+00:00,13478.2500000000)
  (2016-05-12 22:00:00+00:00,10280.5000000000)
  (2016-05-13 08:00:00+00:00,13679.5000000000)
  (2016-05-13 18:00:00+00:00,12219.2500000000)
  (2016-05-14 04:00:00+00:00,9087.0000000000)
  (2016-05-14 14:00:00+00:00,10348.7500000000)
  (2016-05-15 00:00:00+00:00,8161.2500000000)
  (2016-05-15 10:00:00+00:00,10368.2500000000)
  (2016-05-15 20:00:00+00:00,9772.7500000000)
  (2016-05-16 06:00:00+00:00,9835.2500000000)
  (2016-05-16 16:00:00+00:00,10626.2500000000)
  (2016-05-17 02:00:00+00:00,9692.7500000000)
  (2016-05-17 12:00:00+00:00,13581.0000000000)
  (2016-05-17 22:00:00+00:00,10636.5000000000)
  (2016-05-18 08:00:00+00:00,14110.7500000000)
  (2016-05-18 18:00:00+00:00,13312.5000000000)
  (2016-05-19 04:00:00+00:00,12352.0000000000)
  (2016-05-19 14:00:00+00:00,13165.0000000000)
  (2016-05-20 00:00:00+00:00,10114.2500000000)
  (2016-05-20 10:00:00+00:00,14072.5000000000)
  (2016-05-20 20:00:00+00:00,11739.7500000000)
  (2016-05-21 06:00:00+00:00,11240.7500000000)
  (2016-05-21 16:00:00+00:00,11049.2500000000)
  (2016-05-22 02:00:00+00:00,8068.7500000000)
  (2016-05-22 12:00:00+00:00,10195.7500000000)
  (2016-05-22 22:00:00+00:00,10207.2500000000)
  (2016-05-23 08:00:00+00:00,14276.7500000000)
  (2016-05-23 18:00:00+00:00,13468.2500000000)
  (2016-05-24 04:00:00+00:00,12445.2500000000)
  (2016-05-24 14:00:00+00:00,13619.2500000000)
  (2016-05-25 00:00:00+00:00,10207.2500000000)
  (2016-05-25 10:00:00+00:00,14417.0000000000)
  (2016-05-25 20:00:00+00:00,12098.5000000000)
  (2016-05-26 06:00:00+00:00,13922.5000000000)
  (2016-05-26 16:00:00+00:00,13587.2500000000)
  (2016-05-27 02:00:00+00:00,10401.0000000000)
  (2016-05-27 12:00:00+00:00,13775.0000000000)
  (2016-05-27 22:00:00+00:00,10231.7500000000)
  (2016-05-28 08:00:00+00:00,12691.0000000000)
  (2016-05-28 18:00:00+00:00,11123.2500000000)
  (2016-05-29 04:00:00+00:00,8258.7500000000)
  (2016-05-29 14:00:00+00:00,10245.2500000000)
  (2016-05-30 00:00:00+00:00,8919.5000000000)
  (2016-05-30 10:00:00+00:00,14261.0000000000)
  (2016-05-30 20:00:00+00:00,12051.5000000000)
  (2016-05-31 06:00:00+00:00,14152.7500000000)
  (2016-05-31 16:00:00+00:00,13765.7500000000)
  (2016-06-01 02:00:00+00:00,9315.5000000000)
  (2016-06-01 12:00:00+00:00,13357.5000000000)
  (2016-06-01 22:00:00+00:00,9704.0000000000)
  (2016-06-02 08:00:00+00:00,13638.0000000000)
  (2016-06-02 18:00:00+00:00,12184.2500000000)
  (2016-06-03 04:00:00+00:00,11585.5000000000)
  (2016-06-03 14:00:00+00:00,12704.7500000000)
  (2016-06-04 00:00:00+00:00,8914.2500000000)
  (2016-06-04 10:00:00+00:00,11860.5000000000)
  (2016-06-04 20:00:00+00:00,9919.5000000000)
  (2016-06-05 06:00:00+00:00,9445.2500000000)
  (2016-06-05 16:00:00+00:00,10160.5000000000)
  (2016-06-06 02:00:00+00:00,8596.2500000000)
  (2016-06-06 12:00:00+00:00,13078.2500000000)
  (2016-06-06 22:00:00+00:00,9725.2500000000)
  (2016-06-07 08:00:00+00:00,13262.5000000000)
  (2016-06-07 18:00:00+00:00,12211.7500000000)
  (2016-06-08 04:00:00+00:00,11291.2500000000)
  (2016-06-08 14:00:00+00:00,12886.7500000000)
  (2016-06-09 00:00:00+00:00,9196.0000000000)
  (2016-06-09 10:00:00+00:00,13275.5000000000)
  (2016-06-09 20:00:00+00:00,11252.5000000000)
  (2016-06-10 06:00:00+00:00,13134.2500000000)
  (2016-06-10 16:00:00+00:00,12203.2500000000)
  (2016-06-11 02:00:00+00:00,8541.0000000000)
  (2016-06-11 12:00:00+00:00,10781.0000000000)
  (2016-06-11 22:00:00+00:00,8610.5000000000)
  (2016-06-12 08:00:00+00:00,10457.7500000000)
  (2016-06-12 18:00:00+00:00,10283.7500000000)
  (2016-06-13 04:00:00+00:00,10996.7500000000)
  (2016-06-13 14:00:00+00:00,12530.7500000000)
  (2016-06-14 00:00:00+00:00,9108.2500000000)
  (2016-06-14 10:00:00+00:00,13347.2500000000)
  (2016-06-14 20:00:00+00:00,11082.7500000000)
  (2016-06-15 06:00:00+00:00,13040.0000000000)
  (2016-06-15 16:00:00+00:00,12710.2500000000)
  (2016-06-16 02:00:00+00:00,9446.5000000000)
  (2016-06-16 12:00:00+00:00,12897.7500000000)
  (2016-06-16 22:00:00+00:00,9747.2500000000)
  (2016-06-17 08:00:00+00:00,13782.0000000000)
  (2016-06-17 18:00:00+00:00,11879.2500000000)
  (2016-06-18 04:00:00+00:00,8958.2500000000)
  (2016-06-18 14:00:00+00:00,10728.0000000000)
  (2016-06-19 00:00:00+00:00,7873.5000000000)
  (2016-06-19 10:00:00+00:00,10569.2500000000)
  (2016-06-19 20:00:00+00:00,9910.7500000000)
  (2016-06-20 06:00:00+00:00,12794.2500000000)
  (2016-06-20 16:00:00+00:00,12526.7500000000)
  (2016-06-21 02:00:00+00:00,9578.0000000000)
  (2016-06-21 12:00:00+00:00,13194.7500000000)
  (2016-06-21 22:00:00+00:00,9472.7500000000)
  (2016-06-22 08:00:00+00:00,13480.0000000000)
  (2016-06-22 18:00:00+00:00,12205.2500000000)
  (2016-06-23 04:00:00+00:00,11466.2500000000)
  (2016-06-23 14:00:00+00:00,13108.5000000000)
  (2016-06-24 00:00:00+00:00,9600.0000000000)
  (2016-06-24 10:00:00+00:00,14012.5000000000)
  (2016-06-24 20:00:00+00:00,11341.2500000000)
  (2016-06-25 06:00:00+00:00,10987.5000000000)
  (2016-06-25 16:00:00+00:00,11316.2500000000)
  (2016-06-26 02:00:00+00:00,8151.0000000000)
  (2016-06-26 12:00:00+00:00,9813.2500000000)
  (2016-06-26 22:00:00+00:00,9250.5000000000)
  (2016-06-27 08:00:00+00:00,13301.7500000000)
  (2016-06-27 18:00:00+00:00,12354.7500000000)
  (2016-06-28 04:00:00+00:00,11691.0000000000)
  (2016-06-28 14:00:00+00:00,12926.2500000000)
  (2016-06-29 00:00:00+00:00,9239.5000000000)
  (2016-06-29 10:00:00+00:00,14046.2500000000)
  (2016-06-29 20:00:00+00:00,11377.0000000000)
  (2016-06-30 06:00:00+00:00,13105.5000000000)
  (2016-06-30 16:00:00+00:00,12487.2500000000)
  (2016-07-01 02:00:00+00:00,9451.5000000000)
  (2016-07-01 12:00:00+00:00,12794.5000000000)
  (2016-07-01 22:00:00+00:00,9662.7500000000)
  (2016-07-02 08:00:00+00:00,12041.7500000000)
  (2016-07-02 18:00:00+00:00,10274.2500000000)
  (2016-07-03 04:00:00+00:00,7518.7500000000)
  (2016-07-03 14:00:00+00:00,9463.0000000000)
  (2016-07-04 00:00:00+00:00,8643.2500000000)
  (2016-07-04 10:00:00+00:00,13391.7500000000)
  (2016-07-04 20:00:00+00:00,11125.7500000000)
  (2016-07-05 06:00:00+00:00,13047.7500000000)
  (2016-07-05 16:00:00+00:00,12890.7500000000)
  (2016-07-06 02:00:00+00:00,9503.5000000000)
  (2016-07-06 12:00:00+00:00,13133.2500000000)
  (2016-07-06 22:00:00+00:00,9870.2500000000)
  (2016-07-07 08:00:00+00:00,13432.2500000000)
  (2016-07-07 18:00:00+00:00,11995.0000000000)
  (2016-07-08 04:00:00+00:00,11508.0000000000)
  (2016-07-08 14:00:00+00:00,12825.2500000000)
  (2016-07-09 00:00:00+00:00,8927.5000000000)
  (2016-07-09 10:00:00+00:00,11583.0000000000)
  (2016-07-09 20:00:00+00:00,9984.2500000000)
  (2016-07-10 06:00:00+00:00,9327.2500000000)
  (2016-07-10 16:00:00+00:00,10714.5000000000)
  (2016-07-11 02:00:00+00:00,9399.0000000000)
  (2016-07-11 12:00:00+00:00,13270.2500000000)
  (2016-07-11 22:00:00+00:00,10079.7500000000)
  (2016-07-12 08:00:00+00:00,13856.7500000000)
  (2016-07-12 18:00:00+00:00,12150.2500000000)
  (2016-07-13 04:00:00+00:00,11184.2500000000)
  (2016-07-13 14:00:00+00:00,12788.7500000000)
  (2016-07-14 00:00:00+00:00,9199.0000000000)
  (2016-07-14 10:00:00+00:00,13779.7500000000)
  (2016-07-14 20:00:00+00:00,11356.7500000000)
  (2016-07-15 06:00:00+00:00,13375.2500000000)
  (2016-07-15 16:00:00+00:00,12501.2500000000)
  (2016-07-16 02:00:00+00:00,8846.0000000000)
  (2016-07-16 12:00:00+00:00,11069.2500000000)
  (2016-07-16 22:00:00+00:00,8796.2500000000)
  (2016-07-17 08:00:00+00:00,10678.0000000000)
  (2016-07-17 18:00:00+00:00,10260.2500000000)
  (2016-07-18 04:00:00+00:00,11004.0000000000)
  (2016-07-18 14:00:00+00:00,12592.7500000000)
  (2016-07-19 00:00:00+00:00,9053.2500000000)
  (2016-07-19 10:00:00+00:00,13466.7500000000)
  (2016-07-19 20:00:00+00:00,11054.7500000000)
  (2016-07-20 06:00:00+00:00,12961.2500000000)
  (2016-07-20 16:00:00+00:00,12669.0000000000)
  (2016-07-21 02:00:00+00:00,9595.5000000000)
  (2016-07-21 12:00:00+00:00,13339.0000000000)
  (2016-07-21 22:00:00+00:00,9819.2500000000)
  (2016-07-22 08:00:00+00:00,13724.0000000000)
  (2016-07-22 18:00:00+00:00,11920.0000000000)
  (2016-07-23 04:00:00+00:00,9178.0000000000)
  (2016-07-23 14:00:00+00:00,10964.5000000000)
  (2016-07-24 00:00:00+00:00,8309.0000000000)
  (2016-07-24 10:00:00+00:00,10901.7500000000)
  (2016-07-24 20:00:00+00:00,10067.2500000000)
  (2016-07-25 06:00:00+00:00,12763.5000000000)
  (2016-07-25 16:00:00+00:00,12565.7500000000)
  (2016-07-26 02:00:00+00:00,9273.0000000000)
  (2016-07-26 12:00:00+00:00,13203.7500000000)
  (2016-07-26 22:00:00+00:00,9723.7500000000)
  (2016-07-27 08:00:00+00:00,13508.7500000000)
  (2016-07-27 18:00:00+00:00,11840.0000000000)
  (2016-07-28 04:00:00+00:00,10734.0000000000)
  (2016-07-28 14:00:00+00:00,12485.5000000000)
  (2016-07-29 00:00:00+00:00,9112.2500000000)
  (2016-07-29 10:00:00+00:00,13345.0000000000)
  (2016-07-29 20:00:00+00:00,10677.5000000000)
  (2016-07-30 06:00:00+00:00,10251.5000000000)
  (2016-07-30 16:00:00+00:00,10713.7500000000)
  (2016-07-31 02:00:00+00:00,8057.5000000000)
  (2016-07-31 12:00:00+00:00,10073.2500000000)
  (2016-07-31 22:00:00+00:00,8794.0000000000)
  (2016-08-01 08:00:00+00:00,12940.5000000000)
  (2016-08-01 18:00:00+00:00,11663.2500000000)
  (2016-08-02 04:00:00+00:00,10813.7500000000)
  (2016-08-02 14:00:00+00:00,12398.7500000000)
  (2016-08-03 00:00:00+00:00,8977.2500000000)
  (2016-08-03 10:00:00+00:00,13498.5000000000)
  (2016-08-03 20:00:00+00:00,11007.2500000000)
  (2016-08-04 06:00:00+00:00,12894.7500000000)
  (2016-08-04 16:00:00+00:00,12624.5000000000)
  (2016-08-05 02:00:00+00:00,9356.0000000000)
  (2016-08-05 12:00:00+00:00,12783.2500000000)
  (2016-08-05 22:00:00+00:00,9343.2500000000)
  (2016-08-06 08:00:00+00:00,11410.7500000000)
  (2016-08-06 18:00:00+00:00,10137.5000000000)
  (2016-08-07 04:00:00+00:00,7585.7500000000)
  (2016-08-07 14:00:00+00:00,9124.2500000000)
  (2016-08-08 00:00:00+00:00,8322.7500000000)
  (2016-08-08 10:00:00+00:00,13611.0000000000)
  (2016-08-08 20:00:00+00:00,11052.7500000000)
  (2016-08-09 06:00:00+00:00,12757.0000000000)
  (2016-08-09 16:00:00+00:00,12580.5000000000)
  (2016-08-10 02:00:00+00:00,9379.7500000000)
  (2016-08-10 12:00:00+00:00,12891.2500000000)
  (2016-08-10 22:00:00+00:00,9584.7500000000)
  (2016-08-11 08:00:00+00:00,13077.2500000000)
  (2016-08-11 18:00:00+00:00,12254.7500000000)
  (2016-08-12 04:00:00+00:00,11207.5000000000)
  (2016-08-12 14:00:00+00:00,12369.7500000000)
  (2016-08-13 00:00:00+00:00,8573.5000000000)
  (2016-08-13 10:00:00+00:00,11751.0000000000)
  (2016-08-13 20:00:00+00:00,9797.2500000000)
  (2016-08-14 06:00:00+00:00,9280.2500000000)
  (2016-08-14 16:00:00+00:00,9999.5000000000)
  (2016-08-15 02:00:00+00:00,8497.0000000000)
  (2016-08-15 12:00:00+00:00,12806.0000000000)
  (2016-08-15 22:00:00+00:00,9230.2500000000)
  (2016-08-16 08:00:00+00:00,13231.5000000000)
  (2016-08-16 18:00:00+00:00,12017.2500000000)
  (2016-08-17 04:00:00+00:00,10824.5000000000)
  (2016-08-17 14:00:00+00:00,12139.5000000000)
  (2016-08-18 00:00:00+00:00,8668.5000000000)
  (2016-08-18 10:00:00+00:00,13319.5000000000)
  (2016-08-18 20:00:00+00:00,10926.5000000000)
  (2016-08-19 06:00:00+00:00,12944.7500000000)
  (2016-08-19 16:00:00+00:00,12350.0000000000)
  (2016-08-20 02:00:00+00:00,8588.5000000000)
  (2016-08-20 12:00:00+00:00,10891.5000000000)
  (2016-08-20 22:00:00+00:00,8431.0000000000)
  (2016-08-21 08:00:00+00:00,9794.0000000000)
  (2016-08-21 18:00:00+00:00,10194.7500000000)
  (2016-08-22 04:00:00+00:00,10951.0000000000)
  (2016-08-22 14:00:00+00:00,12446.7500000000)
  (2016-08-23 00:00:00+00:00,9060.7500000000)
  (2016-08-23 10:00:00+00:00,13395.7500000000)
  (2016-08-23 20:00:00+00:00,10975.5000000000)
  (2016-08-24 06:00:00+00:00,13059.0000000000)
  (2016-08-24 16:00:00+00:00,12794.0000000000)
  (2016-08-25 02:00:00+00:00,9325.7500000000)
  (2016-08-25 12:00:00+00:00,13371.0000000000)
  (2016-08-25 22:00:00+00:00,10020.7500000000)
  (2016-08-26 08:00:00+00:00,13706.2500000000)
  (2016-08-26 18:00:00+00:00,12234.2500000000)
  (2016-08-27 04:00:00+00:00,9206.5000000000)
  (2016-08-27 14:00:00+00:00,10791.7500000000)
  (2016-08-28 00:00:00+00:00,8316.0000000000)
  (2016-08-28 10:00:00+00:00,10650.7500000000)
  (2016-08-28 20:00:00+00:00,9929.5000000000)
  (2016-08-29 06:00:00+00:00,13057.5000000000)
  (2016-08-29 16:00:00+00:00,12508.5000000000)
  (2016-08-30 02:00:00+00:00,9247.7500000000)
  (2016-08-30 12:00:00+00:00,12670.5000000000)
  (2016-08-30 22:00:00+00:00,9326.2500000000)
  (2016-08-31 08:00:00+00:00,13450.7500000000)
  (2016-08-31 18:00:00+00:00,12767.0000000000)
  (2016-09-01 04:00:00+00:00,11386.2500000000)
  (2016-09-01 14:00:00+00:00,12849.5000000000)
  (2016-09-02 00:00:00+00:00,9202.0000000000)
  (2016-09-02 10:00:00+00:00,13576.7500000000)
  (2016-09-02 20:00:00+00:00,10620.2500000000)
  (2016-09-03 06:00:00+00:00,10770.7500000000)
  (2016-09-03 16:00:00+00:00,10652.0000000000)
  (2016-09-04 02:00:00+00:00,7935.0000000000)
  (2016-09-04 12:00:00+00:00,10268.0000000000)
  (2016-09-04 22:00:00+00:00,9048.2500000000)
  (2016-09-05 08:00:00+00:00,13648.2500000000)
  (2016-09-05 18:00:00+00:00,12902.7500000000)
  (2016-09-06 04:00:00+00:00,11959.5000000000)
  (2016-09-06 14:00:00+00:00,12404.2500000000)
  (2016-09-07 00:00:00+00:00,9222.7500000000)
  (2016-09-07 10:00:00+00:00,13371.0000000000)
  (2016-09-07 20:00:00+00:00,10991.7500000000)
  (2016-09-08 06:00:00+00:00,12947.2500000000)
  (2016-09-08 16:00:00+00:00,13017.7500000000)
  (2016-09-09 02:00:00+00:00,9905.7500000000)
  (2016-09-09 12:00:00+00:00,13348.0000000000)
  (2016-09-09 22:00:00+00:00,9526.2500000000)
  (2016-09-10 08:00:00+00:00,11975.5000000000)
  (2016-09-10 18:00:00+00:00,11192.5000000000)
  (2016-09-11 04:00:00+00:00,8618.7500000000)
  (2016-09-11 14:00:00+00:00,10019.0000000000)
  (2016-09-12 00:00:00+00:00,8989.0000000000)
  (2016-09-12 10:00:00+00:00,13775.2500000000)
  (2016-09-12 20:00:00+00:00,11416.5000000000)
  (2016-09-13 06:00:00+00:00,13383.0000000000)
  (2016-09-13 16:00:00+00:00,13075.2500000000)
  (2016-09-14 02:00:00+00:00,9301.7500000000)
  (2016-09-14 12:00:00+00:00,13096.0000000000)
  (2016-09-14 22:00:00+00:00,9736.2500000000)
  (2016-09-15 08:00:00+00:00,13301.2500000000)
  (2016-09-15 18:00:00+00:00,12769.7500000000)
  (2016-09-16 04:00:00+00:00,11674.0000000000)
  (2016-09-16 14:00:00+00:00,12584.5000000000)
  (2016-09-17 00:00:00+00:00,9075.2500000000)
  (2016-09-17 10:00:00+00:00,12217.5000000000)
  (2016-09-17 20:00:00+00:00,9743.2500000000)
  (2016-09-18 06:00:00+00:00,9807.2500000000)
  (2016-09-18 16:00:00+00:00,10897.0000000000)
  (2016-09-19 02:00:00+00:00,9082.2500000000)
  (2016-09-19 12:00:00+00:00,13205.7500000000)
  (2016-09-19 22:00:00+00:00,9395.0000000000)
  (2016-09-20 08:00:00+00:00,13558.0000000000)
  (2016-09-20 18:00:00+00:00,12968.5000000000)
  (2016-09-21 04:00:00+00:00,12058.7500000000)
  (2016-09-21 14:00:00+00:00,12751.0000000000)
  (2016-09-22 00:00:00+00:00,9265.5000000000)
  (2016-09-22 10:00:00+00:00,13444.2500000000)
  (2016-09-22 20:00:00+00:00,10990.2500000000)
  (2016-09-23 06:00:00+00:00,13325.7500000000)
  (2016-09-23 16:00:00+00:00,12614.7500000000)
  (2016-09-24 02:00:00+00:00,8901.5000000000)
  (2016-09-24 12:00:00+00:00,10903.2500000000)
  (2016-09-24 22:00:00+00:00,8654.0000000000)
  (2016-09-25 08:00:00+00:00,10619.2500000000)
  (2016-09-25 18:00:00+00:00,11125.2500000000)
  (2016-09-26 04:00:00+00:00,11912.7500000000)
  (2016-09-26 14:00:00+00:00,12748.2500000000)
  (2016-09-27 00:00:00+00:00,9331.5000000000)
  (2016-09-27 10:00:00+00:00,13516.5000000000)
  (2016-09-27 20:00:00+00:00,11120.2500000000)
  (2016-09-28 06:00:00+00:00,13110.2500000000)
  (2016-09-28 16:00:00+00:00,12874.7500000000)
  (2016-09-29 02:00:00+00:00,9317.2500000000)
  (2016-09-29 12:00:00+00:00,12956.2500000000)
  (2016-09-29 22:00:00+00:00,9498.5000000000)
  (2016-09-30 08:00:00+00:00,13761.0000000000)
  (2016-09-30 18:00:00+00:00,12726.2500000000)
  (2016-10-01 04:00:00+00:00,9365.0000000000)
  (2016-10-01 14:00:00+00:00,10971.5000000000)
  (2016-10-02 00:00:00+00:00,8271.7500000000)
  (2016-10-02 10:00:00+00:00,10547.5000000000)
  (2016-10-02 20:00:00+00:00,9664.0000000000)
  (2016-10-03 06:00:00+00:00,9702.0000000000)
  (2016-10-03 16:00:00+00:00,10939.7500000000)
  (2016-10-04 02:00:00+00:00,9264.7500000000)
  (2016-10-04 12:00:00+00:00,13421.2500000000)
  (2016-10-04 22:00:00+00:00,9841.7500000000)
  (2016-10-05 08:00:00+00:00,13911.0000000000)
  (2016-10-05 18:00:00+00:00,13450.5000000000)
  (2016-10-06 04:00:00+00:00,12563.7500000000)
  (2016-10-06 14:00:00+00:00,13722.0000000000)
  (2016-10-07 00:00:00+00:00,9652.2500000000)
  (2016-10-07 10:00:00+00:00,13999.7500000000)
  (2016-10-07 20:00:00+00:00,11315.7500000000)
  (2016-10-08 06:00:00+00:00,11278.7500000000)
  (2016-10-08 16:00:00+00:00,12154.5000000000)
  (2016-10-09 02:00:00+00:00,8848.7500000000)
  (2016-10-09 12:00:00+00:00,10593.7500000000)
  (2016-10-09 22:00:00+00:00,9388.0000000000)
  (2016-10-10 08:00:00+00:00,13877.7500000000)
  (2016-10-10 18:00:00+00:00,13334.0000000000)
  (2016-10-11 04:00:00+00:00,12537.0000000000)
  (2016-10-11 14:00:00+00:00,13472.7500000000)
  (2016-10-12 00:00:00+00:00,9700.0000000000)
  (2016-10-12 10:00:00+00:00,14196.0000000000)
  (2016-10-12 20:00:00+00:00,11559.2500000000)
  (2016-10-13 06:00:00+00:00,14068.0000000000)
  (2016-10-13 16:00:00+00:00,14310.0000000000)
  (2016-10-14 02:00:00+00:00,10564.0000000000)
  (2016-10-14 12:00:00+00:00,13267.2500000000)
  (2016-10-14 22:00:00+00:00,10272.0000000000)
  (2016-10-15 08:00:00+00:00,12389.2500000000)
  (2016-10-15 18:00:00+00:00,11476.2500000000)
  (2016-10-16 04:00:00+00:00,8795.2500000000)
  (2016-10-16 14:00:00+00:00,10331.5000000000)
  (2016-10-17 00:00:00+00:00,9145.7500000000)
  (2016-10-17 10:00:00+00:00,13984.2500000000)
  (2016-10-17 20:00:00+00:00,11532.5000000000)
  (2016-10-18 06:00:00+00:00,13865.7500000000)
  (2016-10-18 16:00:00+00:00,13993.7500000000)
  (2016-10-19 02:00:00+00:00,10100.0000000000)
  (2016-10-19 12:00:00+00:00,13521.5000000000)
  (2016-10-19 22:00:00+00:00,10314.7500000000)
  (2016-10-20 08:00:00+00:00,13988.7500000000)
  (2016-10-20 18:00:00+00:00,13378.2500000000)
  (2016-10-21 04:00:00+00:00,12372.2500000000)
  (2016-10-21 14:00:00+00:00,13198.7500000000)
  (2016-10-22 00:00:00+00:00,9017.2500000000)
  (2016-10-22 10:00:00+00:00,11908.2500000000)
  (2016-10-22 20:00:00+00:00,9986.5000000000)
  (2016-10-23 06:00:00+00:00,9560.7500000000)
  (2016-10-23 16:00:00+00:00,11629.7500000000)
  (2016-10-24 02:00:00+00:00,9475.5000000000)
  (2016-10-24 12:00:00+00:00,13309.0000000000)
  (2016-10-24 22:00:00+00:00,9817.0000000000)
  (2016-10-25 08:00:00+00:00,13831.7500000000)
  (2016-10-25 18:00:00+00:00,12670.5000000000)
  (2016-10-26 04:00:00+00:00,11888.5000000000)
  (2016-10-26 14:00:00+00:00,12911.5000000000)
  (2016-10-27 00:00:00+00:00,9375.2500000000)
  (2016-10-27 10:00:00+00:00,13821.5000000000)
  (2016-10-27 20:00:00+00:00,11381.5000000000)
  (2016-10-28 06:00:00+00:00,13845.7500000000)
  (2016-10-28 16:00:00+00:00,13697.5000000000)
  (2016-10-29 02:00:00+00:00,9403.2500000000)
  (2016-10-29 12:00:00+00:00,11165.2500000000)
  (2016-10-29 22:00:00+00:00,9081.2500000000)
  (2016-10-30 08:00:00+00:00,10628.2500000000)
  (2016-10-30 18:00:00+00:00,11206.0000000000)
  (2016-10-31 04:00:00+00:00,9053.2500000000)
  (2016-10-31 14:00:00+00:00,11201.5000000000)
  (2016-11-01 00:00:00+00:00,9084.0000000000)
  (2016-11-01 10:00:00+00:00,13950.7500000000)
  (2016-11-01 20:00:00+00:00,12439.2500000000)
  (2016-11-02 06:00:00+00:00,13609.2500000000)
  (2016-11-02 16:00:00+00:00,14736.5000000000)
  (2016-11-03 02:00:00+00:00,10507.7500000000)
  (2016-11-03 12:00:00+00:00,13796.7500000000)
  (2016-11-03 22:00:00+00:00,11091.5000000000)
  (2016-11-04 08:00:00+00:00,14322.5000000000)
  (2016-11-04 18:00:00+00:00,14222.5000000000)
  (2016-11-05 04:00:00+00:00,9960.5000000000)
  (2016-11-05 14:00:00+00:00,12076.7500000000)
  (2016-11-06 00:00:00+00:00,9001.7500000000)
  (2016-11-06 10:00:00+00:00,12203.0000000000)
  (2016-11-06 20:00:00+00:00,11069.0000000000)
  (2016-11-07 06:00:00+00:00,13875.0000000000)
  (2016-11-07 16:00:00+00:00,14971.0000000000)
  (2016-11-08 02:00:00+00:00,10601.7500000000)
  (2016-11-08 12:00:00+00:00,14529.0000000000)
  (2016-11-08 22:00:00+00:00,11454.5000000000)
  (2016-11-09 08:00:00+00:00,14283.2500000000)
  (2016-11-09 18:00:00+00:00,14749.5000000000)
  (2016-11-10 04:00:00+00:00,11524.0000000000)
  (2016-11-10 14:00:00+00:00,14329.2500000000)
  (2016-11-11 00:00:00+00:00,10644.5000000000)
  (2016-11-11 10:00:00+00:00,14919.5000000000)
  (2016-11-11 20:00:00+00:00,12730.5000000000)
  (2016-11-12 06:00:00+00:00,11353.2500000000)
  (2016-11-12 16:00:00+00:00,13484.0000000000)
  (2016-11-13 02:00:00+00:00,10197.7500000000)
  (2016-11-13 12:00:00+00:00,11594.5000000000)
  (2016-11-13 22:00:00+00:00,11259.7500000000)
  (2016-11-14 08:00:00+00:00,14645.5000000000)
  (2016-11-14 18:00:00+00:00,15173.5000000000)
  (2016-11-15 04:00:00+00:00,12066.2500000000)
  (2016-11-15 14:00:00+00:00,14685.5000000000)
  (2016-11-16 00:00:00+00:00,10079.2500000000)
  (2016-11-16 10:00:00+00:00,14177.2500000000)
  (2016-11-16 20:00:00+00:00,12369.2500000000)
  (2016-11-17 06:00:00+00:00,13972.5000000000)
  (2016-11-17 16:00:00+00:00,15168.2500000000)
  (2016-11-18 02:00:00+00:00,10054.7500000000)
  (2016-11-18 12:00:00+00:00,13854.5000000000)
  (2016-11-18 22:00:00+00:00,11120.2500000000)
  (2016-11-19 08:00:00+00:00,12505.2500000000)
  (2016-11-19 18:00:00+00:00,12763.2500000000)
  (2016-11-20 04:00:00+00:00,8910.2500000000)
  (2016-11-20 14:00:00+00:00,10974.2500000000)
  (2016-11-21 00:00:00+00:00,9270.0000000000)
  (2016-11-21 10:00:00+00:00,14307.2500000000)
  (2016-11-21 20:00:00+00:00,12678.7500000000)
  (2016-11-22 06:00:00+00:00,14079.5000000000)
  (2016-11-22 16:00:00+00:00,14935.7500000000)
  (2016-11-23 02:00:00+00:00,10238.2500000000)
  (2016-11-23 12:00:00+00:00,13857.2500000000)
  (2016-11-23 22:00:00+00:00,11284.7500000000)
  (2016-11-24 08:00:00+00:00,14435.2500000000)
  (2016-11-24 18:00:00+00:00,14842.7500000000)
  (2016-11-25 04:00:00+00:00,11640.2500000000)
  (2016-11-25 14:00:00+00:00,13980.0000000000)
  (2016-11-26 00:00:00+00:00,10363.2500000000)
  (2016-11-26 10:00:00+00:00,13225.7500000000)
  (2016-11-26 20:00:00+00:00,11600.5000000000)
  (2016-11-27 06:00:00+00:00,10089.2500000000)
  (2016-11-27 16:00:00+00:00,12279.0000000000)
  (2016-11-28 02:00:00+00:00,10279.2500000000)
  (2016-11-28 12:00:00+00:00,14491.2500000000)
  (2016-11-28 22:00:00+00:00,12024.5000000000)
  (2016-11-29 08:00:00+00:00,14687.5000000000)
  (2016-11-29 18:00:00+00:00,15582.5000000000)
  (2016-11-30 04:00:00+00:00,12225.0000000000)
  (2016-11-30 14:00:00+00:00,14911.2500000000)
  (2016-12-01 00:00:00+00:00,10804.0000000000)
  (2016-12-01 10:00:00+00:00,15295.7500000000)
  (2016-12-01 20:00:00+00:00,13433.2500000000)
  (2016-12-02 06:00:00+00:00,14716.0000000000)
  (2016-12-02 16:00:00+00:00,15327.2500000000)
  (2016-12-03 02:00:00+00:00,10874.2500000000)
  (2016-12-03 12:00:00+00:00,12489.5000000000)
  (2016-12-03 22:00:00+00:00,11254.2500000000)
  (2016-12-04 08:00:00+00:00,11808.7500000000)
  (2016-12-04 18:00:00+00:00,13065.7500000000)
  (2016-12-05 04:00:00+00:00,11992.5000000000)
  (2016-12-05 14:00:00+00:00,14587.5000000000)
  (2016-12-06 00:00:00+00:00,11788.2500000000)
  (2016-12-06 10:00:00+00:00,15330.2500000000)
  (2016-12-06 20:00:00+00:00,13669.2500000000)
  (2016-12-07 06:00:00+00:00,15176.5000000000)
  (2016-12-07 16:00:00+00:00,15986.0000000000)
  (2016-12-08 02:00:00+00:00,11293.5000000000)
  (2016-12-08 12:00:00+00:00,14305.5000000000)
  (2016-12-08 22:00:00+00:00,11496.5000000000)
  (2016-12-09 08:00:00+00:00,14859.7500000000)
  (2016-12-09 18:00:00+00:00,14161.5000000000)
  (2016-12-10 04:00:00+00:00,10094.5000000000)
  (2016-12-10 14:00:00+00:00,11928.5000000000)
  (2016-12-11 00:00:00+00:00,8890.5000000000)
  (2016-12-11 10:00:00+00:00,12683.0000000000)
  (2016-12-11 20:00:00+00:00,11275.7500000000)
  (2016-12-12 06:00:00+00:00,14055.0000000000)
  (2016-12-12 16:00:00+00:00,15390.0000000000)
  (2016-12-13 02:00:00+00:00,11173.5000000000)
  (2016-12-13 12:00:00+00:00,14582.2500000000)
  (2016-12-13 22:00:00+00:00,11815.0000000000)
  (2016-12-14 08:00:00+00:00,14862.0000000000)
  (2016-12-14 18:00:00+00:00,15016.0000000000)
  (2016-12-15 04:00:00+00:00,12035.7500000000)
  (2016-12-15 14:00:00+00:00,14699.2500000000)
  (2016-12-16 00:00:00+00:00,11307.7500000000)
  (2016-12-16 10:00:00+00:00,14859.2500000000)
  (2016-12-16 20:00:00+00:00,12968.2500000000)
  (2016-12-17 06:00:00+00:00,11491.2500000000)
  (2016-12-17 16:00:00+00:00,13589.7500000000)
  (2016-12-18 02:00:00+00:00,9692.0000000000)
  (2016-12-18 12:00:00+00:00,12030.2500000000)
  (2016-12-18 22:00:00+00:00,10773.0000000000)
  (2016-12-19 08:00:00+00:00,14394.2500000000)
  (2016-12-19 18:00:00+00:00,14546.7500000000)
  (2016-12-20 04:00:00+00:00,11648.7500000000)
  (2016-12-20 14:00:00+00:00,14315.0000000000)
  (2016-12-21 00:00:00+00:00,10900.7500000000)
  (2016-12-21 10:00:00+00:00,14678.5000000000)
  (2016-12-21 20:00:00+00:00,13202.2500000000)
  (2016-12-22 06:00:00+00:00,13933.5000000000)
  (2016-12-22 16:00:00+00:00,14674.5000000000)
  (2016-12-23 02:00:00+00:00,10334.2500000000)
  (2016-12-23 12:00:00+00:00,13268.7500000000)
  (2016-12-23 22:00:00+00:00,10161.5000000000)
  (2016-12-24 08:00:00+00:00,11427.5000000000)
  (2016-12-24 18:00:00+00:00,10656.7500000000)
  (2016-12-25 04:00:00+00:00,8131.0000000000)
  (2016-12-25 14:00:00+00:00,9962.7500000000)
  (2016-12-26 00:00:00+00:00,8126.7500000000)
  (2016-12-26 10:00:00+00:00,11399.0000000000)
  (2016-12-26 20:00:00+00:00,10303.7500000000)
  (2016-12-27 06:00:00+00:00,10937.2500000000)
  (2016-12-27 16:00:00+00:00,13074.2500000000)
  (2016-12-28 02:00:00+00:00,8999.2500000000)
  (2016-12-28 12:00:00+00:00,12520.2500000000)
  (2016-12-28 22:00:00+00:00,10157.0000000000)
  (2016-12-29 08:00:00+00:00,12033.7500000000)
  (2016-12-29 18:00:00+00:00,12945.2500000000)
  (2016-12-30 04:00:00+00:00,9603.7500000000)
  (2016-12-30 14:00:00+00:00,11301.2500000000)
  (2016-12-31 00:00:00+00:00,9102.5000000000)
  (2016-12-31 10:00:00+00:00,11406.0000000000)
  (2016-12-31 20:00:00+00:00,10026.2500000000)
  (2017-01-01 06:00:00+00:00,8384.0000000000)
  (2017-01-01 16:00:00+00:00,11144.2500000000)
  (2017-01-02 02:00:00+00:00,8890.0000000000)
  (2017-01-02 12:00:00+00:00,13551.2500000000)
  (2017-01-02 22:00:00+00:00,11304.7500000000)
  (2017-01-03 08:00:00+00:00,14267.7500000000)
  (2017-01-03 18:00:00+00:00,14800.0000000000)
  (2017-01-04 04:00:00+00:00,11640.5000000000)
  (2017-01-04 14:00:00+00:00,14322.7500000000)
  (2017-01-05 00:00:00+00:00,10828.5000000000)
  (2017-01-05 10:00:00+00:00,14585.2500000000)
  (2017-01-05 20:00:00+00:00,13517.7500000000)
  (2017-01-06 06:00:00+00:00,14279.0000000000)
  (2017-01-06 16:00:00+00:00,15170.2500000000)
  (2017-01-07 02:00:00+00:00,11357.5000000000)
  (2017-01-07 12:00:00+00:00,12967.0000000000)
  (2017-01-07 22:00:00+00:00,11104.5000000000)
  (2017-01-08 08:00:00+00:00,11918.5000000000)
  (2017-01-08 18:00:00+00:00,13011.0000000000)
  (2017-01-09 04:00:00+00:00,11820.2500000000)
  (2017-01-09 14:00:00+00:00,14633.0000000000)
  (2017-01-10 00:00:00+00:00,11631.0000000000)
  (2017-01-10 10:00:00+00:00,15113.5000000000)
  (2017-01-10 20:00:00+00:00,13744.2500000000)
  (2017-01-11 06:00:00+00:00,14658.0000000000)
  (2017-01-11 16:00:00+00:00,15400.5000000000)
  (2017-01-12 02:00:00+00:00,11203.2500000000)
  (2017-01-12 12:00:00+00:00,14490.2500000000)
  (2017-01-12 22:00:00+00:00,11526.5000000000)
  (2017-01-13 08:00:00+00:00,14738.2500000000)
  (2017-01-13 18:00:00+00:00,14291.7500000000)
  (2017-01-14 04:00:00+00:00,10407.5000000000)
  (2017-01-14 14:00:00+00:00,12221.7500000000)
  (2017-01-15 00:00:00+00:00,9955.5000000000)
  (2017-01-15 10:00:00+00:00,12690.0000000000)
  (2017-01-15 20:00:00+00:00,11475.2500000000)
  (2017-01-16 06:00:00+00:00,14280.7500000000)
  (2017-01-16 16:00:00+00:00,14910.2500000000)
  (2017-01-17 02:00:00+00:00,11167.2500000000)
  (2017-01-17 12:00:00+00:00,14389.2500000000)
  (2017-01-17 22:00:00+00:00,12100.7500000000)
  (2017-01-18 08:00:00+00:00,14802.2500000000)
  (2017-01-18 18:00:00+00:00,15317.7500000000)
  (2017-01-19 04:00:00+00:00,12686.0000000000)
  (2017-01-19 14:00:00+00:00,14610.5000000000)
  (2017-01-20 00:00:00+00:00,11886.5000000000)
  (2017-01-20 10:00:00+00:00,15376.0000000000)
  (2017-01-20 20:00:00+00:00,13168.7500000000)
  (2017-01-21 06:00:00+00:00,11539.0000000000)
  (2017-01-21 16:00:00+00:00,13520.2500000000)
  (2017-01-22 02:00:00+00:00,10319.7500000000)
  (2017-01-22 12:00:00+00:00,12002.7500000000)
  (2017-01-22 22:00:00+00:00,11646.0000000000)
  (2017-01-23 08:00:00+00:00,15237.7500000000)
  (2017-01-23 18:00:00+00:00,15355.0000000000)
  (2017-01-24 04:00:00+00:00,12596.7500000000)
  (2017-01-24 14:00:00+00:00,14933.0000000000)
  (2017-01-25 00:00:00+00:00,11611.2500000000)
  (2017-01-25 10:00:00+00:00,15361.7500000000)
  (2017-01-25 20:00:00+00:00,13401.7500000000)
  (2017-01-26 06:00:00+00:00,14810.2500000000)
  (2017-01-26 16:00:00+00:00,15366.0000000000)
  (2017-01-27 02:00:00+00:00,11847.2500000000)
  (2017-01-27 12:00:00+00:00,14521.5000000000)
  (2017-01-27 22:00:00+00:00,11972.2500000000)
  (2017-01-28 08:00:00+00:00,12737.7500000000)
  (2017-01-28 18:00:00+00:00,13095.2500000000)
  (2017-01-29 04:00:00+00:00,10008.7500000000)
  (2017-01-29 14:00:00+00:00,11208.5000000000)
  (2017-01-30 00:00:00+00:00,10084.5000000000)
  (2017-01-30 10:00:00+00:00,14841.0000000000)
  (2017-01-30 20:00:00+00:00,13080.2500000000)
  (2017-01-31 06:00:00+00:00,14653.5000000000)
  (2017-01-31 16:00:00+00:00,15166.5000000000)
  (2017-02-01 02:00:00+00:00,11535.0000000000)
  (2017-02-01 12:00:00+00:00,14939.5000000000)
  (2017-02-01 22:00:00+00:00,12472.5000000000)
  (2017-02-02 08:00:00+00:00,15116.7500000000)
  (2017-02-02 18:00:00+00:00,15163.7500000000)
  (2017-02-03 04:00:00+00:00,12202.0000000000)
  (2017-02-03 14:00:00+00:00,14222.7500000000)
  (2017-02-04 00:00:00+00:00,10404.7500000000)
  (2017-02-04 10:00:00+00:00,12924.5000000000)
  (2017-02-04 20:00:00+00:00,11178.7500000000)
  (2017-02-05 06:00:00+00:00,9633.2500000000)
  (2017-02-05 16:00:00+00:00,12023.0000000000)
  (2017-02-06 02:00:00+00:00,10213.0000000000)
  (2017-02-06 12:00:00+00:00,14895.2500000000)
  (2017-02-06 22:00:00+00:00,12165.5000000000)
  (2017-02-07 08:00:00+00:00,15199.7500000000)
  (2017-02-07 18:00:00+00:00,15235.7500000000)
  (2017-02-08 04:00:00+00:00,12207.0000000000)
  (2017-02-08 14:00:00+00:00,14622.7500000000)
  (2017-02-09 00:00:00+00:00,11724.5000000000)
  (2017-02-09 10:00:00+00:00,15215.5000000000)
  (2017-02-09 20:00:00+00:00,13749.7500000000)
  (2017-02-10 06:00:00+00:00,14927.2500000000)
  (2017-02-10 16:00:00+00:00,15104.2500000000)
  (2017-02-11 02:00:00+00:00,10772.2500000000)
  (2017-02-11 12:00:00+00:00,12746.7500000000)
  (2017-02-11 22:00:00+00:00,11004.2500000000)
  (2017-02-12 08:00:00+00:00,11684.0000000000)
  (2017-02-12 18:00:00+00:00,12935.0000000000)
  (2017-02-13 04:00:00+00:00,12242.0000000000)
  (2017-02-13 14:00:00+00:00,14024.0000000000)
  (2017-02-14 00:00:00+00:00,11671.5000000000)
  (2017-02-14 10:00:00+00:00,14954.2500000000)
  (2017-02-14 20:00:00+00:00,13519.5000000000)
  (2017-02-15 06:00:00+00:00,14438.5000000000)
  (2017-02-15 16:00:00+00:00,14252.2500000000)
  (2017-02-16 02:00:00+00:00,11330.2500000000)
  (2017-02-16 12:00:00+00:00,14076.0000000000)
  (2017-02-16 22:00:00+00:00,11865.2500000000)
  (2017-02-17 08:00:00+00:00,14527.0000000000)
  (2017-02-17 18:00:00+00:00,14121.0000000000)
  (2017-02-18 04:00:00+00:00,10140.0000000000)
  (2017-02-18 14:00:00+00:00,11816.2500000000)
  (2017-02-19 00:00:00+00:00,9753.7500000000)
  (2017-02-19 10:00:00+00:00,12352.2500000000)
  (2017-02-19 20:00:00+00:00,11219.2500000000)
  (2017-02-20 06:00:00+00:00,13485.5000000000)
  (2017-02-20 16:00:00+00:00,14487.7500000000)
  (2017-02-21 02:00:00+00:00,10363.0000000000)
  (2017-02-21 12:00:00+00:00,14159.2500000000)
  (2017-02-21 22:00:00+00:00,11652.5000000000)
  (2017-02-22 08:00:00+00:00,14306.2500000000)
  (2017-02-22 18:00:00+00:00,14576.2500000000)
  (2017-02-23 04:00:00+00:00,11524.5000000000)
  (2017-02-23 14:00:00+00:00,14048.0000000000)
  (2017-02-24 00:00:00+00:00,10539.7500000000)
  (2017-02-24 10:00:00+00:00,14489.0000000000)
  (2017-02-24 20:00:00+00:00,12749.0000000000)
  (2017-02-25 06:00:00+00:00,11202.0000000000)
  (2017-02-25 16:00:00+00:00,12485.2500000000)
  (2017-02-26 02:00:00+00:00,9159.2500000000)
  (2017-02-26 12:00:00+00:00,11123.0000000000)
  (2017-02-26 22:00:00+00:00,10544.5000000000)
  (2017-02-27 08:00:00+00:00,14079.5000000000)
  (2017-02-27 18:00:00+00:00,14146.2500000000)
  (2017-02-28 04:00:00+00:00,11368.7500000000)
  (2017-02-28 14:00:00+00:00,13419.5000000000)
  (2017-03-01 00:00:00+00:00,10800.0000000000)
  (2017-03-01 10:00:00+00:00,14572.0000000000)
  (2017-03-01 20:00:00+00:00,13029.5000000000)
  (2017-03-02 06:00:00+00:00,13669.5000000000)
  (2017-03-02 16:00:00+00:00,14041.7500000000)
  (2017-03-03 02:00:00+00:00,10870.5000000000)
  (2017-03-03 12:00:00+00:00,13668.5000000000)
  (2017-03-03 22:00:00+00:00,11106.2500000000)
  (2017-03-04 08:00:00+00:00,11863.2500000000)
  (2017-03-04 18:00:00+00:00,12123.0000000000)
  (2017-03-05 04:00:00+00:00,8901.7500000000)
  (2017-03-05 14:00:00+00:00,10225.5000000000)
  (2017-03-06 00:00:00+00:00,9826.2500000000)
  (2017-03-06 10:00:00+00:00,14287.0000000000)
  (2017-03-06 20:00:00+00:00,12966.7500000000)
  (2017-03-07 06:00:00+00:00,14006.2500000000)
  (2017-03-07 16:00:00+00:00,14078.5000000000)
  (2017-03-08 02:00:00+00:00,10776.2500000000)
  (2017-03-08 12:00:00+00:00,14205.7500000000)
  (2017-03-08 22:00:00+00:00,11421.2500000000)
  (2017-03-09 08:00:00+00:00,14121.7500000000)
  (2017-03-09 18:00:00+00:00,14724.5000000000)
  (2017-03-10 04:00:00+00:00,11977.0000000000)
  (2017-03-10 14:00:00+00:00,13481.7500000000)
  (2017-03-11 00:00:00+00:00,10387.0000000000)
  (2017-03-11 10:00:00+00:00,12793.0000000000)
  (2017-03-11 20:00:00+00:00,11317.0000000000)
  (2017-03-12 06:00:00+00:00,9331.5000000000)
  (2017-03-12 16:00:00+00:00,11282.0000000000)
  (2017-03-13 02:00:00+00:00,10221.7500000000)
  (2017-03-13 12:00:00+00:00,13742.2500000000)
  (2017-03-13 22:00:00+00:00,11331.7500000000)
  (2017-03-14 08:00:00+00:00,14217.2500000000)
  (2017-03-14 18:00:00+00:00,14709.0000000000)
  (2017-03-15 04:00:00+00:00,11191.2500000000)
  (2017-03-15 14:00:00+00:00,12979.0000000000)
  (2017-03-16 00:00:00+00:00,10058.0000000000)
  (2017-03-16 10:00:00+00:00,13976.0000000000)
  (2017-03-16 20:00:00+00:00,12781.5000000000)
  (2017-03-17 06:00:00+00:00,13113.0000000000)
  (2017-03-17 16:00:00+00:00,13285.7500000000)
  (2017-03-18 02:00:00+00:00,9590.7500000000)
  (2017-03-18 12:00:00+00:00,12084.7500000000)
  (2017-03-18 22:00:00+00:00,10030.0000000000)
  (2017-03-19 08:00:00+00:00,10551.2500000000)
  (2017-03-19 18:00:00+00:00,12039.5000000000)
  (2017-03-20 04:00:00+00:00,10629.0000000000)
  (2017-03-20 14:00:00+00:00,13111.0000000000)
  (2017-03-21 00:00:00+00:00,9974.5000000000)
  (2017-03-21 10:00:00+00:00,14036.2500000000)
  (2017-03-21 20:00:00+00:00,12787.7500000000)
  (2017-03-22 06:00:00+00:00,13571.2500000000)
  (2017-03-22 16:00:00+00:00,13197.0000000000)
  (2017-03-23 02:00:00+00:00,10434.5000000000)
  (2017-03-23 12:00:00+00:00,13919.0000000000)
  (2017-03-23 22:00:00+00:00,11115.0000000000)
  (2017-03-24 08:00:00+00:00,13887.0000000000)
  (2017-03-24 18:00:00+00:00,13534.2500000000)
  (2017-03-25 04:00:00+00:00,9651.0000000000)
  (2017-03-25 14:00:00+00:00,10640.0000000000)
  (2017-03-26 00:00:00+00:00,8719.0000000000)
  (2017-03-26 10:00:00+00:00,11045.0000000000)
  (2017-03-26 20:00:00+00:00,10326.2500000000)
  (2017-03-27 06:00:00+00:00,13492.5000000000)
  (2017-03-27 16:00:00+00:00,12577.0000000000)
  (2017-03-28 02:00:00+00:00,10431.0000000000)
  (2017-03-28 12:00:00+00:00,13045.5000000000)
  (2017-03-28 22:00:00+00:00,10225.5000000000)
  (2017-03-29 08:00:00+00:00,13772.2500000000)
  (2017-03-29 18:00:00+00:00,12843.0000000000)
  (2017-03-30 04:00:00+00:00,12418.0000000000)
  (2017-03-30 14:00:00+00:00,12272.2500000000)
  (2017-03-31 00:00:00+00:00,9530.0000000000)
  (2017-03-31 10:00:00+00:00,13037.2500000000)
  (2017-03-31 20:00:00+00:00,10885.5000000000)
  (2017-04-01 06:00:00+00:00,10868.7500000000)
  (2017-04-01 16:00:00+00:00,10641.0000000000)
  (2017-04-02 02:00:00+00:00,8253.0000000000)
  (2017-04-02 12:00:00+00:00,9886.5000000000)
  (2017-04-02 22:00:00+00:00,9028.2500000000)
  (2017-04-03 08:00:00+00:00,13280.0000000000)
  (2017-04-03 18:00:00+00:00,12835.0000000000)
  (2017-04-04 04:00:00+00:00,11951.0000000000)
  (2017-04-04 14:00:00+00:00,12606.7500000000)
  (2017-04-05 00:00:00+00:00,9327.7500000000)
  (2017-04-05 10:00:00+00:00,13597.0000000000)
  (2017-04-05 20:00:00+00:00,11456.7500000000)
  (2017-04-06 06:00:00+00:00,13414.7500000000)
  (2017-04-06 16:00:00+00:00,12975.2500000000)
  (2017-04-07 02:00:00+00:00,10204.2500000000)
  (2017-04-07 12:00:00+00:00,13430.2500000000)
  (2017-04-07 22:00:00+00:00,9759.5000000000)
  (2017-04-08 08:00:00+00:00,11961.5000000000)
  (2017-04-08 18:00:00+00:00,11083.5000000000)
  (2017-04-09 04:00:00+00:00,8415.5000000000)
  (2017-04-09 14:00:00+00:00,9845.5000000000)
  (2017-04-10 00:00:00+00:00,9268.7500000000)
  (2017-04-10 10:00:00+00:00,13778.7500000000)
  (2017-04-10 20:00:00+00:00,11676.7500000000)
  (2017-04-11 06:00:00+00:00,13662.2500000000)
  (2017-04-11 16:00:00+00:00,13227.5000000000)
  (2017-04-12 02:00:00+00:00,10407.0000000000)
  (2017-04-12 12:00:00+00:00,13290.0000000000)
  (2017-04-12 22:00:00+00:00,10155.0000000000)
  (2017-04-13 08:00:00+00:00,13685.0000000000)
  (2017-04-13 18:00:00+00:00,12315.7500000000)
  (2017-04-14 04:00:00+00:00,8944.5000000000)
  (2017-04-14 14:00:00+00:00,9836.2500000000)
  (2017-04-15 00:00:00+00:00,8468.5000000000)
  (2017-04-15 10:00:00+00:00,11536.0000000000)
  (2017-04-15 20:00:00+00:00,9894.5000000000)
  (2017-04-16 06:00:00+00:00,9568.0000000000)
  (2017-04-16 16:00:00+00:00,9369.5000000000)
  (2017-04-17 02:00:00+00:00,8166.7500000000)
  (2017-04-17 12:00:00+00:00,9710.7500000000)
  (2017-04-17 22:00:00+00:00,9344.0000000000)
  (2017-04-18 08:00:00+00:00,13659.0000000000)
  (2017-04-18 18:00:00+00:00,12827.2500000000)
  (2017-04-19 04:00:00+00:00,12372.0000000000)
  (2017-04-19 14:00:00+00:00,13291.2500000000)
  (2017-04-20 00:00:00+00:00,10363.5000000000)
  (2017-04-20 10:00:00+00:00,13988.2500000000)
  (2017-04-20 20:00:00+00:00,12143.5000000000)
  (2017-04-21 06:00:00+00:00,13955.2500000000)
  (2017-04-21 16:00:00+00:00,13482.7500000000)
  (2017-04-22 02:00:00+00:00,9576.0000000000)
  (2017-04-22 12:00:00+00:00,11524.5000000000)
  (2017-04-22 22:00:00+00:00,9361.0000000000)
  (2017-04-23 08:00:00+00:00,11114.5000000000)
  (2017-04-23 18:00:00+00:00,11404.0000000000)
  (2017-04-24 04:00:00+00:00,12446.7500000000)
  (2017-04-24 14:00:00+00:00,12882.0000000000)
  (2017-04-25 00:00:00+00:00,9809.0000000000)
  (2017-04-25 10:00:00+00:00,13941.5000000000)
  (2017-04-25 20:00:00+00:00,12045.2500000000)
  (2017-04-26 06:00:00+00:00,13746.7500000000)
  (2017-04-26 16:00:00+00:00,13200.0000000000)
  (2017-04-27 02:00:00+00:00,10700.0000000000)
  (2017-04-27 12:00:00+00:00,13342.0000000000)
  (2017-04-27 22:00:00+00:00,10480.7500000000)
  (2017-04-28 08:00:00+00:00,13862.7500000000)
  (2017-04-28 18:00:00+00:00,12683.7500000000)
  (2017-04-29 04:00:00+00:00,9893.0000000000)
  (2017-04-29 14:00:00+00:00,10639.7500000000)
  (2017-04-30 00:00:00+00:00,8447.5000000000)
  (2017-04-30 10:00:00+00:00,10529.7500000000)
  (2017-04-30 20:00:00+00:00,9814.0000000000)
  (2017-05-01 06:00:00+00:00,9352.5000000000)
  (2017-05-01 16:00:00+00:00,10438.7500000000)
  (2017-05-02 02:00:00+00:00,9300.7500000000)
  (2017-05-02 12:00:00+00:00,13978.7500000000)
  (2017-05-02 22:00:00+00:00,10528.7500000000)
  (2017-05-03 08:00:00+00:00,13868.7500000000)
  (2017-05-03 18:00:00+00:00,13086.5000000000)
  (2017-05-04 04:00:00+00:00,12486.0000000000)
  (2017-05-04 14:00:00+00:00,13649.0000000000)
  (2017-05-05 00:00:00+00:00,10133.7500000000)
  (2017-05-05 10:00:00+00:00,14036.7500000000)
  (2017-05-05 20:00:00+00:00,11371.5000000000)
  (2017-05-06 06:00:00+00:00,11230.2500000000)
  (2017-05-06 16:00:00+00:00,11084.5000000000)
  (2017-05-07 02:00:00+00:00,8723.7500000000)
  (2017-05-07 12:00:00+00:00,10461.0000000000)
  (2017-05-07 22:00:00+00:00,9064.7500000000)
  (2017-05-08 08:00:00+00:00,13763.0000000000)
  (2017-05-08 18:00:00+00:00,12438.0000000000)
  (2017-05-09 04:00:00+00:00,11808.2500000000)
  (2017-05-09 14:00:00+00:00,12566.7500000000)
  (2017-05-10 00:00:00+00:00,9490.7500000000)
  (2017-05-10 10:00:00+00:00,13359.5000000000)
  (2017-05-10 20:00:00+00:00,11541.2500000000)
  (2017-05-11 06:00:00+00:00,12913.2500000000)
  (2017-05-11 16:00:00+00:00,12372.2500000000)
  (2017-05-12 02:00:00+00:00,9674.2500000000)
  (2017-05-12 12:00:00+00:00,12779.0000000000)
  (2017-05-12 22:00:00+00:00,9377.7500000000)
  (2017-05-13 08:00:00+00:00,11468.5000000000)
  (2017-05-13 18:00:00+00:00,10200.7500000000)
  (2017-05-14 04:00:00+00:00,7779.0000000000)
  (2017-05-14 14:00:00+00:00,9354.2500000000)
  (2017-05-15 00:00:00+00:00,8663.5000000000)
  (2017-05-15 10:00:00+00:00,13294.7500000000)
  (2017-05-15 20:00:00+00:00,11328.7500000000)
  (2017-05-16 06:00:00+00:00,13174.7500000000)
  (2017-05-16 16:00:00+00:00,12829.2500000000)
  (2017-05-17 02:00:00+00:00,9725.2500000000)
  (2017-05-17 12:00:00+00:00,13043.2500000000)
  (2017-05-17 22:00:00+00:00,10193.5000000000)
  (2017-05-18 08:00:00+00:00,13592.5000000000)
  (2017-05-18 18:00:00+00:00,12290.7500000000)
  (2017-05-19 04:00:00+00:00,11572.5000000000)
  (2017-05-19 14:00:00+00:00,12727.5000000000)
  (2017-05-20 00:00:00+00:00,9235.7500000000)
  (2017-05-20 10:00:00+00:00,11790.2500000000)
  (2017-05-20 20:00:00+00:00,9696.5000000000)
  (2017-05-21 06:00:00+00:00,9506.5000000000)
  (2017-05-21 16:00:00+00:00,10345.0000000000)
  (2017-05-22 02:00:00+00:00,9223.2500000000)
  (2017-05-22 12:00:00+00:00,12839.5000000000)
  (2017-05-22 22:00:00+00:00,9766.0000000000)
  (2017-05-23 08:00:00+00:00,13181.2500000000)
  (2017-05-23 18:00:00+00:00,12612.2500000000)
  (2017-05-24 04:00:00+00:00,11743.5000000000)
  (2017-05-24 14:00:00+00:00,12712.5000000000)
  (2017-05-25 00:00:00+00:00,8753.2500000000)
  (2017-05-25 10:00:00+00:00,10356.5000000000)
  (2017-05-25 20:00:00+00:00,9769.5000000000)
  (2017-05-26 06:00:00+00:00,11204.2500000000)
  (2017-05-26 16:00:00+00:00,11298.5000000000)
  (2017-05-27 02:00:00+00:00,8295.7500000000)
  (2017-05-27 12:00:00+00:00,10473.5000000000)
  (2017-05-27 22:00:00+00:00,8811.5000000000)
  (2017-05-28 08:00:00+00:00,10081.2500000000)
  (2017-05-28 18:00:00+00:00,10367.0000000000)
  (2017-05-29 04:00:00+00:00,11096.7500000000)
  (2017-05-29 14:00:00+00:00,12663.0000000000)
  (2017-05-30 00:00:00+00:00,9354.7500000000)
  (2017-05-30 10:00:00+00:00,14000.2500000000)
  (2017-05-30 20:00:00+00:00,11553.0000000000)
  (2017-05-31 06:00:00+00:00,13596.0000000000)
  (2017-05-31 16:00:00+00:00,13043.5000000000)
  (2017-06-01 02:00:00+00:00,9684.7500000000)
  (2017-06-01 12:00:00+00:00,12964.5000000000)
  (2017-06-01 22:00:00+00:00,9864.5000000000)
  (2017-06-02 08:00:00+00:00,13391.7500000000)
  (2017-06-02 18:00:00+00:00,11854.7500000000)
  (2017-06-03 04:00:00+00:00,9196.0000000000)
  (2017-06-03 14:00:00+00:00,10664.2500000000)
  (2017-06-04 00:00:00+00:00,7890.7500000000)
  (2017-06-04 10:00:00+00:00,10473.2500000000)
  (2017-06-04 20:00:00+00:00,9287.7500000000)
  (2017-06-05 06:00:00+00:00,8973.7500000000)
  (2017-06-05 16:00:00+00:00,10044.7500000000)
  (2017-06-06 02:00:00+00:00,8487.5000000000)
  (2017-06-06 12:00:00+00:00,13153.2500000000)
  (2017-06-06 22:00:00+00:00,9593.2500000000)
  (2017-06-07 08:00:00+00:00,13139.2500000000)
  (2017-06-07 18:00:00+00:00,12581.2500000000)
  (2017-06-08 04:00:00+00:00,11695.2500000000)
  (2017-06-08 14:00:00+00:00,12695.2500000000)
  (2017-06-09 00:00:00+00:00,9440.5000000000)
  (2017-06-09 10:00:00+00:00,13890.7500000000)
  (2017-06-09 20:00:00+00:00,10985.5000000000)
  (2017-06-10 06:00:00+00:00,10978.2500000000)
  (2017-06-10 16:00:00+00:00,10755.7500000000)
  (2017-06-11 02:00:00+00:00,7506.5000000000)
  (2017-06-11 12:00:00+00:00,9614.2500000000)
  (2017-06-11 22:00:00+00:00,8960.5000000000)
  (2017-06-12 08:00:00+00:00,13364.7500000000)
  (2017-06-12 18:00:00+00:00,12218.5000000000)
  (2017-06-13 04:00:00+00:00,11519.0000000000)
  (2017-06-13 14:00:00+00:00,12737.2500000000)
  (2017-06-14 00:00:00+00:00,9045.0000000000)
  (2017-06-14 10:00:00+00:00,13042.7500000000)
  (2017-06-14 20:00:00+00:00,10883.2500000000)
  (2017-06-15 06:00:00+00:00,12632.0000000000)
  (2017-06-15 16:00:00+00:00,12578.0000000000)
  (2017-06-16 02:00:00+00:00,9384.5000000000)
  (2017-06-16 12:00:00+00:00,12656.7500000000)
  (2017-06-16 22:00:00+00:00,9614.5000000000)
  (2017-06-17 08:00:00+00:00,11506.0000000000)
  (2017-06-17 18:00:00+00:00,9806.7500000000)
  (2017-06-18 04:00:00+00:00,7685.2500000000)
  (2017-06-18 14:00:00+00:00,9445.2500000000)
  (2017-06-19 00:00:00+00:00,8432.2500000000)
  (2017-06-19 10:00:00+00:00,13305.5000000000)
  (2017-06-19 20:00:00+00:00,11090.7500000000)
  (2017-06-20 06:00:00+00:00,13093.2500000000)
  (2017-06-20 16:00:00+00:00,12805.7500000000)
  (2017-06-21 02:00:00+00:00,9179.5000000000)
  (2017-06-21 12:00:00+00:00,12928.2500000000)
  (2017-06-21 22:00:00+00:00,9630.5000000000)
  (2017-06-22 08:00:00+00:00,13176.7500000000)
  (2017-06-22 18:00:00+00:00,12532.0000000000)
  (2017-06-23 04:00:00+00:00,11721.0000000000)
  (2017-06-23 14:00:00+00:00,12536.5000000000)
  (2017-06-24 00:00:00+00:00,8730.5000000000)
  (2017-06-24 10:00:00+00:00,11630.0000000000)
  (2017-06-24 20:00:00+00:00,9670.0000000000)
  (2017-06-25 06:00:00+00:00,9207.7500000000)
  (2017-06-25 16:00:00+00:00,10231.5000000000)
  (2017-06-26 02:00:00+00:00,8851.2500000000)
  (2017-06-26 12:00:00+00:00,12775.0000000000)
  (2017-06-26 22:00:00+00:00,9552.0000000000)
  (2017-06-27 08:00:00+00:00,12978.7500000000)
  (2017-06-27 18:00:00+00:00,11773.5000000000)
  (2017-06-28 04:00:00+00:00,11150.7500000000)
  (2017-06-28 14:00:00+00:00,12591.0000000000)
  (2017-06-29 00:00:00+00:00,9062.7500000000)
  (2017-06-29 10:00:00+00:00,13644.2500000000)
  (2017-06-29 20:00:00+00:00,11382.7500000000)
  (2017-06-30 06:00:00+00:00,13344.0000000000)
  (2017-06-30 16:00:00+00:00,12407.5000000000)
  (2017-07-01 02:00:00+00:00,8510.0000000000)
  (2017-07-01 12:00:00+00:00,10944.0000000000)
  (2017-07-01 22:00:00+00:00,8548.5000000000)
  (2017-07-02 08:00:00+00:00,10414.2500000000)
  (2017-07-02 18:00:00+00:00,10049.7500000000)
  (2017-07-03 04:00:00+00:00,10816.5000000000)
  (2017-07-03 14:00:00+00:00,12261.5000000000)
  (2017-07-04 00:00:00+00:00,9030.5000000000)
  (2017-07-04 10:00:00+00:00,13193.0000000000)
  (2017-07-04 20:00:00+00:00,10820.5000000000)
  (2017-07-05 06:00:00+00:00,12745.5000000000)
  (2017-07-05 16:00:00+00:00,12164.5000000000)
  (2017-07-06 02:00:00+00:00,9166.2500000000)
  (2017-07-06 12:00:00+00:00,12666.0000000000)
  (2017-07-06 22:00:00+00:00,9629.7500000000)
  (2017-07-07 08:00:00+00:00,13212.7500000000)
  (2017-07-07 18:00:00+00:00,11461.0000000000)
  (2017-07-08 04:00:00+00:00,8893.2500000000)
  (2017-07-08 14:00:00+00:00,10582.5000000000)
  (2017-07-09 00:00:00+00:00,7877.7500000000)
  (2017-07-09 10:00:00+00:00,10365.0000000000)
  (2017-07-09 20:00:00+00:00,9733.5000000000)
  (2017-07-10 06:00:00+00:00,12601.2500000000)
  (2017-07-10 16:00:00+00:00,12502.7500000000)
  (2017-07-11 02:00:00+00:00,9177.2500000000)
  (2017-07-11 12:00:00+00:00,12927.7500000000)
  (2017-07-11 22:00:00+00:00,9459.7500000000)
  (2017-07-12 08:00:00+00:00,13017.2500000000)
  (2017-07-12 18:00:00+00:00,12237.5000000000)
  (2017-07-13 04:00:00+00:00,11161.5000000000)
  (2017-07-13 14:00:00+00:00,12121.5000000000)
  (2017-07-14 00:00:00+00:00,8708.7500000000)
  (2017-07-14 10:00:00+00:00,12774.5000000000)
  (2017-07-14 20:00:00+00:00,10268.2500000000)
  (2017-07-15 06:00:00+00:00,10131.5000000000)
  (2017-07-15 16:00:00+00:00,10114.7500000000)
  (2017-07-16 02:00:00+00:00,7549.2500000000)
  (2017-07-16 12:00:00+00:00,9505.0000000000)
  (2017-07-16 22:00:00+00:00,8515.7500000000)
  (2017-07-17 08:00:00+00:00,12774.2500000000)
  (2017-07-17 18:00:00+00:00,11455.7500000000)
  (2017-07-18 04:00:00+00:00,10788.0000000000)
  (2017-07-18 14:00:00+00:00,12034.2500000000)
  (2017-07-19 00:00:00+00:00,8783.5000000000)
  (2017-07-19 10:00:00+00:00,13096.5000000000)
  (2017-07-19 20:00:00+00:00,11095.5000000000)
  (2017-07-20 06:00:00+00:00,12488.2500000000)
  (2017-07-20 16:00:00+00:00,12479.2500000000)
  (2017-07-21 02:00:00+00:00,9228.7500000000)
  (2017-07-21 12:00:00+00:00,12440.2500000000)
  (2017-07-21 22:00:00+00:00,9258.5000000000)
  (2017-07-22 08:00:00+00:00,11269.0000000000)
  (2017-07-22 18:00:00+00:00,9940.7500000000)
  (2017-07-23 04:00:00+00:00,7710.5000000000)
  (2017-07-23 14:00:00+00:00,9302.5000000000)
  (2017-07-24 00:00:00+00:00,8013.5000000000)
  (2017-07-24 10:00:00+00:00,12694.7500000000)
  (2017-07-24 20:00:00+00:00,10386.2500000000)
  (2017-07-25 06:00:00+00:00,12498.7500000000)
  (2017-07-25 16:00:00+00:00,12219.5000000000)
  (2017-07-26 02:00:00+00:00,8985.5000000000)
  (2017-07-26 12:00:00+00:00,12849.5000000000)
  (2017-07-26 22:00:00+00:00,9086.2500000000)
  (2017-07-27 08:00:00+00:00,12690.2500000000)
  (2017-07-27 18:00:00+00:00,11353.0000000000)
  (2017-07-28 04:00:00+00:00,10517.5000000000)
  (2017-07-28 14:00:00+00:00,12029.2500000000)
  (2017-07-29 00:00:00+00:00,8040.5000000000)
  (2017-07-29 10:00:00+00:00,11128.7500000000)
  (2017-07-29 20:00:00+00:00,9150.0000000000)
  (2017-07-30 06:00:00+00:00,8969.0000000000)
  (2017-07-30 16:00:00+00:00,10131.2500000000)
  (2017-07-31 02:00:00+00:00,8556.0000000000)
  (2017-07-31 12:00:00+00:00,12662.2500000000)
  (2017-07-31 22:00:00+00:00,9220.2500000000)
  (2017-08-01 08:00:00+00:00,13028.2500000000)
  (2017-08-01 18:00:00+00:00,11843.2500000000)
  (2017-08-02 04:00:00+00:00,10786.0000000000)
  (2017-08-02 14:00:00+00:00,12307.7500000000)
  (2017-08-03 00:00:00+00:00,8864.5000000000)
  (2017-08-03 10:00:00+00:00,13265.5000000000)
  (2017-08-03 20:00:00+00:00,10828.0000000000)
  (2017-08-04 06:00:00+00:00,12567.5000000000)
  (2017-08-04 16:00:00+00:00,11991.2500000000)
  (2017-08-05 02:00:00+00:00,8339.5000000000)
  (2017-08-05 12:00:00+00:00,10342.7500000000)
  (2017-08-05 22:00:00+00:00,8241.0000000000)
  (2017-08-06 08:00:00+00:00,9833.2500000000)
  (2017-08-06 18:00:00+00:00,9489.5000000000)
  (2017-08-07 04:00:00+00:00,10209.2500000000)
  (2017-08-07 14:00:00+00:00,11937.5000000000)
  (2017-08-08 00:00:00+00:00,8727.2500000000)
  (2017-08-08 10:00:00+00:00,13337.7500000000)
  (2017-08-08 20:00:00+00:00,10813.0000000000)
  (2017-08-09 06:00:00+00:00,12567.2500000000)
  (2017-08-09 16:00:00+00:00,12307.7500000000)
  (2017-08-10 02:00:00+00:00,9171.0000000000)
  (2017-08-10 12:00:00+00:00,12710.5000000000)
  (2017-08-10 22:00:00+00:00,9513.7500000000)
  (2017-08-11 08:00:00+00:00,13279.2500000000)
  (2017-08-11 18:00:00+00:00,11586.5000000000)
  (2017-08-12 04:00:00+00:00,8713.2500000000)
  (2017-08-12 14:00:00+00:00,9979.5000000000)
  (2017-08-13 00:00:00+00:00,7550.5000000000)
  (2017-08-13 10:00:00+00:00,10116.7500000000)
  (2017-08-13 20:00:00+00:00,9355.7500000000)
  (2017-08-14 06:00:00+00:00,12207.2500000000)
  (2017-08-14 16:00:00+00:00,12249.5000000000)
  (2017-08-15 02:00:00+00:00,9351.2500000000)
  (2017-08-15 12:00:00+00:00,13162.5000000000)
  (2017-08-15 22:00:00+00:00,9990.0000000000)
  (2017-08-16 08:00:00+00:00,13628.2500000000)
  (2017-08-16 18:00:00+00:00,12403.7500000000)
  (2017-08-17 04:00:00+00:00,11341.0000000000)
  (2017-08-17 14:00:00+00:00,12893.2500000000)
  (2017-08-18 00:00:00+00:00,9429.0000000000)
  (2017-08-18 10:00:00+00:00,13665.2500000000)
  (2017-08-18 20:00:00+00:00,10789.0000000000)
  (2017-08-19 06:00:00+00:00,10225.0000000000)
  (2017-08-19 16:00:00+00:00,10201.5000000000)
  (2017-08-20 02:00:00+00:00,7534.5000000000)
  (2017-08-20 12:00:00+00:00,9582.7500000000)
  (2017-08-20 22:00:00+00:00,8844.7500000000)
  (2017-08-21 08:00:00+00:00,13430.7500000000)
  (2017-08-21 18:00:00+00:00,12424.7500000000)
  (2017-08-22 04:00:00+00:00,11410.2500000000)
  (2017-08-22 14:00:00+00:00,12552.0000000000)
  (2017-08-23 00:00:00+00:00,8809.7500000000)
  (2017-08-23 10:00:00+00:00,12982.5000000000)
  (2017-08-23 20:00:00+00:00,10837.5000000000)
  (2017-08-24 06:00:00+00:00,12513.7500000000)
  (2017-08-24 16:00:00+00:00,12338.5000000000)
  (2017-08-25 02:00:00+00:00,9570.7500000000)
  (2017-08-25 12:00:00+00:00,12895.5000000000)
  (2017-08-25 22:00:00+00:00,9335.7500000000)
  (2017-08-26 08:00:00+00:00,11488.2500000000)
  (2017-08-26 18:00:00+00:00,10706.7500000000)
  (2017-08-27 04:00:00+00:00,8263.7500000000)
  (2017-08-27 14:00:00+00:00,9889.0000000000)
  (2017-08-28 00:00:00+00:00,8570.5000000000)
  (2017-08-28 10:00:00+00:00,13299.7500000000)
  (2017-08-28 20:00:00+00:00,10992.7500000000)
  (2017-08-29 06:00:00+00:00,13035.7500000000)
  (2017-08-29 16:00:00+00:00,12858.7500000000)
  (2017-08-30 02:00:00+00:00,9649.5000000000)
  (2017-08-30 12:00:00+00:00,13604.7500000000)
  (2017-08-30 22:00:00+00:00,9941.0000000000)
  (2017-08-31 08:00:00+00:00,13521.7500000000)
  (2017-08-31 18:00:00+00:00,12812.0000000000)
  (2017-09-01 04:00:00+00:00,11669.0000000000)
  (2017-09-01 14:00:00+00:00,12403.7500000000)
  (2017-09-02 00:00:00+00:00,8790.0000000000)
  (2017-09-02 10:00:00+00:00,11574.0000000000)
  (2017-09-02 20:00:00+00:00,9595.7500000000)
  (2017-09-03 06:00:00+00:00,9514.2500000000)
  (2017-09-03 16:00:00+00:00,10438.5000000000)
  (2017-09-04 02:00:00+00:00,9065.2500000000)
  (2017-09-04 12:00:00+00:00,13186.0000000000)
  (2017-09-04 22:00:00+00:00,9657.5000000000)
  (2017-09-05 08:00:00+00:00,13260.2500000000)
  (2017-09-05 18:00:00+00:00,12717.2500000000)
  (2017-09-06 04:00:00+00:00,11954.7500000000)
  (2017-09-06 14:00:00+00:00,12879.5000000000)
  (2017-09-07 00:00:00+00:00,9318.5000000000)
  (2017-09-07 10:00:00+00:00,13563.5000000000)
  (2017-09-07 20:00:00+00:00,11229.0000000000)
  (2017-09-08 06:00:00+00:00,13622.2500000000)
  (2017-09-08 16:00:00+00:00,13029.2500000000)
  (2017-09-09 02:00:00+00:00,8859.7500000000)
  (2017-09-09 12:00:00+00:00,11361.0000000000)
  (2017-09-09 22:00:00+00:00,8696.7500000000)
  (2017-09-10 08:00:00+00:00,10714.0000000000)
  (2017-09-10 18:00:00+00:00,10882.2500000000)
  (2017-09-11 04:00:00+00:00,11408.5000000000)
  (2017-09-11 14:00:00+00:00,12467.5000000000)
  (2017-09-12 00:00:00+00:00,9048.5000000000)
  (2017-09-12 10:00:00+00:00,13218.7500000000)
  (2017-09-12 20:00:00+00:00,11257.5000000000)
  (2017-09-13 06:00:00+00:00,13253.2500000000)
  (2017-09-13 16:00:00+00:00,12760.5000000000)
  (2017-09-14 02:00:00+00:00,9515.7500000000)
  (2017-09-14 12:00:00+00:00,13091.2500000000)
  (2017-09-14 22:00:00+00:00,9691.0000000000)
  (2017-09-15 08:00:00+00:00,13997.7500000000)
  (2017-09-15 18:00:00+00:00,12766.0000000000)
  (2017-09-16 04:00:00+00:00,9590.0000000000)
  (2017-09-16 14:00:00+00:00,10620.0000000000)
  (2017-09-17 00:00:00+00:00,8229.7500000000)
  (2017-09-17 10:00:00+00:00,10894.7500000000)
  (2017-09-17 20:00:00+00:00,10099.0000000000)
  (2017-09-18 06:00:00+00:00,13207.2500000000)
  (2017-09-18 16:00:00+00:00,12821.5000000000)
  (2017-09-19 02:00:00+00:00,9966.5000000000)
  (2017-09-19 12:00:00+00:00,13099.2500000000)
  (2017-09-19 22:00:00+00:00,10065.7500000000)
  (2017-09-20 08:00:00+00:00,13750.7500000000)
  (2017-09-20 18:00:00+00:00,13148.7500000000)
  (2017-09-21 04:00:00+00:00,12397.5000000000)
  (2017-09-21 14:00:00+00:00,12715.0000000000)
  (2017-09-22 00:00:00+00:00,9574.5000000000)
  (2017-09-22 10:00:00+00:00,13538.5000000000)
  (2017-09-22 20:00:00+00:00,11035.2500000000)
  (2017-09-23 06:00:00+00:00,11125.7500000000)
  (2017-09-23 16:00:00+00:00,11310.2500000000)
  (2017-09-24 02:00:00+00:00,8225.7500000000)
  (2017-09-24 12:00:00+00:00,10327.5000000000)
  (2017-09-24 22:00:00+00:00,9233.5000000000)
  (2017-09-25 08:00:00+00:00,13849.0000000000)
  (2017-09-25 18:00:00+00:00,13086.5000000000)
  (2017-09-26 04:00:00+00:00,12265.5000000000)
  (2017-09-26 14:00:00+00:00,12799.2500000000)
  (2017-09-27 00:00:00+00:00,9471.5000000000)
  (2017-09-27 10:00:00+00:00,13699.2500000000)
  (2017-09-27 20:00:00+00:00,11525.0000000000)
  (2017-09-28 06:00:00+00:00,13709.7500000000)
  (2017-09-28 16:00:00+00:00,13178.2500000000)
  (2017-09-29 02:00:00+00:00,10014.5000000000)
  (2017-09-29 12:00:00+00:00,13088.5000000000)
  (2017-09-29 22:00:00+00:00,9822.5000000000)
  (2017-09-30 08:00:00+00:00,11930.5000000000)
  (2017-09-30 18:00:00+00:00,11510.0000000000)
  (2017-10-01 04:00:00+00:00,8848.5000000000)
  (2017-10-01 14:00:00+00:00,10154.0000000000)
  (2017-10-02 00:00:00+00:00,8385.0000000000)
  (2017-10-02 10:00:00+00:00,12775.7500000000)
  (2017-10-02 20:00:00+00:00,10309.0000000000)
  (2017-10-03 06:00:00+00:00,10229.2500000000)
  (2017-10-03 16:00:00+00:00,10787.2500000000)
  (2017-10-04 02:00:00+00:00,9796.0000000000)
  (2017-10-04 12:00:00+00:00,13278.2500000000)
  (2017-10-04 22:00:00+00:00,10219.2500000000)
  (2017-10-05 08:00:00+00:00,14396.0000000000)
  (2017-10-05 18:00:00+00:00,13373.2500000000)
  (2017-10-06 04:00:00+00:00,12633.5000000000)
  (2017-10-06 14:00:00+00:00,13154.5000000000)
  (2017-10-07 00:00:00+00:00,8937.0000000000)
  (2017-10-07 10:00:00+00:00,11891.2500000000)
  (2017-10-07 20:00:00+00:00,10095.5000000000)
  (2017-10-08 06:00:00+00:00,10047.2500000000)
  (2017-10-08 16:00:00+00:00,11470.0000000000)
  (2017-10-09 02:00:00+00:00,10049.7500000000)
  (2017-10-09 12:00:00+00:00,13626.5000000000)
  (2017-10-09 22:00:00+00:00,10307.0000000000)
  (2017-10-10 08:00:00+00:00,14214.7500000000)
  (2017-10-10 18:00:00+00:00,13718.0000000000)
  (2017-10-11 04:00:00+00:00,12817.5000000000)
  (2017-10-11 14:00:00+00:00,13277.7500000000)
  (2017-10-12 00:00:00+00:00,9712.5000000000)
  (2017-10-12 10:00:00+00:00,13885.2500000000)
  (2017-10-12 20:00:00+00:00,12059.2500000000)
  (2017-10-13 06:00:00+00:00,14044.5000000000)
  (2017-10-13 16:00:00+00:00,13491.2500000000)
  (2017-10-14 02:00:00+00:00,9582.0000000000)
  (2017-10-14 12:00:00+00:00,11135.5000000000)
  (2017-10-14 22:00:00+00:00,9110.5000000000)
  (2017-10-15 08:00:00+00:00,11061.0000000000)
  (2017-10-15 18:00:00+00:00,10951.2500000000)
  (2017-10-16 04:00:00+00:00,11976.5000000000)
  (2017-10-16 14:00:00+00:00,12272.7500000000)
  (2017-10-17 00:00:00+00:00,9359.5000000000)
  (2017-10-17 10:00:00+00:00,13284.0000000000)
  (2017-10-17 20:00:00+00:00,11216.7500000000)
  (2017-10-18 06:00:00+00:00,13724.0000000000)
  (2017-10-18 16:00:00+00:00,13929.5000000000)
  (2017-10-19 02:00:00+00:00,10197.5000000000)
  (2017-10-19 12:00:00+00:00,12860.2500000000)
  (2017-10-19 22:00:00+00:00,10325.7500000000)
  (2017-10-20 08:00:00+00:00,14008.7500000000)
  (2017-10-20 18:00:00+00:00,13115.7500000000)
  (2017-10-21 04:00:00+00:00,10099.5000000000)
  (2017-10-21 14:00:00+00:00,11294.2500000000)
  (2017-10-22 00:00:00+00:00,8124.7500000000)
  (2017-10-22 10:00:00+00:00,11051.0000000000)
  (2017-10-22 20:00:00+00:00,10535.5000000000)
  (2017-10-23 06:00:00+00:00,13923.0000000000)
  (2017-10-23 16:00:00+00:00,14196.5000000000)
  (2017-10-24 02:00:00+00:00,10386.0000000000)
  (2017-10-24 12:00:00+00:00,13598.0000000000)
  (2017-10-24 22:00:00+00:00,10077.2500000000)
  (2017-10-25 08:00:00+00:00,14101.0000000000)
  (2017-10-25 18:00:00+00:00,13214.2500000000)
  (2017-10-26 04:00:00+00:00,12269.0000000000)
  (2017-10-26 14:00:00+00:00,12935.5000000000)
  (2017-10-27 00:00:00+00:00,9474.2500000000)
  (2017-10-27 10:00:00+00:00,14088.5000000000)
  (2017-10-27 20:00:00+00:00,11183.2500000000)
  (2017-10-28 06:00:00+00:00,11355.5000000000)
  (2017-10-28 16:00:00+00:00,12632.5000000000)
  (2017-10-29 02:00:00+00:00,8559.7500000000)
  (2017-10-29 12:00:00+00:00,10738.5000000000)
  (2017-10-29 22:00:00+00:00,9458.7500000000)
  (2017-10-30 08:00:00+00:00,12985.0000000000)
  (2017-10-30 18:00:00+00:00,13194.0000000000)
  (2017-10-31 04:00:00+00:00,9522.0000000000)
  (2017-10-31 14:00:00+00:00,10851.2500000000)
  (2017-11-01 00:00:00+00:00,9557.2500000000)
  (2017-11-01 10:00:00+00:00,14757.0000000000)
  (2017-11-01 20:00:00+00:00,12979.5000000000)
  (2017-11-02 06:00:00+00:00,13863.7500000000)
  (2017-11-02 16:00:00+00:00,14725.5000000000)
  (2017-11-03 02:00:00+00:00,10021.5000000000)
  (2017-11-03 12:00:00+00:00,13570.7500000000)
  (2017-11-03 22:00:00+00:00,10890.2500000000)
  (2017-11-04 08:00:00+00:00,12244.2500000000)
  (2017-11-04 18:00:00+00:00,12634.2500000000)
  (2017-11-05 04:00:00+00:00,8844.5000000000)
  (2017-11-05 14:00:00+00:00,11202.2500000000)
  (2017-11-06 00:00:00+00:00,9440.5000000000)
  (2017-11-06 10:00:00+00:00,14160.0000000000)
  (2017-11-06 20:00:00+00:00,12500.2500000000)
  (2017-11-07 06:00:00+00:00,13812.0000000000)
  (2017-11-07 16:00:00+00:00,14899.7500000000)
  (2017-11-08 02:00:00+00:00,10373.2500000000)
  (2017-11-08 12:00:00+00:00,14033.7500000000)
  (2017-11-08 22:00:00+00:00,11187.0000000000)
  (2017-11-09 08:00:00+00:00,14515.0000000000)
  (2017-11-09 18:00:00+00:00,14687.2500000000)
  (2017-11-10 04:00:00+00:00,11685.2500000000)
  (2017-11-10 14:00:00+00:00,14114.7500000000)
  (2017-11-11 00:00:00+00:00,10170.0000000000)
  (2017-11-11 10:00:00+00:00,13486.5000000000)
  (2017-11-11 20:00:00+00:00,11437.0000000000)
  (2017-11-12 06:00:00+00:00,9916.0000000000)
  (2017-11-12 16:00:00+00:00,12273.0000000000)
  (2017-11-13 02:00:00+00:00,9817.5000000000)
  (2017-11-13 12:00:00+00:00,14031.7500000000)
  (2017-11-13 22:00:00+00:00,11223.2500000000)
  (2017-11-14 08:00:00+00:00,14434.5000000000)
  (2017-11-14 18:00:00+00:00,14832.5000000000)
  (2017-11-15 04:00:00+00:00,11466.5000000000)
  (2017-11-15 14:00:00+00:00,13800.5000000000)
  (2017-11-16 00:00:00+00:00,10214.0000000000)
  (2017-11-16 10:00:00+00:00,14435.7500000000)
  (2017-11-16 20:00:00+00:00,12902.5000000000)
  (2017-11-17 06:00:00+00:00,13900.5000000000)
  (2017-11-17 16:00:00+00:00,14590.5000000000)
  (2017-11-18 02:00:00+00:00,10245.7500000000)
  (2017-11-18 12:00:00+00:00,12751.7500000000)
  (2017-11-18 22:00:00+00:00,10444.0000000000)
  (2017-11-19 08:00:00+00:00,11280.5000000000)
  (2017-11-19 18:00:00+00:00,12393.0000000000)
  (2017-11-20 04:00:00+00:00,11265.7500000000)
  (2017-11-20 14:00:00+00:00,14122.5000000000)
  (2017-11-21 00:00:00+00:00,10174.2500000000)
  (2017-11-21 10:00:00+00:00,14422.2500000000)
  (2017-11-21 20:00:00+00:00,12577.7500000000)
  (2017-11-22 06:00:00+00:00,13051.5000000000)
  (2017-11-22 16:00:00+00:00,13976.0000000000)
  (2017-11-23 02:00:00+00:00,10380.7500000000)
  (2017-11-23 12:00:00+00:00,13991.0000000000)
  (2017-11-23 22:00:00+00:00,11189.5000000000)
  (2017-11-24 08:00:00+00:00,14187.2500000000)
  (2017-11-24 18:00:00+00:00,13669.0000000000)
  (2017-11-25 04:00:00+00:00,9672.5000000000)
  (2017-11-25 14:00:00+00:00,12067.5000000000)
  (2017-11-26 00:00:00+00:00,9335.7500000000)
  (2017-11-26 10:00:00+00:00,12354.7500000000)
  (2017-11-26 20:00:00+00:00,11156.2500000000)
  (2017-11-27 06:00:00+00:00,14284.7500000000)
  (2017-11-27 16:00:00+00:00,15383.7500000000)
  (2017-11-28 02:00:00+00:00,10745.0000000000)
  (2017-11-28 12:00:00+00:00,14555.2500000000)
  (2017-11-28 22:00:00+00:00,11756.0000000000)
  (2017-11-29 08:00:00+00:00,14847.2500000000)
  (2017-11-29 18:00:00+00:00,14895.2500000000)
  (2017-11-30 04:00:00+00:00,12027.0000000000)
  (2017-11-30 14:00:00+00:00,14380.0000000000)
  (2017-12-01 00:00:00+00:00,11026.0000000000)
  (2017-12-01 10:00:00+00:00,14973.2500000000)
  (2017-12-01 20:00:00+00:00,12899.5000000000)
  (2017-12-02 06:00:00+00:00,11371.0000000000)
  (2017-12-02 16:00:00+00:00,13481.5000000000)
  (2017-12-03 02:00:00+00:00,10052.5000000000)
  (2017-12-03 12:00:00+00:00,12460.7500000000)
  (2017-12-03 22:00:00+00:00,11285.0000000000)
  (2017-12-04 08:00:00+00:00,15187.0000000000)
  (2017-12-04 18:00:00+00:00,15308.5000000000)
  (2017-12-05 04:00:00+00:00,12371.0000000000)
  (2017-12-05 14:00:00+00:00,15252.0000000000)
  (2017-12-06 00:00:00+00:00,10926.7500000000)
  (2017-12-06 10:00:00+00:00,15387.7500000000)
  (2017-12-06 20:00:00+00:00,13661.7500000000)
  (2017-12-07 06:00:00+00:00,14497.5000000000)
  (2017-12-07 16:00:00+00:00,15422.0000000000)
  (2017-12-08 02:00:00+00:00,10848.5000000000)
  (2017-12-08 12:00:00+00:00,14821.0000000000)
  (2017-12-08 22:00:00+00:00,11624.5000000000)
  (2017-12-09 08:00:00+00:00,12586.7500000000)
  (2017-12-09 18:00:00+00:00,12886.7500000000)
  (2017-12-10 04:00:00+00:00,9540.0000000000)
  (2017-12-10 14:00:00+00:00,11653.2500000000)
  (2017-12-11 00:00:00+00:00,10158.5000000000)
  (2017-12-11 10:00:00+00:00,14990.2500000000)
  (2017-12-11 20:00:00+00:00,13403.7500000000)
  (2017-12-12 06:00:00+00:00,15067.5000000000)
  (2017-12-12 16:00:00+00:00,15913.0000000000)
  (2017-12-13 02:00:00+00:00,11077.0000000000)
  (2017-12-13 12:00:00+00:00,14896.0000000000)
  (2017-12-13 22:00:00+00:00,12050.5000000000)
  (2017-12-14 08:00:00+00:00,15180.2500000000)
  (2017-12-14 18:00:00+00:00,15223.0000000000)
  (2017-12-15 04:00:00+00:00,11768.0000000000)
  (2017-12-15 14:00:00+00:00,14364.7500000000)
  (2017-12-16 00:00:00+00:00,10734.5000000000)
  (2017-12-16 10:00:00+00:00,13672.5000000000)
  (2017-12-16 20:00:00+00:00,11975.0000000000)
  (2017-12-17 06:00:00+00:00,10218.2500000000)
  (2017-12-17 16:00:00+00:00,12958.0000000000)
  (2017-12-18 02:00:00+00:00,10515.2500000000)
  (2017-12-18 12:00:00+00:00,14249.7500000000)
  (2017-12-18 22:00:00+00:00,11937.5000000000)
  (2017-12-19 08:00:00+00:00,14860.5000000000)
  (2017-12-19 18:00:00+00:00,14760.0000000000)
  (2017-12-20 04:00:00+00:00,11617.2500000000)
  (2017-12-20 14:00:00+00:00,14179.7500000000)
  (2017-12-21 00:00:00+00:00,10573.5000000000)
  (2017-12-21 10:00:00+00:00,14609.7500000000)
  (2017-12-21 20:00:00+00:00,12801.2500000000)
  (2017-12-22 06:00:00+00:00,13108.2500000000)
  (2017-12-22 16:00:00+00:00,14062.7500000000)
  (2017-12-23 02:00:00+00:00,9143.2500000000)
  (2017-12-23 12:00:00+00:00,12223.2500000000)
  (2017-12-23 22:00:00+00:00,9951.5000000000)
  (2017-12-24 08:00:00+00:00,10603.5000000000)
  (2017-12-24 18:00:00+00:00,10277.5000000000)
  (2017-12-25 04:00:00+00:00,8020.7500000000)
  (2017-12-25 14:00:00+00:00,9652.7500000000)
  (2017-12-26 00:00:00+00:00,8164.5000000000)
  (2017-12-26 10:00:00+00:00,11311.2500000000)
  (2017-12-26 20:00:00+00:00,10104.0000000000)
  (2017-12-27 06:00:00+00:00,10462.5000000000)
  (2017-12-27 16:00:00+00:00,13014.0000000000)
  (2017-12-28 02:00:00+00:00,8703.0000000000)
  (2017-12-28 12:00:00+00:00,12371.7500000000)
  (2017-12-28 22:00:00+00:00,10057.2500000000)
  (2017-12-29 08:00:00+00:00,12061.5000000000)
  (2017-12-29 18:00:00+00:00,12826.7500000000)
  (2017-12-30 04:00:00+00:00,8853.7500000000)
  (2017-12-30 14:00:00+00:00,11221.5000000000)
  (2017-12-31 00:00:00+00:00,8474.0000000000)
  (2017-12-31 10:00:00+00:00,10802.0000000000)
  (2017-12-31 20:00:00+00:00,9478.2500000000)
  (2018-01-01 06:00:00+00:00,9240.7500000000)
  (2018-01-01 16:00:00+00:00,11761.5000000000)
  (2018-01-02 02:00:00+00:00,8454.5000000000)
  (2018-01-02 12:00:00+00:00,14115.0000000000)
  (2018-01-02 22:00:00+00:00,11251.7500000000)
  (2018-01-03 08:00:00+00:00,15470.2500000000)
  (2018-01-03 18:00:00+00:00,15720.0000000000)
  (2018-01-04 04:00:00+00:00,12243.0000000000)
  (2018-01-04 14:00:00+00:00,14463.7500000000)
  (2018-01-05 00:00:00+00:00,10709.2500000000)
  (2018-01-05 10:00:00+00:00,15782.7500000000)
  (2018-01-05 20:00:00+00:00,12680.0000000000)
  (2018-01-06 06:00:00+00:00,11109.2500000000)
  (2018-01-06 16:00:00+00:00,13169.0000000000)
  (2018-01-07 02:00:00+00:00,9541.5000000000)
  (2018-01-07 12:00:00+00:00,11834.2500000000)
  (2018-01-07 22:00:00+00:00,11168.2500000000)
  (2018-01-08 08:00:00+00:00,15061.7500000000)
  (2018-01-08 18:00:00+00:00,15899.0000000000)
  (2018-01-09 04:00:00+00:00,13350.5000000000)
  (2018-01-09 14:00:00+00:00,15618.5000000000)
  (2018-01-10 00:00:00+00:00,11308.0000000000)
  (2018-01-10 10:00:00+00:00,14696.0000000000)
  (2018-01-10 20:00:00+00:00,12960.2500000000)
  (2018-01-11 06:00:00+00:00,14475.0000000000)
  (2018-01-11 16:00:00+00:00,15280.0000000000)
  (2018-01-12 02:00:00+00:00,11195.2500000000)
  (2018-01-12 12:00:00+00:00,14992.7500000000)
  (2018-01-12 22:00:00+00:00,12086.0000000000)
  (2018-01-13 08:00:00+00:00,13601.7500000000)
  (2018-01-13 18:00:00+00:00,13984.2500000000)
  (2018-01-14 04:00:00+00:00,10137.5000000000)
  (2018-01-14 14:00:00+00:00,11948.0000000000)
  (2018-01-15 00:00:00+00:00,10825.0000000000)
  (2018-01-15 10:00:00+00:00,16457.0000000000)
  (2018-01-15 20:00:00+00:00,14571.5000000000)
  (2018-01-16 06:00:00+00:00,15516.7500000000)
  (2018-01-16 16:00:00+00:00,16096.2500000000)
  (2018-01-17 02:00:00+00:00,11186.2500000000)
  (2018-01-17 12:00:00+00:00,15456.2500000000)
  (2018-01-17 22:00:00+00:00,12437.7500000000)
  (2018-01-18 08:00:00+00:00,15919.0000000000)
  (2018-01-18 18:00:00+00:00,16147.0000000000)
  (2018-01-19 04:00:00+00:00,11762.0000000000)
  (2018-01-19 14:00:00+00:00,14342.0000000000)
  (2018-01-20 00:00:00+00:00,10921.7500000000)
  (2018-01-20 10:00:00+00:00,13613.7500000000)
  (2018-01-20 20:00:00+00:00,11661.0000000000)
  (2018-01-21 06:00:00+00:00,10106.2500000000)
  (2018-01-21 16:00:00+00:00,12763.5000000000)
  (2018-01-22 02:00:00+00:00,10650.2500000000)
  (2018-01-22 12:00:00+00:00,15470.7500000000)
  (2018-01-22 22:00:00+00:00,12070.7500000000)
  (2018-01-23 08:00:00+00:00,14934.7500000000)
  (2018-01-23 18:00:00+00:00,15575.0000000000)
  (2018-01-24 04:00:00+00:00,12676.0000000000)
  (2018-01-24 14:00:00+00:00,15279.0000000000)
  (2018-01-25 00:00:00+00:00,11441.7500000000)
  (2018-01-25 10:00:00+00:00,15092.0000000000)
  (2018-01-25 20:00:00+00:00,13025.0000000000)
  (2018-01-26 06:00:00+00:00,14207.0000000000)
  (2018-01-26 16:00:00+00:00,14634.7500000000)
  (2018-01-27 02:00:00+00:00,10259.2500000000)
  (2018-01-27 12:00:00+00:00,12698.5000000000)
  (2018-01-27 22:00:00+00:00,11551.7500000000)
  (2018-01-28 08:00:00+00:00,12238.7500000000)
  (2018-01-28 18:00:00+00:00,13222.7500000000)
  (2018-01-29 04:00:00+00:00,12314.2500000000)
  (2018-01-29 14:00:00+00:00,15949.0000000000)
  (2018-01-30 00:00:00+00:00,12074.5000000000)
  (2018-01-30 10:00:00+00:00,15822.7500000000)
  (2018-01-30 20:00:00+00:00,13737.5000000000)
  (2018-01-31 06:00:00+00:00,15336.0000000000)
  (2018-01-31 16:00:00+00:00,15937.2500000000)
  (2018-02-01 02:00:00+00:00,11035.2500000000)
  (2018-02-01 12:00:00+00:00,15204.2500000000)
  (2018-02-01 22:00:00+00:00,13210.7500000000)
  (2018-02-02 08:00:00+00:00,15435.0000000000)
  (2018-02-02 18:00:00+00:00,14599.0000000000)
  (2018-02-03 04:00:00+00:00,10706.2500000000)
  (2018-02-03 14:00:00+00:00,12337.2500000000)
  (2018-02-04 00:00:00+00:00,10264.0000000000)
  (2018-02-04 10:00:00+00:00,13288.2500000000)
  (2018-02-04 20:00:00+00:00,12367.0000000000)
  (2018-02-05 06:00:00+00:00,14379.2500000000)
  (2018-02-05 16:00:00+00:00,14784.0000000000)
  (2018-02-06 02:00:00+00:00,11597.0000000000)
  (2018-02-06 12:00:00+00:00,14494.0000000000)
  (2018-02-06 22:00:00+00:00,12329.2500000000)
  (2018-02-07 08:00:00+00:00,14948.0000000000)
  (2018-02-07 18:00:00+00:00,15070.0000000000)
  (2018-02-08 04:00:00+00:00,12417.7500000000)
  (2018-02-08 14:00:00+00:00,14366.2500000000)
  (2018-02-09 00:00:00+00:00,11849.0000000000)
  (2018-02-09 10:00:00+00:00,15351.5000000000)
  (2018-02-09 20:00:00+00:00,13342.5000000000)
  (2018-02-10 06:00:00+00:00,11532.7500000000)
  (2018-02-10 16:00:00+00:00,13098.0000000000)
  (2018-02-11 02:00:00+00:00,9686.7500000000)
  (2018-02-11 12:00:00+00:00,12808.0000000000)
  (2018-02-11 22:00:00+00:00,11795.2500000000)
  (2018-02-12 08:00:00+00:00,15825.0000000000)
  (2018-02-12 18:00:00+00:00,15826.7500000000)
  (2018-02-13 04:00:00+00:00,12816.0000000000)
  (2018-02-13 14:00:00+00:00,13583.7500000000)
  (2018-02-14 00:00:00+00:00,11207.0000000000)
  (2018-02-14 10:00:00+00:00,14542.2500000000)
  (2018-02-14 20:00:00+00:00,13603.2500000000)
  (2018-02-15 06:00:00+00:00,14990.0000000000)
  (2018-02-15 16:00:00+00:00,15762.5000000000)
  (2018-02-16 02:00:00+00:00,11391.7500000000)
  (2018-02-16 12:00:00+00:00,15013.0000000000)
  (2018-02-16 22:00:00+00:00,12036.0000000000)
  (2018-02-17 08:00:00+00:00,12870.7500000000)
  (2018-02-17 18:00:00+00:00,12971.5000000000)
  (2018-02-18 04:00:00+00:00,9670.5000000000)
  (2018-02-18 14:00:00+00:00,11396.5000000000)
  (2018-02-19 00:00:00+00:00,10952.7500000000)
  (2018-02-19 10:00:00+00:00,14724.2500000000)
  (2018-02-19 20:00:00+00:00,13218.5000000000)
  (2018-02-20 06:00:00+00:00,14534.7500000000)
  (2018-02-20 16:00:00+00:00,14667.0000000000)
  (2018-02-21 02:00:00+00:00,11741.0000000000)
  (2018-02-21 12:00:00+00:00,14432.5000000000)
  (2018-02-21 22:00:00+00:00,12260.2500000000)
  (2018-02-22 08:00:00+00:00,14741.2500000000)
  (2018-02-22 18:00:00+00:00,15156.0000000000)
  (2018-02-23 04:00:00+00:00,12656.5000000000)
  (2018-02-23 14:00:00+00:00,14052.2500000000)
  (2018-02-24 00:00:00+00:00,12000.2500000000)
  (2018-02-24 10:00:00+00:00,14623.0000000000)
  (2018-02-24 20:00:00+00:00,12676.0000000000)
  (2018-02-25 06:00:00+00:00,11027.0000000000)
  (2018-02-25 16:00:00+00:00,13052.5000000000)
  (2018-02-26 02:00:00+00:00,12003.5000000000)
  (2018-02-26 12:00:00+00:00,15523.5000000000)
  (2018-02-26 22:00:00+00:00,13700.2500000000)
  (2018-02-27 08:00:00+00:00,15825.7500000000)
  (2018-02-27 18:00:00+00:00,16903.5000000000)
  (2018-02-28 04:00:00+00:00,13987.7500000000)
  (2018-02-28 14:00:00+00:00,15349.7500000000)
  (2018-03-01 00:00:00+00:00,13443.5000000000)
  (2018-03-01 10:00:00+00:00,16065.0000000000)
  (2018-03-01 20:00:00+00:00,15107.5000000000)
  (2018-03-02 06:00:00+00:00,15790.5000000000)
  (2018-03-02 16:00:00+00:00,15653.0000000000)
  (2018-03-03 02:00:00+00:00,12223.5000000000)
  (2018-03-03 12:00:00+00:00,13707.0000000000)
  (2018-03-03 22:00:00+00:00,12653.5000000000)
  (2018-03-04 08:00:00+00:00,12996.7500000000)
  (2018-03-04 18:00:00+00:00,14296.7500000000)
  (2018-03-05 04:00:00+00:00,12981.7500000000)
  (2018-03-05 14:00:00+00:00,13922.0000000000)
  (2018-03-06 00:00:00+00:00,11615.0000000000)
  (2018-03-06 10:00:00+00:00,14971.0000000000)
  (2018-03-06 20:00:00+00:00,13858.7500000000)
  (2018-03-07 06:00:00+00:00,14701.0000000000)
  (2018-03-07 16:00:00+00:00,14696.0000000000)
  (2018-03-08 02:00:00+00:00,11320.5000000000)
  (2018-03-08 12:00:00+00:00,14475.5000000000)
  (2018-03-08 22:00:00+00:00,12923.2500000000)
  (2018-03-09 08:00:00+00:00,16019.0000000000)
  (2018-03-09 18:00:00+00:00,14574.5000000000)
  (2018-03-10 04:00:00+00:00,10499.0000000000)
  (2018-03-10 14:00:00+00:00,12138.2500000000)
  (2018-03-11 00:00:00+00:00,10042.7500000000)
  (2018-03-11 10:00:00+00:00,11998.2500000000)
  (2018-03-11 20:00:00+00:00,11959.2500000000)
  (2018-03-12 06:00:00+00:00,13525.0000000000)
  (2018-03-12 16:00:00+00:00,14105.2500000000)
  (2018-03-13 02:00:00+00:00,11025.5000000000)
  (2018-03-13 12:00:00+00:00,14667.0000000000)
  (2018-03-13 22:00:00+00:00,12037.2500000000)
  (2018-03-14 08:00:00+00:00,14435.0000000000)
  (2018-03-14 18:00:00+00:00,14834.2500000000)
  (2018-03-15 04:00:00+00:00,12288.5000000000)
  (2018-03-15 14:00:00+00:00,14916.7500000000)
  (2018-03-16 00:00:00+00:00,12017.7500000000)
  (2018-03-16 10:00:00+00:00,16703.7500000000)
  (2018-03-16 20:00:00+00:00,14267.5000000000)
  (2018-03-17 06:00:00+00:00,12751.2500000000)
  (2018-03-17 16:00:00+00:00,13834.7500000000)
  (2018-03-18 02:00:00+00:00,11705.2500000000)
  (2018-03-18 12:00:00+00:00,13366.0000000000)
  (2018-03-18 22:00:00+00:00,12229.0000000000)
  (2018-03-19 08:00:00+00:00,14720.7500000000)
  (2018-03-19 18:00:00+00:00,15152.7500000000)
  (2018-03-20 04:00:00+00:00,12936.0000000000)
  (2018-03-20 14:00:00+00:00,15018.2500000000)
  (2018-03-21 00:00:00+00:00,11035.0000000000)
  (2018-03-21 10:00:00+00:00,14472.2500000000)
  (2018-03-21 20:00:00+00:00,13838.2500000000)
  (2018-03-22 06:00:00+00:00,15171.0000000000)
  (2018-03-22 16:00:00+00:00,14516.7500000000)
  (2018-03-23 02:00:00+00:00,11291.7500000000)
  (2018-03-23 12:00:00+00:00,14420.5000000000)
  (2018-03-23 22:00:00+00:00,11792.7500000000)
  (2018-03-24 08:00:00+00:00,12724.5000000000)
  (2018-03-24 18:00:00+00:00,13049.7500000000)
  (2018-03-25 04:00:00+00:00,9663.7500000000)
  (2018-03-25 14:00:00+00:00,10670.2500000000)
  (2018-03-26 00:00:00+00:00,10034.7500000000)
  (2018-03-26 10:00:00+00:00,14682.0000000000)
  (2018-03-26 20:00:00+00:00,12448.0000000000)
  (2018-03-27 06:00:00+00:00,14462.2500000000)
  (2018-03-27 16:00:00+00:00,13550.0000000000)
  (2018-03-28 02:00:00+00:00,12063.0000000000)
  (2018-03-28 12:00:00+00:00,15618.2500000000)
  (2018-03-28 22:00:00+00:00,12265.2500000000)
  (2018-03-29 08:00:00+00:00,16272.0000000000)
  (2018-03-29 18:00:00+00:00,13957.7500000000)
  (2018-03-30 04:00:00+00:00,9800.5000000000)
  (2018-03-30 14:00:00+00:00,10045.7500000000)
  (2018-03-31 00:00:00+00:00,9564.2500000000)
  (2018-03-31 10:00:00+00:00,13213.5000000000)
  (2018-03-31 20:00:00+00:00,11571.5000000000)
  (2018-04-01 06:00:00+00:00,10998.0000000000)
  (2018-04-01 16:00:00+00:00,11176.2500000000)
  (2018-04-02 02:00:00+00:00,9767.2500000000)
  (2018-04-02 12:00:00+00:00,10145.2500000000)
  (2018-04-02 22:00:00+00:00,9607.5000000000)
  (2018-04-03 08:00:00+00:00,14466.7500000000)
  (2018-04-03 18:00:00+00:00,13461.0000000000)
  (2018-04-04 04:00:00+00:00,13073.7500000000)
  (2018-04-04 14:00:00+00:00,12930.0000000000)
  (2018-04-05 00:00:00+00:00,10086.5000000000)
  (2018-04-05 10:00:00+00:00,14123.2500000000)
  (2018-04-05 20:00:00+00:00,13276.2500000000)
  (2018-04-06 06:00:00+00:00,14407.5000000000)
  (2018-04-06 16:00:00+00:00,12787.5000000000)
  (2018-04-07 02:00:00+00:00,10380.5000000000)
  (2018-04-07 12:00:00+00:00,11421.2500000000)
  (2018-04-07 22:00:00+00:00,9622.7500000000)
  (2018-04-08 08:00:00+00:00,11295.5000000000)
  (2018-04-08 18:00:00+00:00,11437.5000000000)
  (2018-04-09 04:00:00+00:00,12188.0000000000)
  (2018-04-09 14:00:00+00:00,12656.0000000000)
  (2018-04-10 00:00:00+00:00,9770.7500000000)
  (2018-04-10 10:00:00+00:00,14672.2500000000)
  (2018-04-10 20:00:00+00:00,12605.5000000000)
  (2018-04-11 06:00:00+00:00,14644.2500000000)
  (2018-04-11 16:00:00+00:00,14189.0000000000)
  (2018-04-12 02:00:00+00:00,10948.0000000000)
  (2018-04-12 12:00:00+00:00,14315.0000000000)
  (2018-04-12 22:00:00+00:00,11559.2500000000)
  (2018-04-13 08:00:00+00:00,15049.7500000000)
  (2018-04-13 18:00:00+00:00,13534.7500000000)
  (2018-04-14 04:00:00+00:00,10230.5000000000)
  (2018-04-14 14:00:00+00:00,10955.7500000000)
  (2018-04-15 00:00:00+00:00,8509.2500000000)
  (2018-04-15 10:00:00+00:00,10839.5000000000)
  (2018-04-15 20:00:00+00:00,10242.0000000000)
  (2018-04-16 06:00:00+00:00,13513.2500000000)
  (2018-04-16 16:00:00+00:00,12921.0000000000)
  (2018-04-17 02:00:00+00:00,10222.0000000000)
  (2018-04-17 12:00:00+00:00,13084.5000000000)
  (2018-04-17 22:00:00+00:00,9880.5000000000)
  (2018-04-18 08:00:00+00:00,13615.2500000000)
  (2018-04-18 18:00:00+00:00,12973.5000000000)
  (2018-04-19 04:00:00+00:00,11533.7500000000)
  (2018-04-19 14:00:00+00:00,12515.2500000000)
  (2018-04-20 00:00:00+00:00,9662.2500000000)
  (2018-04-20 10:00:00+00:00,13706.2500000000)
  (2018-04-20 20:00:00+00:00,11838.0000000000)
  (2018-04-21 06:00:00+00:00,11633.5000000000)
  (2018-04-21 16:00:00+00:00,10963.5000000000)
  (2018-04-22 02:00:00+00:00,7866.5000000000)
  (2018-04-22 12:00:00+00:00,9477.2500000000)
  (2018-04-22 22:00:00+00:00,9501.7500000000)
  (2018-04-23 08:00:00+00:00,13928.5000000000)
  (2018-04-23 18:00:00+00:00,13007.7500000000)
  (2018-04-24 04:00:00+00:00,12317.7500000000)
  (2018-04-24 14:00:00+00:00,13809.0000000000)
  (2018-04-25 00:00:00+00:00,9777.5000000000)
  (2018-04-25 10:00:00+00:00,14874.2500000000)
  (2018-04-25 20:00:00+00:00,12081.2500000000)
  (2018-04-26 06:00:00+00:00,14041.5000000000)
  (2018-04-26 16:00:00+00:00,13386.2500000000)
  (2018-04-27 02:00:00+00:00,10899.7500000000)
  (2018-04-27 12:00:00+00:00,12767.0000000000)
  (2018-04-27 22:00:00+00:00,9806.5000000000)
  (2018-04-28 08:00:00+00:00,11542.7500000000)
  (2018-04-28 18:00:00+00:00,10883.7500000000)
  (2018-04-29 04:00:00+00:00,8172.0000000000)
  (2018-04-29 14:00:00+00:00,9189.2500000000)
  (2018-04-30 00:00:00+00:00,8993.2500000000)
  (2018-04-30 10:00:00+00:00,12862.7500000000)
  (2018-04-30 20:00:00+00:00,9954.0000000000)
  (2018-05-01 06:00:00+00:00,10796.2500000000)
  (2018-05-01 16:00:00+00:00,11217.7500000000)
  (2018-05-02 02:00:00+00:00,9263.7500000000)
  (2018-05-02 12:00:00+00:00,12887.2500000000)
  (2018-05-02 22:00:00+00:00,9907.7500000000)
  (2018-05-03 08:00:00+00:00,13477.2500000000)
  (2018-05-03 18:00:00+00:00,12572.0000000000)
  (2018-05-04 04:00:00+00:00,11834.7500000000)
  (2018-05-04 14:00:00+00:00,12200.5000000000)
  (2018-05-05 00:00:00+00:00,9102.0000000000)
  (2018-05-05 10:00:00+00:00,11871.7500000000)
  (2018-05-05 20:00:00+00:00,10438.2500000000)
  (2018-05-06 06:00:00+00:00,9691.2500000000)
  (2018-05-06 16:00:00+00:00,10570.2500000000)
  (2018-05-07 02:00:00+00:00,9264.0000000000)
  (2018-05-07 12:00:00+00:00,13101.0000000000)
  (2018-05-07 22:00:00+00:00,10175.2500000000)
  (2018-05-08 08:00:00+00:00,13626.2500000000)
  (2018-05-08 18:00:00+00:00,12889.2500000000)
  (2018-05-09 04:00:00+00:00,12007.2500000000)
  (2018-05-09 14:00:00+00:00,13228.5000000000)
  (2018-05-10 00:00:00+00:00,8909.0000000000)
  (2018-05-10 10:00:00+00:00,10477.5000000000)
  (2018-05-10 20:00:00+00:00,9864.7500000000)
  (2018-05-11 06:00:00+00:00,11866.0000000000)
  (2018-05-11 16:00:00+00:00,11519.5000000000)
  (2018-05-12 02:00:00+00:00,8433.7500000000)
  (2018-05-12 12:00:00+00:00,10571.2500000000)
  (2018-05-12 22:00:00+00:00,8614.5000000000)
  (2018-05-13 08:00:00+00:00,10538.5000000000)
  (2018-05-13 18:00:00+00:00,10722.2500000000)
  (2018-05-14 04:00:00+00:00,11888.7500000000)
  (2018-05-14 14:00:00+00:00,13452.0000000000)
  (2018-05-15 00:00:00+00:00,9549.5000000000)
  (2018-05-15 10:00:00+00:00,13895.2500000000)
  (2018-05-15 20:00:00+00:00,11344.0000000000)
  (2018-05-16 06:00:00+00:00,13161.2500000000)
  (2018-05-16 16:00:00+00:00,13243.5000000000)
  (2018-05-17 02:00:00+00:00,10031.2500000000)
  (2018-05-17 12:00:00+00:00,13593.0000000000)
  (2018-05-17 22:00:00+00:00,10447.7500000000)
  (2018-05-18 08:00:00+00:00,13842.2500000000)
  (2018-05-18 18:00:00+00:00,12364.7500000000)
  (2018-05-19 04:00:00+00:00,9466.0000000000)
  (2018-05-19 14:00:00+00:00,10606.5000000000)
  (2018-05-20 00:00:00+00:00,7809.2500000000)
  (2018-05-20 10:00:00+00:00,9886.7500000000)
  (2018-05-20 20:00:00+00:00,9481.7500000000)
  (2018-05-21 06:00:00+00:00,9062.0000000000)
  (2018-05-21 16:00:00+00:00,10591.2500000000)
  (2018-05-22 02:00:00+00:00,8980.5000000000)
  (2018-05-22 12:00:00+00:00,13257.7500000000)
  (2018-05-22 22:00:00+00:00,10346.0000000000)
  (2018-05-23 08:00:00+00:00,13889.2500000000)
  (2018-05-23 18:00:00+00:00,12998.7500000000)
  (2018-05-24 04:00:00+00:00,12389.5000000000)
  (2018-05-24 14:00:00+00:00,14007.0000000000)
  (2018-05-25 00:00:00+00:00,10547.2500000000)
  (2018-05-25 10:00:00+00:00,14374.2500000000)
  (2018-05-25 20:00:00+00:00,11619.5000000000)
  (2018-05-26 06:00:00+00:00,11328.5000000000)
  (2018-05-26 16:00:00+00:00,11386.0000000000)
  (2018-05-27 02:00:00+00:00,8122.2500000000)
  (2018-05-27 12:00:00+00:00,10372.7500000000)
  (2018-05-27 22:00:00+00:00,10215.0000000000)
  (2018-05-28 08:00:00+00:00,14168.7500000000)
  (2018-05-28 18:00:00+00:00,13527.7500000000)
  (2018-05-29 04:00:00+00:00,12240.0000000000)
  (2018-05-29 14:00:00+00:00,14208.0000000000)
  (2018-05-30 00:00:00+00:00,10337.7500000000)
  (2018-05-30 10:00:00+00:00,14357.2500000000)
  (2018-05-30 20:00:00+00:00,12396.5000000000)
  (2018-05-31 06:00:00+00:00,14109.5000000000)
  (2018-05-31 16:00:00+00:00,14275.2500000000)
  (2018-06-01 02:00:00+00:00,10470.5000000000)
  (2018-06-01 12:00:00+00:00,13552.5000000000)
  (2018-06-01 22:00:00+00:00,10141.2500000000)
  (2018-06-02 08:00:00+00:00,12167.7500000000)
  (2018-06-02 18:00:00+00:00,10817.0000000000)
  (2018-06-03 04:00:00+00:00,8405.5000000000)
  (2018-06-03 14:00:00+00:00,10294.5000000000)
  (2018-06-04 00:00:00+00:00,9077.7500000000)
  (2018-06-04 10:00:00+00:00,13955.2500000000)
  (2018-06-04 20:00:00+00:00,12231.0000000000)
  (2018-06-05 06:00:00+00:00,13775.2500000000)
  (2018-06-05 16:00:00+00:00,13133.2500000000)
  (2018-06-06 02:00:00+00:00,10419.7500000000)
  (2018-06-06 12:00:00+00:00,13896.5000000000)
  (2018-06-06 22:00:00+00:00,10758.7500000000)
  (2018-06-07 08:00:00+00:00,13968.5000000000)
  (2018-06-07 18:00:00+00:00,13064.7500000000)
  (2018-06-08 04:00:00+00:00,11872.2500000000)
  (2018-06-08 14:00:00+00:00,13108.5000000000)
  (2018-06-09 00:00:00+00:00,9693.0000000000)
  (2018-06-09 10:00:00+00:00,12446.0000000000)
  (2018-06-09 20:00:00+00:00,10474.7500000000)
  (2018-06-10 06:00:00+00:00,10057.5000000000)
  (2018-06-10 16:00:00+00:00,11097.0000000000)
  (2018-06-11 02:00:00+00:00,9776.0000000000)
  (2018-06-11 12:00:00+00:00,13921.0000000000)
  (2018-06-11 22:00:00+00:00,10933.0000000000)
  (2018-06-12 08:00:00+00:00,13779.2500000000)
  (2018-06-12 18:00:00+00:00,12666.2500000000)
  (2018-06-13 04:00:00+00:00,12036.0000000000)
  (2018-06-13 14:00:00+00:00,13325.5000000000)
  (2018-06-14 00:00:00+00:00,9828.5000000000)
  (2018-06-14 10:00:00+00:00,13707.5000000000)
  (2018-06-14 20:00:00+00:00,11843.7500000000)
  (2018-06-15 06:00:00+00:00,13392.0000000000)
  (2018-06-15 16:00:00+00:00,12655.5000000000)
  (2018-06-16 02:00:00+00:00,9080.2500000000)
  (2018-06-16 12:00:00+00:00,11527.5000000000)
  (2018-06-16 22:00:00+00:00,9405.5000000000)
  (2018-06-17 08:00:00+00:00,10471.7500000000)
  (2018-06-17 18:00:00+00:00,10607.7500000000)
  (2018-06-18 04:00:00+00:00,11450.7500000000)
  (2018-06-18 14:00:00+00:00,13115.7500000000)
  (2018-06-19 00:00:00+00:00,10024.5000000000)
  (2018-06-19 10:00:00+00:00,14721.2500000000)
  (2018-06-19 20:00:00+00:00,12102.2500000000)
  (2018-06-20 06:00:00+00:00,13561.2500000000)
  (2018-06-20 16:00:00+00:00,13812.2500000000)
  (2018-06-21 02:00:00+00:00,10331.2500000000)
  (2018-06-21 12:00:00+00:00,14261.7500000000)
  (2018-06-21 22:00:00+00:00,10944.0000000000)
  (2018-06-22 08:00:00+00:00,14225.5000000000)
  (2018-06-22 18:00:00+00:00,12784.0000000000)
  (2018-06-23 04:00:00+00:00,9899.7500000000)
  (2018-06-23 14:00:00+00:00,11687.5000000000)
  (2018-06-24 00:00:00+00:00,8605.7500000000)
  (2018-06-24 10:00:00+00:00,10745.2500000000)
  (2018-06-24 20:00:00+00:00,10269.2500000000)
  (2018-06-25 06:00:00+00:00,13136.2500000000)
  (2018-06-25 16:00:00+00:00,13341.2500000000)
  (2018-06-26 02:00:00+00:00,9780.0000000000)
  (2018-06-26 12:00:00+00:00,13037.0000000000)
  (2018-06-26 22:00:00+00:00,10287.2500000000)
  (2018-06-27 08:00:00+00:00,13424.0000000000)
  (2018-06-27 18:00:00+00:00,12460.2500000000)
  (2018-06-28 04:00:00+00:00,11831.2500000000)
  (2018-06-28 14:00:00+00:00,13710.0000000000)
  (2018-06-29 00:00:00+00:00,9882.5000000000)
  (2018-06-29 10:00:00+00:00,14374.5000000000)
  (2018-06-29 20:00:00+00:00,11569.2500000000)
  (2018-06-30 06:00:00+00:00,11102.7500000000)
  (2018-06-30 16:00:00+00:00,11195.5000000000)
  (2018-07-01 02:00:00+00:00,7824.0000000000)
  (2018-07-01 12:00:00+00:00,10056.0000000000)
  (2018-07-01 22:00:00+00:00,9325.0000000000)
  (2018-07-02 08:00:00+00:00,13541.0000000000)
  (2018-07-02 18:00:00+00:00,12439.5000000000)
  (2018-07-03 04:00:00+00:00,11742.0000000000)
  (2018-07-03 14:00:00+00:00,13386.7500000000)
  (2018-07-04 00:00:00+00:00,9526.7500000000)
  (2018-07-04 10:00:00+00:00,13949.5000000000)
  (2018-07-04 20:00:00+00:00,11810.5000000000)
  (2018-07-05 06:00:00+00:00,13327.7500000000)
  (2018-07-05 16:00:00+00:00,13828.5000000000)
  (2018-07-06 02:00:00+00:00,10071.5000000000)
  (2018-07-06 12:00:00+00:00,13889.5000000000)
  (2018-07-06 22:00:00+00:00,10661.5000000000)
  (2018-07-07 08:00:00+00:00,11973.0000000000)
  (2018-07-07 18:00:00+00:00,11034.2500000000)
  (2018-07-08 04:00:00+00:00,8226.2500000000)
  (2018-07-08 14:00:00+00:00,9872.5000000000)
  (2018-07-09 00:00:00+00:00,9125.0000000000)
  (2018-07-09 10:00:00+00:00,14128.2500000000)
  (2018-07-09 20:00:00+00:00,11051.2500000000)
  (2018-07-10 06:00:00+00:00,12902.2500000000)
  (2018-07-10 16:00:00+00:00,13013.0000000000)
  (2018-07-11 02:00:00+00:00,10030.5000000000)
  (2018-07-11 12:00:00+00:00,13154.7500000000)
  (2018-07-11 22:00:00+00:00,9794.7500000000)
  (2018-07-12 08:00:00+00:00,14273.0000000000)
  (2018-07-12 18:00:00+00:00,12232.7500000000)
  (2018-07-13 04:00:00+00:00,11314.5000000000)
  (2018-07-13 14:00:00+00:00,12903.0000000000)
  (2018-07-14 00:00:00+00:00,9456.2500000000)
  (2018-07-14 10:00:00+00:00,11979.5000000000)
  (2018-07-14 20:00:00+00:00,9942.0000000000)
  (2018-07-15 06:00:00+00:00,9565.0000000000)
  (2018-07-15 16:00:00+00:00,10224.2500000000)
  (2018-07-16 02:00:00+00:00,8888.2500000000)
  (2018-07-16 12:00:00+00:00,13103.5000000000)
  (2018-07-16 22:00:00+00:00,9834.0000000000)
  (2018-07-17 08:00:00+00:00,13526.5000000000)
  (2018-07-17 18:00:00+00:00,12508.2500000000)
  (2018-07-18 04:00:00+00:00,11983.5000000000)
  (2018-07-18 14:00:00+00:00,13571.0000000000)
  (2018-07-19 00:00:00+00:00,10106.0000000000)
  (2018-07-19 10:00:00+00:00,14424.2500000000)
  (2018-07-19 20:00:00+00:00,11472.7500000000)
  (2018-07-20 06:00:00+00:00,12891.7500000000)
  (2018-07-20 16:00:00+00:00,12677.2500000000)
  (2018-07-21 02:00:00+00:00,8965.7500000000)
  (2018-07-21 12:00:00+00:00,11202.0000000000)
  (2018-07-21 22:00:00+00:00,8972.0000000000)
  (2018-07-22 08:00:00+00:00,10209.2500000000)
  (2018-07-22 18:00:00+00:00,9892.0000000000)
  (2018-07-23 04:00:00+00:00,10550.7500000000)
  (2018-07-23 14:00:00+00:00,12544.5000000000)
  (2018-07-24 00:00:00+00:00,9311.7500000000)
  (2018-07-24 10:00:00+00:00,13744.2500000000)
  (2018-07-24 20:00:00+00:00,11416.5000000000)
  (2018-07-25 06:00:00+00:00,13029.0000000000)
  (2018-07-25 16:00:00+00:00,12900.2500000000)
  (2018-07-26 02:00:00+00:00,9635.7500000000)
  (2018-07-26 12:00:00+00:00,13422.7500000000)
  (2018-07-26 22:00:00+00:00,10216.5000000000)
  (2018-07-27 08:00:00+00:00,13740.5000000000)
  (2018-07-27 18:00:00+00:00,12523.0000000000)
  (2018-07-28 04:00:00+00:00,9390.7500000000)
  (2018-07-28 14:00:00+00:00,11005.7500000000)
  (2018-07-29 00:00:00+00:00,8443.2500000000)
  (2018-07-29 10:00:00+00:00,10316.2500000000)
  (2018-07-29 20:00:00+00:00,10083.2500000000)
  (2018-07-30 06:00:00+00:00,12634.2500000000)
  (2018-07-30 16:00:00+00:00,12624.5000000000)
  (2018-07-31 02:00:00+00:00,9600.7500000000)
  (2018-07-31 12:00:00+00:00,13276.5000000000)
  (2018-07-31 22:00:00+00:00,10403.2500000000)
  (2018-08-01 08:00:00+00:00,13373.7500000000)
  (2018-08-01 18:00:00+00:00,12285.2500000000)
  (2018-08-02 04:00:00+00:00,11204.5000000000)
  (2018-08-02 14:00:00+00:00,12863.2500000000)
  (2018-08-03 00:00:00+00:00,9365.2500000000)
  (2018-08-03 10:00:00+00:00,13488.2500000000)
  (2018-08-03 20:00:00+00:00,10982.0000000000)
  (2018-08-04 06:00:00+00:00,10560.2500000000)
  (2018-08-04 16:00:00+00:00,11099.2500000000)
  (2018-08-05 02:00:00+00:00,8779.2500000000)
  (2018-08-05 12:00:00+00:00,10902.0000000000)
  (2018-08-05 22:00:00+00:00,9406.7500000000)
  (2018-08-06 08:00:00+00:00,13024.7500000000)
  (2018-08-06 18:00:00+00:00,11887.7500000000)
  (2018-08-07 04:00:00+00:00,11154.0000000000)
  (2018-08-07 14:00:00+00:00,12902.5000000000)
  (2018-08-08 00:00:00+00:00,10458.2500000000)
  (2018-08-08 10:00:00+00:00,14184.5000000000)
  (2018-08-08 20:00:00+00:00,11719.7500000000)
  (2018-08-09 06:00:00+00:00,12934.7500000000)
  (2018-08-09 16:00:00+00:00,13500.7500000000)
  (2018-08-10 02:00:00+00:00,10918.5000000000)
  (2018-08-10 12:00:00+00:00,13592.0000000000)
  (2018-08-10 22:00:00+00:00,9631.7500000000)
  (2018-08-11 08:00:00+00:00,12206.5000000000)
  (2018-08-11 18:00:00+00:00,11303.2500000000)
  (2018-08-12 04:00:00+00:00,8559.2500000000)
  (2018-08-12 14:00:00+00:00,9834.0000000000)
  (2018-08-13 00:00:00+00:00,8704.0000000000)
  (2018-08-13 10:00:00+00:00,13858.2500000000)
  (2018-08-13 20:00:00+00:00,11841.0000000000)
  (2018-08-14 06:00:00+00:00,13336.5000000000)
  (2018-08-14 16:00:00+00:00,13795.5000000000)
  (2018-08-15 02:00:00+00:00,10249.5000000000)
  (2018-08-15 12:00:00+00:00,13362.5000000000)
  (2018-08-15 22:00:00+00:00,9958.5000000000)
  (2018-08-16 08:00:00+00:00,13589.7500000000)
  (2018-08-16 18:00:00+00:00,13133.0000000000)
  (2018-08-17 04:00:00+00:00,11789.0000000000)
  (2018-08-17 14:00:00+00:00,12956.2500000000)
  (2018-08-18 00:00:00+00:00,9166.2500000000)
  (2018-08-18 10:00:00+00:00,11921.2500000000)
  (2018-08-18 20:00:00+00:00,9632.0000000000)
  (2018-08-19 06:00:00+00:00,9974.7500000000)
  (2018-08-19 16:00:00+00:00,10861.0000000000)
  (2018-08-20 02:00:00+00:00,9888.5000000000)
  (2018-08-20 12:00:00+00:00,13819.5000000000)
  (2018-08-20 22:00:00+00:00,10377.0000000000)
  (2018-08-21 08:00:00+00:00,13440.5000000000)
  (2018-08-21 18:00:00+00:00,12403.0000000000)
  (2018-08-22 04:00:00+00:00,11570.0000000000)
  (2018-08-22 14:00:00+00:00,12819.0000000000)
  (2018-08-23 00:00:00+00:00,9472.0000000000)
  (2018-08-23 10:00:00+00:00,13588.0000000000)
  (2018-08-23 20:00:00+00:00,11163.5000000000)
  (2018-08-24 06:00:00+00:00,13569.2500000000)
  (2018-08-24 16:00:00+00:00,13097.2500000000)
  (2018-08-25 02:00:00+00:00,9353.0000000000)
  (2018-08-25 12:00:00+00:00,11534.7500000000)
  (2018-08-25 22:00:00+00:00,9400.0000000000)
  (2018-08-26 08:00:00+00:00,11322.5000000000)
  (2018-08-26 18:00:00+00:00,10746.2500000000)
  (2018-08-27 04:00:00+00:00,12574.0000000000)
  (2018-08-27 14:00:00+00:00,13550.5000000000)
  (2018-08-28 00:00:00+00:00,9889.5000000000)
  (2018-08-28 10:00:00+00:00,13553.2500000000)
  (2018-08-28 20:00:00+00:00,11023.0000000000)
  (2018-08-29 06:00:00+00:00,13293.0000000000)
  (2018-08-29 16:00:00+00:00,13137.7500000000)
  (2018-08-30 02:00:00+00:00,10023.7500000000)
  (2018-08-30 12:00:00+00:00,13872.2500000000)
  (2018-08-30 22:00:00+00:00,10645.2500000000)
  (2018-08-31 08:00:00+00:00,13587.0000000000)
  (2018-08-31 18:00:00+00:00,12186.0000000000)
  (2018-09-01 04:00:00+00:00,9311.5000000000)
  (2018-09-01 14:00:00+00:00,10782.7500000000)
  (2018-09-02 00:00:00+00:00,8254.7500000000)
  (2018-09-02 10:00:00+00:00,11308.0000000000)
  (2018-09-02 20:00:00+00:00,10773.2500000000)
  (2018-09-03 06:00:00+00:00,13837.7500000000)
  (2018-09-03 16:00:00+00:00,13439.7500000000)
  (2018-09-04 02:00:00+00:00,9890.7500000000)
  (2018-09-04 12:00:00+00:00,13382.7500000000)
  (2018-09-04 22:00:00+00:00,10094.7500000000)
  (2018-09-05 08:00:00+00:00,13305.2500000000)
  (2018-09-05 18:00:00+00:00,12735.0000000000)
  (2018-09-06 04:00:00+00:00,11728.5000000000)
  (2018-09-06 14:00:00+00:00,12790.2500000000)
  (2018-09-07 00:00:00+00:00,9865.5000000000)
  (2018-09-07 10:00:00+00:00,13891.0000000000)
  (2018-09-07 20:00:00+00:00,11517.0000000000)
  (2018-09-08 06:00:00+00:00,11669.2500000000)
  (2018-09-08 16:00:00+00:00,10884.7500000000)
  (2018-09-09 02:00:00+00:00,7895.7500000000)
  (2018-09-09 12:00:00+00:00,9805.5000000000)
  (2018-09-09 22:00:00+00:00,9069.5000000000)
  (2018-09-10 08:00:00+00:00,13401.0000000000)
  (2018-09-10 18:00:00+00:00,12815.2500000000)
  (2018-09-11 04:00:00+00:00,12268.2500000000)
  (2018-09-11 14:00:00+00:00,14050.7500000000)
  (2018-09-12 00:00:00+00:00,9866.2500000000)
  (2018-09-12 10:00:00+00:00,13819.0000000000)
  (2018-09-12 20:00:00+00:00,11073.0000000000)
  (2018-09-13 06:00:00+00:00,12759.0000000000)
  (2018-09-13 16:00:00+00:00,12494.5000000000)
  (2018-09-14 02:00:00+00:00,9843.2500000000)
  (2018-09-14 12:00:00+00:00,13027.5000000000)
  (2018-09-14 22:00:00+00:00,9666.0000000000)
  (2018-09-15 08:00:00+00:00,12611.2500000000)
  (2018-09-15 18:00:00+00:00,11385.0000000000)
  (2018-09-16 04:00:00+00:00,8439.2500000000)
  (2018-09-16 14:00:00+00:00,9939.7500000000)
  (2018-09-17 00:00:00+00:00,9223.5000000000)
  (2018-09-17 10:00:00+00:00,13568.0000000000)
  (2018-09-17 20:00:00+00:00,10876.7500000000)
  (2018-09-18 06:00:00+00:00,13575.5000000000)
  (2018-09-18 16:00:00+00:00,13620.7500000000)
  (2018-09-19 02:00:00+00:00,10340.2500000000)
  (2018-09-19 12:00:00+00:00,13548.7500000000)
  (2018-09-19 22:00:00+00:00,10853.7500000000)
  (2018-09-20 08:00:00+00:00,14034.7500000000)
  (2018-09-20 18:00:00+00:00,13280.2500000000)
  (2018-09-21 04:00:00+00:00,13846.7500000000)
  (2018-09-21 14:00:00+00:00,14382.7500000000)
  (2018-09-22 00:00:00+00:00,10108.5000000000)
  (2018-09-22 10:00:00+00:00,12755.5000000000)
  (2018-09-22 20:00:00+00:00,10905.0000000000)
  (2018-09-23 06:00:00+00:00,10606.7500000000)
  (2018-09-23 16:00:00+00:00,11945.0000000000)
  (2018-09-24 02:00:00+00:00,10725.0000000000)
  (2018-09-24 12:00:00+00:00,14485.0000000000)
  (2018-09-24 22:00:00+00:00,10863.0000000000)
  (2018-09-25 08:00:00+00:00,14245.7500000000)
  (2018-09-25 18:00:00+00:00,13586.0000000000)
  (2018-09-26 04:00:00+00:00,13866.0000000000)
  (2018-09-26 14:00:00+00:00,14024.5000000000)
  (2018-09-27 00:00:00+00:00,10986.0000000000)
  (2018-09-27 10:00:00+00:00,14541.0000000000)
  (2018-09-27 20:00:00+00:00,12603.7500000000)
  (2018-09-28 06:00:00+00:00,14965.0000000000)
  (2018-09-28 16:00:00+00:00,13875.5000000000)
  (2018-09-29 02:00:00+00:00,9871.0000000000)
  (2018-09-29 12:00:00+00:00,10997.0000000000)
  (2018-09-29 22:00:00+00:00,9703.0000000000)
  (2018-09-30 08:00:00+00:00,11146.2500000000)
  (2018-09-30 18:00:00+00:00,10915.5000000000)
  (2018-10-01 04:00:00+00:00,11606.7500000000)
  (2018-10-01 14:00:00+00:00,12740.2500000000)
  (2018-10-02 00:00:00+00:00,10063.0000000000)
  (2018-10-02 10:00:00+00:00,14948.2500000000)
  (2018-10-02 20:00:00+00:00,12230.7500000000)
  (2018-10-03 06:00:00+00:00,11768.2500000000)
  (2018-10-03 16:00:00+00:00,12499.5000000000)
  (2018-10-04 02:00:00+00:00,10554.7500000000)
  (2018-10-04 12:00:00+00:00,13317.2500000000)
  (2018-10-04 22:00:00+00:00,11048.5000000000)
  (2018-10-05 08:00:00+00:00,14133.5000000000)
  (2018-10-05 18:00:00+00:00,13351.7500000000)
  (2018-10-06 04:00:00+00:00,10128.7500000000)
  (2018-10-06 14:00:00+00:00,10652.7500000000)
  (2018-10-07 00:00:00+00:00,8152.5000000000)
  (2018-10-07 10:00:00+00:00,11610.0000000000)
  (2018-10-07 20:00:00+00:00,9906.0000000000)
  (2018-10-08 06:00:00+00:00,13797.0000000000)
  (2018-10-08 16:00:00+00:00,13147.2500000000)
  (2018-10-09 02:00:00+00:00,10058.2500000000)
  (2018-10-09 12:00:00+00:00,13023.2500000000)
  (2018-10-09 22:00:00+00:00,10032.2500000000)
  (2018-10-10 08:00:00+00:00,13239.0000000000)
  (2018-10-10 18:00:00+00:00,13746.2500000000)
  (2018-10-11 04:00:00+00:00,13409.7500000000)
  (2018-10-11 14:00:00+00:00,13543.0000000000)
  (2018-10-12 00:00:00+00:00,10568.5000000000)
  (2018-10-12 10:00:00+00:00,13473.5000000000)
  (2018-10-12 20:00:00+00:00,11277.5000000000)
  (2018-10-13 06:00:00+00:00,10980.5000000000)
  (2018-10-13 16:00:00+00:00,11921.5000000000)
  (2018-10-14 02:00:00+00:00,8933.0000000000)
  (2018-10-14 12:00:00+00:00,10809.0000000000)
  (2018-10-14 22:00:00+00:00,10036.7500000000)
  (2018-10-15 08:00:00+00:00,13602.0000000000)
  (2018-10-15 18:00:00+00:00,13272.0000000000)
  (2018-10-16 04:00:00+00:00,11997.0000000000)
  (2018-10-16 14:00:00+00:00,12475.5000000000)
  (2018-10-17 00:00:00+00:00,9341.2500000000)
  (2018-10-17 10:00:00+00:00,13313.7500000000)
  (2018-10-17 20:00:00+00:00,11366.5000000000)
  (2018-10-18 06:00:00+00:00,13705.0000000000)
  (2018-10-18 16:00:00+00:00,14364.0000000000)
  (2018-10-19 02:00:00+00:00,10210.2500000000)
  (2018-10-19 12:00:00+00:00,13037.0000000000)
  (2018-10-19 22:00:00+00:00,9957.0000000000)
  (2018-10-20 08:00:00+00:00,12608.2500000000)
  (2018-10-20 18:00:00+00:00,11739.2500000000)
  (2018-10-21 04:00:00+00:00,8639.7500000000)
  (2018-10-21 14:00:00+00:00,10147.5000000000)
  (2018-10-22 00:00:00+00:00,9733.7500000000)
  (2018-10-22 10:00:00+00:00,14536.0000000000)
  (2018-10-22 20:00:00+00:00,12468.5000000000)
  (2018-10-23 06:00:00+00:00,15171.0000000000)
  (2018-10-23 16:00:00+00:00,15479.2500000000)
  (2018-10-24 02:00:00+00:00,11496.7500000000)
  (2018-10-24 12:00:00+00:00,14846.2500000000)
  (2018-10-24 22:00:00+00:00,11169.2500000000)
  (2018-10-25 08:00:00+00:00,15585.0000000000)
  (2018-10-25 18:00:00+00:00,14561.7500000000)
  (2018-10-26 04:00:00+00:00,13006.0000000000)
  (2018-10-26 14:00:00+00:00,13455.2500000000)
  (2018-10-27 00:00:00+00:00,10203.0000000000)
  (2018-10-27 10:00:00+00:00,13602.0000000000)
  (2018-10-27 20:00:00+00:00,10626.0000000000)
  (2018-10-28 06:00:00+00:00,9796.7500000000)
  (2018-10-28 16:00:00+00:00,12764.5000000000)
  (2018-10-29 02:00:00+00:00,10306.5000000000)
  (2018-10-29 12:00:00+00:00,14808.5000000000)
  (2018-10-29 22:00:00+00:00,11954.5000000000)
  (2018-10-30 08:00:00+00:00,14782.7500000000)
  (2018-10-30 18:00:00+00:00,14779.7500000000)
  (2018-10-31 04:00:00+00:00,11140.2500000000)
  (2018-10-31 14:00:00+00:00,10978.5000000000)
  (2018-11-01 00:00:00+00:00,9566.2500000000)
  (2018-11-01 10:00:00+00:00,13942.2500000000)
  (2018-11-01 20:00:00+00:00,12244.2500000000)
  (2018-11-02 06:00:00+00:00,13755.2500000000)
  (2018-11-02 16:00:00+00:00,14502.0000000000)
  (2018-11-03 02:00:00+00:00,9742.2500000000)
  (2018-11-03 12:00:00+00:00,11612.2500000000)
  (2018-11-03 22:00:00+00:00,9854.5000000000)
  (2018-11-04 08:00:00+00:00,11167.2500000000)
  (2018-11-04 18:00:00+00:00,11843.7500000000)
  (2018-11-05 04:00:00+00:00,10890.2500000000)
  (2018-11-05 14:00:00+00:00,13319.0000000000)
  (2018-11-06 00:00:00+00:00,10335.2500000000)
  (2018-11-06 10:00:00+00:00,13632.2500000000)
  (2018-11-06 20:00:00+00:00,12349.5000000000)
  (2018-11-07 06:00:00+00:00,13696.5000000000)
  (2018-11-07 16:00:00+00:00,14414.0000000000)
  (2018-11-08 02:00:00+00:00,9974.0000000000)
  (2018-11-08 12:00:00+00:00,13546.2500000000)
  (2018-11-08 22:00:00+00:00,10833.7500000000)
  (2018-11-09 08:00:00+00:00,14167.5000000000)
  (2018-11-09 18:00:00+00:00,13938.2500000000)
  (2018-11-10 04:00:00+00:00,9781.0000000000)
  (2018-11-10 14:00:00+00:00,11776.0000000000)
  (2018-11-11 00:00:00+00:00,8941.2500000000)
  (2018-11-11 10:00:00+00:00,11968.2500000000)
  (2018-11-11 20:00:00+00:00,10932.5000000000)
  (2018-11-12 06:00:00+00:00,13164.7500000000)
  (2018-11-12 16:00:00+00:00,14464.0000000000)
  (2018-11-13 02:00:00+00:00,10014.0000000000)
  (2018-11-13 12:00:00+00:00,13932.2500000000)
  (2018-11-13 22:00:00+00:00,11614.2500000000)
  (2018-11-14 08:00:00+00:00,14102.0000000000)
  (2018-11-14 18:00:00+00:00,14342.2500000000)
  (2018-11-15 04:00:00+00:00,11442.5000000000)
  (2018-11-15 14:00:00+00:00,13451.7500000000)
  (2018-11-16 00:00:00+00:00,10867.0000000000)
  (2018-11-16 10:00:00+00:00,14243.2500000000)
  (2018-11-16 20:00:00+00:00,12679.2500000000)
  (2018-11-17 06:00:00+00:00,11008.2500000000)
  (2018-11-17 16:00:00+00:00,12862.2500000000)
  (2018-11-18 02:00:00+00:00,9284.5000000000)
  (2018-11-18 12:00:00+00:00,11636.0000000000)
  (2018-11-18 22:00:00+00:00,10803.7500000000)
  (2018-11-19 08:00:00+00:00,14306.2500000000)
  (2018-11-19 18:00:00+00:00,14581.2500000000)
  (2018-11-20 04:00:00+00:00,11745.0000000000)
  (2018-11-20 14:00:00+00:00,14444.7500000000)
  (2018-11-21 00:00:00+00:00,10615.2500000000)
  (2018-11-21 10:00:00+00:00,14536.0000000000)
  (2018-11-21 20:00:00+00:00,12993.0000000000)
  (2018-11-22 06:00:00+00:00,14305.7500000000)
  (2018-11-22 16:00:00+00:00,14763.0000000000)
  (2018-11-23 02:00:00+00:00,10607.2500000000)
  (2018-11-23 12:00:00+00:00,14230.7500000000)
  (2018-11-23 22:00:00+00:00,11401.0000000000)
  (2018-11-24 08:00:00+00:00,12514.5000000000)
  (2018-11-24 18:00:00+00:00,12720.5000000000)
  (2018-11-25 04:00:00+00:00,9507.5000000000)
  (2018-11-25 14:00:00+00:00,11894.0000000000)
  (2018-11-26 00:00:00+00:00,10374.7500000000)
  (2018-11-26 10:00:00+00:00,14850.2500000000)
  (2018-11-26 20:00:00+00:00,12924.2500000000)
  (2018-11-27 06:00:00+00:00,14554.5000000000)
  (2018-11-27 16:00:00+00:00,15258.7500000000)
  (2018-11-28 02:00:00+00:00,11435.5000000000)
  (2018-11-28 12:00:00+00:00,14426.2500000000)
  (2018-11-28 22:00:00+00:00,12722.0000000000)
  (2018-11-29 08:00:00+00:00,14646.0000000000)
  (2018-11-29 18:00:00+00:00,14795.2500000000)
  (2018-11-30 04:00:00+00:00,11883.0000000000)
  (2018-11-30 14:00:00+00:00,14527.2500000000)
  (2018-12-01 00:00:00+00:00,10970.7500000000)
  (2018-12-01 10:00:00+00:00,13040.7500000000)
  (2018-12-01 20:00:00+00:00,11467.2500000000)
  (2018-12-02 06:00:00+00:00,9603.0000000000)
  (2018-12-02 16:00:00+00:00,12737.5000000000)
  (2018-12-03 02:00:00+00:00,9289.0000000000)
  (2018-12-03 12:00:00+00:00,14247.2500000000)
  (2018-12-03 22:00:00+00:00,11370.0000000000)
  (2018-12-04 08:00:00+00:00,14769.2500000000)
  (2018-12-04 18:00:00+00:00,15662.2500000000)
  (2018-12-05 04:00:00+00:00,12485.0000000000)
  (2018-12-05 14:00:00+00:00,14109.2500000000)
  (2018-12-06 00:00:00+00:00,11493.7500000000)
  (2018-12-06 10:00:00+00:00,15054.5000000000)
  (2018-12-06 20:00:00+00:00,13462.5000000000)
  (2018-12-07 06:00:00+00:00,14327.7500000000)
  (2018-12-07 16:00:00+00:00,14906.7500000000)
  (2018-12-08 02:00:00+00:00,10006.2500000000)
  (2018-12-08 12:00:00+00:00,12257.5000000000)
  (2018-12-08 22:00:00+00:00,10288.5000000000)
  (2018-12-09 08:00:00+00:00,10735.5000000000)
  (2018-12-09 18:00:00+00:00,12179.2500000000)
  (2018-12-10 04:00:00+00:00,11298.5000000000)
  (2018-12-10 14:00:00+00:00,14421.5000000000)
  (2018-12-11 00:00:00+00:00,10829.7500000000)
  (2018-12-11 10:00:00+00:00,14801.7500000000)
  (2018-12-11 20:00:00+00:00,13358.2500000000)
  (2018-12-12 06:00:00+00:00,14532.0000000000)
  (2018-12-12 16:00:00+00:00,15353.0000000000)
  (2018-12-13 02:00:00+00:00,10884.2500000000)
  (2018-12-13 12:00:00+00:00,14456.0000000000)
  (2018-12-13 22:00:00+00:00,11924.0000000000)
  (2018-12-14 08:00:00+00:00,14923.2500000000)
  (2018-12-14 18:00:00+00:00,14286.2500000000)
  (2018-12-15 04:00:00+00:00,10379.5000000000)
  (2018-12-15 14:00:00+00:00,12705.7500000000)
  (2018-12-16 00:00:00+00:00,10086.5000000000)
  (2018-12-16 10:00:00+00:00,13230.2500000000)
  (2018-12-16 20:00:00+00:00,11741.0000000000)
  (2018-12-17 06:00:00+00:00,14186.5000000000)
  (2018-12-17 16:00:00+00:00,14738.7500000000)
  (2018-12-18 02:00:00+00:00,10420.5000000000)
  (2018-12-18 12:00:00+00:00,14047.7500000000)
  (2018-12-18 22:00:00+00:00,11607.2500000000)
  (2018-12-19 08:00:00+00:00,15293.5000000000)
  (2018-12-19 18:00:00+00:00,14627.7500000000)
  (2018-12-20 04:00:00+00:00,11156.0000000000)
  (2018-12-20 14:00:00+00:00,13489.7500000000)
  (2018-12-21 00:00:00+00:00,9692.7500000000)
  (2018-12-21 10:00:00+00:00,13365.5000000000)
  (2018-12-21 20:00:00+00:00,11284.7500000000)
  (2018-12-22 06:00:00+00:00,10711.0000000000)
  (2018-12-22 16:00:00+00:00,12469.7500000000)
  (2018-12-23 02:00:00+00:00,8038.7500000000)
  (2018-12-23 12:00:00+00:00,11490.5000000000)
  (2018-12-23 22:00:00+00:00,9796.2500000000)
  (2018-12-24 08:00:00+00:00,11458.5000000000)
  (2018-12-24 18:00:00+00:00,10471.0000000000)
  (2018-12-25 04:00:00+00:00,8410.7500000000)
  (2018-12-25 14:00:00+00:00,9948.0000000000)
  (2018-12-26 00:00:00+00:00,8427.7500000000)
  (2018-12-26 10:00:00+00:00,11563.7500000000)
  (2018-12-26 20:00:00+00:00,10437.0000000000)
  (2018-12-27 06:00:00+00:00,11139.7500000000)
  (2018-12-27 16:00:00+00:00,13229.0000000000)
  (2018-12-28 02:00:00+00:00,9344.7500000000)
  (2018-12-28 12:00:00+00:00,12099.7500000000)
  (2018-12-28 22:00:00+00:00,9854.7500000000)
  (2018-12-29 08:00:00+00:00,11391.7500000000)
  (2018-12-29 18:00:00+00:00,11911.5000000000)
  (2018-12-30 04:00:00+00:00,8501.5000000000)
  (2018-12-30 14:00:00+00:00,10872.7500000000)
  (2018-12-31 00:00:00+00:00,8593.0000000000)
  (2018-12-31 10:00:00+00:00,11212.0000000000)
  (2018-12-31 20:00:00+00:00,10432.7500000000)
  (2019-01-01 06:00:00+00:00,7357.5000000000)
  (2019-01-01 16:00:00+00:00,10476.0000000000)
  (2019-01-02 02:00:00+00:00,7735.2500000000)
  (2019-01-02 12:00:00+00:00,12792.5000000000)
  (2019-01-02 22:00:00+00:00,11044.7500000000)
  (2019-01-03 08:00:00+00:00,13702.7500000000)
  (2019-01-03 18:00:00+00:00,13898.7500000000)
  (2019-01-04 04:00:00+00:00,11590.0000000000)
  (2019-01-04 14:00:00+00:00,12981.0000000000)
  (2019-01-05 00:00:00+00:00,9469.5000000000)
  (2019-01-05 10:00:00+00:00,12326.5000000000)
  (2019-01-05 20:00:00+00:00,10956.7500000000)
  (2019-01-06 06:00:00+00:00,9506.5000000000)
  (2019-01-06 16:00:00+00:00,12400.5000000000)
  (2019-01-07 02:00:00+00:00,9896.7500000000)
  (2019-01-07 12:00:00+00:00,14040.2500000000)
  (2019-01-07 22:00:00+00:00,10988.7500000000)
  (2019-01-08 08:00:00+00:00,13711.7500000000)
  (2019-01-08 18:00:00+00:00,13865.7500000000)
  (2019-01-09 04:00:00+00:00,11453.7500000000)
  (2019-01-09 14:00:00+00:00,14610.7500000000)
  (2019-01-10 00:00:00+00:00,10855.7500000000)
  (2019-01-10 10:00:00+00:00,14359.2500000000)
  (2019-01-10 20:00:00+00:00,13006.2500000000)
  (2019-01-11 06:00:00+00:00,14864.7500000000)
  (2019-01-11 16:00:00+00:00,14798.0000000000)
  (2019-01-12 02:00:00+00:00,9538.7500000000)
  (2019-01-12 12:00:00+00:00,12161.2500000000)
  (2019-01-12 22:00:00+00:00,9776.7500000000)
  (2019-01-13 08:00:00+00:00,10369.5000000000)
  (2019-01-13 18:00:00+00:00,11471.5000000000)
  (2019-01-14 04:00:00+00:00,10651.5000000000)
  (2019-01-14 14:00:00+00:00,13327.7500000000)
  (2019-01-15 00:00:00+00:00,10363.0000000000)
  (2019-01-15 10:00:00+00:00,13972.0000000000)
  (2019-01-15 20:00:00+00:00,12247.2500000000)
  (2019-01-16 06:00:00+00:00,13888.0000000000)
  (2019-01-16 16:00:00+00:00,14623.7500000000)
  (2019-01-17 02:00:00+00:00,10368.0000000000)
  (2019-01-17 12:00:00+00:00,14575.0000000000)
  (2019-01-17 22:00:00+00:00,11200.7500000000)
  (2019-01-18 08:00:00+00:00,14695.5000000000)
  (2019-01-18 18:00:00+00:00,14322.7500000000)
  (2019-01-19 04:00:00+00:00,10672.7500000000)
  (2019-01-19 14:00:00+00:00,11956.0000000000)
  (2019-01-20 00:00:00+00:00,10262.5000000000)
  (2019-01-20 10:00:00+00:00,12389.2500000000)
  (2019-01-20 20:00:00+00:00,11849.7500000000)
  (2019-01-21 06:00:00+00:00,15043.2500000000)
  (2019-01-21 16:00:00+00:00,15738.2500000000)
  (2019-01-22 02:00:00+00:00,11483.5000000000)
  (2019-01-22 12:00:00+00:00,14298.5000000000)
  (2019-01-22 22:00:00+00:00,12746.0000000000)
  (2019-01-23 08:00:00+00:00,15519.0000000000)
  (2019-01-23 18:00:00+00:00,15312.7500000000)
  (2019-01-24 04:00:00+00:00,12035.2500000000)
  (2019-01-24 14:00:00+00:00,14952.2500000000)
  (2019-01-25 00:00:00+00:00,11626.0000000000)
  (2019-01-25 10:00:00+00:00,15236.7500000000)
  (2019-01-25 20:00:00+00:00,12856.0000000000)
  (2019-01-26 06:00:00+00:00,11108.7500000000)
  (2019-01-26 16:00:00+00:00,12845.5000000000)
  (2019-01-27 02:00:00+00:00,8936.5000000000)
  (2019-01-27 12:00:00+00:00,11143.5000000000)
  (2019-01-27 22:00:00+00:00,10612.5000000000)
  (2019-01-28 08:00:00+00:00,14314.7500000000)
  (2019-01-28 18:00:00+00:00,14261.2500000000)
  (2019-01-29 04:00:00+00:00,12341.5000000000)
  (2019-01-29 14:00:00+00:00,14294.2500000000)
  (2019-01-30 00:00:00+00:00,11120.5000000000)
  (2019-01-30 10:00:00+00:00,14533.0000000000)
  (2019-01-30 20:00:00+00:00,13586.5000000000)
  (2019-01-31 06:00:00+00:00,14929.5000000000)
  (2019-01-31 16:00:00+00:00,14636.2500000000)
  (2019-02-01 02:00:00+00:00,11087.7500000000)
  (2019-02-01 12:00:00+00:00,14696.5000000000)
  (2019-02-01 22:00:00+00:00,11663.2500000000)
  (2019-02-02 08:00:00+00:00,12899.2500000000)
  (2019-02-02 18:00:00+00:00,13055.7500000000)
  (2019-02-03 04:00:00+00:00,9760.2500000000)
  (2019-02-03 14:00:00+00:00,11566.0000000000)
  (2019-02-04 00:00:00+00:00,10242.5000000000)
  (2019-02-04 10:00:00+00:00,14355.7500000000)
  (2019-02-04 20:00:00+00:00,13232.2500000000)
  (2019-02-05 06:00:00+00:00,14627.5000000000)
  (2019-02-05 16:00:00+00:00,15269.5000000000)
  (2019-02-06 02:00:00+00:00,11238.7500000000)
  (2019-02-06 12:00:00+00:00,14112.0000000000)
  (2019-02-06 22:00:00+00:00,11716.7500000000)
  (2019-02-07 08:00:00+00:00,14228.2500000000)
  (2019-02-07 18:00:00+00:00,14151.5000000000)
  (2019-02-08 04:00:00+00:00,11853.7500000000)
  (2019-02-08 14:00:00+00:00,12936.0000000000)
  (2019-02-09 00:00:00+00:00,9649.7500000000)
  (2019-02-09 10:00:00+00:00,11842.7500000000)
  (2019-02-09 20:00:00+00:00,10387.5000000000)
  (2019-02-10 06:00:00+00:00,9008.2500000000)
  (2019-02-10 16:00:00+00:00,10961.2500000000)
  (2019-02-11 02:00:00+00:00,9002.0000000000)
  (2019-02-11 12:00:00+00:00,13338.0000000000)
  (2019-02-11 22:00:00+00:00,11175.0000000000)
  (2019-02-12 08:00:00+00:00,14347.5000000000)
  (2019-02-12 18:00:00+00:00,14254.5000000000)
  (2019-02-13 04:00:00+00:00,11536.5000000000)
  (2019-02-13 14:00:00+00:00,14193.5000000000)
  (2019-02-14 00:00:00+00:00,10952.0000000000)
  (2019-02-14 10:00:00+00:00,14081.2500000000)
  (2019-02-14 20:00:00+00:00,12779.5000000000)
  (2019-02-15 06:00:00+00:00,13865.7500000000)
  (2019-02-15 16:00:00+00:00,13243.5000000000)
  (2019-02-16 02:00:00+00:00,9800.2500000000)
  (2019-02-16 12:00:00+00:00,11155.0000000000)
  (2019-02-16 22:00:00+00:00,10398.0000000000)
  (2019-02-17 08:00:00+00:00,10913.5000000000)
  (2019-02-17 18:00:00+00:00,11815.7500000000)
  (2019-02-18 04:00:00+00:00,10620.0000000000)
  (2019-02-18 14:00:00+00:00,12207.2500000000)
  (2019-02-19 00:00:00+00:00,9556.2500000000)
  (2019-02-19 10:00:00+00:00,14116.5000000000)
  (2019-02-19 20:00:00+00:00,12748.0000000000)
  (2019-02-20 06:00:00+00:00,14164.5000000000)
  (2019-02-20 16:00:00+00:00,13861.2500000000)
  (2019-02-21 02:00:00+00:00,10405.0000000000)
  (2019-02-21 12:00:00+00:00,13568.7500000000)
  (2019-02-21 22:00:00+00:00,11248.7500000000)
  (2019-02-22 08:00:00+00:00,14399.5000000000)
  (2019-02-22 18:00:00+00:00,13973.5000000000)
  (2019-02-23 04:00:00+00:00,10196.2500000000)
  (2019-02-23 14:00:00+00:00,10903.2500000000)
  (2019-02-24 00:00:00+00:00,9535.2500000000)
  (2019-02-24 10:00:00+00:00,11295.0000000000)
  (2019-02-24 20:00:00+00:00,10953.0000000000)
  (2019-02-25 06:00:00+00:00,13183.5000000000)
  (2019-02-25 16:00:00+00:00,13227.5000000000)
  (2019-02-26 02:00:00+00:00,10044.0000000000)
  (2019-02-26 12:00:00+00:00,12676.0000000000)
  (2019-02-26 22:00:00+00:00,10811.0000000000)
  (2019-02-27 08:00:00+00:00,13757.2500000000)
  (2019-02-27 18:00:00+00:00,14206.0000000000)
  (2019-02-28 04:00:00+00:00,10861.0000000000)
  (2019-02-28 14:00:00+00:00,12588.7500000000)
  (2019-03-01 00:00:00+00:00,9637.2500000000)
  (2019-03-01 10:00:00+00:00,13985.5000000000)
  (2019-03-01 20:00:00+00:00,12019.0000000000)
  (2019-03-02 06:00:00+00:00,10307.2500000000)
  (2019-03-02 16:00:00+00:00,12057.0000000000)
  (2019-03-03 02:00:00+00:00,8362.0000000000)
  (2019-03-03 12:00:00+00:00,9933.0000000000)
  (2019-03-03 22:00:00+00:00,8748.2500000000)
  (2019-03-04 08:00:00+00:00,12637.0000000000)
  (2019-03-04 18:00:00+00:00,13087.5000000000)
  (2019-03-05 04:00:00+00:00,10241.7500000000)
  (2019-03-05 14:00:00+00:00,12717.5000000000)
  (2019-03-06 00:00:00+00:00,9372.0000000000)
  (2019-03-06 10:00:00+00:00,12502.2500000000)
  (2019-03-06 20:00:00+00:00,11915.2500000000)
  (2019-03-07 06:00:00+00:00,12656.2500000000)
  (2019-03-07 16:00:00+00:00,12496.2500000000)
  (2019-03-08 02:00:00+00:00,9059.2500000000)
  (2019-03-08 12:00:00+00:00,12514.5000000000)
  (2019-03-08 22:00:00+00:00,10118.5000000000)
  (2019-03-09 08:00:00+00:00,10822.0000000000)
  (2019-03-09 18:00:00+00:00,11139.2500000000)
  (2019-03-10 04:00:00+00:00,8513.7500000000)
  (2019-03-10 14:00:00+00:00,10044.0000000000)
  (2019-03-11 00:00:00+00:00,8507.0000000000)
  (2019-03-11 10:00:00+00:00,12628.5000000000)
  (2019-03-11 20:00:00+00:00,12026.5000000000)
  (2019-03-12 06:00:00+00:00,12993.7500000000)
  (2019-03-12 16:00:00+00:00,12747.0000000000)
  (2019-03-13 02:00:00+00:00,9480.7500000000)
  (2019-03-13 12:00:00+00:00,12745.7500000000)
  (2019-03-13 22:00:00+00:00,10451.0000000000)
  (2019-03-14 08:00:00+00:00,12936.2500000000)
  (2019-03-14 18:00:00+00:00,12903.7500000000)
  (2019-03-15 04:00:00+00:00,10640.0000000000)
  (2019-03-15 14:00:00+00:00,12467.0000000000)
  (2019-03-16 00:00:00+00:00,9048.7500000000)
  (2019-03-16 10:00:00+00:00,12057.7500000000)
  (2019-03-16 20:00:00+00:00,10193.5000000000)
  (2019-03-17 06:00:00+00:00,8522.7500000000)
  (2019-03-17 16:00:00+00:00,10428.5000000000)
  (2019-03-18 02:00:00+00:00,9129.0000000000)
  (2019-03-18 12:00:00+00:00,12895.2500000000)
  (2019-03-18 22:00:00+00:00,10709.0000000000)
  (2019-03-19 08:00:00+00:00,12753.0000000000)
  (2019-03-19 18:00:00+00:00,13969.2500000000)
  (2019-03-20 04:00:00+00:00,11633.2500000000)
  (2019-03-20 14:00:00+00:00,12781.7500000000)
  (2019-03-21 00:00:00+00:00,9791.5000000000)
  (2019-03-21 10:00:00+00:00,13611.0000000000)
  (2019-03-21 20:00:00+00:00,13039.5000000000)
  (2019-03-22 06:00:00+00:00,13337.5000000000)
  (2019-03-22 16:00:00+00:00,12499.5000000000)
  (2019-03-23 02:00:00+00:00,8791.0000000000)
  (2019-03-23 12:00:00+00:00,10360.0000000000)
  (2019-03-23 22:00:00+00:00,9410.2500000000)
  (2019-03-24 08:00:00+00:00,10133.7500000000)
  (2019-03-24 18:00:00+00:00,11692.7500000000)
  (2019-03-25 04:00:00+00:00,9952.2500000000)
  (2019-03-25 14:00:00+00:00,12076.2500000000)
  (2019-03-26 00:00:00+00:00,9841.5000000000)
  (2019-03-26 10:00:00+00:00,13449.7500000000)
  (2019-03-26 20:00:00+00:00,12105.7500000000)
  (2019-03-27 06:00:00+00:00,13573.0000000000)
  (2019-03-27 16:00:00+00:00,13584.2500000000)
  (2019-03-28 02:00:00+00:00,9561.0000000000)
  (2019-03-28 12:00:00+00:00,13374.5000000000)
  (2019-03-28 22:00:00+00:00,11325.0000000000)
  (2019-03-29 08:00:00+00:00,13865.5000000000)
  (2019-03-29 18:00:00+00:00,13939.7500000000)
  (2019-03-30 04:00:00+00:00,9293.2500000000)
  (2019-03-30 14:00:00+00:00,9552.0000000000)
  (2019-03-31 00:00:00+00:00,8379.5000000000)
  (2019-03-31 10:00:00+00:00,10803.7500000000)
  (2019-03-31 20:00:00+00:00,9720.5000000000)
  (2019-04-01 06:00:00+00:00,12591.7500000000)
  (2019-04-01 16:00:00+00:00,11853.0000000000)
  (2019-04-02 02:00:00+00:00,10354.2500000000)
  (2019-04-02 12:00:00+00:00,12491.2500000000)
  (2019-04-02 22:00:00+00:00,9931.0000000000)
  (2019-04-03 08:00:00+00:00,13484.7500000000)
  (2019-04-03 18:00:00+00:00,13264.0000000000)
  (2019-04-04 04:00:00+00:00,12752.0000000000)
  (2019-04-04 14:00:00+00:00,12300.7500000000)
  (2019-04-05 00:00:00+00:00,9471.2500000000)
  (2019-04-05 10:00:00+00:00,13499.7500000000)
  (2019-04-05 20:00:00+00:00,10728.0000000000)
  (2019-04-06 06:00:00+00:00,10603.2500000000)
  (2019-04-06 16:00:00+00:00,10101.7500000000)
  (2019-04-07 02:00:00+00:00,7962.5000000000)
  (2019-04-07 12:00:00+00:00,8565.0000000000)
  (2019-04-07 22:00:00+00:00,8320.5000000000)
  (2019-04-08 08:00:00+00:00,12905.7500000000)
  (2019-04-08 18:00:00+00:00,12567.5000000000)
  (2019-04-09 04:00:00+00:00,11880.7500000000)
  (2019-04-09 14:00:00+00:00,11742.7500000000)
  (2019-04-10 00:00:00+00:00,9184.5000000000)
  (2019-04-10 10:00:00+00:00,12357.5000000000)
  (2019-04-10 20:00:00+00:00,10830.7500000000)
  (2019-04-11 06:00:00+00:00,13254.2500000000)
  (2019-04-11 16:00:00+00:00,12904.7500000000)
  (2019-04-12 02:00:00+00:00,10119.7500000000)
  (2019-04-12 12:00:00+00:00,12876.0000000000)
  (2019-04-12 22:00:00+00:00,10062.5000000000)
  (2019-04-13 08:00:00+00:00,11794.0000000000)
  (2019-04-13 18:00:00+00:00,11239.5000000000)
  (2019-04-14 04:00:00+00:00,8742.7500000000)
  (2019-04-14 14:00:00+00:00,10087.5000000000)
  (2019-04-15 00:00:00+00:00,9028.2500000000)
  (2019-04-15 10:00:00+00:00,12241.5000000000)
  (2019-04-15 20:00:00+00:00,10978.5000000000)
  (2019-04-16 06:00:00+00:00,13083.5000000000)
  (2019-04-16 16:00:00+00:00,12214.7500000000)
  (2019-04-17 02:00:00+00:00,9540.2500000000)
  (2019-04-17 12:00:00+00:00,11370.7500000000)
  (2019-04-17 22:00:00+00:00,9588.5000000000)
  (2019-04-18 08:00:00+00:00,12258.5000000000)
  (2019-04-18 18:00:00+00:00,11150.0000000000)
  (2019-04-19 04:00:00+00:00,7593.5000000000)
  (2019-04-19 14:00:00+00:00,8035.5000000000)
  (2019-04-20 00:00:00+00:00,7108.2500000000)
  (2019-04-20 10:00:00+00:00,9465.2500000000)
  (2019-04-20 20:00:00+00:00,8103.7500000000)
  (2019-04-21 06:00:00+00:00,8297.5000000000)
  (2019-04-21 16:00:00+00:00,8204.5000000000)
  (2019-04-22 02:00:00+00:00,6650.5000000000)
  (2019-04-22 12:00:00+00:00,7459.5000000000)
  (2019-04-22 22:00:00+00:00,7433.7500000000)
  (2019-04-23 08:00:00+00:00,11603.2500000000)
  (2019-04-23 18:00:00+00:00,11398.0000000000)
  (2019-04-24 04:00:00+00:00,10951.5000000000)
  (2019-04-24 14:00:00+00:00,11348.5000000000)
  (2019-04-25 00:00:00+00:00,8772.2500000000)
  (2019-04-25 10:00:00+00:00,12178.0000000000)
  (2019-04-25 20:00:00+00:00,9981.7500000000)
  (2019-04-26 06:00:00+00:00,12326.7500000000)
  (2019-04-26 16:00:00+00:00,11281.5000000000)
  (2019-04-27 02:00:00+00:00,8708.0000000000)
  (2019-04-27 12:00:00+00:00,9796.5000000000)
  (2019-04-27 22:00:00+00:00,7927.0000000000)
  (2019-04-28 08:00:00+00:00,9905.2500000000)
  (2019-04-28 18:00:00+00:00,10310.2500000000)
  (2019-04-29 04:00:00+00:00,11094.2500000000)
  (2019-04-29 14:00:00+00:00,12101.5000000000)
  (2019-04-30 00:00:00+00:00,8477.7500000000)
  (2019-04-30 10:00:00+00:00,11647.7500000000)
  (2019-04-30 20:00:00+00:00,9885.7500000000)
  (2019-05-01 06:00:00+00:00,8909.2500000000)
  (2019-05-01 16:00:00+00:00,9451.7500000000)
  (2019-05-02 02:00:00+00:00,8541.5000000000)
  (2019-05-02 12:00:00+00:00,12041.2500000000)
  (2019-05-02 22:00:00+00:00,9435.5000000000)
  (2019-05-03 08:00:00+00:00,12283.5000000000)
  (2019-05-03 18:00:00+00:00,11898.0000000000)
  (2019-05-04 04:00:00+00:00,9344.7500000000)
  (2019-05-04 14:00:00+00:00,9942.7500000000)
  (2019-05-05 00:00:00+00:00,7998.0000000000)
  (2019-05-05 10:00:00+00:00,10186.7500000000)
  (2019-05-05 20:00:00+00:00,10115.5000000000)
  (2019-05-06 06:00:00+00:00,12714.5000000000)
  (2019-05-06 16:00:00+00:00,12476.2500000000)
  (2019-05-07 02:00:00+00:00,10365.0000000000)
  (2019-05-07 12:00:00+00:00,13105.0000000000)
  (2019-05-07 22:00:00+00:00,9805.2500000000)
  (2019-05-08 08:00:00+00:00,12526.7500000000)
  (2019-05-08 18:00:00+00:00,12196.5000000000)
  (2019-05-09 04:00:00+00:00,11668.0000000000)
  (2019-05-09 14:00:00+00:00,11871.7500000000)
  (2019-05-10 00:00:00+00:00,8808.5000000000)
  (2019-05-10 10:00:00+00:00,12623.0000000000)
  (2019-05-10 20:00:00+00:00,11193.0000000000)
  (2019-05-11 06:00:00+00:00,10133.0000000000)
  (2019-05-11 16:00:00+00:00,10149.0000000000)
  (2019-05-12 02:00:00+00:00,7722.2500000000)
  (2019-05-12 12:00:00+00:00,8525.2500000000)
  (2019-05-12 22:00:00+00:00,8295.0000000000)
  (2019-05-13 08:00:00+00:00,12159.7500000000)
  (2019-05-13 18:00:00+00:00,11833.7500000000)
  (2019-05-14 04:00:00+00:00,11093.2500000000)
  (2019-05-14 14:00:00+00:00,11255.7500000000)
  (2019-05-15 00:00:00+00:00,8720.0000000000)
  (2019-05-15 10:00:00+00:00,12658.0000000000)
  (2019-05-15 20:00:00+00:00,11095.0000000000)
  (2019-05-16 06:00:00+00:00,13831.7500000000)
  (2019-05-16 16:00:00+00:00,13350.2500000000)
  (2019-05-17 02:00:00+00:00,9944.2500000000)
  (2019-05-17 12:00:00+00:00,12003.2500000000)
  (2019-05-17 22:00:00+00:00,8696.0000000000)
  (2019-05-18 08:00:00+00:00,10482.7500000000)
  (2019-05-18 18:00:00+00:00,10038.5000000000)
  (2019-05-19 04:00:00+00:00,7774.7500000000)
  (2019-05-19 14:00:00+00:00,9101.2500000000)
  (2019-05-20 00:00:00+00:00,8023.5000000000)
  (2019-05-20 10:00:00+00:00,12679.0000000000)
  (2019-05-20 20:00:00+00:00,10537.2500000000)
  (2019-05-21 06:00:00+00:00,12502.2500000000)
  (2019-05-21 16:00:00+00:00,11690.0000000000)
  (2019-05-22 02:00:00+00:00,8438.0000000000)
  (2019-05-22 12:00:00+00:00,11954.2500000000)
  (2019-05-22 22:00:00+00:00,9346.0000000000)
  (2019-05-23 08:00:00+00:00,12267.0000000000)
  (2019-05-23 18:00:00+00:00,12118.7500000000)
  (2019-05-24 04:00:00+00:00,10943.2500000000)
  (2019-05-24 14:00:00+00:00,11284.2500000000)
  (2019-05-25 00:00:00+00:00,8262.2500000000)
  (2019-05-25 10:00:00+00:00,10753.5000000000)
  (2019-05-25 20:00:00+00:00,9657.2500000000)
  (2019-05-26 06:00:00+00:00,9017.5000000000)
  (2019-05-26 16:00:00+00:00,9644.7500000000)
  (2019-05-27 02:00:00+00:00,8541.7500000000)
  (2019-05-27 12:00:00+00:00,12513.5000000000)
  (2019-05-27 22:00:00+00:00,9297.0000000000)
  (2019-05-28 08:00:00+00:00,12850.2500000000)
  (2019-05-28 18:00:00+00:00,11701.2500000000)
  (2019-05-29 04:00:00+00:00,11077.0000000000)
  (2019-05-29 14:00:00+00:00,11049.5000000000)
  (2019-05-30 00:00:00+00:00,7989.7500000000)
  (2019-05-30 10:00:00+00:00,9022.7500000000)
  (2019-05-30 20:00:00+00:00,9171.2500000000)
  (2019-05-31 06:00:00+00:00,10659.7500000000)
  (2019-05-31 16:00:00+00:00,10472.2500000000)
  (2019-06-01 02:00:00+00:00,7704.2500000000)
  (2019-06-01 12:00:00+00:00,9552.2500000000)
  (2019-06-01 22:00:00+00:00,8189.7500000000)
  (2019-06-02 08:00:00+00:00,9150.7500000000)
  (2019-06-02 18:00:00+00:00,9813.0000000000)
  (2019-06-03 04:00:00+00:00,10328.5000000000)
  (2019-06-03 14:00:00+00:00,11561.0000000000)
  (2019-06-04 00:00:00+00:00,8999.7500000000)
  (2019-06-04 10:00:00+00:00,12490.7500000000)
  (2019-06-04 20:00:00+00:00,10416.5000000000)
  (2019-06-05 06:00:00+00:00,12479.0000000000)
  (2019-06-05 16:00:00+00:00,12492.0000000000)
  (2019-06-06 02:00:00+00:00,9272.5000000000)
  (2019-06-06 12:00:00+00:00,11744.7500000000)
  (2019-06-06 22:00:00+00:00,9536.0000000000)
  (2019-06-07 08:00:00+00:00,12101.0000000000)
  (2019-06-07 18:00:00+00:00,10650.2500000000)
  (2019-06-08 04:00:00+00:00,7719.7500000000)
  (2019-06-08 14:00:00+00:00,8303.0000000000)
  (2019-06-09 00:00:00+00:00,7128.2500000000)
  (2019-06-09 10:00:00+00:00,8671.0000000000)
  (2019-06-09 20:00:00+00:00,7924.5000000000)
  (2019-06-10 06:00:00+00:00,8591.7500000000)
  (2019-06-10 16:00:00+00:00,9873.5000000000)
  (2019-06-11 02:00:00+00:00,8531.5000000000)
  (2019-06-11 12:00:00+00:00,13050.2500000000)
  (2019-06-11 22:00:00+00:00,9313.2500000000)
  (2019-06-12 08:00:00+00:00,13649.0000000000)
  (2019-06-12 18:00:00+00:00,12180.5000000000)
  (2019-06-13 04:00:00+00:00,10880.7500000000)
  (2019-06-13 14:00:00+00:00,11509.5000000000)
  (2019-06-14 00:00:00+00:00,9086.2500000000)
  (2019-06-14 10:00:00+00:00,11819.0000000000)
  (2019-06-14 20:00:00+00:00,10846.2500000000)
  (2019-06-15 06:00:00+00:00,10422.7500000000)
  (2019-06-15 16:00:00+00:00,10122.0000000000)
  (2019-06-16 02:00:00+00:00,7735.2500000000)
  (2019-06-16 12:00:00+00:00,9186.7500000000)
  (2019-06-16 22:00:00+00:00,8728.2500000000)
  (2019-06-17 08:00:00+00:00,12300.0000000000)
  (2019-06-17 18:00:00+00:00,11770.5000000000)
  (2019-06-18 04:00:00+00:00,10856.7500000000)
  (2019-06-18 14:00:00+00:00,11985.5000000000)
  (2019-06-19 00:00:00+00:00,8810.2500000000)
  (2019-06-19 10:00:00+00:00,12733.5000000000)
  (2019-06-19 20:00:00+00:00,10821.5000000000)
  (2019-06-20 06:00:00+00:00,12259.0000000000)
  (2019-06-20 16:00:00+00:00,12193.0000000000)
  (2019-06-21 02:00:00+00:00,8820.0000000000)
  (2019-06-21 12:00:00+00:00,11720.2500000000)
  (2019-06-21 22:00:00+00:00,8783.5000000000)
  (2019-06-22 08:00:00+00:00,10080.2500000000)
  (2019-06-22 18:00:00+00:00,9096.2500000000)
  (2019-06-23 04:00:00+00:00,7021.7500000000)
  (2019-06-23 14:00:00+00:00,8367.7500000000)
  (2019-06-24 00:00:00+00:00,7674.0000000000)
  (2019-06-24 10:00:00+00:00,11928.2500000000)
  (2019-06-24 20:00:00+00:00,10641.0000000000)
  (2019-06-25 06:00:00+00:00,12221.7500000000)
  (2019-06-25 16:00:00+00:00,12757.7500000000)
  (2019-06-26 02:00:00+00:00,8947.0000000000)
  (2019-06-26 12:00:00+00:00,12484.2500000000)
  (2019-06-26 22:00:00+00:00,9640.7500000000)
  (2019-06-27 08:00:00+00:00,12301.7500000000)
  (2019-06-27 18:00:00+00:00,11021.0000000000)
  (2019-06-28 04:00:00+00:00,10747.0000000000)
  (2019-06-28 14:00:00+00:00,11281.7500000000)
  (2019-06-29 00:00:00+00:00,8191.7500000000)
  (2019-06-29 10:00:00+00:00,10378.2500000000)
  (2019-06-29 20:00:00+00:00,9458.0000000000)
  (2019-06-30 06:00:00+00:00,8490.7500000000)
  (2019-06-30 16:00:00+00:00,9820.7500000000)
  (2019-07-01 02:00:00+00:00,9508.2500000000)
  (2019-07-01 12:00:00+00:00,14268.0000000000)
  (2019-07-01 22:00:00+00:00,11039.2500000000)
  (2019-07-02 08:00:00+00:00,14321.7500000000)
  (2019-07-02 18:00:00+00:00,13061.5000000000)
  (2019-07-03 04:00:00+00:00,11979.5000000000)
  (2019-07-03 14:00:00+00:00,13264.2500000000)
  (2019-07-04 00:00:00+00:00,10003.7500000000)
  (2019-07-04 10:00:00+00:00,14087.2500000000)
  (2019-07-04 20:00:00+00:00,11985.2500000000)
  (2019-07-05 06:00:00+00:00,13670.5000000000)
  (2019-07-05 16:00:00+00:00,12247.7500000000)
  (2019-07-06 02:00:00+00:00,8382.0000000000)
  (2019-07-06 12:00:00+00:00,11244.5000000000)
  (2019-07-06 22:00:00+00:00,9568.0000000000)
  (2019-07-07 08:00:00+00:00,11125.2500000000)
  (2019-07-07 18:00:00+00:00,10458.7500000000)
  (2019-07-08 04:00:00+00:00,11069.2500000000)
  (2019-07-08 14:00:00+00:00,12581.5000000000)
  (2019-07-09 00:00:00+00:00,9437.0000000000)
  (2019-07-09 10:00:00+00:00,13262.5000000000)
  (2019-07-09 20:00:00+00:00,10672.7500000000)
  (2019-07-10 06:00:00+00:00,12352.2500000000)
  (2019-07-10 16:00:00+00:00,12782.0000000000)
  (2019-07-11 02:00:00+00:00,8996.7500000000)
  (2019-07-11 12:00:00+00:00,12398.7500000000)
  (2019-07-11 22:00:00+00:00,9039.2500000000)
  (2019-07-12 08:00:00+00:00,12636.2500000000)
  (2019-07-12 18:00:00+00:00,10719.7500000000)
  (2019-07-13 04:00:00+00:00,8246.0000000000)
  (2019-07-13 14:00:00+00:00,9982.2500000000)
  (2019-07-14 00:00:00+00:00,7736.2500000000)
  (2019-07-14 10:00:00+00:00,9948.0000000000)
  (2019-07-14 20:00:00+00:00,9268.7500000000)
  (2019-07-15 06:00:00+00:00,12247.2500000000)
  (2019-07-15 16:00:00+00:00,12837.0000000000)
  (2019-07-16 02:00:00+00:00,9202.0000000000)
  (2019-07-16 12:00:00+00:00,12686.0000000000)
  (2019-07-16 22:00:00+00:00,10035.5000000000)
  (2019-07-17 08:00:00+00:00,12603.7500000000)
  (2019-07-17 18:00:00+00:00,11487.2500000000)
  (2019-07-18 04:00:00+00:00,10523.2500000000)
  (2019-07-18 14:00:00+00:00,11835.7500000000)
  (2019-07-19 00:00:00+00:00,8994.7500000000)
  (2019-07-19 10:00:00+00:00,13073.2500000000)
  (2019-07-19 20:00:00+00:00,10068.2500000000)
  (2019-07-20 06:00:00+00:00,9922.2500000000)
  (2019-07-20 16:00:00+00:00,11131.7500000000)
  (2019-07-21 02:00:00+00:00,7627.7500000000)
  (2019-07-21 12:00:00+00:00,10598.2500000000)
  (2019-07-21 22:00:00+00:00,9254.2500000000)
  (2019-07-22 08:00:00+00:00,12728.5000000000)
  (2019-07-22 18:00:00+00:00,11660.0000000000)
  (2019-07-23 04:00:00+00:00,10921.5000000000)
  (2019-07-23 14:00:00+00:00,12583.2500000000)
  (2019-07-24 00:00:00+00:00,8940.7500000000)
  (2019-07-24 10:00:00+00:00,13350.2500000000)
  (2019-07-24 20:00:00+00:00,11513.2500000000)
  (2019-07-25 06:00:00+00:00,12728.5000000000)
  (2019-07-25 16:00:00+00:00,12820.5000000000)
  (2019-07-26 02:00:00+00:00,9617.7500000000)
  (2019-07-26 12:00:00+00:00,13055.2500000000)
  (2019-07-26 22:00:00+00:00,10151.5000000000)
  (2019-07-27 08:00:00+00:00,11868.5000000000)
  (2019-07-27 18:00:00+00:00,11190.0000000000)
  (2019-07-28 04:00:00+00:00,8837.0000000000)
  (2019-07-28 14:00:00+00:00,10352.5000000000)
  (2019-07-29 00:00:00+00:00,8643.0000000000)
  (2019-07-29 10:00:00+00:00,13270.7500000000)
  (2019-07-29 20:00:00+00:00,10687.0000000000)
  (2019-07-30 06:00:00+00:00,12361.5000000000)
  (2019-07-30 16:00:00+00:00,12913.2500000000)
  (2019-07-31 02:00:00+00:00,9831.7500000000)
  (2019-07-31 12:00:00+00:00,12916.5000000000)
  (2019-07-31 22:00:00+00:00,9448.5000000000)
  (2019-08-01 08:00:00+00:00,12549.5000000000)
  (2019-08-01 18:00:00+00:00,11387.0000000000)
  (2019-08-02 04:00:00+00:00,10539.2500000000)
  (2019-08-02 14:00:00+00:00,11766.0000000000)
  (2019-08-03 00:00:00+00:00,8046.0000000000)
  (2019-08-03 10:00:00+00:00,10842.5000000000)
  (2019-08-03 20:00:00+00:00,9089.0000000000)
  (2019-08-04 06:00:00+00:00,8673.0000000000)
  (2019-08-04 16:00:00+00:00,9877.5000000000)
  (2019-08-05 02:00:00+00:00,8288.0000000000)
  (2019-08-05 12:00:00+00:00,12418.7500000000)
  (2019-08-05 22:00:00+00:00,9292.2500000000)
  (2019-08-06 08:00:00+00:00,13321.7500000000)
  (2019-08-06 18:00:00+00:00,11692.0000000000)
  (2019-08-07 04:00:00+00:00,10817.0000000000)
  (2019-08-07 14:00:00+00:00,12389.0000000000)
  (2019-08-08 00:00:00+00:00,9031.2500000000)
  (2019-08-08 10:00:00+00:00,13575.5000000000)
  (2019-08-08 20:00:00+00:00,11580.0000000000)
  (2019-08-09 06:00:00+00:00,12992.0000000000)
  (2019-08-09 16:00:00+00:00,12056.0000000000)
  (2019-08-10 02:00:00+00:00,9280.0000000000)
  (2019-08-10 12:00:00+00:00,11646.2500000000)
  (2019-08-10 22:00:00+00:00,9553.0000000000)
  (2019-08-11 08:00:00+00:00,11375.5000000000)
  (2019-08-11 18:00:00+00:00,10020.0000000000)
  (2019-08-12 04:00:00+00:00,10802.0000000000)
  (2019-08-12 14:00:00+00:00,12441.2500000000)
  (2019-08-13 00:00:00+00:00,8838.0000000000)
  (2019-08-13 10:00:00+00:00,13909.2500000000)
  (2019-08-13 20:00:00+00:00,11333.2500000000)
  (2019-08-14 06:00:00+00:00,12561.0000000000)
  (2019-08-14 16:00:00+00:00,12637.0000000000)
  (2019-08-15 02:00:00+00:00,9334.5000000000)
  (2019-08-15 12:00:00+00:00,13198.2500000000)
  (2019-08-15 22:00:00+00:00,9815.2500000000)
  (2019-08-16 08:00:00+00:00,13388.0000000000)
  (2019-08-16 18:00:00+00:00,11688.7500000000)
  (2019-08-17 04:00:00+00:00,9423.2500000000)
  (2019-08-17 14:00:00+00:00,10974.0000000000)
  (2019-08-18 00:00:00+00:00,7983.0000000000)
  (2019-08-18 10:00:00+00:00,10443.0000000000)
  (2019-08-18 20:00:00+00:00,10442.7500000000)
  (2019-08-19 06:00:00+00:00,12636.5000000000)
  (2019-08-19 16:00:00+00:00,12604.7500000000)
  (2019-08-20 02:00:00+00:00,9295.2500000000)
  (2019-08-20 12:00:00+00:00,12677.2500000000)
  (2019-08-20 22:00:00+00:00,9308.5000000000)
  (2019-08-21 08:00:00+00:00,13009.5000000000)
  (2019-08-21 18:00:00+00:00,12314.0000000000)
  (2019-08-22 04:00:00+00:00,11042.5000000000)
  (2019-08-22 14:00:00+00:00,12150.0000000000)
  (2019-08-23 00:00:00+00:00,9429.2500000000)
  (2019-08-23 10:00:00+00:00,12878.0000000000)
  (2019-08-23 20:00:00+00:00,10142.0000000000)
  (2019-08-24 06:00:00+00:00,10154.2500000000)
  (2019-08-24 16:00:00+00:00,10625.7500000000)
  (2019-08-25 02:00:00+00:00,7984.0000000000)
  (2019-08-25 12:00:00+00:00,10076.7500000000)
  (2019-08-25 22:00:00+00:00,9146.5000000000)
  (2019-08-26 08:00:00+00:00,13224.2500000000)
  (2019-08-26 18:00:00+00:00,12851.0000000000)
  (2019-08-27 04:00:00+00:00,12013.7500000000)
  (2019-08-27 14:00:00+00:00,13971.2500000000)
  (2019-08-28 00:00:00+00:00,9687.2500000000)
  (2019-08-28 10:00:00+00:00,14327.7500000000)
  (2019-08-28 20:00:00+00:00,11844.7500000000)
  (2019-08-29 06:00:00+00:00,13036.0000000000)
  (2019-08-29 16:00:00+00:00,13075.7500000000)
  (2019-08-30 02:00:00+00:00,9472.5000000000)
  (2019-08-30 12:00:00+00:00,12944.2500000000)
  (2019-08-30 22:00:00+00:00,9262.2500000000)
  (2019-08-31 08:00:00+00:00,11450.5000000000)
  (2019-08-31 18:00:00+00:00,11496.0000000000)
  (2019-09-01 04:00:00+00:00,8860.2500000000)
  (2019-09-01 14:00:00+00:00,10336.0000000000)
  (2019-09-02 00:00:00+00:00,8685.5000000000)
  (2019-09-02 10:00:00+00:00,13565.0000000000)
  (2019-09-02 20:00:00+00:00,10899.7500000000)
  (2019-09-03 06:00:00+00:00,13797.5000000000)
  (2019-09-03 16:00:00+00:00,13419.2500000000)
  (2019-09-04 02:00:00+00:00,9504.0000000000)
  (2019-09-04 12:00:00+00:00,13154.5000000000)
  (2019-09-04 22:00:00+00:00,10157.5000000000)
  (2019-09-05 08:00:00+00:00,14158.5000000000)
  (2019-09-05 18:00:00+00:00,13884.5000000000)
  (2019-09-06 04:00:00+00:00,12532.5000000000)
  (2019-09-06 14:00:00+00:00,12404.5000000000)
  (2019-09-07 00:00:00+00:00,9329.5000000000)
  (2019-09-07 10:00:00+00:00,11008.7500000000)
  (2019-09-07 20:00:00+00:00,9093.2500000000)
  (2019-09-08 06:00:00+00:00,9147.0000000000)
  (2019-09-08 16:00:00+00:00,10243.7500000000)
  (2019-09-09 02:00:00+00:00,8793.5000000000)
  (2019-09-09 12:00:00+00:00,13015.0000000000)
  (2019-09-09 22:00:00+00:00,9558.7500000000)
  (2019-09-10 08:00:00+00:00,14607.5000000000)
  (2019-09-10 18:00:00+00:00,13630.5000000000)
  (2019-09-11 04:00:00+00:00,12365.0000000000)
  (2019-09-11 14:00:00+00:00,13069.7500000000)
  (2019-09-12 00:00:00+00:00,9992.7500000000)
  (2019-09-12 10:00:00+00:00,13780.5000000000)
  (2019-09-12 20:00:00+00:00,11391.7500000000)
  (2019-09-13 06:00:00+00:00,14740.5000000000)
  (2019-09-13 16:00:00+00:00,13428.7500000000)
  (2019-09-14 02:00:00+00:00,9497.0000000000)
  (2019-09-14 12:00:00+00:00,10944.0000000000)
  (2019-09-14 22:00:00+00:00,9392.2500000000)
  (2019-09-15 08:00:00+00:00,11921.7500000000)
  (2019-09-15 18:00:00+00:00,12134.5000000000)
  (2019-09-16 04:00:00+00:00,12758.5000000000)
  (2019-09-16 14:00:00+00:00,13501.0000000000)
  (2019-09-17 00:00:00+00:00,10487.7500000000)
  (2019-09-17 10:00:00+00:00,15117.2500000000)
  (2019-09-17 20:00:00+00:00,13225.7500000000)
  (2019-09-18 06:00:00+00:00,14663.7500000000)
  (2019-09-18 16:00:00+00:00,13796.5000000000)
  (2019-09-19 02:00:00+00:00,10250.0000000000)
  (2019-09-19 12:00:00+00:00,12996.7500000000)
  (2019-09-19 22:00:00+00:00,9676.7500000000)
  (2019-09-20 08:00:00+00:00,13001.7500000000)
  (2019-09-20 18:00:00+00:00,12606.5000000000)
  (2019-09-21 04:00:00+00:00,9635.7500000000)
  (2019-09-21 14:00:00+00:00,10119.2500000000)
  (2019-09-22 00:00:00+00:00,7895.5000000000)
  (2019-09-22 10:00:00+00:00,10439.0000000000)
  (2019-09-22 20:00:00+00:00,10981.5000000000)
  (2019-09-23 06:00:00+00:00,13885.5000000000)
  (2019-09-23 16:00:00+00:00,13193.0000000000)
  (2019-09-24 02:00:00+00:00,9484.5000000000)
  (2019-09-24 12:00:00+00:00,12648.0000000000)
  (2019-09-24 22:00:00+00:00,9630.0000000000)
  (2019-09-25 08:00:00+00:00,13584.2500000000)
  (2019-09-25 18:00:00+00:00,12972.0000000000)
  (2019-09-26 04:00:00+00:00,12036.2500000000)
  (2019-09-26 14:00:00+00:00,12371.5000000000)
  (2019-09-27 00:00:00+00:00,9991.7500000000)
  (2019-09-27 10:00:00+00:00,13965.2500000000)
  (2019-09-27 20:00:00+00:00,10743.7500000000)
  (2019-09-28 06:00:00+00:00,11587.7500000000)
  (2019-09-28 16:00:00+00:00,11546.0000000000)
  (2019-09-29 02:00:00+00:00,9671.0000000000)
  (2019-09-29 12:00:00+00:00,11381.0000000000)
  (2019-09-29 22:00:00+00:00,9949.5000000000)
  (2019-09-30 08:00:00+00:00,15272.2500000000)
  (2019-09-30 18:00:00+00:00,14423.2500000000)
  (2019-10-01 04:00:00+00:00,12577.7500000000)
  (2019-10-01 14:00:00+00:00,13209.7500000000)
  (2019-10-02 00:00:00+00:00,9376.5000000000)
  (2019-10-02 10:00:00+00:00,14044.7500000000)
  (2019-10-02 20:00:00+00:00,11405.7500000000)
  (2019-10-03 06:00:00+00:00,10754.0000000000)
  (2019-10-03 16:00:00+00:00,10566.2500000000)
  (2019-10-04 02:00:00+00:00,8495.0000000000)
  (2019-10-04 12:00:00+00:00,11500.0000000000)
  (2019-10-04 22:00:00+00:00,10035.0000000000)
  (2019-10-05 08:00:00+00:00,12675.5000000000)
  (2019-10-05 18:00:00+00:00,11461.5000000000)
  (2019-10-06 04:00:00+00:00,8473.2500000000)
  (2019-10-06 14:00:00+00:00,9912.2500000000)
  (2019-10-07 00:00:00+00:00,8706.0000000000)
  (2019-10-07 10:00:00+00:00,13099.5000000000)
  (2019-10-07 20:00:00+00:00,11214.7500000000)
  (2019-10-08 06:00:00+00:00,14443.5000000000)
  (2019-10-08 16:00:00+00:00,14321.5000000000)
  (2019-10-09 02:00:00+00:00,11062.5000000000)
  (2019-10-09 12:00:00+00:00,13757.7500000000)
  (2019-10-09 22:00:00+00:00,10352.5000000000)
  (2019-10-10 08:00:00+00:00,14616.7500000000)
  (2019-10-10 18:00:00+00:00,14538.7500000000)
  (2019-10-11 04:00:00+00:00,13932.5000000000)
  (2019-10-11 14:00:00+00:00,14289.7500000000)
  (2019-10-12 00:00:00+00:00,11021.7500000000)
  (2019-10-12 10:00:00+00:00,12899.2500000000)
  (2019-10-12 20:00:00+00:00,10647.2500000000)
  (2019-10-13 06:00:00+00:00,9723.0000000000)
  (2019-10-13 16:00:00+00:00,12112.0000000000)
  (2019-10-14 02:00:00+00:00,10762.0000000000)
  (2019-10-14 12:00:00+00:00,12903.0000000000)
  (2019-10-14 22:00:00+00:00,10266.0000000000)
  (2019-10-15 08:00:00+00:00,13661.7500000000)
  (2019-10-15 18:00:00+00:00,12543.7500000000)
  (2019-10-16 04:00:00+00:00,12238.0000000000)
  (2019-10-16 14:00:00+00:00,13230.2500000000)
  (2019-10-17 00:00:00+00:00,9811.5000000000)
  (2019-10-17 10:00:00+00:00,13889.5000000000)
  (2019-10-17 20:00:00+00:00,11233.7500000000)
  (2019-10-18 06:00:00+00:00,13664.5000000000)
  (2019-10-18 16:00:00+00:00,14695.5000000000)
  (2019-10-19 02:00:00+00:00,9967.2500000000)
  (2019-10-19 12:00:00+00:00,11403.7500000000)
  (2019-10-19 22:00:00+00:00,9571.7500000000)
  (2019-10-20 08:00:00+00:00,10433.7500000000)
  (2019-10-20 18:00:00+00:00,10736.5000000000)
  (2019-10-21 04:00:00+00:00,12095.7500000000)
  (2019-10-21 14:00:00+00:00,12441.2500000000)
  (2019-10-22 00:00:00+00:00,9119.7500000000)
  (2019-10-22 10:00:00+00:00,13450.0000000000)
  (2019-10-22 20:00:00+00:00,11684.7500000000)
  (2019-10-23 06:00:00+00:00,13368.0000000000)
  (2019-10-23 16:00:00+00:00,13610.0000000000)
  (2019-10-24 02:00:00+00:00,10051.2500000000)
  (2019-10-24 12:00:00+00:00,13182.5000000000)
  (2019-10-24 22:00:00+00:00,10444.5000000000)
  (2019-10-25 08:00:00+00:00,14587.0000000000)
  (2019-10-25 18:00:00+00:00,13647.7500000000)
  (2019-10-26 04:00:00+00:00,11632.0000000000)
  (2019-10-26 14:00:00+00:00,12343.2500000000)
  (2019-10-27 00:00:00+00:00,9878.0000000000)
  (2019-10-27 10:00:00+00:00,13019.2500000000)
  (2019-10-27 20:00:00+00:00,11921.0000000000)
  (2019-10-28 06:00:00+00:00,14842.5000000000)
  (2019-10-28 16:00:00+00:00,14342.0000000000)
  (2019-10-29 02:00:00+00:00,10112.5000000000)
  (2019-10-29 12:00:00+00:00,13735.7500000000)
  (2019-10-29 22:00:00+00:00,10844.5000000000)
  (2019-10-30 08:00:00+00:00,13773.5000000000)
  (2019-10-30 18:00:00+00:00,13712.0000000000)
  (2019-10-31 04:00:00+00:00,9895.5000000000)
  (2019-10-31 14:00:00+00:00,10924.2500000000)
  (2019-11-01 00:00:00+00:00,9303.2500000000)
  (2019-11-01 10:00:00+00:00,13259.0000000000)
  (2019-11-01 20:00:00+00:00,12202.5000000000)
  (2019-11-02 06:00:00+00:00,10462.0000000000)
  (2019-11-02 16:00:00+00:00,12598.7500000000)
  (2019-11-03 02:00:00+00:00,8457.5000000000)
  (2019-11-03 12:00:00+00:00,10316.5000000000)
  (2019-11-03 22:00:00+00:00,9127.7500000000)
  (2019-11-04 08:00:00+00:00,13693.7500000000)
  (2019-11-04 18:00:00+00:00,14211.0000000000)
  (2019-11-05 04:00:00+00:00,11302.5000000000)
  (2019-11-05 14:00:00+00:00,13975.2500000000)
  (2019-11-06 00:00:00+00:00,10166.5000000000)
  (2019-11-06 10:00:00+00:00,14105.7500000000)
  (2019-11-06 20:00:00+00:00,12214.5000000000)
  (2019-11-07 06:00:00+00:00,13499.7500000000)
  (2019-11-07 16:00:00+00:00,14610.7500000000)
  (2019-11-08 02:00:00+00:00,10413.0000000000)
  (2019-11-08 12:00:00+00:00,13131.2500000000)
  (2019-11-08 22:00:00+00:00,10330.2500000000)
  (2019-11-09 08:00:00+00:00,11663.0000000000)
  (2019-11-09 18:00:00+00:00,12421.2500000000)
  (2019-11-10 04:00:00+00:00,8887.0000000000)
  (2019-11-10 14:00:00+00:00,10172.7500000000)
  (2019-11-11 00:00:00+00:00,9482.7500000000)
  (2019-11-11 10:00:00+00:00,14505.5000000000)
  (2019-11-11 20:00:00+00:00,13652.7500000000)
  (2019-11-12 06:00:00+00:00,14290.0000000000)
  (2019-11-12 16:00:00+00:00,15216.5000000000)
  (2019-11-13 02:00:00+00:00,10132.2500000000)
  (2019-11-13 12:00:00+00:00,13712.2500000000)
  (2019-11-13 22:00:00+00:00,10576.2500000000)
  (2019-11-14 08:00:00+00:00,13589.0000000000)
  (2019-11-14 18:00:00+00:00,14602.7500000000)
  (2019-11-15 04:00:00+00:00,12783.0000000000)
  (2019-11-15 14:00:00+00:00,14537.7500000000)
  (2019-11-16 00:00:00+00:00,10276.2500000000)
  (2019-11-16 10:00:00+00:00,13494.2500000000)
  (2019-11-16 20:00:00+00:00,12118.2500000000)
  (2019-11-17 06:00:00+00:00,9538.7500000000)
  (2019-11-17 16:00:00+00:00,13087.5000000000)
  (2019-11-18 02:00:00+00:00,10693.7500000000)
  (2019-11-18 12:00:00+00:00,14608.0000000000)
  (2019-11-18 22:00:00+00:00,12196.2500000000)
  (2019-11-19 08:00:00+00:00,15836.7500000000)
  (2019-11-19 18:00:00+00:00,15123.7500000000)
  (2019-11-20 04:00:00+00:00,10980.2500000000)
  (2019-11-20 14:00:00+00:00,13725.2500000000)
  (2019-11-21 00:00:00+00:00,9971.0000000000)
  (2019-11-21 10:00:00+00:00,14881.5000000000)
  (2019-11-21 20:00:00+00:00,13496.2500000000)
  (2019-11-22 06:00:00+00:00,14711.7500000000)
  (2019-11-22 16:00:00+00:00,14969.2500000000)
  (2019-11-23 02:00:00+00:00,10338.0000000000)
  (2019-11-23 12:00:00+00:00,12744.0000000000)
  (2019-11-23 22:00:00+00:00,11040.0000000000)
  (2019-11-24 08:00:00+00:00,11619.2500000000)
  (2019-11-24 18:00:00+00:00,12324.5000000000)
  (2019-11-25 04:00:00+00:00,10848.2500000000)
  (2019-11-25 14:00:00+00:00,14285.5000000000)
  (2019-11-26 00:00:00+00:00,9990.2500000000)
  (2019-11-26 10:00:00+00:00,14670.5000000000)
  (2019-11-26 20:00:00+00:00,13081.0000000000)
  (2019-11-27 06:00:00+00:00,14263.5000000000)
  (2019-11-27 16:00:00+00:00,15723.7500000000)
  (2019-11-28 02:00:00+00:00,11394.7500000000)
  (2019-11-28 12:00:00+00:00,15713.5000000000)
  (2019-11-28 22:00:00+00:00,12322.5000000000)
  (2019-11-29 08:00:00+00:00,15564.5000000000)
  (2019-11-29 18:00:00+00:00,15215.2500000000)
  (2019-11-30 04:00:00+00:00,10297.0000000000)
  (2019-11-30 14:00:00+00:00,11900.2500000000)
  (2019-12-01 00:00:00+00:00,9275.0000000000)
  (2019-12-01 10:00:00+00:00,12191.5000000000)
  (2019-12-01 20:00:00+00:00,11085.0000000000)
  (2019-12-02 06:00:00+00:00,14261.2500000000)
  (2019-12-02 16:00:00+00:00,15684.5000000000)
  (2019-12-03 02:00:00+00:00,10392.7500000000)
  (2019-12-03 12:00:00+00:00,14319.0000000000)
  (2019-12-03 22:00:00+00:00,12058.7500000000)
  (2019-12-04 08:00:00+00:00,14809.5000000000)
  (2019-12-04 18:00:00+00:00,15083.2500000000)
  (2019-12-05 04:00:00+00:00,11994.2500000000)
  (2019-12-05 14:00:00+00:00,14468.7500000000)
  (2019-12-06 00:00:00+00:00,11999.7500000000)
  (2019-12-06 10:00:00+00:00,15792.5000000000)
  (2019-12-06 20:00:00+00:00,13621.5000000000)
  (2019-12-07 06:00:00+00:00,12414.2500000000)
  (2019-12-07 16:00:00+00:00,13785.7500000000)
  (2019-12-08 02:00:00+00:00,9804.0000000000)
  (2019-12-08 12:00:00+00:00,12975.2500000000)
  (2019-12-08 22:00:00+00:00,10943.7500000000)
  (2019-12-09 08:00:00+00:00,14868.0000000000)
  (2019-12-09 18:00:00+00:00,14955.5000000000)
  (2019-12-10 04:00:00+00:00,12552.7500000000)
  (2019-12-10 14:00:00+00:00,13828.7500000000)
  (2019-12-11 00:00:00+00:00,11840.7500000000)
  (2019-12-11 10:00:00+00:00,15404.5000000000)
  (2019-12-11 20:00:00+00:00,13001.2500000000)
  (2019-12-12 06:00:00+00:00,14135.0000000000)
  (2019-12-12 16:00:00+00:00,15197.0000000000)
  (2019-12-13 02:00:00+00:00,11567.5000000000)
  (2019-12-13 12:00:00+00:00,14656.7500000000)
  (2019-12-13 22:00:00+00:00,10995.2500000000)
  (2019-12-14 08:00:00+00:00,12558.0000000000)
  (2019-12-14 18:00:00+00:00,13825.5000000000)
  (2019-12-15 04:00:00+00:00,10113.2500000000)
  (2019-12-15 14:00:00+00:00,12398.5000000000)
  (2019-12-16 00:00:00+00:00,10599.2500000000)
  (2019-12-16 10:00:00+00:00,14684.7500000000)
  (2019-12-16 20:00:00+00:00,12413.2500000000)
  (2019-12-17 06:00:00+00:00,14002.7500000000)
  (2019-12-17 16:00:00+00:00,15083.5000000000)
  (2019-12-18 02:00:00+00:00,10496.5000000000)
  (2019-12-18 12:00:00+00:00,13753.5000000000)
  (2019-12-18 22:00:00+00:00,10754.7500000000)
  (2019-12-19 08:00:00+00:00,14203.2500000000)
  (2019-12-19 18:00:00+00:00,13921.2500000000)
  (2019-12-20 04:00:00+00:00,10537.0000000000)
  (2019-12-20 14:00:00+00:00,12408.2500000000)
  (2019-12-21 00:00:00+00:00,9100.5000000000)
  (2019-12-21 10:00:00+00:00,11929.5000000000)
  (2019-12-21 20:00:00+00:00,10539.7500000000)
  (2019-12-22 06:00:00+00:00,9270.5000000000)
  (2019-12-22 16:00:00+00:00,11858.0000000000)
  (2019-12-23 02:00:00+00:00,8513.7500000000)
  (2019-12-23 12:00:00+00:00,11986.0000000000)
  (2019-12-23 22:00:00+00:00,9601.7500000000)
  (2019-12-24 08:00:00+00:00,10957.0000000000)
  (2019-12-24 18:00:00+00:00,9866.5000000000)
  (2019-12-25 04:00:00+00:00,7439.7500000000)
  (2019-12-25 14:00:00+00:00,9455.5000000000)
  (2019-12-26 00:00:00+00:00,8022.2500000000)
  (2019-12-26 10:00:00+00:00,10799.5000000000)
  (2019-12-26 20:00:00+00:00,9495.5000000000)
  (2019-12-27 06:00:00+00:00,9686.7500000000)
  (2019-12-27 16:00:00+00:00,11871.0000000000)
  (2019-12-28 02:00:00+00:00,8245.7500000000)
  (2019-12-28 12:00:00+00:00,10458.0000000000)
  (2019-12-28 22:00:00+00:00,9292.2500000000)
  (2019-12-29 08:00:00+00:00,9975.5000000000)
  (2019-12-29 18:00:00+00:00,11411.5000000000)
  (2019-12-30 04:00:00+00:00,9920.5000000000)
  (2019-12-30 14:00:00+00:00,12106.7500000000)
  (2019-12-31 00:00:00+00:00,9609.0000000000)
  (2019-12-31 10:00:00+00:00,12158.7500000000)
  (2019-12-31 20:00:00+00:00,9949.2500000000)
  (2020-01-01 06:00:00+00:00,8703.5000000000)
  (2020-01-01 16:00:00+00:00,11639.5000000000)
  (2020-01-02 02:00:00+00:00,9239.5000000000)
  (2020-01-02 12:00:00+00:00,12708.2500000000)
  (2020-01-02 22:00:00+00:00,11203.0000000000)
  (2020-01-03 08:00:00+00:00,13736.0000000000)
  (2020-01-03 18:00:00+00:00,13542.7500000000)
  (2020-01-04 04:00:00+00:00,10767.2500000000)
  (2020-01-04 14:00:00+00:00,13178.2500000000)
  (2020-01-05 00:00:00+00:00,10191.0000000000)
  (2020-01-05 10:00:00+00:00,11527.2500000000)
  (2020-01-05 20:00:00+00:00,11412.0000000000)
  (2020-01-06 06:00:00+00:00,13194.2500000000)
  (2020-01-06 16:00:00+00:00,14054.7500000000)
  (2020-01-07 02:00:00+00:00,10349.7500000000)
  (2020-01-07 12:00:00+00:00,14242.2500000000)
  (2020-01-07 22:00:00+00:00,12057.7500000000)
  (2020-01-08 08:00:00+00:00,15325.7500000000)
  (2020-01-08 18:00:00+00:00,14600.5000000000)
  (2020-01-09 04:00:00+00:00,10761.7500000000)
  (2020-01-09 14:00:00+00:00,14559.7500000000)
  (2020-01-10 00:00:00+00:00,10259.2500000000)
  (2020-01-10 10:00:00+00:00,14512.2500000000)
  (2020-01-10 20:00:00+00:00,12655.0000000000)
  (2020-01-11 06:00:00+00:00,10338.2500000000)
  (2020-01-11 16:00:00+00:00,13480.0000000000)
  (2020-01-12 02:00:00+00:00,10355.2500000000)
  (2020-01-12 12:00:00+00:00,12143.5000000000)
  (2020-01-12 22:00:00+00:00,11308.2500000000)
  (2020-01-13 08:00:00+00:00,14389.7500000000)
  (2020-01-13 18:00:00+00:00,14370.5000000000)
  (2020-01-14 04:00:00+00:00,12479.7500000000)
  (2020-01-14 14:00:00+00:00,14704.5000000000)
  (2020-01-15 00:00:00+00:00,11273.2500000000)
  (2020-01-15 10:00:00+00:00,15041.2500000000)
  (2020-01-15 20:00:00+00:00,13716.2500000000)
  (2020-01-16 06:00:00+00:00,14353.5000000000)
  (2020-01-16 16:00:00+00:00,14426.0000000000)
  (2020-01-17 02:00:00+00:00,11020.2500000000)
  (2020-01-17 12:00:00+00:00,13633.5000000000)
  (2020-01-17 22:00:00+00:00,11182.2500000000)
  (2020-01-18 08:00:00+00:00,12716.2500000000)
  (2020-01-18 18:00:00+00:00,12756.7500000000)
  (2020-01-19 04:00:00+00:00,9243.2500000000)
  (2020-01-19 14:00:00+00:00,10491.0000000000)
  (2020-01-20 00:00:00+00:00,9430.2500000000)
  (2020-01-20 10:00:00+00:00,14475.0000000000)
  (2020-01-20 20:00:00+00:00,13335.0000000000)
  (2020-01-21 06:00:00+00:00,15253.0000000000)
  (2020-01-21 16:00:00+00:00,15282.7500000000)
  (2020-01-22 02:00:00+00:00,11499.7500000000)
  (2020-01-22 12:00:00+00:00,15001.0000000000)
  (2020-01-22 22:00:00+00:00,11243.2500000000)
  (2020-01-23 08:00:00+00:00,14789.5000000000)
  (2020-01-23 18:00:00+00:00,15165.7500000000)
  (2020-01-24 04:00:00+00:00,12435.7500000000)
  (2020-01-24 14:00:00+00:00,14033.0000000000)
  (2020-01-25 00:00:00+00:00,10724.2500000000)
  (2020-01-25 10:00:00+00:00,13483.0000000000)
  (2020-01-25 20:00:00+00:00,11303.7500000000)
  (2020-01-26 06:00:00+00:00,9765.0000000000)
  (2020-01-26 16:00:00+00:00,12159.2500000000)
  (2020-01-27 02:00:00+00:00,10211.0000000000)
  (2020-01-27 12:00:00+00:00,14614.0000000000)
  (2020-01-27 22:00:00+00:00,11300.5000000000)
  (2020-01-28 08:00:00+00:00,14824.5000000000)
  (2020-01-28 18:00:00+00:00,15393.2500000000)
  (2020-01-29 04:00:00+00:00,13281.5000000000)
  (2020-01-29 14:00:00+00:00,15753.2500000000)
  (2020-01-30 00:00:00+00:00,11913.5000000000)
  (2020-01-30 10:00:00+00:00,16378.7500000000)
  (2020-01-30 20:00:00+00:00,13879.2500000000)
  (2020-01-31 06:00:00+00:00,15420.5000000000)
  (2020-01-31 16:00:00+00:00,14802.5000000000)
  (2020-02-01 02:00:00+00:00,11067.2500000000)
  (2020-02-01 12:00:00+00:00,13274.2500000000)
  (2020-02-01 22:00:00+00:00,11725.7500000000)
  (2020-02-02 08:00:00+00:00,12765.0000000000)
  (2020-02-02 18:00:00+00:00,11919.5000000000)
  (2020-02-03 04:00:00+00:00,11117.0000000000)
  (2020-02-03 14:00:00+00:00,14175.5000000000)
  (2020-02-04 00:00:00+00:00,10940.7500000000)
  (2020-02-04 10:00:00+00:00,16045.7500000000)
  (2020-02-04 20:00:00+00:00,14551.2500000000)
  (2020-02-05 06:00:00+00:00,14545.0000000000)
  (2020-02-05 16:00:00+00:00,15255.0000000000)
  (2020-02-06 02:00:00+00:00,12253.0000000000)
  (2020-02-06 12:00:00+00:00,15049.7500000000)
  (2020-02-06 22:00:00+00:00,11746.5000000000)
  (2020-02-07 08:00:00+00:00,13834.5000000000)
  (2020-02-07 18:00:00+00:00,14119.2500000000)
  (2020-02-08 04:00:00+00:00,10887.5000000000)
  (2020-02-08 14:00:00+00:00,11461.2500000000)
  (2020-02-09 00:00:00+00:00,10128.0000000000)
  (2020-02-09 10:00:00+00:00,13128.5000000000)
  (2020-02-09 20:00:00+00:00,12662.2500000000)
  (2020-02-10 06:00:00+00:00,15249.2500000000)
  (2020-02-10 16:00:00+00:00,16202.2500000000)
  (2020-02-11 02:00:00+00:00,12383.0000000000)
  (2020-02-11 12:00:00+00:00,15807.2500000000)
  (2020-02-11 22:00:00+00:00,12926.2500000000)
  (2020-02-12 08:00:00+00:00,15908.0000000000)
  (2020-02-12 18:00:00+00:00,15928.0000000000)
  (2020-02-13 04:00:00+00:00,13285.7500000000)
  (2020-02-13 14:00:00+00:00,14329.7500000000)
  (2020-02-14 00:00:00+00:00,10439.0000000000)
  (2020-02-14 10:00:00+00:00,14126.0000000000)
  (2020-02-14 20:00:00+00:00,11925.7500000000)
  (2020-02-15 06:00:00+00:00,11744.0000000000)
  (2020-02-15 16:00:00+00:00,12951.7500000000)
  (2020-02-16 02:00:00+00:00,10588.7500000000)
  (2020-02-16 12:00:00+00:00,12616.0000000000)
  (2020-02-16 22:00:00+00:00,11797.7500000000)
  (2020-02-17 08:00:00+00:00,15296.5000000000)
  (2020-02-17 18:00:00+00:00,15345.2500000000)
  (2020-02-18 04:00:00+00:00,12955.0000000000)
  (2020-02-18 14:00:00+00:00,15394.5000000000)
  (2020-02-19 00:00:00+00:00,11677.5000000000)
  (2020-02-19 10:00:00+00:00,15458.0000000000)
  (2020-02-19 20:00:00+00:00,13549.0000000000)
  (2020-02-20 06:00:00+00:00,13801.7500000000)
  (2020-02-20 16:00:00+00:00,15007.5000000000)
  (2020-02-21 02:00:00+00:00,12167.5000000000)
  (2020-02-21 12:00:00+00:00,15130.5000000000)
  (2020-02-21 22:00:00+00:00,12197.2500000000)
  (2020-02-22 08:00:00+00:00,13360.0000000000)
  (2020-02-22 18:00:00+00:00,13995.7500000000)
  (2020-02-23 04:00:00+00:00,10543.7500000000)
  (2020-02-23 14:00:00+00:00,12030.5000000000)
  (2020-02-24 00:00:00+00:00,11351.2500000000)
  (2020-02-24 10:00:00+00:00,14254.7500000000)
  (2020-02-24 20:00:00+00:00,13279.5000000000)
  (2020-02-25 06:00:00+00:00,14848.7500000000)
  (2020-02-25 16:00:00+00:00,14912.2500000000)
  (2020-02-26 02:00:00+00:00,11128.2500000000)
  (2020-02-26 12:00:00+00:00,14040.5000000000)
  (2020-02-26 22:00:00+00:00,11397.5000000000)
  (2020-02-27 08:00:00+00:00,14033.7500000000)
  (2020-02-27 18:00:00+00:00,13967.0000000000)
  (2020-02-28 04:00:00+00:00,12192.7500000000)
  (2020-02-28 14:00:00+00:00,14481.7500000000)
  (2020-02-29 00:00:00+00:00,10895.7500000000)
  (2020-02-29 10:00:00+00:00,13598.7500000000)
  (2020-02-29 20:00:00+00:00,12464.0000000000)
  (2020-03-01 06:00:00+00:00,11290.2500000000)
  (2020-03-01 16:00:00+00:00,11880.0000000000)
  (2020-03-02 02:00:00+00:00,9730.0000000000)
  (2020-03-02 12:00:00+00:00,13175.5000000000)
  (2020-03-02 22:00:00+00:00,10768.2500000000)
  (2020-03-03 08:00:00+00:00,13886.0000000000)
  (2020-03-03 18:00:00+00:00,14217.0000000000)
  (2020-03-04 04:00:00+00:00,12218.2500000000)
  (2020-03-04 14:00:00+00:00,13597.0000000000)
  (2020-03-05 00:00:00+00:00,10374.5000000000)
  (2020-03-05 10:00:00+00:00,13502.0000000000)
  (2020-03-05 20:00:00+00:00,13223.2500000000)
  (2020-03-06 06:00:00+00:00,14461.5000000000)
  (2020-03-06 16:00:00+00:00,14474.0000000000)
  (2020-03-07 02:00:00+00:00,10319.0000000000)
  (2020-03-07 12:00:00+00:00,12826.0000000000)
  (2020-03-07 22:00:00+00:00,10521.5000000000)
  (2020-03-08 08:00:00+00:00,11869.0000000000)
  (2020-03-08 18:00:00+00:00,13055.7500000000)
  (2020-03-09 04:00:00+00:00,10696.5000000000)
  (2020-03-09 14:00:00+00:00,13260.2500000000)
  (2020-03-10 00:00:00+00:00,10831.5000000000)
  (2020-03-10 10:00:00+00:00,16373.2500000000)
  (2020-03-10 20:00:00+00:00,14544.0000000000)
  (2020-03-11 06:00:00+00:00,14655.2500000000)
  (2020-03-11 16:00:00+00:00,13932.7500000000)
  (2020-03-12 02:00:00+00:00,11703.7500000000)
  (2020-03-12 12:00:00+00:00,15548.5000000000)
  (2020-03-12 22:00:00+00:00,12892.0000000000)
  (2020-03-13 08:00:00+00:00,16002.2500000000)
  (2020-03-13 18:00:00+00:00,15880.0000000000)
  (2020-03-14 04:00:00+00:00,9949.0000000000)
  (2020-03-14 14:00:00+00:00,10786.7500000000)
  (2020-03-15 00:00:00+00:00,11168.0000000000)
  (2020-03-15 10:00:00+00:00,12767.5000000000)
  (2020-03-15 20:00:00+00:00,12499.5000000000)
  (2020-03-16 06:00:00+00:00,13804.7500000000)
  (2020-03-16 16:00:00+00:00,12652.7500000000)
  (2020-03-17 02:00:00+00:00,9513.2500000000)
  (2020-03-17 12:00:00+00:00,13796.0000000000)
  (2020-03-17 22:00:00+00:00,11228.7500000000)
  (2020-03-18 08:00:00+00:00,14193.5000000000)
  (2020-03-18 18:00:00+00:00,14116.0000000000)
  (2020-03-19 04:00:00+00:00,10493.2500000000)
  (2020-03-19 14:00:00+00:00,11998.7500000000)
  (2020-03-20 00:00:00+00:00,9206.7500000000)
  (2020-03-20 10:00:00+00:00,13515.5000000000)
  (2020-03-20 20:00:00+00:00,12228.0000000000)
  (2020-03-21 06:00:00+00:00,10800.7500000000)
  (2020-03-21 16:00:00+00:00,11418.5000000000)
  (2020-03-22 02:00:00+00:00,9810.5000000000)
  (2020-03-22 12:00:00+00:00,10793.5000000000)
  (2020-03-22 22:00:00+00:00,10832.7500000000)
  (2020-03-23 08:00:00+00:00,12583.5000000000)
  (2020-03-23 18:00:00+00:00,13131.5000000000)
  (2020-03-24 04:00:00+00:00,10735.5000000000)
  (2020-03-24 14:00:00+00:00,11871.7500000000)
  (2020-03-25 00:00:00+00:00,10504.5000000000)
  (2020-03-25 10:00:00+00:00,13321.0000000000)
  (2020-03-25 20:00:00+00:00,12219.0000000000)
  (2020-03-26 06:00:00+00:00,12134.2500000000)
  (2020-03-26 16:00:00+00:00,12749.0000000000)
  (2020-03-27 02:00:00+00:00,10034.7500000000)
  (2020-03-27 12:00:00+00:00,11914.0000000000)
  (2020-03-27 22:00:00+00:00,10016.0000000000)
  (2020-03-28 08:00:00+00:00,10679.5000000000)
  (2020-03-28 18:00:00+00:00,12333.0000000000)
  (2020-03-29 04:00:00+00:00,9365.2500000000)
  (2020-03-29 14:00:00+00:00,11941.7500000000)
  (2020-03-30 00:00:00+00:00,10038.5000000000)
  (2020-03-30 10:00:00+00:00,13249.2500000000)
  (2020-03-30 20:00:00+00:00,11316.0000000000)
  (2020-03-31 06:00:00+00:00,12354.2500000000)
  (2020-03-31 16:00:00+00:00,12661.7500000000)
  (2020-04-01 02:00:00+00:00,11576.5000000000)
  (2020-04-01 12:00:00+00:00,13486.5000000000)
  (2020-04-01 22:00:00+00:00,11201.5000000000)
  (2020-04-02 08:00:00+00:00,13807.5000000000)
  (2020-04-02 18:00:00+00:00,14037.0000000000)
  (2020-04-03 04:00:00+00:00,13035.5000000000)
  (2020-04-03 14:00:00+00:00,13361.2500000000)
  (2020-04-04 00:00:00+00:00,9225.2500000000)
  (2020-04-04 10:00:00+00:00,11249.0000000000)
  (2020-04-04 20:00:00+00:00,9838.2500000000)
  (2020-04-05 06:00:00+00:00,9927.2500000000)
  (2020-04-05 16:00:00+00:00,11123.7500000000)
  (2020-04-06 02:00:00+00:00,10468.5000000000)
  (2020-04-06 12:00:00+00:00,12652.0000000000)
  (2020-04-06 22:00:00+00:00,9758.2500000000)
  (2020-04-07 08:00:00+00:00,13055.5000000000)
  (2020-04-07 18:00:00+00:00,11553.0000000000)
  (2020-04-08 04:00:00+00:00,10604.7500000000)
  (2020-04-08 14:00:00+00:00,10714.2500000000)
  (2020-04-09 00:00:00+00:00,9184.0000000000)
  (2020-04-09 10:00:00+00:00,13037.0000000000)
  (2020-04-09 20:00:00+00:00,10007.7500000000)
  (2020-04-10 06:00:00+00:00,8978.7500000000)
  (2020-04-10 16:00:00+00:00,9659.0000000000)
  (2020-04-11 02:00:00+00:00,7727.0000000000)
  (2020-04-11 12:00:00+00:00,9368.7500000000)
  (2020-04-11 22:00:00+00:00,8592.7500000000)
  (2020-04-12 08:00:00+00:00,10028.2500000000)
  (2020-04-12 18:00:00+00:00,10256.0000000000)
  (2020-04-13 04:00:00+00:00,8757.2500000000)
  (2020-04-13 14:00:00+00:00,10087.5000000000)
  (2020-04-14 00:00:00+00:00,9067.5000000000)
  (2020-04-14 10:00:00+00:00,13217.2500000000)
  (2020-04-14 20:00:00+00:00,11186.7500000000)
  (2020-04-15 06:00:00+00:00,12804.2500000000)
  (2020-04-15 16:00:00+00:00,12379.0000000000)
  (2020-04-16 02:00:00+00:00,10346.2500000000)
  (2020-04-16 12:00:00+00:00,11741.7500000000)
  (2020-04-16 22:00:00+00:00,9081.0000000000)
  (2020-04-17 08:00:00+00:00,11688.7500000000)
  (2020-04-17 18:00:00+00:00,10974.2500000000)
  (2020-04-18 04:00:00+00:00,8650.5000000000)
  (2020-04-18 14:00:00+00:00,9500.2500000000)
  (2020-04-19 00:00:00+00:00,8292.7500000000)
  (2020-04-19 10:00:00+00:00,11014.5000000000)
  (2020-04-19 20:00:00+00:00,9871.2500000000)
  (2020-04-20 06:00:00+00:00,11742.0000000000)
  (2020-04-20 16:00:00+00:00,12129.0000000000)
  (2020-04-21 02:00:00+00:00,9962.7500000000)
  (2020-04-21 12:00:00+00:00,12507.2500000000)
  (2020-04-21 22:00:00+00:00,9984.0000000000)
  (2020-04-22 08:00:00+00:00,12479.5000000000)
  (2020-04-22 18:00:00+00:00,11561.2500000000)
  (2020-04-23 04:00:00+00:00,10495.0000000000)
  (2020-04-23 14:00:00+00:00,11106.5000000000)
  (2020-04-24 00:00:00+00:00,8625.0000000000)
  (2020-04-24 10:00:00+00:00,12665.5000000000)
  (2020-04-24 20:00:00+00:00,10972.2500000000)
  (2020-04-25 06:00:00+00:00,10287.2500000000)
  (2020-04-25 16:00:00+00:00,10970.0000000000)
  (2020-04-26 02:00:00+00:00,8195.0000000000)
  (2020-04-26 12:00:00+00:00,9470.7500000000)
  (2020-04-26 22:00:00+00:00,8543.7500000000)
  (2020-04-27 08:00:00+00:00,11885.2500000000)
  (2020-04-27 18:00:00+00:00,11377.5000000000)
  (2020-04-28 04:00:00+00:00,10323.0000000000)
  (2020-04-28 14:00:00+00:00,11173.7500000000)
  (2020-04-29 00:00:00+00:00,8262.0000000000)
  (2020-04-29 10:00:00+00:00,13184.2500000000)
  (2020-04-29 20:00:00+00:00,10982.5000000000)
  (2020-04-30 06:00:00+00:00,12420.7500000000)
  (2020-04-30 16:00:00+00:00,11146.5000000000)
  (2020-05-01 02:00:00+00:00,8769.2500000000)
  (2020-05-01 12:00:00+00:00,9633.5000000000)
  (2020-05-01 22:00:00+00:00,8009.2500000000)
  (2020-05-02 08:00:00+00:00,10653.0000000000)
  (2020-05-02 18:00:00+00:00,10534.5000000000)
  (2020-05-03 04:00:00+00:00,8201.7500000000)
  (2020-05-03 14:00:00+00:00,9784.7500000000)
  (2020-05-04 00:00:00+00:00,8090.0000000000)
  (2020-05-04 10:00:00+00:00,12503.5000000000)
  (2020-05-04 20:00:00+00:00,11032.2500000000)
  (2020-05-05 06:00:00+00:00,12314.0000000000)
  (2020-05-05 16:00:00+00:00,11901.7500000000)
  (2020-05-06 02:00:00+00:00,9687.5000000000)
  (2020-05-06 12:00:00+00:00,13016.5000000000)
  (2020-05-06 22:00:00+00:00,10471.0000000000)
  (2020-05-07 08:00:00+00:00,12576.2500000000)
  (2020-05-07 18:00:00+00:00,11876.5000000000)
  (2020-05-08 04:00:00+00:00,10255.7500000000)
  (2020-05-08 14:00:00+00:00,11278.7500000000)
  (2020-05-09 00:00:00+00:00,7957.0000000000)
  (2020-05-09 10:00:00+00:00,10700.5000000000)
  (2020-05-09 20:00:00+00:00,9246.5000000000)
  (2020-05-10 06:00:00+00:00,9238.0000000000)
  (2020-05-10 16:00:00+00:00,9867.7500000000)
  (2020-05-11 02:00:00+00:00,9582.0000000000)
  (2020-05-11 12:00:00+00:00,13749.7500000000)
  (2020-05-11 22:00:00+00:00,10414.7500000000)
  (2020-05-12 08:00:00+00:00,13253.7500000000)
  (2020-05-12 18:00:00+00:00,11920.2500000000)
  (2020-05-13 04:00:00+00:00,10960.7500000000)
  (2020-05-13 14:00:00+00:00,11509.0000000000)
  (2020-05-14 00:00:00+00:00,8773.0000000000)
  (2020-05-14 10:00:00+00:00,13192.0000000000)
  (2020-05-14 20:00:00+00:00,10667.2500000000)
  (2020-05-15 06:00:00+00:00,12515.0000000000)
  (2020-05-15 16:00:00+00:00,12825.2500000000)
  (2020-05-16 02:00:00+00:00,9699.0000000000)
  (2020-05-16 12:00:00+00:00,11637.5000000000)
  (2020-05-16 22:00:00+00:00,9256.5000000000)
  (2020-05-17 08:00:00+00:00,11599.5000000000)
  (2020-05-17 18:00:00+00:00,11019.7500000000)
  (2020-05-18 04:00:00+00:00,11226.5000000000)
  (2020-05-18 14:00:00+00:00,12282.0000000000)
  (2020-05-19 00:00:00+00:00,9457.0000000000)
  (2020-05-19 10:00:00+00:00,13658.7500000000)
  (2020-05-19 20:00:00+00:00,10743.7500000000)
  (2020-05-20 06:00:00+00:00,11710.7500000000)
  (2020-05-20 16:00:00+00:00,11307.5000000000)
  (2020-05-21 02:00:00+00:00,8063.0000000000)
  (2020-05-21 12:00:00+00:00,9130.7500000000)
  (2020-05-21 22:00:00+00:00,8662.5000000000)
  (2020-05-22 08:00:00+00:00,11578.0000000000)
  (2020-05-22 18:00:00+00:00,11765.5000000000)
  (2020-05-23 04:00:00+00:00,8833.0000000000)
  (2020-05-23 14:00:00+00:00,11164.7500000000)
  (2020-05-24 00:00:00+00:00,8445.0000000000)
  (2020-05-24 10:00:00+00:00,11367.0000000000)
  (2020-05-24 20:00:00+00:00,10524.0000000000)
  (2020-05-25 06:00:00+00:00,12666.0000000000)
  (2020-05-25 16:00:00+00:00,11941.2500000000)
  (2020-05-26 02:00:00+00:00,8829.2500000000)
  (2020-05-26 12:00:00+00:00,11826.7500000000)
  (2020-05-26 22:00:00+00:00,8666.5000000000)
  (2020-05-27 08:00:00+00:00,12437.7500000000)
  (2020-05-27 18:00:00+00:00,12042.0000000000)
  (2020-05-28 04:00:00+00:00,10996.5000000000)
  (2020-05-28 14:00:00+00:00,11930.7500000000)
  (2020-05-29 00:00:00+00:00,8475.7500000000)
  (2020-05-29 10:00:00+00:00,12202.7500000000)
  (2020-05-29 20:00:00+00:00,9851.2500000000)
  (2020-05-30 06:00:00+00:00,9901.7500000000)
  (2020-05-30 16:00:00+00:00,10481.0000000000)
  (2020-05-31 02:00:00+00:00,8209.5000000000)
  (2020-05-31 12:00:00+00:00,10208.7500000000)
  (2020-05-31 22:00:00+00:00,8775.0000000000)
  (2020-06-01 08:00:00+00:00,9683.0000000000)
  (2020-06-01 18:00:00+00:00,9298.7500000000)
  (2020-06-02 04:00:00+00:00,10264.0000000000)
  (2020-06-02 14:00:00+00:00,12251.2500000000)
  (2020-06-03 00:00:00+00:00,8562.0000000000)
  (2020-06-03 10:00:00+00:00,12716.0000000000)
  (2020-06-03 20:00:00+00:00,10820.5000000000)
  (2020-06-04 06:00:00+00:00,12848.0000000000)
  (2020-06-04 16:00:00+00:00,13537.0000000000)
  (2020-06-05 02:00:00+00:00,9540.5000000000)
  (2020-06-05 12:00:00+00:00,12214.5000000000)
  (2020-06-05 22:00:00+00:00,10526.2500000000)
  (2020-06-06 08:00:00+00:00,12728.5000000000)
  (2020-06-06 18:00:00+00:00,9776.0000000000)
  (2020-06-07 04:00:00+00:00,7822.0000000000)
  (2020-06-07 14:00:00+00:00,9804.0000000000)
  (2020-06-08 00:00:00+00:00,8400.5000000000)
  (2020-06-08 10:00:00+00:00,12744.5000000000)
  (2020-06-08 20:00:00+00:00,10389.7500000000)
  (2020-06-09 06:00:00+00:00,12269.7500000000)
  (2020-06-09 16:00:00+00:00,12132.7500000000)
  (2020-06-10 02:00:00+00:00,8793.5000000000)
  (2020-06-10 12:00:00+00:00,12610.2500000000)
  (2020-06-10 22:00:00+00:00,9485.7500000000)
  (2020-06-11 08:00:00+00:00,13523.2500000000)
  (2020-06-11 18:00:00+00:00,12490.7500000000)
  (2020-06-12 04:00:00+00:00,11015.5000000000)
  (2020-06-12 14:00:00+00:00,12845.2500000000)
  (2020-06-13 00:00:00+00:00,9084.7500000000)
  (2020-06-13 10:00:00+00:00,12329.2500000000)
  (2020-06-13 20:00:00+00:00,10576.2500000000)
  (2020-06-14 06:00:00+00:00,10191.0000000000)
  (2020-06-14 16:00:00+00:00,11011.2500000000)
  (2020-06-15 02:00:00+00:00,8577.7500000000)
  (2020-06-15 12:00:00+00:00,12174.0000000000)
  (2020-06-15 22:00:00+00:00,9073.7500000000)
  (2020-06-16 08:00:00+00:00,12996.7500000000)
  (2020-06-16 18:00:00+00:00,11439.0000000000)
  (2020-06-17 04:00:00+00:00,10507.5000000000)
  (2020-06-17 14:00:00+00:00,11984.0000000000)
  (2020-06-18 00:00:00+00:00,9004.0000000000)
  (2020-06-18 10:00:00+00:00,13184.2500000000)
  (2020-06-18 20:00:00+00:00,10501.5000000000)
  (2020-06-19 06:00:00+00:00,12131.7500000000)
  (2020-06-19 16:00:00+00:00,11762.0000000000)
  (2020-06-20 02:00:00+00:00,8544.2500000000)
  (2020-06-20 12:00:00+00:00,11105.2500000000)
  (2020-06-20 22:00:00+00:00,9117.0000000000)
  (2020-06-21 08:00:00+00:00,10411.5000000000)
  (2020-06-21 18:00:00+00:00,9594.5000000000)
  (2020-06-22 04:00:00+00:00,10542.7500000000)
  (2020-06-22 14:00:00+00:00,12653.5000000000)
  (2020-06-23 00:00:00+00:00,8924.2500000000)
  (2020-06-23 10:00:00+00:00,13261.7500000000)
  (2020-06-23 20:00:00+00:00,10691.0000000000)
  (2020-06-24 06:00:00+00:00,12006.2500000000)
  (2020-06-24 16:00:00+00:00,12017.7500000000)
  (2020-06-25 02:00:00+00:00,9208.7500000000)
  (2020-06-25 12:00:00+00:00,13349.2500000000)
  (2020-06-25 22:00:00+00:00,9756.0000000000)
  (2020-06-26 08:00:00+00:00,12875.0000000000)
  (2020-06-26 18:00:00+00:00,11689.0000000000)
  (2020-06-27 04:00:00+00:00,8574.0000000000)
  (2020-06-27 14:00:00+00:00,10346.2500000000)
  (2020-06-28 00:00:00+00:00,7909.5000000000)
  (2020-06-28 10:00:00+00:00,10783.0000000000)
  (2020-06-28 20:00:00+00:00,9420.0000000000)
  (2020-06-29 06:00:00+00:00,11810.2500000000)
  (2020-06-29 16:00:00+00:00,11831.2500000000)
  (2020-06-30 02:00:00+00:00,9585.0000000000)
  (2020-06-30 12:00:00+00:00,14149.0000000000)
  (2020-06-30 22:00:00+00:00,11399.7500000000)
  (2020-07-01 08:00:00+00:00,13214.7500000000)
  (2020-07-01 18:00:00+00:00,12356.7500000000)
  (2020-07-02 04:00:00+00:00,10211.7500000000)
  (2020-07-02 14:00:00+00:00,12776.2500000000)
  (2020-07-03 00:00:00+00:00,9070.7500000000)
  (2020-07-03 10:00:00+00:00,12504.7500000000)
  (2020-07-03 20:00:00+00:00,10966.0000000000)
  (2020-07-04 06:00:00+00:00,9869.7500000000)
  (2020-07-04 16:00:00+00:00,11277.2500000000)
  (2020-07-05 02:00:00+00:00,8870.5000000000)
  (2020-07-05 12:00:00+00:00,11878.7500000000)
  (2020-07-05 22:00:00+00:00,10735.2500000000)
  (2020-07-06 08:00:00+00:00,14092.7500000000)
  (2020-07-06 18:00:00+00:00,13582.5000000000)
  (2020-07-07 04:00:00+00:00,10206.2500000000)
  (2020-07-07 14:00:00+00:00,13416.7500000000)
  (2020-07-08 00:00:00+00:00,9729.0000000000)
  (2020-07-08 10:00:00+00:00,13252.5000000000)
  (2020-07-08 20:00:00+00:00,11669.7500000000)
  (2020-07-09 06:00:00+00:00,10362.2500000000)
  (2020-07-09 16:00:00+00:00,12219.0000000000)
  (2020-07-10 02:00:00+00:00,8544.2500000000)
  (2020-07-10 12:00:00+00:00,13261.2500000000)
  (2020-07-10 22:00:00+00:00,10940.2500000000)
  (2020-07-11 08:00:00+00:00,10887.7500000000)
  (2020-07-11 18:00:00+00:00,10746.7500000000)
  (2020-07-12 04:00:00+00:00,8015.0000000000)
  (2020-07-12 14:00:00+00:00,10216.2500000000)
  (2020-07-13 00:00:00+00:00,8153.5000000000)
  (2020-07-13 10:00:00+00:00,12852.0000000000)
  (2020-07-13 20:00:00+00:00,10165.7500000000)
  (2020-07-14 06:00:00+00:00,12010.0000000000)
  (2020-07-14 16:00:00+00:00,12141.5000000000)
  (2020-07-15 02:00:00+00:00,8662.0000000000)
  (2020-07-15 12:00:00+00:00,12288.2500000000)
  (2020-07-15 22:00:00+00:00,8976.7500000000)
  (2020-07-16 08:00:00+00:00,12714.2500000000)
  (2020-07-16 18:00:00+00:00,11372.0000000000)
  (2020-07-17 04:00:00+00:00,10580.7500000000)
  (2020-07-17 14:00:00+00:00,11484.5000000000)
  (2020-07-18 00:00:00+00:00,7792.5000000000)
  (2020-07-18 10:00:00+00:00,10541.5000000000)
  (2020-07-18 20:00:00+00:00,8789.5000000000)
  (2020-07-19 06:00:00+00:00,8601.7500000000)
  (2020-07-19 16:00:00+00:00,9483.5000000000)
  (2020-07-20 02:00:00+00:00,8042.0000000000)
  (2020-07-20 12:00:00+00:00,12660.5000000000)
  (2020-07-20 22:00:00+00:00,9340.0000000000)
  (2020-07-21 08:00:00+00:00,12852.5000000000)
  (2020-07-21 18:00:00+00:00,11680.5000000000)
  (2020-07-22 04:00:00+00:00,11015.5000000000)
  (2020-07-22 14:00:00+00:00,11876.5000000000)
  (2020-07-23 00:00:00+00:00,8269.7500000000)
  (2020-07-23 10:00:00+00:00,12299.2500000000)
  (2020-07-23 20:00:00+00:00,9937.0000000000)
  (2020-07-24 06:00:00+00:00,12144.7500000000)
  (2020-07-24 16:00:00+00:00,12528.7500000000)
  (2020-07-25 02:00:00+00:00,7970.7500000000)
  (2020-07-25 12:00:00+00:00,9884.5000000000)
  (2020-07-25 22:00:00+00:00,8252.5000000000)
  (2020-07-26 08:00:00+00:00,10094.7500000000)
  (2020-07-26 18:00:00+00:00,9220.2500000000)
  (2020-07-27 04:00:00+00:00,9773.0000000000)
  (2020-07-27 14:00:00+00:00,11630.5000000000)
  (2020-07-28 00:00:00+00:00,9281.0000000000)
  (2020-07-28 10:00:00+00:00,13870.0000000000)
  (2020-07-28 20:00:00+00:00,11023.0000000000)
  (2020-07-29 06:00:00+00:00,13005.2500000000)
  (2020-07-29 16:00:00+00:00,13202.0000000000)
  (2020-07-30 02:00:00+00:00,9781.2500000000)
  (2020-07-30 12:00:00+00:00,12833.0000000000)
  (2020-07-30 22:00:00+00:00,9607.7500000000)
  (2020-07-31 08:00:00+00:00,12874.0000000000)
  (2020-07-31 18:00:00+00:00,10539.2500000000)
  (2020-08-01 04:00:00+00:00,8359.7500000000)
  (2020-08-01 14:00:00+00:00,10064.2500000000)
  (2020-08-02 00:00:00+00:00,8320.0000000000)
  (2020-08-02 10:00:00+00:00,10168.0000000000)
  (2020-08-02 20:00:00+00:00,9040.7500000000)
  (2020-08-03 06:00:00+00:00,11700.7500000000)
  (2020-08-03 16:00:00+00:00,11596.7500000000)
  (2020-08-04 02:00:00+00:00,8581.7500000000)
  (2020-08-04 12:00:00+00:00,12066.5000000000)
  (2020-08-04 22:00:00+00:00,8724.0000000000)
  (2020-08-05 08:00:00+00:00,12520.0000000000)
  (2020-08-05 18:00:00+00:00,11341.5000000000)
  (2020-08-06 04:00:00+00:00,10408.5000000000)
  (2020-08-06 14:00:00+00:00,11861.7500000000)
  (2020-08-07 00:00:00+00:00,8750.0000000000)
  (2020-08-07 10:00:00+00:00,12662.0000000000)
  (2020-08-07 20:00:00+00:00,10459.0000000000)
  (2020-08-08 06:00:00+00:00,9922.0000000000)
  (2020-08-08 16:00:00+00:00,10457.2500000000)
  (2020-08-09 02:00:00+00:00,7683.5000000000)
  (2020-08-09 12:00:00+00:00,9642.5000000000)
  (2020-08-09 22:00:00+00:00,8464.0000000000)
  (2020-08-10 08:00:00+00:00,12979.0000000000)
  (2020-08-10 18:00:00+00:00,12128.7500000000)
  (2020-08-11 04:00:00+00:00,11139.5000000000)
  (2020-08-11 14:00:00+00:00,12890.7500000000)
  (2020-08-12 00:00:00+00:00,9478.0000000000)
  (2020-08-12 10:00:00+00:00,13568.5000000000)
  (2020-08-12 20:00:00+00:00,11386.0000000000)
  (2020-08-13 06:00:00+00:00,12588.7500000000)
  (2020-08-13 16:00:00+00:00,12792.5000000000)
  (2020-08-14 02:00:00+00:00,9258.5000000000)
  (2020-08-14 12:00:00+00:00,12317.0000000000)
  (2020-08-14 22:00:00+00:00,8901.5000000000)
  (2020-08-15 08:00:00+00:00,10931.7500000000)
  (2020-08-15 18:00:00+00:00,10023.2500000000)
  (2020-08-16 04:00:00+00:00,7907.0000000000)
  (2020-08-16 14:00:00+00:00,9746.7500000000)
  (2020-08-17 00:00:00+00:00,8525.0000000000)
  (2020-08-17 10:00:00+00:00,13490.5000000000)
  (2020-08-17 20:00:00+00:00,11242.0000000000)
  (2020-08-18 06:00:00+00:00,12387.0000000000)
  (2020-08-18 16:00:00+00:00,12209.7500000000)
  (2020-08-19 02:00:00+00:00,9081.2500000000)
  (2020-08-19 12:00:00+00:00,12520.0000000000)
  (2020-08-19 22:00:00+00:00,9461.0000000000)
  (2020-08-20 08:00:00+00:00,12936.0000000000)
  (2020-08-20 18:00:00+00:00,12582.7500000000)
  (2020-08-21 04:00:00+00:00,11021.5000000000)
  (2020-08-21 14:00:00+00:00,12455.7500000000)
  (2020-08-22 00:00:00+00:00,8694.0000000000)
  (2020-08-22 10:00:00+00:00,11841.0000000000)
  (2020-08-22 20:00:00+00:00,9969.7500000000)
  (2020-08-23 06:00:00+00:00,10046.7500000000)
  (2020-08-23 16:00:00+00:00,11288.5000000000)
  (2020-08-24 02:00:00+00:00,9347.5000000000)
  (2020-08-24 12:00:00+00:00,12575.7500000000)
  (2020-08-24 22:00:00+00:00,9537.5000000000)
  (2020-08-25 08:00:00+00:00,12966.5000000000)
  (2020-08-25 18:00:00+00:00,12119.7500000000)
  (2020-08-26 04:00:00+00:00,12091.7500000000)
  (2020-08-26 14:00:00+00:00,14039.0000000000)
  (2020-08-27 00:00:00+00:00,9821.5000000000)
  (2020-08-27 10:00:00+00:00,13877.0000000000)
  (2020-08-27 20:00:00+00:00,10498.7500000000)
  (2020-08-28 06:00:00+00:00,12974.7500000000)
  (2020-08-28 16:00:00+00:00,12534.7500000000)
  (2020-08-29 02:00:00+00:00,8248.2500000000)
  (2020-08-29 12:00:00+00:00,10947.5000000000)
  (2020-08-29 22:00:00+00:00,8177.7500000000)
  (2020-08-30 08:00:00+00:00,10102.0000000000)
  (2020-08-30 18:00:00+00:00,10418.5000000000)
  (2020-08-31 04:00:00+00:00,11504.7500000000)
  (2020-08-31 14:00:00+00:00,12515.5000000000)
  (2020-09-01 00:00:00+00:00,8934.5000000000)
  (2020-09-01 10:00:00+00:00,13463.5000000000)
  (2020-09-01 20:00:00+00:00,11374.5000000000)
  (2020-09-02 06:00:00+00:00,13381.2500000000)
  (2020-09-02 16:00:00+00:00,13251.7500000000)
  (2020-09-03 02:00:00+00:00,9410.5000000000)
  (2020-09-03 12:00:00+00:00,12684.0000000000)
  (2020-09-03 22:00:00+00:00,10647.0000000000)
  (2020-09-04 08:00:00+00:00,13874.7500000000)
  (2020-09-04 18:00:00+00:00,12265.5000000000)
  (2020-09-05 04:00:00+00:00,9190.7500000000)
  (2020-09-05 14:00:00+00:00,11072.0000000000)
  (2020-09-06 00:00:00+00:00,8507.2500000000)
  (2020-09-06 10:00:00+00:00,11303.2500000000)
  (2020-09-06 20:00:00+00:00,9966.5000000000)
  (2020-09-07 06:00:00+00:00,13443.7500000000)
  (2020-09-07 16:00:00+00:00,12856.5000000000)
  (2020-09-08 02:00:00+00:00,10534.7500000000)
  (2020-09-08 12:00:00+00:00,14112.2500000000)
  (2020-09-08 22:00:00+00:00,10308.2500000000)
  (2020-09-09 08:00:00+00:00,13910.5000000000)
  (2020-09-09 18:00:00+00:00,13866.2500000000)
  (2020-09-10 04:00:00+00:00,12723.2500000000)
  (2020-09-10 14:00:00+00:00,12636.5000000000)
  (2020-09-11 00:00:00+00:00,8822.7500000000)
  (2020-09-11 10:00:00+00:00,12831.7500000000)
  (2020-09-11 20:00:00+00:00,10261.2500000000)
  (2020-09-12 06:00:00+00:00,11068.0000000000)
  (2020-09-12 16:00:00+00:00,12116.5000000000)
  (2020-09-13 02:00:00+00:00,8633.7500000000)
  (2020-09-13 12:00:00+00:00,10763.2500000000)
  (2020-09-13 22:00:00+00:00,9115.2500000000)
  (2020-09-14 08:00:00+00:00,13155.5000000000)
  (2020-09-14 18:00:00+00:00,12874.5000000000)
  (2020-09-15 04:00:00+00:00,12051.7500000000)
  (2020-09-15 14:00:00+00:00,12912.2500000000)
  (2020-09-16 00:00:00+00:00,8925.5000000000)
  (2020-09-16 10:00:00+00:00,13434.2500000000)
  (2020-09-16 20:00:00+00:00,12383.2500000000)
  (2020-09-17 06:00:00+00:00,13122.5000000000)
  (2020-09-17 16:00:00+00:00,12663.2500000000)
  (2020-09-18 02:00:00+00:00,9072.5000000000)
  (2020-09-18 12:00:00+00:00,12038.5000000000)
  (2020-09-18 22:00:00+00:00,8869.2500000000)
  (2020-09-19 08:00:00+00:00,10894.0000000000)
  (2020-09-19 18:00:00+00:00,10683.5000000000)
  (2020-09-20 04:00:00+00:00,8078.2500000000)
  (2020-09-20 14:00:00+00:00,9268.5000000000)
  (2020-09-21 00:00:00+00:00,8130.7500000000)
  (2020-09-21 10:00:00+00:00,12784.0000000000)
  (2020-09-21 20:00:00+00:00,10590.0000000000)
  (2020-09-22 06:00:00+00:00,13107.7500000000)
  (2020-09-22 16:00:00+00:00,12878.5000000000)
  (2020-09-23 02:00:00+00:00,9092.2500000000)
  (2020-09-23 12:00:00+00:00,12761.5000000000)
  (2020-09-23 22:00:00+00:00,9489.5000000000)
  (2020-09-24 08:00:00+00:00,13945.0000000000)
  (2020-09-24 18:00:00+00:00,12158.7500000000)
  (2020-09-25 04:00:00+00:00,12119.7500000000)
  (2020-09-25 14:00:00+00:00,12375.5000000000)
  (2020-09-26 00:00:00+00:00,8447.5000000000)
  (2020-09-26 10:00:00+00:00,12340.2500000000)
  (2020-09-26 20:00:00+00:00,11242.2500000000)
  (2020-09-27 06:00:00+00:00,9287.5000000000)
  (2020-09-27 16:00:00+00:00,11347.2500000000)
  (2020-09-28 02:00:00+00:00,8986.0000000000)
  (2020-09-28 12:00:00+00:00,13227.2500000000)
  (2020-09-28 22:00:00+00:00,9416.0000000000)
  (2020-09-29 08:00:00+00:00,13661.2500000000)
  (2020-09-29 18:00:00+00:00,12925.2500000000)
  (2020-09-30 04:00:00+00:00,12120.7500000000)
  (2020-09-30 14:00:00+00:00,12564.5000000000)
  (2020-10-01 00:00:00+00:00,9187.0000000000)
  (2020-10-01 10:00:00+00:00,13512.0000000000)
  (2020-10-01 20:00:00+00:00,11475.2500000000)
  (2020-10-02 06:00:00+00:00,13756.2500000000)
  (2020-10-02 16:00:00+00:00,13873.7500000000)
  (2020-10-03 02:00:00+00:00,10164.5000000000)
  (2020-10-03 12:00:00+00:00,11623.5000000000)
  (2020-10-03 22:00:00+00:00,9333.5000000000)
  (2020-10-04 08:00:00+00:00,10639.2500000000)
  (2020-10-04 18:00:00+00:00,10983.2500000000)
  (2020-10-05 04:00:00+00:00,12086.5000000000)
  (2020-10-05 14:00:00+00:00,12573.7500000000)
  (2020-10-06 00:00:00+00:00,10234.7500000000)
  (2020-10-06 10:00:00+00:00,14466.5000000000)
  (2020-10-06 20:00:00+00:00,11853.5000000000)
  (2020-10-07 06:00:00+00:00,13517.5000000000)
  (2020-10-07 16:00:00+00:00,13484.2500000000)
  (2020-10-08 02:00:00+00:00,10628.0000000000)
  (2020-10-08 12:00:00+00:00,14209.0000000000)
  (2020-10-08 22:00:00+00:00,11638.0000000000)
  (2020-10-09 08:00:00+00:00,13739.7500000000)
  (2020-10-09 18:00:00+00:00,12396.7500000000)
  (2020-10-10 04:00:00+00:00,9855.7500000000)
  (2020-10-10 14:00:00+00:00,11821.5000000000)
  (2020-10-11 00:00:00+00:00,9097.7500000000)
  (2020-10-11 10:00:00+00:00,11212.0000000000)
  (2020-10-11 20:00:00+00:00,10506.2500000000)
  (2020-10-12 06:00:00+00:00,13351.2500000000)
  (2020-10-12 16:00:00+00:00,13700.7500000000)
  (2020-10-13 02:00:00+00:00,9907.7500000000)
  (2020-10-13 12:00:00+00:00,13039.5000000000)
  (2020-10-13 22:00:00+00:00,11280.7500000000)
  (2020-10-14 08:00:00+00:00,16136.7500000000)
  (2020-10-14 18:00:00+00:00,14729.7500000000)
  (2020-10-15 04:00:00+00:00,13801.0000000000)
  (2020-10-15 14:00:00+00:00,14428.0000000000)
  (2020-10-16 00:00:00+00:00,9752.5000000000)
  (2020-10-16 10:00:00+00:00,14005.0000000000)
  (2020-10-16 20:00:00+00:00,11452.7500000000)
  (2020-10-17 06:00:00+00:00,10609.0000000000)
  (2020-10-17 16:00:00+00:00,11666.2500000000)
  (2020-10-18 02:00:00+00:00,9976.7500000000)
  (2020-10-18 12:00:00+00:00,11824.7500000000)
  (2020-10-18 22:00:00+00:00,9818.2500000000)
  (2020-10-19 08:00:00+00:00,14570.0000000000)
  (2020-10-19 18:00:00+00:00,13270.2500000000)
  (2020-10-20 04:00:00+00:00,12596.0000000000)
  (2020-10-20 14:00:00+00:00,13716.0000000000)
  (2020-10-21 00:00:00+00:00,9675.5000000000)
  (2020-10-21 10:00:00+00:00,13956.0000000000)
  (2020-10-21 20:00:00+00:00,13305.7500000000)
  (2020-10-22 06:00:00+00:00,14688.2500000000)
  (2020-10-22 16:00:00+00:00,14092.7500000000)
  (2020-10-23 02:00:00+00:00,10939.0000000000)
  (2020-10-23 12:00:00+00:00,13553.2500000000)
  (2020-10-23 22:00:00+00:00,9328.5000000000)
  (2020-10-24 08:00:00+00:00,12834.5000000000)
  (2020-10-24 18:00:00+00:00,11352.5000000000)
  (2020-10-25 04:00:00+00:00,9885.7500000000)
  (2020-10-25 14:00:00+00:00,11152.2500000000)
  (2020-10-26 00:00:00+00:00,9039.7500000000)
  (2020-10-26 10:00:00+00:00,13623.0000000000)
  (2020-10-26 20:00:00+00:00,11747.2500000000)
  (2020-10-27 06:00:00+00:00,13882.5000000000)
  (2020-10-27 16:00:00+00:00,14843.7500000000)
  (2020-10-28 02:00:00+00:00,11311.7500000000)
  (2020-10-28 12:00:00+00:00,13897.2500000000)
  (2020-10-28 22:00:00+00:00,11460.5000000000)
  (2020-10-29 08:00:00+00:00,15009.7500000000)
  (2020-10-29 18:00:00+00:00,15223.7500000000)
  (2020-10-30 04:00:00+00:00,10986.5000000000)
  (2020-10-30 14:00:00+00:00,13658.0000000000)
  (2020-10-31 00:00:00+00:00,8933.5000000000)
  (2020-10-31 10:00:00+00:00,11371.2500000000)
  (2020-10-31 20:00:00+00:00,11165.5000000000)
  (2020-11-01 06:00:00+00:00,9483.0000000000)
  (2020-11-01 16:00:00+00:00,12944.2500000000)
  (2020-11-02 02:00:00+00:00,10343.5000000000)
  (2020-11-02 12:00:00+00:00,14833.5000000000)
  (2020-11-02 22:00:00+00:00,11860.2500000000)
  (2020-11-03 08:00:00+00:00,14494.2500000000)
  (2020-11-03 18:00:00+00:00,13409.2500000000)
  (2020-11-04 04:00:00+00:00,12049.0000000000)
  (2020-11-04 14:00:00+00:00,13698.0000000000)
  (2020-11-05 00:00:00+00:00,11221.5000000000)
  (2020-11-05 10:00:00+00:00,14584.5000000000)
  (2020-11-05 20:00:00+00:00,13198.7500000000)
  (2020-11-06 06:00:00+00:00,13756.5000000000)
  (2020-11-06 16:00:00+00:00,13617.5000000000)
  (2020-11-07 02:00:00+00:00,8982.2500000000)
  (2020-11-07 12:00:00+00:00,10957.7500000000)
  (2020-11-07 22:00:00+00:00,9305.0000000000)
  (2020-11-08 08:00:00+00:00,10376.2500000000)
  (2020-11-08 18:00:00+00:00,11327.7500000000)
  (2020-11-09 04:00:00+00:00,10429.2500000000)
  (2020-11-09 14:00:00+00:00,13202.0000000000)
  (2020-11-10 00:00:00+00:00,9959.7500000000)
  (2020-11-10 10:00:00+00:00,14523.5000000000)
  (2020-11-10 20:00:00+00:00,11940.7500000000)
  (2020-11-11 06:00:00+00:00,13462.0000000000)
  (2020-11-11 16:00:00+00:00,14540.0000000000)
  (2020-11-12 02:00:00+00:00,9816.5000000000)
  (2020-11-12 12:00:00+00:00,13640.7500000000)
  (2020-11-12 22:00:00+00:00,11154.5000000000)
  (2020-11-13 08:00:00+00:00,13428.5000000000)
  (2020-11-13 18:00:00+00:00,13603.7500000000)
  (2020-11-14 04:00:00+00:00,9209.2500000000)
  (2020-11-14 14:00:00+00:00,10709.0000000000)
  (2020-11-15 00:00:00+00:00,9376.0000000000)
  (2020-11-15 10:00:00+00:00,12237.0000000000)
  (2020-11-15 20:00:00+00:00,12444.0000000000)
  (2020-11-16 06:00:00+00:00,14105.5000000000)
  (2020-11-16 16:00:00+00:00,15279.2500000000)
  (2020-11-17 02:00:00+00:00,10943.7500000000)
  (2020-11-17 12:00:00+00:00,14890.0000000000)
  (2020-11-17 22:00:00+00:00,11800.7500000000)
  (2020-11-18 08:00:00+00:00,13799.2500000000)
  (2020-11-18 18:00:00+00:00,14883.5000000000)
  (2020-11-19 04:00:00+00:00,12588.0000000000)
  (2020-11-19 14:00:00+00:00,15436.7500000000)
  (2020-11-20 00:00:00+00:00,11717.0000000000)
  (2020-11-20 10:00:00+00:00,15044.2500000000)
  (2020-11-20 20:00:00+00:00,12140.2500000000)
  (2020-11-21 06:00:00+00:00,12093.5000000000)
  (2020-11-21 16:00:00+00:00,14781.0000000000)
  (2020-11-22 02:00:00+00:00,10605.2500000000)
  (2020-11-22 12:00:00+00:00,12674.2500000000)
  (2020-11-22 22:00:00+00:00,11557.0000000000)
  (2020-11-23 08:00:00+00:00,15327.0000000000)
  (2020-11-23 18:00:00+00:00,15612.2500000000)
  (2020-11-24 04:00:00+00:00,12917.5000000000)
  (2020-11-24 14:00:00+00:00,14386.2500000000)
  (2020-11-25 00:00:00+00:00,10624.2500000000)
  (2020-11-25 10:00:00+00:00,14681.2500000000)
  (2020-11-25 20:00:00+00:00,13042.5000000000)
  (2020-11-26 06:00:00+00:00,14495.0000000000)
  (2020-11-26 16:00:00+00:00,15347.5000000000)
  (2020-11-27 02:00:00+00:00,10317.2500000000)
  (2020-11-27 12:00:00+00:00,13789.7500000000)
  (2020-11-27 22:00:00+00:00,10822.0000000000)
  (2020-11-28 08:00:00+00:00,11846.5000000000)
  (2020-11-28 18:00:00+00:00,12412.7500000000)
  (2020-11-29 04:00:00+00:00,9229.7500000000)
  (2020-11-29 14:00:00+00:00,11302.5000000000)
  (2020-11-30 00:00:00+00:00,9787.0000000000)
  (2020-11-30 10:00:00+00:00,14978.2500000000)
  (2020-11-30 20:00:00+00:00,13448.0000000000)
  (2020-12-01 06:00:00+00:00,14212.2500000000)
  (2020-12-01 16:00:00+00:00,14997.7500000000)
  (2020-12-02 02:00:00+00:00,10469.5000000000)
  (2020-12-02 12:00:00+00:00,14372.0000000000)
  (2020-12-02 22:00:00+00:00,12441.5000000000)
  (2020-12-03 08:00:00+00:00,15872.2500000000)
  (2020-12-03 18:00:00+00:00,16272.2500000000)
  (2020-12-04 04:00:00+00:00,13267.0000000000)
  (2020-12-04 14:00:00+00:00,14945.2500000000)
  (2020-12-05 00:00:00+00:00,10742.0000000000)
  (2020-12-05 10:00:00+00:00,13709.7500000000)
  (2020-12-05 20:00:00+00:00,11003.7500000000)
  (2020-12-06 06:00:00+00:00,10421.5000000000)
  (2020-12-06 16:00:00+00:00,13765.7500000000)
  (2020-12-07 02:00:00+00:00,11449.5000000000)
  (2020-12-07 12:00:00+00:00,15823.2500000000)
  (2020-12-07 22:00:00+00:00,11299.0000000000)
  (2020-12-08 08:00:00+00:00,14589.0000000000)
  (2020-12-08 18:00:00+00:00,15030.2500000000)
  (2020-12-09 04:00:00+00:00,11794.5000000000)
  (2020-12-09 14:00:00+00:00,15078.2500000000)
  (2020-12-10 00:00:00+00:00,10734.5000000000)
  (2020-12-10 10:00:00+00:00,15239.0000000000)
  (2020-12-10 20:00:00+00:00,12915.5000000000)
  (2020-12-11 06:00:00+00:00,14186.7500000000)
  (2020-12-11 16:00:00+00:00,15529.0000000000)
  (2020-12-12 02:00:00+00:00,10541.2500000000)
  (2020-12-12 12:00:00+00:00,12535.0000000000)
  (2020-12-12 22:00:00+00:00,10131.5000000000)
  (2020-12-13 08:00:00+00:00,10844.7500000000)
  (2020-12-13 18:00:00+00:00,11944.2500000000)
  (2020-12-14 04:00:00+00:00,11245.0000000000)
  (2020-12-14 14:00:00+00:00,14863.7500000000)
  (2020-12-15 00:00:00+00:00,10442.5000000000)
  (2020-12-15 10:00:00+00:00,15327.0000000000)
  (2020-12-15 20:00:00+00:00,13067.5000000000)
  (2020-12-16 06:00:00+00:00,13634.0000000000)
  (2020-12-16 16:00:00+00:00,14764.0000000000)
  (2020-12-17 02:00:00+00:00,11548.0000000000)
  (2020-12-17 12:00:00+00:00,14892.0000000000)
  (2020-12-17 22:00:00+00:00,12221.0000000000)
  (2020-12-18 08:00:00+00:00,14015.2500000000)
  (2020-12-18 18:00:00+00:00,14687.5000000000)
  (2020-12-19 04:00:00+00:00,10475.2500000000)
  (2020-12-19 14:00:00+00:00,12501.2500000000)
  (2020-12-20 00:00:00+00:00,10203.7500000000)
  (2020-12-20 10:00:00+00:00,12314.7500000000)
  (2020-12-20 20:00:00+00:00,11099.7500000000)
  (2020-12-21 06:00:00+00:00,12464.0000000000)
  (2020-12-21 16:00:00+00:00,13825.5000000000)
  (2020-12-22 02:00:00+00:00,10938.2500000000)
  (2020-12-22 12:00:00+00:00,14003.5000000000)
  (2020-12-22 22:00:00+00:00,10520.2500000000)
  (2020-12-23 08:00:00+00:00,11992.0000000000)
  (2020-12-23 18:00:00+00:00,12127.0000000000)
  (2020-12-24 04:00:00+00:00,8965.2500000000)
  (2020-12-24 14:00:00+00:00,11133.2500000000)
  (2020-12-25 00:00:00+00:00,8707.0000000000)
  (2020-12-25 10:00:00+00:00,11975.7500000000)
  (2020-12-25 20:00:00+00:00,10498.5000000000)
  (2020-12-26 06:00:00+00:00,8829.0000000000)
  (2020-12-26 16:00:00+00:00,12353.2500000000)
  (2020-12-27 02:00:00+00:00,10267.0000000000)
  (2020-12-27 12:00:00+00:00,12924.5000000000)
  (2020-12-27 22:00:00+00:00,11448.7500000000)
  (2020-12-28 08:00:00+00:00,12572.0000000000)
  (2020-12-28 18:00:00+00:00,12296.0000000000)
  (2020-12-29 04:00:00+00:00,9629.5000000000)
  (2020-12-29 14:00:00+00:00,11384.7500000000)
  (2020-12-30 00:00:00+00:00,10280.0000000000)
  (2020-12-30 10:00:00+00:00,12442.5000000000)
  (2020-12-30 20:00:00+00:00,11690.5000000000)
  (2020-12-31 06:00:00+00:00,10168.5000000000)
  (2020-12-31 16:00:00+00:00,12175.0000000000)
  (2021-01-01 02:00:00+00:00,8897.5000000000)
  (2021-01-01 12:00:00+00:00,10249.2500000000)
  (2021-01-01 22:00:00+00:00,9671.2500000000)
  (2021-01-02 08:00:00+00:00,10648.0000000000)
  (2021-01-02 18:00:00+00:00,11973.0000000000)
  (2021-01-03 04:00:00+00:00,9824.5000000000)
  (2021-01-03 14:00:00+00:00,12642.2500000000)
  (2021-01-04 00:00:00+00:00,10666.7500000000)
  (2021-01-04 10:00:00+00:00,15794.5000000000)
  (2021-01-04 20:00:00+00:00,13856.7500000000)
  (2021-01-05 06:00:00+00:00,14696.5000000000)
  (2021-01-05 16:00:00+00:00,16231.7500000000)
  (2021-01-06 02:00:00+00:00,11589.2500000000)
  (2021-01-06 12:00:00+00:00,15094.5000000000)
  (2021-01-06 22:00:00+00:00,11949.7500000000)
  (2021-01-07 08:00:00+00:00,14756.5000000000)
  (2021-01-07 18:00:00+00:00,14987.5000000000)
  (2021-01-08 04:00:00+00:00,11709.5000000000)
  (2021-01-08 14:00:00+00:00,14266.2500000000)
  (2021-01-09 00:00:00+00:00,10650.5000000000)
  (2021-01-09 10:00:00+00:00,13906.7500000000)
  (2021-01-09 20:00:00+00:00,11782.5000000000)
  (2021-01-10 06:00:00+00:00,10174.7500000000)
  (2021-01-10 16:00:00+00:00,14023.0000000000)
  (2021-01-11 02:00:00+00:00,11437.5000000000)
  (2021-01-11 12:00:00+00:00,16476.0000000000)
  (2021-01-11 22:00:00+00:00,13980.7500000000)
  (2021-01-12 08:00:00+00:00,15647.5000000000)
  (2021-01-12 18:00:00+00:00,15787.5000000000)
  (2021-01-13 04:00:00+00:00,13623.0000000000)
  (2021-01-13 14:00:00+00:00,16095.7500000000)
  (2021-01-14 00:00:00+00:00,12881.0000000000)
  (2021-01-14 10:00:00+00:00,16532.5000000000)
  (2021-01-14 20:00:00+00:00,13992.5000000000)
  (2021-01-15 06:00:00+00:00,15126.5000000000)
  (2021-01-15 16:00:00+00:00,15820.0000000000)
  (2021-01-16 02:00:00+00:00,11121.5000000000)
  (2021-01-16 12:00:00+00:00,13879.5000000000)
  (2021-01-16 22:00:00+00:00,11358.2500000000)
  (2021-01-17 08:00:00+00:00,12233.0000000000)
  (2021-01-17 18:00:00+00:00,13401.7500000000)
  (2021-01-18 04:00:00+00:00,12012.0000000000)
  (2021-01-18 14:00:00+00:00,15583.2500000000)
  (2021-01-19 00:00:00+00:00,11739.0000000000)
  (2021-01-19 10:00:00+00:00,16209.2500000000)
  (2021-01-19 20:00:00+00:00,13688.5000000000)
  (2021-01-20 06:00:00+00:00,15522.7500000000)
  (2021-01-20 16:00:00+00:00,16475.2500000000)
  (2021-01-21 02:00:00+00:00,12520.5000000000)
  (2021-01-21 12:00:00+00:00,15906.7500000000)
  (2021-01-21 22:00:00+00:00,13065.5000000000)
  (2021-01-22 08:00:00+00:00,16194.7500000000)
  (2021-01-22 18:00:00+00:00,15433.5000000000)
  (2021-01-23 04:00:00+00:00,10132.5000000000)
  (2021-01-23 14:00:00+00:00,11790.7500000000)
  (2021-01-24 00:00:00+00:00,9783.0000000000)
  (2021-01-24 10:00:00+00:00,12717.0000000000)
  (2021-01-24 20:00:00+00:00,11418.2500000000)
  (2021-01-25 06:00:00+00:00,13904.2500000000)
  (2021-01-25 16:00:00+00:00,14993.5000000000)
  (2021-01-26 02:00:00+00:00,11439.2500000000)
  (2021-01-26 12:00:00+00:00,14887.2500000000)
  (2021-01-26 22:00:00+00:00,12528.7500000000)
  (2021-01-27 08:00:00+00:00,14992.2500000000)
  (2021-01-27 18:00:00+00:00,15043.5000000000)
  (2021-01-28 04:00:00+00:00,12980.2500000000)
  (2021-01-28 14:00:00+00:00,14862.5000000000)
  (2021-01-29 00:00:00+00:00,11748.5000000000)
  (2021-01-29 10:00:00+00:00,15248.7500000000)
  (2021-01-29 20:00:00+00:00,13107.2500000000)
  (2021-01-30 06:00:00+00:00,11384.2500000000)
  (2021-01-30 16:00:00+00:00,13526.0000000000)
  (2021-01-31 02:00:00+00:00,10971.7500000000)
  (2021-01-31 12:00:00+00:00,12893.2500000000)
  (2021-01-31 22:00:00+00:00,11876.7500000000)
  (2021-02-01 08:00:00+00:00,15572.7500000000)
  (2021-02-01 18:00:00+00:00,15884.0000000000)
  (2021-02-02 04:00:00+00:00,11803.2500000000)
  (2021-02-02 14:00:00+00:00,14250.0000000000)
  (2021-02-03 00:00:00+00:00,10400.7500000000)
  (2021-02-03 10:00:00+00:00,16079.7500000000)
  (2021-02-03 20:00:00+00:00,14392.7500000000)
  (2021-02-04 06:00:00+00:00,14828.7500000000)
  (2021-02-04 16:00:00+00:00,14955.5000000000)
  (2021-02-05 02:00:00+00:00,11464.0000000000)
  (2021-02-05 12:00:00+00:00,15172.7500000000)
  (2021-02-05 22:00:00+00:00,12955.0000000000)
  (2021-02-06 08:00:00+00:00,14216.7500000000)
  (2021-02-06 18:00:00+00:00,15313.0000000000)
  (2021-02-07 04:00:00+00:00,11930.0000000000)
  (2021-02-07 14:00:00+00:00,14616.7500000000)
  (2021-02-08 00:00:00+00:00,13553.0000000000)
  (2021-02-08 10:00:00+00:00,17058.0000000000)
  (2021-02-08 20:00:00+00:00,14477.7500000000)
  (2021-02-09 06:00:00+00:00,14998.2500000000)
  (2021-02-09 16:00:00+00:00,15864.0000000000)
  (2021-02-10 02:00:00+00:00,12886.7500000000)
  (2021-02-10 12:00:00+00:00,15881.0000000000)
  (2021-02-10 22:00:00+00:00,14034.2500000000)
  (2021-02-11 08:00:00+00:00,16154.5000000000)
  (2021-02-11 18:00:00+00:00,16243.5000000000)
  (2021-02-12 04:00:00+00:00,13617.0000000000)
  (2021-02-12 14:00:00+00:00,14686.7500000000)
  (2021-02-13 00:00:00+00:00,12385.5000000000)
  (2021-02-13 10:00:00+00:00,14344.5000000000)
  (2021-02-13 20:00:00+00:00,13301.5000000000)
  (2021-02-14 06:00:00+00:00,11732.0000000000)
  (2021-02-14 16:00:00+00:00,13192.5000000000)
  (2021-02-15 02:00:00+00:00,12192.2500000000)
  (2021-02-15 12:00:00+00:00,16407.2500000000)
  (2021-02-15 22:00:00+00:00,13875.7500000000)
  (2021-02-16 08:00:00+00:00,14918.7500000000)
  (2021-02-16 18:00:00+00:00,15130.7500000000)
  (2021-02-17 04:00:00+00:00,13067.2500000000)
  (2021-02-17 14:00:00+00:00,14956.2500000000)
  (2021-02-18 00:00:00+00:00,10936.2500000000)
  (2021-02-18 10:00:00+00:00,14603.5000000000)
  (2021-02-18 20:00:00+00:00,14847.7500000000)
  (2021-02-19 06:00:00+00:00,15335.7500000000)
  (2021-02-19 16:00:00+00:00,13907.2500000000)
  (2021-02-20 02:00:00+00:00,10609.0000000000)
  (2021-02-20 12:00:00+00:00,11974.2500000000)
  (2021-02-20 22:00:00+00:00,11458.2500000000)
  (2021-02-21 08:00:00+00:00,11688.5000000000)
  (2021-02-21 18:00:00+00:00,12584.5000000000)
  (2021-02-22 04:00:00+00:00,11620.0000000000)
  (2021-02-22 14:00:00+00:00,12990.0000000000)
  (2021-02-23 00:00:00+00:00,11236.5000000000)
  (2021-02-23 10:00:00+00:00,14524.5000000000)
  (2021-02-23 20:00:00+00:00,13169.7500000000)
  (2021-02-24 06:00:00+00:00,14768.2500000000)
  (2021-02-24 16:00:00+00:00,13831.7500000000)
  (2021-02-25 02:00:00+00:00,11086.2500000000)
  (2021-02-25 12:00:00+00:00,13800.5000000000)
  (2021-02-25 22:00:00+00:00,11774.0000000000)
  (2021-02-26 08:00:00+00:00,14863.7500000000)
  (2021-02-26 18:00:00+00:00,15394.7500000000)
  (2021-02-27 04:00:00+00:00,10584.5000000000)
  (2021-02-27 14:00:00+00:00,11577.0000000000)
  (2021-02-28 00:00:00+00:00,10262.0000000000)
  (2021-02-28 10:00:00+00:00,12610.7500000000)
  (2021-02-28 20:00:00+00:00,11043.7500000000)
  (2021-03-01 06:00:00+00:00,13405.2500000000)
  (2021-03-01 16:00:00+00:00,13768.7500000000)
  (2021-03-02 02:00:00+00:00,10793.5000000000)
  (2021-03-02 12:00:00+00:00,13359.0000000000)
  (2021-03-02 22:00:00+00:00,11112.5000000000)
  (2021-03-03 08:00:00+00:00,13843.2500000000)
  (2021-03-03 18:00:00+00:00,14769.0000000000)
  (2021-03-04 04:00:00+00:00,12126.0000000000)
  (2021-03-04 14:00:00+00:00,13976.2500000000)
  (2021-03-05 00:00:00+00:00,10977.5000000000)
  (2021-03-05 10:00:00+00:00,14503.5000000000)
  (2021-03-05 20:00:00+00:00,12720.0000000000)
  (2021-03-06 06:00:00+00:00,11653.2500000000)
  (2021-03-06 16:00:00+00:00,13147.0000000000)
  (2021-03-07 02:00:00+00:00,11131.5000000000)
  (2021-03-07 12:00:00+00:00,13385.2500000000)
  (2021-03-07 22:00:00+00:00,11113.0000000000)
  (2021-03-08 08:00:00+00:00,14060.5000000000)
  (2021-03-08 18:00:00+00:00,14535.2500000000)
  (2021-03-09 04:00:00+00:00,11919.5000000000)
  (2021-03-09 14:00:00+00:00,13523.2500000000)
  (2021-03-10 00:00:00+00:00,10383.5000000000)
  (2021-03-10 10:00:00+00:00,14167.7500000000)
  (2021-03-10 20:00:00+00:00,14406.0000000000)
  (2021-03-11 06:00:00+00:00,15836.0000000000)
  (2021-03-11 16:00:00+00:00,15872.7500000000)
  (2021-03-12 02:00:00+00:00,12263.0000000000)
  (2021-03-12 12:00:00+00:00,15558.5000000000)
  (2021-03-12 22:00:00+00:00,12821.7500000000)
  (2021-03-13 08:00:00+00:00,13559.0000000000)
  (2021-03-13 18:00:00+00:00,14459.7500000000)
  (2021-03-14 04:00:00+00:00,11509.0000000000)
  (2021-03-14 14:00:00+00:00,11707.2500000000)
  (2021-03-15 00:00:00+00:00,11254.5000000000)
  (2021-03-15 10:00:00+00:00,15295.2500000000)
  (2021-03-15 20:00:00+00:00,13277.7500000000)
  (2021-03-16 06:00:00+00:00,14592.7500000000)
  (2021-03-16 16:00:00+00:00,13526.2500000000)
  (2021-03-17 02:00:00+00:00,10625.7500000000)
  (2021-03-17 12:00:00+00:00,13886.5000000000)
  (2021-03-17 22:00:00+00:00,11408.2500000000)
  (2021-03-18 08:00:00+00:00,13886.2500000000)
  (2021-03-18 18:00:00+00:00,14532.7500000000)
  (2021-03-19 04:00:00+00:00,12646.2500000000)
  (2021-03-19 14:00:00+00:00,14756.5000000000)
  (2021-03-20 00:00:00+00:00,10578.0000000000)
  (2021-03-20 10:00:00+00:00,13334.0000000000)
  (2021-03-20 20:00:00+00:00,12991.5000000000)
  (2021-03-21 06:00:00+00:00,11691.7500000000)
  (2021-03-21 16:00:00+00:00,12791.2500000000)
  (2021-03-22 02:00:00+00:00,10560.0000000000)
  (2021-03-22 12:00:00+00:00,13767.2500000000)
  (2021-03-22 22:00:00+00:00,12113.2500000000)
  (2021-03-23 08:00:00+00:00,14249.7500000000)
  (2021-03-23 18:00:00+00:00,14465.2500000000)
  (2021-03-24 04:00:00+00:00,11941.5000000000)
  (2021-03-24 14:00:00+00:00,12991.7500000000)
  (2021-03-25 00:00:00+00:00,10622.7500000000)
  (2021-03-25 10:00:00+00:00,13443.5000000000)
  (2021-03-25 20:00:00+00:00,12142.2500000000)
  (2021-03-26 06:00:00+00:00,13484.5000000000)
  (2021-03-26 16:00:00+00:00,12594.7500000000)
  (2021-03-27 02:00:00+00:00,10651.2500000000)
  (2021-03-27 12:00:00+00:00,13475.5000000000)
  (2021-03-27 22:00:00+00:00,11794.5000000000)
  (2021-03-28 08:00:00+00:00,12138.0000000000)
  (2021-03-28 18:00:00+00:00,13054.0000000000)
  (2021-03-29 04:00:00+00:00,13039.0000000000)
  (2021-03-29 14:00:00+00:00,13586.7500000000)
  (2021-03-30 00:00:00+00:00,10592.7500000000)
  (2021-03-30 10:00:00+00:00,13863.7500000000)
  (2021-03-30 20:00:00+00:00,11781.0000000000)
  (2021-03-31 06:00:00+00:00,13000.2500000000)
  (2021-03-31 16:00:00+00:00,12041.7500000000)
  (2021-04-01 02:00:00+00:00,9767.2500000000)
  (2021-04-01 12:00:00+00:00,13372.0000000000)
  (2021-04-01 22:00:00+00:00,9813.0000000000)
  (2021-04-02 08:00:00+00:00,11713.5000000000)
  (2021-04-02 18:00:00+00:00,11377.7500000000)
  (2021-04-03 04:00:00+00:00,10267.0000000000)
  (2021-04-03 14:00:00+00:00,10918.7500000000)
  (2021-04-04 00:00:00+00:00,8700.5000000000)
  (2021-04-04 10:00:00+00:00,11276.0000000000)
  (2021-04-04 20:00:00+00:00,11003.2500000000)
  (2021-04-05 06:00:00+00:00,10867.2500000000)
  (2021-04-05 16:00:00+00:00,12684.2500000000)
  (2021-04-06 02:00:00+00:00,11562.0000000000)
  (2021-04-06 12:00:00+00:00,15106.7500000000)
  (2021-04-06 22:00:00+00:00,11804.5000000000)
  (2021-04-07 08:00:00+00:00,15937.5000000000)
  (2021-04-07 18:00:00+00:00,14612.7500000000)
  (2021-04-08 04:00:00+00:00,14060.5000000000)
  (2021-04-08 14:00:00+00:00,14389.0000000000)
  (2021-04-09 00:00:00+00:00,12155.5000000000)
  (2021-04-09 10:00:00+00:00,15700.2500000000)
  (2021-04-09 20:00:00+00:00,12079.7500000000)
  (2021-04-10 06:00:00+00:00,11100.5000000000)
  (2021-04-10 16:00:00+00:00,11364.5000000000)
  (2021-04-11 02:00:00+00:00,10181.7500000000)
  (2021-04-11 12:00:00+00:00,10581.0000000000)
  (2021-04-11 22:00:00+00:00,10668.0000000000)
  (2021-04-12 08:00:00+00:00,14427.7500000000)
  (2021-04-12 18:00:00+00:00,13556.7500000000)
  (2021-04-13 04:00:00+00:00,13455.7500000000)
  (2021-04-13 14:00:00+00:00,13767.0000000000)
  (2021-04-14 00:00:00+00:00,10641.5000000000)
  (2021-04-14 10:00:00+00:00,13911.2500000000)
  (2021-04-14 20:00:00+00:00,12402.0000000000)
  (2021-04-15 06:00:00+00:00,14318.2500000000)
  (2021-04-15 16:00:00+00:00,14428.5000000000)
  (2021-04-16 02:00:00+00:00,12040.2500000000)
  (2021-04-16 12:00:00+00:00,13750.2500000000)
  (2021-04-16 22:00:00+00:00,11008.0000000000)
  (2021-04-17 08:00:00+00:00,12477.2500000000)
  (2021-04-17 18:00:00+00:00,12018.7500000000)
  (2021-04-18 04:00:00+00:00,8797.0000000000)
  (2021-04-18 14:00:00+00:00,10584.0000000000)
  (2021-04-19 00:00:00+00:00,9576.7500000000)
  (2021-04-19 10:00:00+00:00,14515.2500000000)
  (2021-04-19 20:00:00+00:00,12598.0000000000)
  (2021-04-20 06:00:00+00:00,14002.2500000000)
  (2021-04-20 16:00:00+00:00,12777.5000000000)
  (2021-04-21 02:00:00+00:00,10375.7500000000)
  (2021-04-21 12:00:00+00:00,13483.5000000000)
  (2021-04-21 22:00:00+00:00,11764.7500000000)
  (2021-04-22 08:00:00+00:00,15280.2500000000)
  (2021-04-22 18:00:00+00:00,14143.2500000000)
  (2021-04-23 04:00:00+00:00,12807.7500000000)
  (2021-04-23 14:00:00+00:00,13601.5000000000)
  (2021-04-24 00:00:00+00:00,9848.7500000000)
  (2021-04-24 10:00:00+00:00,11668.5000000000)
  (2021-04-24 20:00:00+00:00,10618.0000000000)
  (2021-04-25 06:00:00+00:00,10036.0000000000)
  (2021-04-25 16:00:00+00:00,11806.2500000000)
  (2021-04-26 02:00:00+00:00,10190.2500000000)
  (2021-04-26 12:00:00+00:00,12499.2500000000)
  (2021-04-26 22:00:00+00:00,10383.7500000000)
  (2021-04-27 08:00:00+00:00,12954.5000000000)
  (2021-04-27 18:00:00+00:00,12421.2500000000)
  (2021-04-28 04:00:00+00:00,12140.2500000000)
  (2021-04-28 14:00:00+00:00,12541.0000000000)
  (2021-04-29 00:00:00+00:00,10220.7500000000)
  (2021-04-29 10:00:00+00:00,13865.5000000000)
  (2021-04-29 20:00:00+00:00,12792.7500000000)
  (2021-04-30 06:00:00+00:00,14378.2500000000)
  (2021-04-30 16:00:00+00:00,12370.7500000000)
  (2021-05-01 02:00:00+00:00,8755.7500000000)
  (2021-05-01 12:00:00+00:00,10062.5000000000)
  (2021-05-01 22:00:00+00:00,9237.5000000000)
  (2021-05-02 08:00:00+00:00,11432.0000000000)
  (2021-05-02 18:00:00+00:00,11018.5000000000)
  (2021-05-03 04:00:00+00:00,12351.5000000000)
  (2021-05-03 14:00:00+00:00,13438.7500000000)
  (2021-05-04 00:00:00+00:00,10478.7500000000)
  (2021-05-04 10:00:00+00:00,15782.5000000000)
  (2021-05-04 20:00:00+00:00,13384.0000000000)
  (2021-05-05 06:00:00+00:00,15525.2500000000)
  (2021-05-05 16:00:00+00:00,15392.0000000000)
  (2021-05-06 02:00:00+00:00,12290.7500000000)
  (2021-05-06 12:00:00+00:00,14431.2500000000)
  (2021-05-06 22:00:00+00:00,10382.0000000000)
  (2021-05-07 08:00:00+00:00,14531.5000000000)
  (2021-05-07 18:00:00+00:00,13021.2500000000)
  (2021-05-08 04:00:00+00:00,10549.5000000000)
  (2021-05-08 14:00:00+00:00,10373.7500000000)
  (2021-05-09 00:00:00+00:00,9693.2500000000)
  (2021-05-09 10:00:00+00:00,11168.2500000000)
  (2021-05-09 20:00:00+00:00,11957.7500000000)
  (2021-05-10 06:00:00+00:00,13716.7500000000)
  (2021-05-10 16:00:00+00:00,13015.7500000000)
  (2021-05-11 02:00:00+00:00,10165.0000000000)
  (2021-05-11 12:00:00+00:00,12522.7500000000)
  (2021-05-11 22:00:00+00:00,10335.2500000000)
  (2021-05-12 08:00:00+00:00,13370.5000000000)
  (2021-05-12 18:00:00+00:00,12363.7500000000)
  (2021-05-13 04:00:00+00:00,9089.2500000000)
  (2021-05-13 14:00:00+00:00,10260.5000000000)
  (2021-05-14 00:00:00+00:00,8733.2500000000)
  (2021-05-14 10:00:00+00:00,11996.0000000000)
  (2021-05-14 20:00:00+00:00,10321.2500000000)
  (2021-05-15 06:00:00+00:00,10797.2500000000)
  (2021-05-15 16:00:00+00:00,10449.7500000000)
  (2021-05-16 02:00:00+00:00,8216.0000000000)
  (2021-05-16 12:00:00+00:00,9811.0000000000)
  (2021-05-16 22:00:00+00:00,9041.0000000000)
  (2021-05-17 08:00:00+00:00,12735.5000000000)
  (2021-05-17 18:00:00+00:00,11632.0000000000)
  (2021-05-18 04:00:00+00:00,11377.0000000000)
  (2021-05-18 14:00:00+00:00,12293.5000000000)
  (2021-05-19 00:00:00+00:00,9313.7500000000)
  (2021-05-19 10:00:00+00:00,13298.5000000000)
  (2021-05-19 20:00:00+00:00,11455.7500000000)
  (2021-05-20 06:00:00+00:00,13142.0000000000)
  (2021-05-20 16:00:00+00:00,12715.7500000000)
  (2021-05-21 02:00:00+00:00,9972.7500000000)
  (2021-05-21 12:00:00+00:00,13497.0000000000)
  (2021-05-21 22:00:00+00:00,9561.2500000000)
  (2021-05-22 08:00:00+00:00,12017.0000000000)
  (2021-05-22 18:00:00+00:00,10983.5000000000)
  (2021-05-23 04:00:00+00:00,8687.0000000000)
  (2021-05-23 14:00:00+00:00,10277.2500000000)
  (2021-05-24 00:00:00+00:00,8125.5000000000)
  (2021-05-24 10:00:00+00:00,10987.5000000000)
  (2021-05-24 20:00:00+00:00,10831.5000000000)
  (2021-05-25 06:00:00+00:00,13037.0000000000)
  (2021-05-25 16:00:00+00:00,12825.7500000000)
  (2021-05-26 02:00:00+00:00,10485.7500000000)
  (2021-05-26 12:00:00+00:00,13322.7500000000)
  (2021-05-26 22:00:00+00:00,10309.0000000000)
  (2021-05-27 08:00:00+00:00,13594.5000000000)
  (2021-05-27 18:00:00+00:00,12115.7500000000)
  (2021-05-28 04:00:00+00:00,11605.5000000000)
  (2021-05-28 14:00:00+00:00,12560.2500000000)
  (2021-05-29 00:00:00+00:00,9514.5000000000)
  (2021-05-29 10:00:00+00:00,11904.5000000000)
  (2021-05-29 20:00:00+00:00,9984.7500000000)
  (2021-05-30 06:00:00+00:00,9536.5000000000)
  (2021-05-30 16:00:00+00:00,10071.0000000000)
  (2021-05-31 02:00:00+00:00,9124.0000000000)
  (2021-05-31 12:00:00+00:00,12372.7500000000)
  (2021-05-31 22:00:00+00:00,9760.5000000000)
  (2021-06-01 08:00:00+00:00,12909.2500000000)
  (2021-06-01 18:00:00+00:00,12063.5000000000)
  (2021-06-02 04:00:00+00:00,11236.5000000000)
  (2021-06-02 14:00:00+00:00,12617.0000000000)
  (2021-06-03 00:00:00+00:00,10387.2500000000)
  (2021-06-03 10:00:00+00:00,13738.0000000000)
  (2021-06-03 20:00:00+00:00,12369.5000000000)
  (2021-06-04 06:00:00+00:00,12805.7500000000)
  (2021-06-04 16:00:00+00:00,12469.0000000000)
  (2021-06-05 02:00:00+00:00,8538.0000000000)
  (2021-06-05 12:00:00+00:00,11004.5000000000)
  (2021-06-05 22:00:00+00:00,8818.7500000000)
  (2021-06-06 08:00:00+00:00,10357.2500000000)
  (2021-06-06 18:00:00+00:00,10119.0000000000)
  (2021-06-07 04:00:00+00:00,10813.2500000000)
  (2021-06-07 14:00:00+00:00,12268.0000000000)
  (2021-06-08 00:00:00+00:00,8884.5000000000)
  (2021-06-08 10:00:00+00:00,13272.0000000000)
  (2021-06-08 20:00:00+00:00,11272.5000000000)
  (2021-06-09 06:00:00+00:00,12808.2500000000)
  (2021-06-09 16:00:00+00:00,12891.0000000000)
  (2021-06-10 02:00:00+00:00,9460.0000000000)
  (2021-06-10 12:00:00+00:00,12760.0000000000)
  (2021-06-10 22:00:00+00:00,10080.2500000000)
  (2021-06-11 08:00:00+00:00,13230.5000000000)
  (2021-06-11 18:00:00+00:00,12130.7500000000)
  (2021-06-12 04:00:00+00:00,9552.7500000000)
  (2021-06-12 14:00:00+00:00,12206.5000000000)
  (2021-06-13 00:00:00+00:00,9622.7500000000)
  (2021-06-13 10:00:00+00:00,11581.0000000000)
  (2021-06-13 20:00:00+00:00,10142.2500000000)
  (2021-06-14 06:00:00+00:00,12222.7500000000)
  (2021-06-14 16:00:00+00:00,12614.0000000000)
  (2021-06-15 02:00:00+00:00,9550.0000000000)
  (2021-06-15 12:00:00+00:00,12846.2500000000)
  (2021-06-15 22:00:00+00:00,9872.7500000000)
  (2021-06-16 08:00:00+00:00,13251.7500000000)
  (2021-06-16 18:00:00+00:00,12859.5000000000)
  (2021-06-17 04:00:00+00:00,12084.5000000000)
  (2021-06-17 14:00:00+00:00,13768.5000000000)
  (2021-06-18 00:00:00+00:00,10529.2500000000)
  (2021-06-18 10:00:00+00:00,14042.7500000000)
  (2021-06-18 20:00:00+00:00,12263.2500000000)
  (2021-06-19 06:00:00+00:00,11061.7500000000)
  (2021-06-19 16:00:00+00:00,11128.2500000000)
  (2021-06-20 02:00:00+00:00,8542.7500000000)
  (2021-06-20 12:00:00+00:00,10807.7500000000)
  (2021-06-20 22:00:00+00:00,9752.5000000000)
  (2021-06-21 08:00:00+00:00,14213.0000000000)
  (2021-06-21 18:00:00+00:00,12972.5000000000)
  (2021-06-22 04:00:00+00:00,12015.5000000000)
  (2021-06-22 14:00:00+00:00,13048.0000000000)
  (2021-06-23 00:00:00+00:00,9322.5000000000)
  (2021-06-23 10:00:00+00:00,13550.2500000000)
  (2021-06-23 20:00:00+00:00,11202.0000000000)
  (2021-06-24 06:00:00+00:00,12985.2500000000)
  (2021-06-24 16:00:00+00:00,13038.2500000000)
  (2021-06-25 02:00:00+00:00,9550.5000000000)
  (2021-06-25 12:00:00+00:00,12773.7500000000)
  (2021-06-25 22:00:00+00:00,9475.0000000000)
  (2021-06-26 08:00:00+00:00,10957.0000000000)
  (2021-06-26 18:00:00+00:00,10093.7500000000)
  (2021-06-27 04:00:00+00:00,7640.2500000000)
  (2021-06-27 14:00:00+00:00,9324.2500000000)
  (2021-06-28 00:00:00+00:00,8351.5000000000)
  (2021-06-28 10:00:00+00:00,13594.2500000000)
  (2021-06-28 20:00:00+00:00,11545.7500000000)
  (2021-06-29 06:00:00+00:00,13153.0000000000)
  (2021-06-29 16:00:00+00:00,12954.7500000000)
  (2021-06-30 02:00:00+00:00,9808.0000000000)
  (2021-06-30 12:00:00+00:00,14294.5000000000)
  (2021-06-30 22:00:00+00:00,11287.0000000000)
  (2021-07-01 08:00:00+00:00,14765.0000000000)
  (2021-07-01 18:00:00+00:00,12353.5000000000)
  (2021-07-02 04:00:00+00:00,11986.2500000000)
  (2021-07-02 14:00:00+00:00,12961.2500000000)
  (2021-07-03 00:00:00+00:00,9019.0000000000)
  (2021-07-03 10:00:00+00:00,11137.2500000000)
  (2021-07-03 20:00:00+00:00,9620.0000000000)
  (2021-07-04 06:00:00+00:00,9037.5000000000)
  (2021-07-04 16:00:00+00:00,10257.5000000000)
  (2021-07-05 02:00:00+00:00,8900.0000000000)
  (2021-07-05 12:00:00+00:00,12810.0000000000)
  (2021-07-05 22:00:00+00:00,9755.2500000000)
  (2021-07-06 08:00:00+00:00,13472.5000000000)
  (2021-07-06 18:00:00+00:00,12375.2500000000)
  (2021-07-07 04:00:00+00:00,11625.5000000000)
  (2021-07-07 14:00:00+00:00,12583.5000000000)
  (2021-07-08 00:00:00+00:00,9712.2500000000)
  (2021-07-08 10:00:00+00:00,13332.5000000000)
  (2021-07-08 20:00:00+00:00,10731.0000000000)
  (2021-07-09 06:00:00+00:00,12659.7500000000)
  (2021-07-09 16:00:00+00:00,13291.2500000000)
  (2021-07-10 02:00:00+00:00,9042.7500000000)
  (2021-07-10 12:00:00+00:00,11379.7500000000)
  (2021-07-10 22:00:00+00:00,8333.0000000000)
  (2021-07-11 08:00:00+00:00,10024.2500000000)
  (2021-07-11 18:00:00+00:00,10445.0000000000)
  (2021-07-12 04:00:00+00:00,11396.5000000000)
  (2021-07-12 14:00:00+00:00,13432.2500000000)
  (2021-07-13 00:00:00+00:00,9715.5000000000)
  (2021-07-13 10:00:00+00:00,13947.2500000000)
  (2021-07-13 20:00:00+00:00,11468.2500000000)
  (2021-07-14 06:00:00+00:00,12720.5000000000)
  (2021-07-14 16:00:00+00:00,12786.2500000000)
  (2021-07-15 02:00:00+00:00,9187.7500000000)
  (2021-07-15 12:00:00+00:00,12591.0000000000)
  (2021-07-15 22:00:00+00:00,9561.7500000000)
  (2021-07-16 08:00:00+00:00,13057.0000000000)
  (2021-07-16 18:00:00+00:00,11807.5000000000)
  (2021-07-17 04:00:00+00:00,9001.7500000000)
  (2021-07-17 14:00:00+00:00,10906.5000000000)
  (2021-07-18 00:00:00+00:00,8866.0000000000)
  (2021-07-18 10:00:00+00:00,11250.0000000000)
  (2021-07-18 20:00:00+00:00,10769.5000000000)
  (2021-07-19 06:00:00+00:00,12590.5000000000)
  (2021-07-19 16:00:00+00:00,12465.5000000000)
  (2021-07-20 02:00:00+00:00,9494.2500000000)
  (2021-07-20 12:00:00+00:00,13358.5000000000)
  (2021-07-20 22:00:00+00:00,10278.2500000000)
  (2021-07-21 08:00:00+00:00,13294.7500000000)
  (2021-07-21 18:00:00+00:00,12146.0000000000)
  (2021-07-22 04:00:00+00:00,10760.2500000000)
  (2021-07-22 14:00:00+00:00,12227.5000000000)
  (2021-07-23 00:00:00+00:00,9044.2500000000)
  (2021-07-23 10:00:00+00:00,12813.0000000000)
  (2021-07-23 20:00:00+00:00,10416.2500000000)
  (2021-07-24 06:00:00+00:00,10482.5000000000)
  (2021-07-24 16:00:00+00:00,11807.2500000000)
  (2021-07-25 02:00:00+00:00,8834.2500000000)
  (2021-07-25 12:00:00+00:00,10672.0000000000)
  (2021-07-25 22:00:00+00:00,9433.2500000000)
  (2021-07-26 08:00:00+00:00,13234.5000000000)
  (2021-07-26 18:00:00+00:00,12263.2500000000)
  (2021-07-27 04:00:00+00:00,11187.0000000000)
  (2021-07-27 14:00:00+00:00,12944.2500000000)
  (2021-07-28 00:00:00+00:00,8938.2500000000)
  (2021-07-28 10:00:00+00:00,13279.7500000000)
  (2021-07-28 20:00:00+00:00,10918.5000000000)
  (2021-07-29 06:00:00+00:00,13137.0000000000)
  (2021-07-29 16:00:00+00:00,13600.2500000000)
  (2021-07-30 02:00:00+00:00,10507.7500000000)
  (2021-07-30 12:00:00+00:00,13358.7500000000)
  (2021-07-30 22:00:00+00:00,8826.5000000000)
  (2021-07-31 08:00:00+00:00,12088.7500000000)
  (2021-07-31 18:00:00+00:00,11341.0000000000)
  (2021-08-01 04:00:00+00:00,8441.0000000000)
  (2021-08-01 14:00:00+00:00,9619.7500000000)
  (2021-08-02 00:00:00+00:00,8817.0000000000)
  (2021-08-02 10:00:00+00:00,14026.7500000000)
  (2021-08-02 20:00:00+00:00,10490.5000000000)
  (2021-08-03 06:00:00+00:00,12441.5000000000)
  (2021-08-03 16:00:00+00:00,12029.5000000000)
  (2021-08-04 02:00:00+00:00,8876.2500000000)
  (2021-08-04 12:00:00+00:00,12252.0000000000)
  (2021-08-04 22:00:00+00:00,9294.5000000000)
  (2021-08-05 08:00:00+00:00,13062.0000000000)
  (2021-08-05 18:00:00+00:00,12504.5000000000)
  (2021-08-06 04:00:00+00:00,10643.0000000000)
  (2021-08-06 14:00:00+00:00,12589.5000000000)
  (2021-08-07 00:00:00+00:00,8352.0000000000)
  (2021-08-07 10:00:00+00:00,10985.7500000000)
  (2021-08-07 20:00:00+00:00,9560.5000000000)
  (2021-08-08 06:00:00+00:00,9757.5000000000)
  (2021-08-08 16:00:00+00:00,10711.0000000000)
  (2021-08-09 02:00:00+00:00,8892.2500000000)
  (2021-08-09 12:00:00+00:00,12731.7500000000)
  (2021-08-09 22:00:00+00:00,9202.0000000000)
  (2021-08-10 08:00:00+00:00,12757.0000000000)
  (2021-08-10 18:00:00+00:00,12055.0000000000)
  (2021-08-11 04:00:00+00:00,11063.0000000000)
  (2021-08-11 14:00:00+00:00,12607.5000000000)
  (2021-08-12 00:00:00+00:00,8847.2500000000)
  (2021-08-12 10:00:00+00:00,13287.2500000000)
  (2021-08-12 20:00:00+00:00,11430.0000000000)
  (2021-08-13 06:00:00+00:00,13042.0000000000)
  (2021-08-13 16:00:00+00:00,12402.5000000000)
  (2021-08-14 02:00:00+00:00,9126.5000000000)
  (2021-08-14 12:00:00+00:00,11394.2500000000)
  (2021-08-14 22:00:00+00:00,9549.5000000000)
  (2021-08-15 08:00:00+00:00,11240.0000000000)
  (2021-08-15 18:00:00+00:00,10333.7500000000)
  (2021-08-16 04:00:00+00:00,11686.5000000000)
  (2021-08-16 14:00:00+00:00,13547.2500000000)
  (2021-08-17 00:00:00+00:00,10909.7500000000)
  (2021-08-17 10:00:00+00:00,15332.0000000000)
  (2021-08-17 20:00:00+00:00,12459.2500000000)
  (2021-08-18 06:00:00+00:00,14181.5000000000)
  (2021-08-18 16:00:00+00:00,14169.7500000000)
  (2021-08-19 02:00:00+00:00,10789.0000000000)
  (2021-08-19 12:00:00+00:00,14012.0000000000)
  (2021-08-19 22:00:00+00:00,10503.2500000000)
  (2021-08-20 08:00:00+00:00,13680.2500000000)
  (2021-08-20 18:00:00+00:00,11998.5000000000)
  (2021-08-21 04:00:00+00:00,8868.2500000000)
  (2021-08-21 14:00:00+00:00,10432.5000000000)
  (2021-08-22 00:00:00+00:00,7987.7500000000)
  (2021-08-22 10:00:00+00:00,10594.7500000000)
  (2021-08-22 20:00:00+00:00,10158.7500000000)
  (2021-08-23 06:00:00+00:00,13686.5000000000)
  (2021-08-23 16:00:00+00:00,13632.2500000000)
  (2021-08-24 02:00:00+00:00,9937.0000000000)
  (2021-08-24 12:00:00+00:00,13098.5000000000)
  (2021-08-24 22:00:00+00:00,9822.5000000000)
  (2021-08-25 08:00:00+00:00,14243.2500000000)
  (2021-08-25 18:00:00+00:00,13673.7500000000)
  (2021-08-26 04:00:00+00:00,11893.2500000000)
  (2021-08-26 14:00:00+00:00,13020.7500000000)
  (2021-08-27 00:00:00+00:00,9420.5000000000)
  (2021-08-27 10:00:00+00:00,13722.7500000000)
  (2021-08-27 20:00:00+00:00,11003.7500000000)
  (2021-08-28 06:00:00+00:00,11058.5000000000)
  (2021-08-28 16:00:00+00:00,10909.5000000000)
  (2021-08-29 02:00:00+00:00,8081.5000000000)
  (2021-08-29 12:00:00+00:00,10952.7500000000)
  (2021-08-29 22:00:00+00:00,9734.0000000000)
  (2021-08-30 08:00:00+00:00,14547.2500000000)
  (2021-08-30 18:00:00+00:00,13518.2500000000)
  (2021-08-31 04:00:00+00:00,12753.2500000000)
  (2021-08-31 14:00:00+00:00,13368.2500000000)
  (2021-09-01 00:00:00+00:00,9591.2500000000)
  (2021-09-01 10:00:00+00:00,13878.5000000000)
  (2021-09-01 20:00:00+00:00,11678.2500000000)
  (2021-09-02 06:00:00+00:00,13162.5000000000)
  (2021-09-02 16:00:00+00:00,12783.7500000000)
  (2021-09-03 02:00:00+00:00,10513.7500000000)
  (2021-09-03 12:00:00+00:00,13047.2500000000)
  (2021-09-03 22:00:00+00:00,9523.5000000000)
  (2021-09-04 08:00:00+00:00,11415.5000000000)
  (2021-09-04 18:00:00+00:00,11474.2500000000)
  (2021-09-05 04:00:00+00:00,8915.2500000000)
  (2021-09-05 14:00:00+00:00,9996.2500000000)
  (2021-09-06 00:00:00+00:00,8884.7500000000)
  (2021-09-06 10:00:00+00:00,13094.7500000000)
  (2021-09-06 20:00:00+00:00,10640.7500000000)
  (2021-09-07 06:00:00+00:00,12857.5000000000)
  (2021-09-07 16:00:00+00:00,12385.7500000000)
  (2021-09-08 02:00:00+00:00,9391.0000000000)
  (2021-09-08 12:00:00+00:00,12638.5000000000)
  (2021-09-08 22:00:00+00:00,9731.0000000000)
  (2021-09-09 08:00:00+00:00,13086.0000000000)
  (2021-09-09 18:00:00+00:00,12666.0000000000)
  (2021-09-10 04:00:00+00:00,11634.7500000000)
  (2021-09-10 14:00:00+00:00,12399.7500000000)
  (2021-09-11 00:00:00+00:00,8790.7500000000)
  (2021-09-11 10:00:00+00:00,11837.5000000000)
  (2021-09-11 20:00:00+00:00,10165.7500000000)
  (2021-09-12 06:00:00+00:00,10352.0000000000)
  (2021-09-12 16:00:00+00:00,11520.7500000000)
  (2021-09-13 02:00:00+00:00,9585.7500000000)
  (2021-09-13 12:00:00+00:00,12674.5000000000)
  (2021-09-13 22:00:00+00:00,9781.5000000000)
  (2021-09-14 08:00:00+00:00,13184.5000000000)
  (2021-09-14 18:00:00+00:00,12718.2500000000)
  (2021-09-15 04:00:00+00:00,11819.2500000000)
  (2021-09-15 14:00:00+00:00,13657.7500000000)
  (2021-09-16 00:00:00+00:00,9814.2500000000)
  (2021-09-16 10:00:00+00:00,14401.2500000000)
  (2021-09-16 20:00:00+00:00,11068.2500000000)
  (2021-09-17 06:00:00+00:00,13519.5000000000)
  (2021-09-17 16:00:00+00:00,12866.5000000000)
  (2021-09-18 02:00:00+00:00,9565.2500000000)
  (2021-09-18 12:00:00+00:00,11622.0000000000)
  (2021-09-18 22:00:00+00:00,8957.7500000000)
  (2021-09-19 08:00:00+00:00,11274.2500000000)
  (2021-09-19 18:00:00+00:00,11109.7500000000)
  (2021-09-20 04:00:00+00:00,11594.2500000000)
  (2021-09-20 14:00:00+00:00,11940.5000000000)
  (2021-09-21 00:00:00+00:00,8881.0000000000)
  (2021-09-21 10:00:00+00:00,13289.7500000000)
  (2021-09-21 20:00:00+00:00,11885.2500000000)
  (2021-09-22 06:00:00+00:00,12952.2500000000)
  (2021-09-22 16:00:00+00:00,13087.2500000000)
  (2021-09-23 02:00:00+00:00,11238.0000000000)
  (2021-09-23 12:00:00+00:00,14506.2500000000)
  (2021-09-23 22:00:00+00:00,11923.2500000000)
  (2021-09-24 08:00:00+00:00,15216.5000000000)
  (2021-09-24 18:00:00+00:00,13709.5000000000)
  (2021-09-25 04:00:00+00:00,10831.7500000000)
  (2021-09-25 14:00:00+00:00,10655.2500000000)
  (2021-09-26 00:00:00+00:00,8046.2500000000)
  (2021-09-26 10:00:00+00:00,11080.7500000000)
  (2021-09-26 20:00:00+00:00,10608.0000000000)
  (2021-09-27 06:00:00+00:00,13590.2500000000)
  (2021-09-27 16:00:00+00:00,13448.5000000000)
  (2021-09-28 02:00:00+00:00,9883.2500000000)
  (2021-09-28 12:00:00+00:00,12957.2500000000)
  (2021-09-28 22:00:00+00:00,9397.7500000000)
  (2021-09-29 08:00:00+00:00,13394.0000000000)
  (2021-09-29 18:00:00+00:00,14147.2500000000)
  (2021-09-30 04:00:00+00:00,14002.2500000000)
  (2021-09-30 14:00:00+00:00,13789.7500000000)
  (2021-10-01 00:00:00+00:00,10621.5000000000)
  (2021-10-01 10:00:00+00:00,14236.2500000000)
  (2021-10-01 20:00:00+00:00,12040.2500000000)
  (2021-10-02 06:00:00+00:00,11196.2500000000)
  (2021-10-02 16:00:00+00:00,11475.0000000000)
  (2021-10-03 02:00:00+00:00,9476.7500000000)
  (2021-10-03 12:00:00+00:00,11686.2500000000)
  (2021-10-03 22:00:00+00:00,10328.0000000000)
  (2021-10-04 08:00:00+00:00,13883.0000000000)
  (2021-10-04 18:00:00+00:00,13230.7500000000)
  (2021-10-05 04:00:00+00:00,12483.5000000000)
  (2021-10-05 14:00:00+00:00,14209.0000000000)
  (2021-10-06 00:00:00+00:00,9773.0000000000)
  (2021-10-06 10:00:00+00:00,14021.2500000000)
  (2021-10-06 20:00:00+00:00,11331.5000000000)
  (2021-10-07 06:00:00+00:00,13635.2500000000)
  (2021-10-07 16:00:00+00:00,13386.5000000000)
  (2021-10-08 02:00:00+00:00,10176.7500000000)
  (2021-10-08 12:00:00+00:00,12984.5000000000)
  (2021-10-08 22:00:00+00:00,9903.7500000000)
  (2021-10-09 08:00:00+00:00,11604.2500000000)
  (2021-10-09 18:00:00+00:00,11571.7500000000)
  (2021-10-10 04:00:00+00:00,8945.7500000000)
  (2021-10-10 14:00:00+00:00,10202.5000000000)
  (2021-10-11 00:00:00+00:00,10234.0000000000)
  (2021-10-11 10:00:00+00:00,15304.5000000000)
  (2021-10-11 20:00:00+00:00,13171.7500000000)
  (2021-10-12 06:00:00+00:00,14960.0000000000)
  (2021-10-12 16:00:00+00:00,14838.2500000000)
  (2021-10-13 02:00:00+00:00,11130.5000000000)
  (2021-10-13 12:00:00+00:00,13628.5000000000)
  (2021-10-13 22:00:00+00:00,10709.2500000000)
  (2021-10-14 08:00:00+00:00,15217.5000000000)
  (2021-10-14 18:00:00+00:00,14760.5000000000)
  (2021-10-15 04:00:00+00:00,14151.7500000000)
  (2021-10-15 14:00:00+00:00,14558.0000000000)
  (2021-10-16 00:00:00+00:00,11371.5000000000)
  (2021-10-16 10:00:00+00:00,13225.0000000000)
  (2021-10-16 20:00:00+00:00,11761.2500000000)
  (2021-10-17 06:00:00+00:00,11580.2500000000)
  (2021-10-17 16:00:00+00:00,12737.5000000000)
  (2021-10-18 02:00:00+00:00,10273.0000000000)
  (2021-10-18 12:00:00+00:00,13659.2500000000)
  (2021-10-18 22:00:00+00:00,10579.2500000000)
  (2021-10-19 08:00:00+00:00,14431.0000000000)
  (2021-10-19 18:00:00+00:00,14494.0000000000)
  (2021-10-20 04:00:00+00:00,13665.7500000000)
  (2021-10-20 14:00:00+00:00,14831.2500000000)
  (2021-10-21 00:00:00+00:00,11202.7500000000)
  (2021-10-21 10:00:00+00:00,16033.2500000000)
  (2021-10-21 20:00:00+00:00,13731.5000000000)
  (2021-10-22 06:00:00+00:00,15757.2500000000)
  (2021-10-22 16:00:00+00:00,15505.2500000000)
  (2021-10-23 02:00:00+00:00,11650.2500000000)
  (2021-10-23 12:00:00+00:00,12821.0000000000)
  (2021-10-23 22:00:00+00:00,9407.5000000000)
  (2021-10-24 08:00:00+00:00,11259.5000000000)
  (2021-10-24 18:00:00+00:00,12696.0000000000)
  (2021-10-25 04:00:00+00:00,13291.5000000000)
  (2021-10-25 14:00:00+00:00,13595.0000000000)
  (2021-10-26 00:00:00+00:00,10713.7500000000)
  (2021-10-26 10:00:00+00:00,15318.2500000000)
  (2021-10-26 20:00:00+00:00,12930.5000000000)
  (2021-10-27 06:00:00+00:00,15841.0000000000)
  (2021-10-27 16:00:00+00:00,15827.7500000000)
  (2021-10-28 02:00:00+00:00,11379.5000000000)
  (2021-10-28 12:00:00+00:00,13433.0000000000)
  (2021-10-28 22:00:00+00:00,11951.7500000000)
  (2021-10-29 08:00:00+00:00,14127.7500000000)
  (2021-10-29 18:00:00+00:00,13728.7500000000)
  (2021-10-30 04:00:00+00:00,10233.7500000000)
  (2021-10-30 14:00:00+00:00,11797.2500000000)
  (2021-10-31 00:00:00+00:00,9530.2500000000)
  (2021-10-31 10:00:00+00:00,11188.2500000000)
  (2021-10-31 20:00:00+00:00,12076.0000000000)
  (2021-11-01 06:00:00+00:00,13781.7500000000)
  (2021-11-01 16:00:00+00:00,14517.7500000000)
  (2021-11-02 02:00:00+00:00,9866.2500000000)
  (2021-11-02 12:00:00+00:00,13887.0000000000)
  (2021-11-02 22:00:00+00:00,10949.2500000000)
  (2021-11-03 08:00:00+00:00,13928.0000000000)
  (2021-11-03 18:00:00+00:00,14125.5000000000)
  (2021-11-04 04:00:00+00:00,11798.0000000000)
  (2021-11-04 14:00:00+00:00,16549.7500000000)
  (2021-11-05 00:00:00+00:00,12181.0000000000)
  (2021-11-05 10:00:00+00:00,15961.2500000000)
  (2021-11-05 20:00:00+00:00,14051.5000000000)
  (2021-11-06 06:00:00+00:00,12274.0000000000)
  (2021-11-06 16:00:00+00:00,13955.5000000000)
  (2021-11-07 02:00:00+00:00,10806.5000000000)
  (2021-11-07 12:00:00+00:00,12938.0000000000)
  (2021-11-07 22:00:00+00:00,11132.0000000000)
  (2021-11-08 08:00:00+00:00,14547.2500000000)
  (2021-11-08 18:00:00+00:00,14564.0000000000)
  (2021-11-09 04:00:00+00:00,11894.7500000000)
  (2021-11-09 14:00:00+00:00,14639.0000000000)
  (2021-11-10 00:00:00+00:00,11699.2500000000)
  (2021-11-10 10:00:00+00:00,14668.5000000000)
  (2021-11-10 20:00:00+00:00,13113.5000000000)
  (2021-11-11 06:00:00+00:00,14412.5000000000)
  (2021-11-11 16:00:00+00:00,15349.0000000000)
  (2021-11-12 02:00:00+00:00,10715.7500000000)
  (2021-11-12 12:00:00+00:00,14500.0000000000)
  (2021-11-12 22:00:00+00:00,11810.5000000000)
  (2021-11-13 08:00:00+00:00,12718.5000000000)
  (2021-11-13 18:00:00+00:00,12820.0000000000)
  (2021-11-14 04:00:00+00:00,9712.5000000000)
  (2021-11-14 14:00:00+00:00,12279.2500000000)
  (2021-11-15 00:00:00+00:00,10484.5000000000)
  (2021-11-15 10:00:00+00:00,14980.5000000000)
  (2021-11-15 20:00:00+00:00,12890.0000000000)
  (2021-11-16 06:00:00+00:00,14004.2500000000)
  (2021-11-16 16:00:00+00:00,14933.7500000000)
  (2021-11-17 02:00:00+00:00,10672.5000000000)
  (2021-11-17 12:00:00+00:00,14793.7500000000)
  (2021-11-17 22:00:00+00:00,13185.0000000000)
  (2021-11-18 08:00:00+00:00,16208.2500000000)
  (2021-11-18 18:00:00+00:00,15842.2500000000)
  (2021-11-19 04:00:00+00:00,13177.2500000000)
  (2021-11-19 14:00:00+00:00,15955.5000000000)
  (2021-11-20 00:00:00+00:00,11955.2500000000)
  (2021-11-20 10:00:00+00:00,14574.0000000000)
  (2021-11-20 20:00:00+00:00,12839.2500000000)
  (2021-11-21 06:00:00+00:00,11448.5000000000)
  (2021-11-21 16:00:00+00:00,13238.2500000000)
  (2021-11-22 02:00:00+00:00,10652.0000000000)
  (2021-11-22 12:00:00+00:00,15027.7500000000)
  (2021-11-22 22:00:00+00:00,13628.5000000000)
  (2021-11-23 08:00:00+00:00,14854.0000000000)
  (2021-11-23 18:00:00+00:00,15093.5000000000)
  (2021-11-24 04:00:00+00:00,12407.2500000000)
  (2021-11-24 14:00:00+00:00,15048.5000000000)
  (2021-11-25 00:00:00+00:00,11778.7500000000)
  (2021-11-25 10:00:00+00:00,16097.7500000000)
  (2021-11-25 20:00:00+00:00,14540.0000000000)
  (2021-11-26 06:00:00+00:00,15124.5000000000)
  (2021-11-26 16:00:00+00:00,15115.2500000000)
  (2021-11-27 02:00:00+00:00,10962.2500000000)
  (2021-11-27 12:00:00+00:00,12999.0000000000)
  (2021-11-27 22:00:00+00:00,11081.5000000000)
  (2021-11-28 08:00:00+00:00,11869.2500000000)
  (2021-11-28 18:00:00+00:00,12836.0000000000)
  (2021-11-29 04:00:00+00:00,11865.0000000000)
  (2021-11-29 14:00:00+00:00,15700.5000000000)
  (2021-11-30 00:00:00+00:00,11600.5000000000)
  (2021-11-30 10:00:00+00:00,18106.0000000000)
  (2021-11-30 20:00:00+00:00,14018.5000000000)
  (2021-12-01 06:00:00+00:00,15723.5000000000)
  (2021-12-01 16:00:00+00:00,17528.5000000000)
  (2021-12-02 02:00:00+00:00,12266.0000000000)
  (2021-12-02 12:00:00+00:00,16977.7500000000)
  (2021-12-02 22:00:00+00:00,13468.5000000000)
  (2021-12-03 08:00:00+00:00,17053.7500000000)
  (2021-12-03 18:00:00+00:00,16960.2500000000)
  (2021-12-04 04:00:00+00:00,11772.0000000000)
  (2021-12-04 14:00:00+00:00,12795.5000000000)
  (2021-12-05 00:00:00+00:00,10503.5000000000)
  (2021-12-05 10:00:00+00:00,13864.2500000000)
  (2021-12-05 20:00:00+00:00,12753.2500000000)
  (2021-12-06 06:00:00+00:00,15410.5000000000)
  (2021-12-06 16:00:00+00:00,16131.5000000000)
  (2021-12-07 02:00:00+00:00,11958.5000000000)
  (2021-12-07 12:00:00+00:00,16260.0000000000)
  (2021-12-07 22:00:00+00:00,13396.2500000000)
  (2021-12-08 08:00:00+00:00,17450.2500000000)
  (2021-12-08 18:00:00+00:00,17074.5000000000)
  (2021-12-09 04:00:00+00:00,12986.7500000000)
  (2021-12-09 14:00:00+00:00,15524.2500000000)
  (2021-12-10 00:00:00+00:00,11510.5000000000)
  (2021-12-10 10:00:00+00:00,15549.5000000000)
  (2021-12-10 20:00:00+00:00,13213.2500000000)
  (2021-12-11 06:00:00+00:00,11887.5000000000)
  (2021-12-11 16:00:00+00:00,14705.2500000000)
  (2021-12-12 02:00:00+00:00,10316.5000000000)
  (2021-12-12 12:00:00+00:00,13269.0000000000)
  (2021-12-12 22:00:00+00:00,11562.7500000000)
  (2021-12-13 08:00:00+00:00,15125.5000000000)
  (2021-12-13 18:00:00+00:00,15543.5000000000)
  (2021-12-14 04:00:00+00:00,13430.2500000000)
  (2021-12-14 14:00:00+00:00,15147.7500000000)
  (2021-12-15 00:00:00+00:00,11423.5000000000)
  (2021-12-15 10:00:00+00:00,16549.0000000000)
  (2021-12-15 20:00:00+00:00,14576.5000000000)
  (2021-12-16 06:00:00+00:00,15624.2500000000)
  (2021-12-16 16:00:00+00:00,16735.7500000000)
  (2021-12-17 02:00:00+00:00,12222.2500000000)
  (2021-12-17 12:00:00+00:00,15208.5000000000)
  (2021-12-17 22:00:00+00:00,11727.0000000000)
  (2021-12-18 08:00:00+00:00,13425.0000000000)
  (2021-12-18 18:00:00+00:00,14681.5000000000)
  (2021-12-19 04:00:00+00:00,11141.7500000000)
  (2021-12-19 14:00:00+00:00,13437.2500000000)
  (2021-12-20 00:00:00+00:00,12064.7500000000)
  (2021-12-20 10:00:00+00:00,15279.7500000000)
  (2021-12-20 20:00:00+00:00,12964.7500000000)
  (2021-12-21 06:00:00+00:00,14085.5000000000)
  (2021-12-21 16:00:00+00:00,14767.0000000000)
  (2021-12-22 02:00:00+00:00,11154.2500000000)
  (2021-12-22 12:00:00+00:00,14720.5000000000)
  (2021-12-22 22:00:00+00:00,12115.7500000000)
  (2021-12-23 08:00:00+00:00,14509.5000000000)
  (2021-12-23 18:00:00+00:00,14574.2500000000)
  (2021-12-24 04:00:00+00:00,9699.2500000000)
  (2021-12-24 14:00:00+00:00,12831.5000000000)
  (2021-12-25 00:00:00+00:00,9627.0000000000)
  (2021-12-25 10:00:00+00:00,13069.7500000000)
  (2021-12-25 20:00:00+00:00,11634.2500000000)
  (2021-12-26 06:00:00+00:00,10620.5000000000)
  (2021-12-26 16:00:00+00:00,12808.7500000000)
  (2021-12-27 02:00:00+00:00,10877.0000000000)
  (2021-12-27 12:00:00+00:00,14780.7500000000)
  (2021-12-27 22:00:00+00:00,11955.5000000000)
  (2021-12-28 08:00:00+00:00,14234.2500000000)
  (2021-12-28 18:00:00+00:00,14376.0000000000)
  (2021-12-29 04:00:00+00:00,9893.2500000000)
  (2021-12-29 14:00:00+00:00,12163.2500000000)
  (2021-12-30 00:00:00+00:00,9230.0000000000)
  (2021-12-30 10:00:00+00:00,13642.7500000000)
  (2021-12-30 20:00:00+00:00,12332.2500000000)
  (2021-12-31 06:00:00+00:00,11156.7500000000)
  (2021-12-31 16:00:00+00:00,13062.7500000000)
  (2022-01-01 02:00:00+00:00,8906.5000000000)
  (2022-01-01 12:00:00+00:00,9713.7500000000)
  (2022-01-01 22:00:00+00:00,10199.0000000000)
  (2022-01-02 08:00:00+00:00,11349.0000000000)
  (2022-01-02 18:00:00+00:00,13391.2500000000)
  (2022-01-03 04:00:00+00:00,10894.2500000000)
  (2022-01-03 14:00:00+00:00,15388.2500000000)
  (2022-01-04 00:00:00+00:00,11631.5000000000)
  (2022-01-04 10:00:00+00:00,15626.5000000000)
  (2022-01-04 20:00:00+00:00,13834.0000000000)
  (2022-01-05 06:00:00+00:00,15683.2500000000)
  (2022-01-05 16:00:00+00:00,16687.2500000000)
  (2022-01-06 02:00:00+00:00,12682.5000000000)
  (2022-01-06 12:00:00+00:00,14784.7500000000)
  (2022-01-06 22:00:00+00:00,12999.7500000000)
  (2022-01-07 08:00:00+00:00,16650.7500000000)
  (2022-01-07 18:00:00+00:00,15462.2500000000)
  (2022-01-08 04:00:00+00:00,11333.2500000000)
  (2022-01-08 14:00:00+00:00,14064.5000000000)
  (2022-01-09 00:00:00+00:00,11720.0000000000)
  (2022-01-09 10:00:00+00:00,14472.5000000000)
  (2022-01-09 20:00:00+00:00,12615.5000000000)
  (2022-01-10 06:00:00+00:00,15198.0000000000)
  (2022-01-10 16:00:00+00:00,16148.5000000000)
  (2022-01-11 02:00:00+00:00,12066.5000000000)
  (2022-01-11 12:00:00+00:00,15559.7500000000)
  (2022-01-11 22:00:00+00:00,12669.7500000000)
  (2022-01-12 08:00:00+00:00,15412.2500000000)
  (2022-01-12 18:00:00+00:00,16235.5000000000)
  (2022-01-13 04:00:00+00:00,14039.5000000000)
  (2022-01-13 14:00:00+00:00,16464.5000000000)
  (2022-01-14 00:00:00+00:00,12949.7500000000)
  (2022-01-14 10:00:00+00:00,17130.0000000000)
  (2022-01-14 20:00:00+00:00,13818.5000000000)
  (2022-01-15 06:00:00+00:00,11798.5000000000)
  (2022-01-15 16:00:00+00:00,13549.7500000000)
  (2022-01-16 02:00:00+00:00,10356.5000000000)
  (2022-01-16 12:00:00+00:00,13796.2500000000)
  (2022-01-16 22:00:00+00:00,12047.2500000000)
  (2022-01-17 08:00:00+00:00,16590.5000000000)
  (2022-01-17 18:00:00+00:00,16938.7500000000)
  (2022-01-18 04:00:00+00:00,13099.2500000000)
  (2022-01-18 14:00:00+00:00,15463.2500000000)
  (2022-01-19 00:00:00+00:00,12609.2500000000)
  (2022-01-19 10:00:00+00:00,16588.0000000000)
  (2022-01-19 20:00:00+00:00,15123.5000000000)
  (2022-01-20 06:00:00+00:00,16622.2500000000)
  (2022-01-20 16:00:00+00:00,17103.7500000000)
  (2022-01-21 02:00:00+00:00,14022.7500000000)
  (2022-01-21 12:00:00+00:00,16717.0000000000)
  (2022-01-21 22:00:00+00:00,13411.7500000000)
  (2022-01-22 08:00:00+00:00,13170.2500000000)
  (2022-01-22 18:00:00+00:00,13564.5000000000)
  (2022-01-23 04:00:00+00:00,10204.5000000000)
  (2022-01-23 14:00:00+00:00,12215.7500000000)
  (2022-01-24 00:00:00+00:00,10670.0000000000)
  (2022-01-24 10:00:00+00:00,15965.7500000000)
  (2022-01-24 20:00:00+00:00,13873.0000000000)
  (2022-01-25 06:00:00+00:00,15629.7500000000)
  (2022-01-25 16:00:00+00:00,15913.2500000000)
  (2022-01-26 02:00:00+00:00,12570.2500000000)
  (2022-01-26 12:00:00+00:00,17045.0000000000)
  (2022-01-26 22:00:00+00:00,13356.2500000000)
  (2022-01-27 08:00:00+00:00,16063.0000000000)
  (2022-01-27 18:00:00+00:00,16385.2500000000)
  (2022-01-28 04:00:00+00:00,13831.7500000000)
  (2022-01-28 14:00:00+00:00,15486.0000000000)
  (2022-01-29 00:00:00+00:00,11939.7500000000)
  (2022-01-29 10:00:00+00:00,15233.2500000000)
  (2022-01-29 20:00:00+00:00,11328.5000000000)
  (2022-01-30 06:00:00+00:00,10395.2500000000)
  (2022-01-30 16:00:00+00:00,13632.2500000000)
  (2022-01-31 02:00:00+00:00,10929.5000000000)
  (2022-01-31 12:00:00+00:00,14859.7500000000)
  (2022-01-31 22:00:00+00:00,12434.5000000000)
  (2022-02-01 08:00:00+00:00,16567.2500000000)
  (2022-02-01 18:00:00+00:00,15851.2500000000)
  (2022-02-02 04:00:00+00:00,13250.2500000000)
  (2022-02-02 14:00:00+00:00,16066.7500000000)
  (2022-02-03 00:00:00+00:00,12519.0000000000)
  (2022-02-03 10:00:00+00:00,15751.2500000000)
  (2022-02-03 20:00:00+00:00,14098.2500000000)
  (2022-02-04 06:00:00+00:00,15174.7500000000)
  (2022-02-04 16:00:00+00:00,16425.0000000000)
  (2022-02-05 02:00:00+00:00,12007.5000000000)
  (2022-02-05 12:00:00+00:00,13501.7500000000)
  (2022-02-05 22:00:00+00:00,12092.2500000000)
  (2022-02-06 08:00:00+00:00,13400.7500000000)
  (2022-02-06 18:00:00+00:00,13878.7500000000)
  (2022-02-07 04:00:00+00:00,13292.0000000000)
  (2022-02-07 14:00:00+00:00,14608.2500000000)
  (2022-02-08 00:00:00+00:00,12321.2500000000)
  (2022-02-08 10:00:00+00:00,17035.0000000000)
  (2022-02-08 20:00:00+00:00,14631.7500000000)
  (2022-02-09 06:00:00+00:00,16003.7500000000)
  (2022-02-09 16:00:00+00:00,16445.5000000000)
  (2022-02-10 02:00:00+00:00,12086.0000000000)
  (2022-02-10 12:00:00+00:00,15354.5000000000)
  (2022-02-10 22:00:00+00:00,13157.5000000000)
  (2022-02-11 08:00:00+00:00,15538.0000000000)
  (2022-02-11 18:00:00+00:00,16367.0000000000)
  (2022-02-12 04:00:00+00:00,11409.0000000000)
  (2022-02-12 14:00:00+00:00,13182.0000000000)
  (2022-02-13 00:00:00+00:00,11480.2500000000)
  (2022-02-13 10:00:00+00:00,13333.5000000000)
  (2022-02-13 20:00:00+00:00,12730.0000000000)
  (2022-02-14 06:00:00+00:00,15098.5000000000)
  (2022-02-14 16:00:00+00:00,15513.7500000000)
  (2022-02-15 02:00:00+00:00,12367.2500000000)
  (2022-02-15 12:00:00+00:00,15370.2500000000)
  (2022-02-15 22:00:00+00:00,13302.5000000000)
  (2022-02-16 08:00:00+00:00,16427.7500000000)
  (2022-02-16 18:00:00+00:00,16251.7500000000)
  (2022-02-17 04:00:00+00:00,13252.7500000000)
  (2022-02-17 14:00:00+00:00,14863.0000000000)
  (2022-02-18 00:00:00+00:00,12315.0000000000)
  (2022-02-18 10:00:00+00:00,15692.2500000000)
  (2022-02-18 20:00:00+00:00,13724.5000000000)
  (2022-02-19 06:00:00+00:00,11577.7500000000)
  (2022-02-19 16:00:00+00:00,13389.0000000000)
  (2022-02-20 02:00:00+00:00,11018.5000000000)
  (2022-02-20 12:00:00+00:00,12772.2500000000)
  (2022-02-20 22:00:00+00:00,11591.7500000000)
  (2022-02-21 08:00:00+00:00,15484.5000000000)
  (2022-02-21 18:00:00+00:00,16203.7500000000)
  (2022-02-22 04:00:00+00:00,13712.2500000000)
  (2022-02-22 14:00:00+00:00,14507.5000000000)
  (2022-02-23 00:00:00+00:00,11351.0000000000)
  (2022-02-23 10:00:00+00:00,14610.0000000000)
  (2022-02-23 20:00:00+00:00,14288.2500000000)
  (2022-02-24 06:00:00+00:00,16047.0000000000)
  (2022-02-24 16:00:00+00:00,15139.0000000000)
  (2022-02-25 02:00:00+00:00,12413.0000000000)
  (2022-02-25 12:00:00+00:00,15750.2500000000)
  (2022-02-25 22:00:00+00:00,12383.5000000000)
  (2022-02-26 08:00:00+00:00,12819.5000000000)
  (2022-02-26 18:00:00+00:00,13073.2500000000)
  (2022-02-27 04:00:00+00:00,9967.5000000000)
  (2022-02-27 14:00:00+00:00,10739.5000000000)
  (2022-02-28 00:00:00+00:00,10762.5000000000)
  (2022-02-28 10:00:00+00:00,14729.5000000000)
  (2022-02-28 20:00:00+00:00,14210.7500000000)
  (2022-03-01 06:00:00+00:00,15484.7500000000)
  (2022-03-01 16:00:00+00:00,14306.2500000000)
  (2022-03-02 02:00:00+00:00,11305.2500000000)
  (2022-03-02 12:00:00+00:00,14309.2500000000)
  (2022-03-02 22:00:00+00:00,12260.7500000000)
  (2022-03-03 08:00:00+00:00,14868.2500000000)
  (2022-03-03 18:00:00+00:00,15157.0000000000)
  (2022-03-04 04:00:00+00:00,12409.5000000000)
  (2022-03-04 14:00:00+00:00,14346.2500000000)
  (2022-03-05 00:00:00+00:00,11269.7500000000)
  (2022-03-05 10:00:00+00:00,13401.0000000000)
  (2022-03-05 20:00:00+00:00,11848.0000000000)
  (2022-03-06 06:00:00+00:00,10561.0000000000)
  (2022-03-06 16:00:00+00:00,12502.7500000000)
  (2022-03-07 02:00:00+00:00,10990.2500000000)
  (2022-03-07 12:00:00+00:00,14538.7500000000)
  (2022-03-07 22:00:00+00:00,12077.0000000000)
  (2022-03-08 08:00:00+00:00,13924.7500000000)
  (2022-03-08 18:00:00+00:00,13974.2500000000)
  (2022-03-09 04:00:00+00:00,11845.7500000000)
  (2022-03-09 14:00:00+00:00,12853.7500000000)
  (2022-03-10 00:00:00+00:00,11329.2500000000)
  (2022-03-10 10:00:00+00:00,14804.2500000000)
  (2022-03-10 20:00:00+00:00,14254.5000000000)
  (2022-03-11 06:00:00+00:00,15343.2500000000)
  (2022-03-11 16:00:00+00:00,14762.0000000000)
  (2022-03-12 02:00:00+00:00,11807.5000000000)
  (2022-03-12 12:00:00+00:00,11898.2500000000)
  (2022-03-12 22:00:00+00:00,11520.5000000000)
  (2022-03-13 08:00:00+00:00,11656.0000000000)
  (2022-03-13 18:00:00+00:00,13805.2500000000)
  (2022-03-14 04:00:00+00:00,12951.7500000000)
  (2022-03-14 14:00:00+00:00,14304.2500000000)
  (2022-03-15 00:00:00+00:00,10995.2500000000)
  (2022-03-15 10:00:00+00:00,14829.0000000000)
  (2022-03-15 20:00:00+00:00,12551.2500000000)
  (2022-03-16 06:00:00+00:00,13676.0000000000)
  (2022-03-16 16:00:00+00:00,13585.5000000000)
  (2022-03-17 02:00:00+00:00,11353.7500000000)
  (2022-03-17 12:00:00+00:00,15215.7500000000)
  (2022-03-17 22:00:00+00:00,11840.7500000000)
  (2022-03-18 08:00:00+00:00,14183.5000000000)
  (2022-03-18 18:00:00+00:00,14927.0000000000)
  (2022-03-19 04:00:00+00:00,11340.7500000000)
  (2022-03-19 14:00:00+00:00,11485.2500000000)
  (2022-03-20 00:00:00+00:00,10225.7500000000)
  (2022-03-20 10:00:00+00:00,11958.7500000000)
  (2022-03-20 20:00:00+00:00,12881.7500000000)
  (2022-03-21 06:00:00+00:00,14755.0000000000)
  (2022-03-21 16:00:00+00:00,14150.5000000000)
  (2022-03-22 02:00:00+00:00,11305.0000000000)
  (2022-03-22 12:00:00+00:00,13169.2500000000)
  (2022-03-22 22:00:00+00:00,11018.7500000000)
  (2022-03-23 08:00:00+00:00,13729.7500000000)
  (2022-03-23 18:00:00+00:00,14432.0000000000)
  (2022-03-24 04:00:00+00:00,11771.0000000000)
  (2022-03-24 14:00:00+00:00,13084.7500000000)
  (2022-03-25 00:00:00+00:00,10457.2500000000)
  (2022-03-25 10:00:00+00:00,13811.2500000000)
  (2022-03-25 20:00:00+00:00,13276.5000000000)
  (2022-03-26 06:00:00+00:00,12169.2500000000)
  (2022-03-26 16:00:00+00:00,12003.0000000000)
  (2022-03-27 02:00:00+00:00,9295.2500000000)
  (2022-03-27 12:00:00+00:00,10722.0000000000)
  (2022-03-27 22:00:00+00:00,11065.0000000000)
  (2022-03-28 08:00:00+00:00,14430.2500000000)
  (2022-03-28 18:00:00+00:00,14530.5000000000)
  (2022-03-29 04:00:00+00:00,13089.0000000000)
  (2022-03-29 14:00:00+00:00,13487.5000000000)
  (2022-03-30 00:00:00+00:00,10354.5000000000)
  (2022-03-30 10:00:00+00:00,14652.5000000000)
  (2022-03-30 20:00:00+00:00,13178.2500000000)
  (2022-03-31 06:00:00+00:00,15077.2500000000)
  (2022-03-31 16:00:00+00:00,15020.7500000000)
  (2022-04-01 02:00:00+00:00,12645.5000000000)
  (2022-04-01 12:00:00+00:00,15474.7500000000)
  (2022-04-01 22:00:00+00:00,12242.2500000000)
  (2022-04-02 08:00:00+00:00,13866.7500000000)
  (2022-04-02 18:00:00+00:00,12561.0000000000)
  (2022-04-03 04:00:00+00:00,9532.5000000000)
  (2022-04-03 14:00:00+00:00,11428.5000000000)
  (2022-04-04 00:00:00+00:00,11205.5000000000)
  (2022-04-04 10:00:00+00:00,15437.7500000000)
  (2022-04-04 20:00:00+00:00,13512.5000000000)
  (2022-04-05 06:00:00+00:00,15240.2500000000)
  (2022-04-05 16:00:00+00:00,14256.7500000000)
  (2022-04-06 02:00:00+00:00,10652.0000000000)
  (2022-04-06 12:00:00+00:00,14967.0000000000)
  (2022-04-06 22:00:00+00:00,11880.0000000000)
  (2022-04-07 08:00:00+00:00,14847.7500000000)
  (2022-04-07 18:00:00+00:00,14505.2500000000)
  (2022-04-08 04:00:00+00:00,14233.7500000000)
  (2022-04-08 14:00:00+00:00,14416.5000000000)
  (2022-04-09 00:00:00+00:00,11006.2500000000)
  (2022-04-09 10:00:00+00:00,13782.7500000000)
  (2022-04-09 20:00:00+00:00,12118.7500000000)
  (2022-04-10 06:00:00+00:00,11771.5000000000)
  (2022-04-10 16:00:00+00:00,12716.0000000000)
  (2022-04-11 02:00:00+00:00,11434.2500000000)
  (2022-04-11 12:00:00+00:00,13394.0000000000)
  (2022-04-11 22:00:00+00:00,10309.0000000000)
  (2022-04-12 08:00:00+00:00,13295.2500000000)
  (2022-04-12 18:00:00+00:00,13835.0000000000)
  (2022-04-13 04:00:00+00:00,13125.2500000000)
  (2022-04-13 14:00:00+00:00,12436.5000000000)
  (2022-04-14 00:00:00+00:00,9691.0000000000)
  (2022-04-14 10:00:00+00:00,13193.2500000000)
  (2022-04-14 20:00:00+00:00,10905.7500000000)
  (2022-04-15 06:00:00+00:00,10774.0000000000)
  (2022-04-15 16:00:00+00:00,11230.0000000000)
  (2022-04-16 02:00:00+00:00,9386.2500000000)
  (2022-04-16 12:00:00+00:00,10245.2500000000)
  (2022-04-16 22:00:00+00:00,9105.7500000000)
  (2022-04-17 08:00:00+00:00,10110.5000000000)
  (2022-04-17 18:00:00+00:00,9810.7500000000)
  (2022-04-18 04:00:00+00:00,8276.7500000000)
  (2022-04-18 14:00:00+00:00,8979.0000000000)
  (2022-04-19 00:00:00+00:00,8571.0000000000)
  (2022-04-19 10:00:00+00:00,13372.5000000000)
  (2022-04-19 20:00:00+00:00,11168.7500000000)
  (2022-04-20 06:00:00+00:00,13418.2500000000)
  (2022-04-20 16:00:00+00:00,13916.0000000000)
  (2022-04-21 02:00:00+00:00,10351.2500000000)
  (2022-04-21 12:00:00+00:00,13517.0000000000)
  (2022-04-21 22:00:00+00:00,11090.2500000000)
  (2022-04-22 08:00:00+00:00,13606.0000000000)
  (2022-04-22 18:00:00+00:00,13375.0000000000)
  (2022-04-23 04:00:00+00:00,9615.5000000000)
  (2022-04-23 14:00:00+00:00,10473.7500000000)
  (2022-04-24 00:00:00+00:00,9086.7500000000)
  (2022-04-24 10:00:00+00:00,12149.0000000000)
  (2022-04-24 20:00:00+00:00,11121.7500000000)
  (2022-04-25 06:00:00+00:00,13779.0000000000)
  (2022-04-25 16:00:00+00:00,12997.7500000000)
  (2022-04-26 02:00:00+00:00,10238.0000000000)
  (2022-04-26 12:00:00+00:00,12730.5000000000)
  (2022-04-26 22:00:00+00:00,9949.0000000000)
  (2022-04-27 08:00:00+00:00,13418.0000000000)
  (2022-04-27 18:00:00+00:00,12858.2500000000)
  (2022-04-28 04:00:00+00:00,12127.0000000000)
  (2022-04-28 14:00:00+00:00,12349.0000000000)
  (2022-04-29 00:00:00+00:00,9611.0000000000)
  (2022-04-29 10:00:00+00:00,13221.0000000000)
  (2022-04-29 20:00:00+00:00,11055.0000000000)
  (2022-04-30 06:00:00+00:00,10539.0000000000)
  (2022-04-30 16:00:00+00:00,10509.2500000000)
  (2022-05-01 02:00:00+00:00,8168.7500000000)
  (2022-05-01 12:00:00+00:00,9495.7500000000)
  (2022-05-01 22:00:00+00:00,8565.0000000000)
  (2022-05-02 08:00:00+00:00,12936.7500000000)
  (2022-05-02 18:00:00+00:00,13034.0000000000)
  (2022-05-03 04:00:00+00:00,12190.0000000000)
  (2022-05-03 14:00:00+00:00,12277.7500000000)
  (2022-05-04 00:00:00+00:00,9166.5000000000)
  (2022-05-04 10:00:00+00:00,12957.7500000000)
  (2022-05-04 20:00:00+00:00,11308.5000000000)
  (2022-05-05 06:00:00+00:00,13201.5000000000)
  (2022-05-05 16:00:00+00:00,13040.0000000000)
  (2022-05-06 02:00:00+00:00,9983.2500000000)
  (2022-05-06 12:00:00+00:00,12524.5000000000)
  (2022-05-06 22:00:00+00:00,10001.7500000000)
  (2022-05-07 08:00:00+00:00,11472.2500000000)
  (2022-05-07 18:00:00+00:00,10429.2500000000)
  (2022-05-08 04:00:00+00:00,8115.2500000000)
  (2022-05-08 14:00:00+00:00,9354.0000000000)
  (2022-05-09 00:00:00+00:00,8719.7500000000)
  (2022-05-09 10:00:00+00:00,12852.0000000000)
  (2022-05-09 20:00:00+00:00,11865.5000000000)
  (2022-05-10 06:00:00+00:00,13403.0000000000)
  (2022-05-10 16:00:00+00:00,12728.2500000000)
  (2022-05-11 02:00:00+00:00,9438.7500000000)
  (2022-05-11 12:00:00+00:00,13336.5000000000)
  (2022-05-11 22:00:00+00:00,9806.2500000000)
  (2022-05-12 08:00:00+00:00,13647.5000000000)
  (2022-05-12 18:00:00+00:00,12949.7500000000)
  (2022-05-13 04:00:00+00:00,11642.0000000000)
  (2022-05-13 14:00:00+00:00,13505.2500000000)
  (2022-05-14 00:00:00+00:00,9780.7500000000)
  (2022-05-14 10:00:00+00:00,11386.5000000000)
  (2022-05-14 20:00:00+00:00,10464.5000000000)
  (2022-05-15 06:00:00+00:00,9563.0000000000)
  (2022-05-15 16:00:00+00:00,10248.5000000000)
  (2022-05-16 02:00:00+00:00,8968.5000000000)
  (2022-05-16 12:00:00+00:00,12691.2500000000)
  (2022-05-16 22:00:00+00:00,10110.5000000000)
  (2022-05-17 08:00:00+00:00,13657.2500000000)
  (2022-05-17 18:00:00+00:00,13151.0000000000)
  (2022-05-18 04:00:00+00:00,12526.7500000000)
  (2022-05-18 14:00:00+00:00,12858.2500000000)
  (2022-05-19 00:00:00+00:00,10506.7500000000)
  (2022-05-19 10:00:00+00:00,13281.2500000000)
  (2022-05-19 20:00:00+00:00,11839.2500000000)
  (2022-05-20 06:00:00+00:00,13419.5000000000)
  (2022-05-20 16:00:00+00:00,12708.0000000000)
  (2022-05-21 02:00:00+00:00,10317.2500000000)
  (2022-05-21 12:00:00+00:00,11784.0000000000)
  (2022-05-21 22:00:00+00:00,8836.2500000000)
  (2022-05-22 08:00:00+00:00,9915.7500000000)
  (2022-05-22 18:00:00+00:00,9934.0000000000)
  (2022-05-23 04:00:00+00:00,10814.7500000000)
  (2022-05-23 14:00:00+00:00,12787.5000000000)
  (2022-05-24 00:00:00+00:00,10193.5000000000)
  (2022-05-24 10:00:00+00:00,13062.7500000000)
  (2022-05-24 20:00:00+00:00,11240.7500000000)
  (2022-05-25 06:00:00+00:00,12813.5000000000)
  (2022-05-25 16:00:00+00:00,11931.5000000000)
  (2022-05-26 02:00:00+00:00,8837.2500000000)
  (2022-05-26 12:00:00+00:00,10526.7500000000)
  (2022-05-26 22:00:00+00:00,9059.5000000000)
  (2022-05-27 08:00:00+00:00,12203.7500000000)
  (2022-05-27 18:00:00+00:00,11526.0000000000)
  (2022-05-28 04:00:00+00:00,9413.2500000000)
  (2022-05-28 14:00:00+00:00,10868.2500000000)
  (2022-05-29 00:00:00+00:00,8574.5000000000)
  (2022-05-29 10:00:00+00:00,10310.7500000000)
  (2022-05-29 20:00:00+00:00,10006.2500000000)
  (2022-05-30 06:00:00+00:00,12850.0000000000)
  (2022-05-30 16:00:00+00:00,12721.5000000000)
  (2022-05-31 02:00:00+00:00,9635.2500000000)
  (2022-05-31 12:00:00+00:00,12557.2500000000)
  (2022-05-31 22:00:00+00:00,9545.5000000000)
  (2022-06-01 08:00:00+00:00,12900.7500000000)
  (2022-06-01 18:00:00+00:00,12057.7500000000)
  (2022-06-02 04:00:00+00:00,11645.5000000000)
  (2022-06-02 14:00:00+00:00,13329.7500000000)
  (2022-06-03 00:00:00+00:00,9023.0000000000)
  (2022-06-03 10:00:00+00:00,12324.2500000000)
  (2022-06-03 20:00:00+00:00,10116.2500000000)
  (2022-06-04 06:00:00+00:00,9831.7500000000)
  (2022-06-04 16:00:00+00:00,9834.7500000000)
  (2022-06-05 02:00:00+00:00,7143.0000000000)
  (2022-06-05 12:00:00+00:00,9084.0000000000)
  (2022-06-05 22:00:00+00:00,8820.0000000000)
  (2022-06-06 08:00:00+00:00,9871.7500000000)
  (2022-06-06 18:00:00+00:00,9533.0000000000)
  (2022-06-07 04:00:00+00:00,10511.7500000000)
  (2022-06-07 14:00:00+00:00,12274.5000000000)
  (2022-06-08 00:00:00+00:00,8516.7500000000)
  (2022-06-08 10:00:00+00:00,12365.7500000000)
  (2022-06-08 20:00:00+00:00,10665.2500000000)
  (2022-06-09 06:00:00+00:00,12300.7500000000)
  (2022-06-09 16:00:00+00:00,12518.7500000000)
  (2022-06-10 02:00:00+00:00,9765.7500000000)
  (2022-06-10 12:00:00+00:00,12044.7500000000)
  (2022-06-10 22:00:00+00:00,9036.7500000000)
  (2022-06-11 08:00:00+00:00,10845.2500000000)
  (2022-06-11 18:00:00+00:00,10148.2500000000)
  (2022-06-12 04:00:00+00:00,7902.0000000000)
  (2022-06-12 14:00:00+00:00,9538.5000000000)
  (2022-06-13 00:00:00+00:00,8504.2500000000)
  (2022-06-13 10:00:00+00:00,13503.7500000000)
  (2022-06-13 20:00:00+00:00,11328.2500000000)
  (2022-06-14 06:00:00+00:00,13088.5000000000)
  (2022-06-14 16:00:00+00:00,12783.0000000000)
  (2022-06-15 02:00:00+00:00,9108.2500000000)
  (2022-06-15 12:00:00+00:00,12454.0000000000)
  (2022-06-15 22:00:00+00:00,9462.5000000000)
  (2022-06-16 08:00:00+00:00,12569.0000000000)
  (2022-06-16 18:00:00+00:00,12093.5000000000)
  (2022-06-17 04:00:00+00:00,10908.5000000000)
  (2022-06-17 14:00:00+00:00,11968.2500000000)
  (2022-06-18 00:00:00+00:00,9020.5000000000)
  (2022-06-18 10:00:00+00:00,11633.5000000000)
  (2022-06-18 20:00:00+00:00,9554.0000000000)
  (2022-06-19 06:00:00+00:00,10010.7500000000)
  (2022-06-19 16:00:00+00:00,10615.5000000000)
  (2022-06-20 02:00:00+00:00,8384.5000000000)
  (2022-06-20 12:00:00+00:00,12375.2500000000)
  (2022-06-20 22:00:00+00:00,10054.2500000000)
  (2022-06-21 08:00:00+00:00,12862.5000000000)
  (2022-06-21 18:00:00+00:00,11757.2500000000)
  (2022-06-22 04:00:00+00:00,10975.5000000000)
  (2022-06-22 14:00:00+00:00,12081.2500000000)
  (2022-06-23 00:00:00+00:00,8717.7500000000)
  (2022-06-23 10:00:00+00:00,12572.5000000000)
  (2022-06-23 20:00:00+00:00,11018.2500000000)
  (2022-06-24 06:00:00+00:00,12631.0000000000)
  (2022-06-24 16:00:00+00:00,13142.2500000000)
  (2022-06-25 02:00:00+00:00,8630.0000000000)
  (2022-06-25 12:00:00+00:00,10352.0000000000)
  (2022-06-25 22:00:00+00:00,8769.2500000000)
  (2022-06-26 08:00:00+00:00,10206.0000000000)
  (2022-06-26 18:00:00+00:00,10261.5000000000)
  (2022-06-27 04:00:00+00:00,10830.5000000000)
  (2022-06-27 14:00:00+00:00,13014.5000000000)
  (2022-06-28 00:00:00+00:00,9340.0000000000)
  (2022-06-28 10:00:00+00:00,12825.7500000000)
  (2022-06-28 20:00:00+00:00,10939.0000000000)
  (2022-06-29 06:00:00+00:00,12433.2500000000)
  (2022-06-29 16:00:00+00:00,12442.2500000000)
  (2022-06-30 02:00:00+00:00,9449.5000000000)
  (2022-06-30 12:00:00+00:00,12768.7500000000)
  (2022-06-30 22:00:00+00:00,9910.7500000000)
  (2022-07-01 08:00:00+00:00,12975.5000000000)
  (2022-07-01 18:00:00+00:00,11505.0000000000)
  (2022-07-02 04:00:00+00:00,8867.0000000000)
  (2022-07-02 14:00:00+00:00,9871.7500000000)
  (2022-07-03 00:00:00+00:00,8521.2500000000)
  (2022-07-03 10:00:00+00:00,10539.2500000000)
  (2022-07-03 20:00:00+00:00,10292.2500000000)
  (2022-07-04 06:00:00+00:00,12082.0000000000)
  (2022-07-04 16:00:00+00:00,12252.2500000000)
  (2022-07-05 02:00:00+00:00,9080.0000000000)
  (2022-07-05 12:00:00+00:00,12851.5000000000)
  (2022-07-05 22:00:00+00:00,9566.5000000000)
  (2022-07-06 08:00:00+00:00,12745.7500000000)
  (2022-07-06 18:00:00+00:00,11918.5000000000)
  (2022-07-07 04:00:00+00:00,11572.2500000000)
  (2022-07-07 14:00:00+00:00,12722.7500000000)
  (2022-07-08 00:00:00+00:00,9492.7500000000)
  (2022-07-08 10:00:00+00:00,13221.7500000000)
  (2022-07-08 20:00:00+00:00,10480.5000000000)
  (2022-07-09 06:00:00+00:00,10189.7500000000)
  (2022-07-09 16:00:00+00:00,11338.5000000000)
  (2022-07-10 02:00:00+00:00,8269.0000000000)
  (2022-07-10 12:00:00+00:00,10256.7500000000)
  (2022-07-10 22:00:00+00:00,9203.2500000000)
  (2022-07-11 08:00:00+00:00,12181.0000000000)
  (2022-07-11 18:00:00+00:00,11170.5000000000)
  (2022-07-12 04:00:00+00:00,10375.2500000000)
  (2022-07-12 14:00:00+00:00,12189.0000000000)
  (2022-07-13 00:00:00+00:00,8735.2500000000)
  (2022-07-13 10:00:00+00:00,13199.2500000000)
  (2022-07-13 20:00:00+00:00,11444.5000000000)
  (2022-07-14 06:00:00+00:00,12346.7500000000)
  (2022-07-14 16:00:00+00:00,12669.5000000000)
  (2022-07-15 02:00:00+00:00,9678.7500000000)
  (2022-07-15 12:00:00+00:00,12733.2500000000)
  (2022-07-15 22:00:00+00:00,8830.2500000000)
  (2022-07-16 08:00:00+00:00,11530.2500000000)
  (2022-07-16 18:00:00+00:00,10049.2500000000)
  (2022-07-17 04:00:00+00:00,7772.7500000000)
  (2022-07-17 14:00:00+00:00,9376.5000000000)
  (2022-07-18 00:00:00+00:00,8005.5000000000)
  (2022-07-18 10:00:00+00:00,12323.5000000000)
  (2022-07-18 20:00:00+00:00,10053.2500000000)
  (2022-07-19 06:00:00+00:00,11824.5000000000)
  (2022-07-19 16:00:00+00:00,11721.5000000000)
  (2022-07-20 02:00:00+00:00,9186.5000000000)
  (2022-07-20 12:00:00+00:00,12448.0000000000)
  (2022-07-20 22:00:00+00:00,10789.2500000000)
  (2022-07-21 08:00:00+00:00,13360.2500000000)
  (2022-07-21 18:00:00+00:00,11273.2500000000)
  (2022-07-22 04:00:00+00:00,10550.5000000000)
  (2022-07-22 14:00:00+00:00,11863.7500000000)
  (2022-07-23 00:00:00+00:00,8133.0000000000)
  (2022-07-23 10:00:00+00:00,10917.2500000000)
  (2022-07-23 20:00:00+00:00,9238.7500000000)
  (2022-07-24 06:00:00+00:00,8720.5000000000)
  (2022-07-24 16:00:00+00:00,9579.0000000000)
  (2022-07-25 02:00:00+00:00,9356.0000000000)
  (2022-07-25 12:00:00+00:00,13025.2500000000)
  (2022-07-25 22:00:00+00:00,9442.2500000000)
  (2022-07-26 08:00:00+00:00,12860.5000000000)
  (2022-07-26 18:00:00+00:00,11008.0000000000)
  (2022-07-27 04:00:00+00:00,10386.0000000000)
  (2022-07-27 14:00:00+00:00,11940.7500000000)
  (2022-07-28 00:00:00+00:00,8801.2500000000)
  (2022-07-28 10:00:00+00:00,12444.2500000000)
  (2022-07-28 20:00:00+00:00,10575.2500000000)
  (2022-07-29 06:00:00+00:00,11704.7500000000)
  (2022-07-29 16:00:00+00:00,11912.2500000000)
  (2022-07-30 02:00:00+00:00,8087.5000000000)
  (2022-07-30 12:00:00+00:00,9850.7500000000)
  (2022-07-30 22:00:00+00:00,7783.0000000000)
  (2022-07-31 08:00:00+00:00,9558.5000000000)
  (2022-07-31 18:00:00+00:00,9126.0000000000)
  (2022-08-01 04:00:00+00:00,10201.2500000000)
  (2022-08-01 14:00:00+00:00,11772.5000000000)
  (2022-08-02 00:00:00+00:00,8856.2500000000)
  (2022-08-02 10:00:00+00:00,12200.7500000000)
  (2022-08-02 20:00:00+00:00,9986.2500000000)
  (2022-08-03 06:00:00+00:00,11760.5000000000)
  (2022-08-03 16:00:00+00:00,11857.0000000000)
  (2022-08-04 02:00:00+00:00,8707.7500000000)
  (2022-08-04 12:00:00+00:00,12138.5000000000)
  (2022-08-04 22:00:00+00:00,9216.0000000000)
  (2022-08-05 08:00:00+00:00,12085.2500000000)
  (2022-08-05 18:00:00+00:00,10787.5000000000)
  (2022-08-06 04:00:00+00:00,9215.0000000000)
  (2022-08-06 14:00:00+00:00,10289.5000000000)
  (2022-08-07 00:00:00+00:00,7482.0000000000)
  (2022-08-07 10:00:00+00:00,9473.7500000000)
  (2022-08-07 20:00:00+00:00,8927.5000000000)
  (2022-08-08 06:00:00+00:00,10959.2500000000)
  (2022-08-08 16:00:00+00:00,11368.2500000000)
  (2022-08-09 02:00:00+00:00,8439.2500000000)
  (2022-08-09 12:00:00+00:00,11395.7500000000)
  (2022-08-09 22:00:00+00:00,8839.7500000000)
  (2022-08-10 08:00:00+00:00,11831.2500000000)
  (2022-08-10 18:00:00+00:00,11065.5000000000)
  (2022-08-11 04:00:00+00:00,10088.5000000000)
  (2022-08-11 14:00:00+00:00,11555.5000000000)
  (2022-08-12 00:00:00+00:00,8308.5000000000)
  (2022-08-12 10:00:00+00:00,12146.7500000000)
  (2022-08-12 20:00:00+00:00,9875.2500000000)
  (2022-08-13 06:00:00+00:00,9408.2500000000)
  (2022-08-13 16:00:00+00:00,10274.0000000000)
  (2022-08-14 02:00:00+00:00,7621.5000000000)
  (2022-08-14 12:00:00+00:00,9612.2500000000)
  (2022-08-14 22:00:00+00:00,8997.0000000000)
  (2022-08-15 08:00:00+00:00,12356.5000000000)
  (2022-08-15 18:00:00+00:00,12283.0000000000)
  (2022-08-16 04:00:00+00:00,10421.7500000000)
  (2022-08-16 14:00:00+00:00,11880.7500000000)
  (2022-08-17 00:00:00+00:00,8447.2500000000)
  (2022-08-17 10:00:00+00:00,12601.7500000000)
  (2022-08-17 20:00:00+00:00,10484.7500000000)
  (2022-08-18 06:00:00+00:00,12150.5000000000)
  (2022-08-18 16:00:00+00:00,12537.5000000000)
  (2022-08-19 02:00:00+00:00,9003.2500000000)
  (2022-08-19 12:00:00+00:00,11975.5000000000)
  (2022-08-19 22:00:00+00:00,8662.2500000000)
  (2022-08-20 08:00:00+00:00,10550.7500000000)
  (2022-08-20 18:00:00+00:00,9798.5000000000)
  (2022-08-21 04:00:00+00:00,7621.0000000000)
  (2022-08-21 14:00:00+00:00,8979.2500000000)
  (2022-08-22 00:00:00+00:00,8212.5000000000)
  (2022-08-22 10:00:00+00:00,12327.0000000000)
  (2022-08-22 20:00:00+00:00,10527.2500000000)
  (2022-08-23 06:00:00+00:00,12271.7500000000)
  (2022-08-23 16:00:00+00:00,12232.0000000000)
  (2022-08-24 02:00:00+00:00,9118.7500000000)
  (2022-08-24 12:00:00+00:00,12187.0000000000)
  (2022-08-24 22:00:00+00:00,9515.5000000000)
  (2022-08-25 08:00:00+00:00,13058.7500000000)
  (2022-08-25 18:00:00+00:00,12693.2500000000)
  (2022-08-26 04:00:00+00:00,10687.0000000000)
  (2022-08-26 14:00:00+00:00,11899.0000000000)
  (2022-08-27 00:00:00+00:00,8000.5000000000)
  (2022-08-27 10:00:00+00:00,11006.5000000000)
  (2022-08-27 20:00:00+00:00,9184.0000000000)
  (2022-08-28 06:00:00+00:00,9059.2500000000)
  (2022-08-28 16:00:00+00:00,10447.2500000000)
  (2022-08-29 02:00:00+00:00,8781.0000000000)
  (2022-08-29 12:00:00+00:00,12100.0000000000)
  (2022-08-29 22:00:00+00:00,8857.5000000000)
  (2022-08-30 08:00:00+00:00,11990.5000000000)
  (2022-08-30 18:00:00+00:00,11369.0000000000)
  (2022-08-31 04:00:00+00:00,10796.0000000000)
  (2022-08-31 14:00:00+00:00,12096.0000000000)
  (2022-09-01 00:00:00+00:00,8895.7500000000)
  (2022-09-01 10:00:00+00:00,12536.2500000000)
  (2022-09-01 20:00:00+00:00,10380.7500000000)
  (2022-09-02 06:00:00+00:00,12246.7500000000)
  (2022-09-02 16:00:00+00:00,11706.5000000000)
  (2022-09-03 02:00:00+00:00,8879.2500000000)
  (2022-09-03 12:00:00+00:00,10640.2500000000)
  (2022-09-03 22:00:00+00:00,9397.5000000000)
  (2022-09-04 08:00:00+00:00,10471.7500000000)
  (2022-09-04 18:00:00+00:00,10955.5000000000)
  (2022-09-05 04:00:00+00:00,11332.5000000000)
  (2022-09-05 14:00:00+00:00,12583.7500000000)
  (2022-09-06 00:00:00+00:00,9980.7500000000)
  (2022-09-06 10:00:00+00:00,13165.5000000000)
  (2022-09-06 20:00:00+00:00,10855.7500000000)
  (2022-09-07 06:00:00+00:00,12007.0000000000)
  (2022-09-07 16:00:00+00:00,12982.7500000000)
  (2022-09-08 02:00:00+00:00,10084.2500000000)
  (2022-09-08 12:00:00+00:00,13464.0000000000)
  (2022-09-08 22:00:00+00:00,9931.2500000000)
  (2022-09-09 08:00:00+00:00,12090.2500000000)
  (2022-09-09 18:00:00+00:00,11247.2500000000)
  (2022-09-10 04:00:00+00:00,8564.5000000000)
  (2022-09-10 14:00:00+00:00,9676.7500000000)
  (2022-09-11 00:00:00+00:00,7405.0000000000)
  (2022-09-11 10:00:00+00:00,10033.0000000000)
  (2022-09-11 20:00:00+00:00,8893.2500000000)
  (2022-09-12 06:00:00+00:00,11734.7500000000)
  (2022-09-12 16:00:00+00:00,11527.5000000000)
  (2022-09-13 02:00:00+00:00,9857.7500000000)
  (2022-09-13 12:00:00+00:00,12700.5000000000)
  (2022-09-13 22:00:00+00:00,9826.7500000000)
  (2022-09-14 08:00:00+00:00,12973.7500000000)
  (2022-09-14 18:00:00+00:00,12165.7500000000)
  (2022-09-15 04:00:00+00:00,11879.2500000000)
  (2022-09-15 14:00:00+00:00,12950.7500000000)
  (2022-09-16 00:00:00+00:00,10079.7500000000)
  (2022-09-16 10:00:00+00:00,13975.7500000000)
  (2022-09-16 20:00:00+00:00,10718.2500000000)
  (2022-09-17 06:00:00+00:00,10580.0000000000)
  (2022-09-17 16:00:00+00:00,10807.0000000000)
  (2022-09-18 02:00:00+00:00,8482.7500000000)
  (2022-09-18 12:00:00+00:00,10045.2500000000)
  (2022-09-18 22:00:00+00:00,8424.5000000000)
  (2022-09-19 08:00:00+00:00,12657.0000000000)
  (2022-09-19 18:00:00+00:00,12443.5000000000)
  (2022-09-20 04:00:00+00:00,11189.7500000000)
  (2022-09-20 14:00:00+00:00,11271.5000000000)
  (2022-09-21 00:00:00+00:00,8373.5000000000)
  (2022-09-21 10:00:00+00:00,11979.5000000000)
  (2022-09-21 20:00:00+00:00,10093.7500000000)
  (2022-09-22 06:00:00+00:00,12214.7500000000)
  (2022-09-22 16:00:00+00:00,11638.5000000000)
  (2022-09-23 02:00:00+00:00,9268.5000000000)
  (2022-09-23 12:00:00+00:00,11573.2500000000)
  (2022-09-23 22:00:00+00:00,8728.5000000000)
  (2022-09-24 08:00:00+00:00,10602.2500000000)
  (2022-09-24 18:00:00+00:00,10142.5000000000)
  (2022-09-25 04:00:00+00:00,7744.5000000000)
  (2022-09-25 14:00:00+00:00,9217.2500000000)
  (2022-09-26 00:00:00+00:00,8467.0000000000)
  (2022-09-26 10:00:00+00:00,13064.0000000000)
  (2022-09-26 20:00:00+00:00,11390.0000000000)
  (2022-09-27 06:00:00+00:00,13128.5000000000)
  (2022-09-27 16:00:00+00:00,12706.5000000000)
  (2022-09-28 02:00:00+00:00,9267.0000000000)
  (2022-09-28 12:00:00+00:00,12130.5000000000)
  (2022-09-28 22:00:00+00:00,9262.7500000000)
  (2022-09-29 08:00:00+00:00,12503.7500000000)
  (2022-09-29 18:00:00+00:00,12036.5000000000)
  (2022-09-30 04:00:00+00:00,11235.0000000000)
  (2022-09-30 14:00:00+00:00,11490.5000000000)
  (2022-10-01 00:00:00+00:00,9208.2500000000)
  (2022-10-01 10:00:00+00:00,12770.0000000000)
  (2022-10-01 20:00:00+00:00,10762.7500000000)
  (2022-10-02 06:00:00+00:00,10673.0000000000)
  (2022-10-02 16:00:00+00:00,11349.2500000000)
  (2022-10-03 02:00:00+00:00,9204.0000000000)
  (2022-10-03 12:00:00+00:00,10589.5000000000)
  (2022-10-03 22:00:00+00:00,9839.5000000000)
  (2022-10-04 08:00:00+00:00,12253.2500000000)
  (2022-10-04 18:00:00+00:00,12344.5000000000)
  (2022-10-05 04:00:00+00:00,12220.2500000000)
  (2022-10-05 14:00:00+00:00,12693.7500000000)
  (2022-10-06 00:00:00+00:00,10645.0000000000)
  (2022-10-06 10:00:00+00:00,13319.5000000000)
  (2022-10-06 20:00:00+00:00,11522.0000000000)
  (2022-10-07 06:00:00+00:00,13345.2500000000)
  (2022-10-07 16:00:00+00:00,12643.5000000000)
  (2022-10-08 02:00:00+00:00,9468.0000000000)
  (2022-10-08 12:00:00+00:00,11424.7500000000)
  (2022-10-08 22:00:00+00:00,9753.0000000000)
  (2022-10-09 08:00:00+00:00,10043.5000000000)
  (2022-10-09 18:00:00+00:00,10602.2500000000)
  (2022-10-10 04:00:00+00:00,12789.7500000000)
  (2022-10-10 14:00:00+00:00,12059.0000000000)
  (2022-10-11 00:00:00+00:00,10160.2500000000)
  (2022-10-11 10:00:00+00:00,13373.2500000000)
  (2022-10-11 20:00:00+00:00,10597.0000000000)
  (2022-10-12 06:00:00+00:00,12349.0000000000)
  (2022-10-12 16:00:00+00:00,12275.7500000000)
  (2022-10-13 02:00:00+00:00,9440.7500000000)
  (2022-10-13 12:00:00+00:00,12062.2500000000)
  (2022-10-13 22:00:00+00:00,9368.5000000000)
  (2022-10-14 08:00:00+00:00,12522.7500000000)
  (2022-10-14 18:00:00+00:00,11808.2500000000)
  (2022-10-15 04:00:00+00:00,9232.2500000000)
  (2022-10-15 14:00:00+00:00,9941.0000000000)
  (2022-10-16 00:00:00+00:00,8927.2500000000)
  (2022-10-16 10:00:00+00:00,11304.7500000000)
  (2022-10-16 20:00:00+00:00,9767.2500000000)
  (2022-10-17 06:00:00+00:00,12734.5000000000)
  (2022-10-17 16:00:00+00:00,13280.0000000000)
  (2022-10-18 02:00:00+00:00,9462.0000000000)
  (2022-10-18 12:00:00+00:00,12857.5000000000)
  (2022-10-18 22:00:00+00:00,10338.5000000000)
  (2022-10-19 08:00:00+00:00,12393.2500000000)
  (2022-10-19 18:00:00+00:00,11708.5000000000)
  (2022-10-20 04:00:00+00:00,11928.0000000000)
  (2022-10-20 14:00:00+00:00,12897.7500000000)
  (2022-10-21 00:00:00+00:00,10116.7500000000)
  (2022-10-21 10:00:00+00:00,12681.0000000000)
  (2022-10-21 20:00:00+00:00,9930.0000000000)
  (2022-10-22 06:00:00+00:00,10200.7500000000)
  (2022-10-22 16:00:00+00:00,10887.2500000000)
  (2022-10-23 02:00:00+00:00,8161.5000000000)
  (2022-10-23 12:00:00+00:00,9303.7500000000)
  (2022-10-23 22:00:00+00:00,9406.7500000000)
  (2022-10-24 08:00:00+00:00,13259.0000000000)
  (2022-10-24 18:00:00+00:00,13073.2500000000)
  (2022-10-25 04:00:00+00:00,11917.2500000000)
  (2022-10-25 14:00:00+00:00,12545.7500000000)
  (2022-10-26 00:00:00+00:00,8869.0000000000)
  (2022-10-26 10:00:00+00:00,12155.0000000000)
  (2022-10-26 20:00:00+00:00,10636.5000000000)
  (2022-10-27 06:00:00+00:00,12984.7500000000)
  (2022-10-27 16:00:00+00:00,12852.0000000000)
  (2022-10-28 02:00:00+00:00,9871.5000000000)
  (2022-10-28 12:00:00+00:00,11664.5000000000)
  (2022-10-28 22:00:00+00:00,10345.5000000000)
  (2022-10-29 08:00:00+00:00,10942.5000000000)
  (2022-10-29 18:00:00+00:00,10028.0000000000)
  (2022-10-30 04:00:00+00:00,7664.0000000000)
  (2022-10-30 14:00:00+00:00,9086.2500000000)
  (2022-10-31 00:00:00+00:00,7733.5000000000)
  (2022-10-31 10:00:00+00:00,10371.7500000000)
  (2022-10-31 20:00:00+00:00,9408.0000000000)
  (2022-11-01 06:00:00+00:00,12420.0000000000)
  (2022-11-01 16:00:00+00:00,13094.7500000000)
  (2022-11-02 02:00:00+00:00,9669.5000000000)
  (2022-11-02 12:00:00+00:00,13191.0000000000)
  (2022-11-02 22:00:00+00:00,10682.2500000000)
  (2022-11-03 08:00:00+00:00,13712.0000000000)
  (2022-11-03 18:00:00+00:00,14299.7500000000)
  (2022-11-04 04:00:00+00:00,10901.5000000000)
  (2022-11-04 14:00:00+00:00,12265.7500000000)
  (2022-11-05 00:00:00+00:00,9392.2500000000)
  (2022-11-05 10:00:00+00:00,12137.5000000000)
  (2022-11-05 20:00:00+00:00,11257.7500000000)
  (2022-11-06 06:00:00+00:00,10538.0000000000)
  (2022-11-06 16:00:00+00:00,12462.7500000000)
  (2022-11-07 02:00:00+00:00,10163.0000000000)
  (2022-11-07 12:00:00+00:00,14101.7500000000)
  (2022-11-07 22:00:00+00:00,11470.5000000000)
  (2022-11-08 08:00:00+00:00,13692.2500000000)
  (2022-11-08 18:00:00+00:00,14703.5000000000)
  (2022-11-09 04:00:00+00:00,11534.7500000000)
  (2022-11-09 14:00:00+00:00,12830.0000000000)
  (2022-11-10 00:00:00+00:00,10910.7500000000)
  (2022-11-10 10:00:00+00:00,14071.5000000000)
  (2022-11-10 20:00:00+00:00,12783.2500000000)
  (2022-11-11 06:00:00+00:00,14005.0000000000)
  (2022-11-11 16:00:00+00:00,14218.2500000000)
  (2022-11-12 02:00:00+00:00,9563.0000000000)
  (2022-11-12 12:00:00+00:00,11251.5000000000)
  (2022-11-12 22:00:00+00:00,9189.0000000000)
  (2022-11-13 08:00:00+00:00,9867.0000000000)
  (2022-11-13 18:00:00+00:00,11280.2500000000)
  (2022-11-14 04:00:00+00:00,10394.2500000000)
  (2022-11-14 14:00:00+00:00,12646.7500000000)
  (2022-11-15 00:00:00+00:00,9562.2500000000)
  (2022-11-15 10:00:00+00:00,12643.5000000000)
  (2022-11-15 20:00:00+00:00,12106.5000000000)
  (2022-11-16 06:00:00+00:00,12955.5000000000)
  (2022-11-16 16:00:00+00:00,14903.2500000000)
  (2022-11-17 02:00:00+00:00,10558.2500000000)
  (2022-11-17 12:00:00+00:00,14706.7500000000)
  (2022-11-17 22:00:00+00:00,11329.7500000000)
  (2022-11-18 08:00:00+00:00,14822.0000000000)
  (2022-11-18 18:00:00+00:00,13967.7500000000)
  (2022-11-19 04:00:00+00:00,10943.0000000000)
  (2022-11-19 14:00:00+00:00,11954.0000000000)
  (2022-11-20 00:00:00+00:00,9894.0000000000)
  (2022-11-20 10:00:00+00:00,11933.7500000000)
  (2022-11-20 20:00:00+00:00,10674.5000000000)
  (2022-11-21 06:00:00+00:00,13180.2500000000)
  (2022-11-21 16:00:00+00:00,14604.5000000000)
  (2022-11-22 02:00:00+00:00,10821.0000000000)
  (2022-11-22 12:00:00+00:00,14827.7500000000)
  (2022-11-22 22:00:00+00:00,12027.5000000000)
  (2022-11-23 08:00:00+00:00,14366.7500000000)
  (2022-11-23 18:00:00+00:00,14769.0000000000)
  (2022-11-24 04:00:00+00:00,11352.0000000000)
  (2022-11-24 14:00:00+00:00,13466.7500000000)
  (2022-11-25 00:00:00+00:00,9727.5000000000)
  (2022-11-25 10:00:00+00:00,14013.5000000000)
  (2022-11-25 20:00:00+00:00,11758.0000000000)
  (2022-11-26 06:00:00+00:00,10148.5000000000)
  (2022-11-26 16:00:00+00:00,11902.7500000000)
  (2022-11-27 02:00:00+00:00,8353.5000000000)
  (2022-11-27 12:00:00+00:00,10443.5000000000)
  (2022-11-27 22:00:00+00:00,10566.5000000000)
  (2022-11-28 08:00:00+00:00,15261.5000000000)
  (2022-11-28 18:00:00+00:00,14753.5000000000)
  (2022-11-29 04:00:00+00:00,11629.5000000000)
  (2022-11-29 14:00:00+00:00,13879.0000000000)
  (2022-11-30 00:00:00+00:00,10045.7500000000)
  (2022-11-30 10:00:00+00:00,13945.0000000000)
  (2022-11-30 20:00:00+00:00,12242.5000000000)
  (2022-12-01 06:00:00+00:00,13248.2500000000)
  (2022-12-01 16:00:00+00:00,14295.2500000000)
  (2022-12-02 02:00:00+00:00,10609.2500000000)
  (2022-12-02 12:00:00+00:00,14308.0000000000)
  (2022-12-02 22:00:00+00:00,11320.2500000000)
  (2022-12-03 08:00:00+00:00,12478.0000000000)
  (2022-12-03 18:00:00+00:00,12747.7500000000)
  (2022-12-04 04:00:00+00:00,8925.5000000000)
  (2022-12-04 14:00:00+00:00,12049.2500000000)
  (2022-12-05 00:00:00+00:00,10002.5000000000)
  (2022-12-05 10:00:00+00:00,14606.5000000000)
  (2022-12-05 20:00:00+00:00,12326.0000000000)
  (2022-12-06 06:00:00+00:00,13242.0000000000)
  (2022-12-06 16:00:00+00:00,14495.2500000000)
  (2022-12-07 02:00:00+00:00,10419.2500000000)
  (2022-12-07 12:00:00+00:00,14628.7500000000)
  (2022-12-07 22:00:00+00:00,11610.5000000000)
  (2022-12-08 08:00:00+00:00,14242.0000000000)
  (2022-12-08 18:00:00+00:00,13903.2500000000)
  (2022-12-09 04:00:00+00:00,11522.0000000000)
  (2022-12-09 14:00:00+00:00,13551.5000000000)
  (2022-12-10 00:00:00+00:00,9882.5000000000)
  (2022-12-10 10:00:00+00:00,12964.5000000000)
  (2022-12-10 20:00:00+00:00,11175.7500000000)
  (2022-12-11 06:00:00+00:00,9933.7500000000)
  (2022-12-11 16:00:00+00:00,12861.7500000000)
  (2022-12-12 02:00:00+00:00,10404.2500000000)
  (2022-12-12 12:00:00+00:00,14602.5000000000)
  (2022-12-12 22:00:00+00:00,12285.2500000000)
  (2022-12-13 08:00:00+00:00,14808.2500000000)
  (2022-12-13 18:00:00+00:00,14816.5000000000)
  (2022-12-14 04:00:00+00:00,12732.0000000000)
  (2022-12-14 14:00:00+00:00,14514.2500000000)
  (2022-12-15 00:00:00+00:00,11406.0000000000)
  (2022-12-15 10:00:00+00:00,14929.2500000000)
  (2022-12-15 20:00:00+00:00,13084.7500000000)
  (2022-12-16 06:00:00+00:00,13872.5000000000)
  (2022-12-16 16:00:00+00:00,14345.2500000000)
  (2022-12-17 02:00:00+00:00,10425.0000000000)
  (2022-12-17 12:00:00+00:00,12771.0000000000)
  (2022-12-17 22:00:00+00:00,11606.7500000000)
  (2022-12-18 08:00:00+00:00,12627.7500000000)
  (2022-12-18 18:00:00+00:00,13263.0000000000)
  (2022-12-19 04:00:00+00:00,12096.0000000000)
  (2022-12-19 14:00:00+00:00,14912.5000000000)
  (2022-12-20 00:00:00+00:00,11069.7500000000)
  (2022-12-20 10:00:00+00:00,14643.5000000000)
  (2022-12-20 20:00:00+00:00,12785.2500000000)
  (2022-12-21 06:00:00+00:00,12339.0000000000)
  (2022-12-21 16:00:00+00:00,12973.7500000000)
  (2022-12-22 02:00:00+00:00,9557.0000000000)
  (2022-12-22 12:00:00+00:00,12371.2500000000)
  (2022-12-22 22:00:00+00:00,10640.5000000000)
  (2022-12-23 08:00:00+00:00,12409.2500000000)
  (2022-12-23 18:00:00+00:00,11579.2500000000)
  (2022-12-24 04:00:00+00:00,7877.2500000000)
  (2022-12-24 14:00:00+00:00,9647.2500000000)
  (2022-12-25 00:00:00+00:00,7820.0000000000)
  (2022-12-25 10:00:00+00:00,10752.5000000000)
  (2022-12-25 20:00:00+00:00,10185.5000000000)
  (2022-12-26 06:00:00+00:00,8526.5000000000)
  (2022-12-26 16:00:00+00:00,11880.5000000000)
  (2022-12-27 02:00:00+00:00,9434.7500000000)
  (2022-12-27 12:00:00+00:00,12533.7500000000)
  (2022-12-27 22:00:00+00:00,10330.0000000000)
  (2022-12-28 08:00:00+00:00,12586.5000000000)
  (2022-12-28 18:00:00+00:00,12141.2500000000)
  (2022-12-29 04:00:00+00:00,9800.2500000000)
  (2022-12-29 14:00:00+00:00,12282.2500000000)
  (2022-12-30 00:00:00+00:00,9447.5000000000)
  (2022-12-30 10:00:00+00:00,11975.2500000000)
  (2022-12-30 20:00:00+00:00,11601.0000000000)
  (2022-12-31 06:00:00+00:00,9648.0000000000)
  (2022-12-31 16:00:00+00:00,12224.0000000000)
  (2023-01-01 02:00:00+00:00,8105.0000000000)
  (2023-01-01 12:00:00+00:00,9750.2500000000)
  (2023-01-01 22:00:00+00:00,8515.0000000000)
  (2023-01-02 08:00:00+00:00,13342.7500000000)
  (2023-01-02 18:00:00+00:00,13316.0000000000)
  (2023-01-03 04:00:00+00:00,10801.0000000000)
  (2023-01-03 14:00:00+00:00,12611.7500000000)
  (2023-01-04 00:00:00+00:00,9797.7500000000)
  (2023-01-04 10:00:00+00:00,15136.0000000000)
  (2023-01-04 20:00:00+00:00,13276.0000000000)
  (2023-01-05 06:00:00+00:00,14346.7500000000)
  (2023-01-05 16:00:00+00:00,13792.0000000000)
  (2023-01-06 02:00:00+00:00,9400.0000000000)
  (2023-01-06 12:00:00+00:00,13939.2500000000)
  (2023-01-06 22:00:00+00:00,10316.0000000000)
  (2023-01-07 08:00:00+00:00,10665.2500000000)
  (2023-01-07 18:00:00+00:00,12042.0000000000)
  (2023-01-08 04:00:00+00:00,9312.7500000000)
  (2023-01-08 14:00:00+00:00,10619.2500000000)
  (2023-01-09 00:00:00+00:00,9276.7500000000)
  (2023-01-09 10:00:00+00:00,14318.7500000000)
  (2023-01-09 20:00:00+00:00,12426.5000000000)
  (2023-01-10 06:00:00+00:00,13703.0000000000)
  (2023-01-10 16:00:00+00:00,14323.7500000000)
  (2023-01-11 02:00:00+00:00,11160.2500000000)
  (2023-01-11 12:00:00+00:00,14456.0000000000)
  (2023-01-11 22:00:00+00:00,11937.0000000000)
  (2023-01-12 08:00:00+00:00,15351.2500000000)
  (2023-01-12 18:00:00+00:00,14992.0000000000)
  (2023-01-13 04:00:00+00:00,12031.0000000000)
  (2023-01-13 14:00:00+00:00,14389.2500000000)
  (2023-01-14 00:00:00+00:00,10949.7500000000)
  (2023-01-14 10:00:00+00:00,12852.2500000000)
  (2023-01-14 20:00:00+00:00,12249.5000000000)
  (2023-01-15 06:00:00+00:00,10217.5000000000)
  (2023-01-15 16:00:00+00:00,13675.2500000000)
  (2023-01-16 02:00:00+00:00,10714.5000000000)
  (2023-01-16 12:00:00+00:00,13686.2500000000)
  (2023-01-16 22:00:00+00:00,12667.7500000000)
  (2023-01-17 08:00:00+00:00,15167.7500000000)
  (2023-01-17 18:00:00+00:00,14903.0000000000)
  (2023-01-18 04:00:00+00:00,11591.7500000000)
  (2023-01-18 14:00:00+00:00,13463.0000000000)
  (2023-01-19 00:00:00+00:00,11179.0000000000)
  (2023-01-19 10:00:00+00:00,15289.2500000000)
  (2023-01-19 20:00:00+00:00,12875.0000000000)
  (2023-01-20 06:00:00+00:00,13815.2500000000)
  (2023-01-20 16:00:00+00:00,14220.2500000000)
  (2023-01-21 02:00:00+00:00,9970.2500000000)
  (2023-01-21 12:00:00+00:00,13712.2500000000)
  (2023-01-21 22:00:00+00:00,10460.2500000000)
  (2023-01-22 08:00:00+00:00,11970.2500000000)
  (2023-01-22 18:00:00+00:00,12325.2500000000)
  (2023-01-23 04:00:00+00:00,10740.7500000000)
  (2023-01-23 14:00:00+00:00,13662.5000000000)
  (2023-01-24 00:00:00+00:00,10498.7500000000)
  (2023-01-24 10:00:00+00:00,14247.0000000000)
  (2023-01-24 20:00:00+00:00,12500.0000000000)
  (2023-01-25 06:00:00+00:00,13897.2500000000)
  (2023-01-25 16:00:00+00:00,14725.5000000000)
  (2023-01-26 02:00:00+00:00,10919.0000000000)
  (2023-01-26 12:00:00+00:00,14392.5000000000)
  (2023-01-26 22:00:00+00:00,12661.2500000000)
  (2023-01-27 08:00:00+00:00,15040.7500000000)
  (2023-01-27 18:00:00+00:00,13989.7500000000)
  (2023-01-28 04:00:00+00:00,10033.5000000000)
  (2023-01-28 14:00:00+00:00,12774.2500000000)
  (2023-01-29 00:00:00+00:00,10809.0000000000)
  (2023-01-29 10:00:00+00:00,13505.7500000000)
  (2023-01-29 20:00:00+00:00,12519.7500000000)
  (2023-01-30 06:00:00+00:00,14424.7500000000)
  (2023-01-30 16:00:00+00:00,14476.7500000000)
  (2023-01-31 02:00:00+00:00,12096.5000000000)
  (2023-01-31 12:00:00+00:00,14616.0000000000)
  (2023-01-31 22:00:00+00:00,11574.0000000000)
  (2023-02-01 08:00:00+00:00,14747.2500000000)
  (2023-02-01 18:00:00+00:00,15264.7500000000)
  (2023-02-02 04:00:00+00:00,13203.5000000000)
  (2023-02-02 14:00:00+00:00,14267.7500000000)
  (2023-02-03 00:00:00+00:00,10973.0000000000)
  (2023-02-03 10:00:00+00:00,15561.7500000000)
  (2023-02-03 20:00:00+00:00,13363.5000000000)
  (2023-02-04 06:00:00+00:00,11271.7500000000)
  (2023-02-04 16:00:00+00:00,12487.5000000000)
  (2023-02-05 02:00:00+00:00,9232.7500000000)
  (2023-02-05 12:00:00+00:00,11333.2500000000)
  (2023-02-05 22:00:00+00:00,10658.0000000000)
  (2023-02-06 08:00:00+00:00,14180.7500000000)
  (2023-02-06 18:00:00+00:00,14162.0000000000)
  (2023-02-07 04:00:00+00:00,12109.7500000000)
  (2023-02-07 14:00:00+00:00,13507.5000000000)
  (2023-02-08 00:00:00+00:00,11244.5000000000)
  (2023-02-08 10:00:00+00:00,14479.7500000000)
  (2023-02-08 20:00:00+00:00,14351.5000000000)
  (2023-02-09 06:00:00+00:00,15270.2500000000)
  (2023-02-09 16:00:00+00:00,15167.5000000000)
  (2023-02-10 02:00:00+00:00,12411.5000000000)
  (2023-02-10 12:00:00+00:00,14538.0000000000)
  (2023-02-10 22:00:00+00:00,11779.7500000000)
  (2023-02-11 08:00:00+00:00,13359.2500000000)
  (2023-02-11 18:00:00+00:00,13266.7500000000)
  (2023-02-12 04:00:00+00:00,9489.2500000000)
  (2023-02-12 14:00:00+00:00,11823.7500000000)
  (2023-02-13 00:00:00+00:00,10872.7500000000)
  (2023-02-13 10:00:00+00:00,14810.7500000000)
  (2023-02-13 20:00:00+00:00,12561.2500000000)
  (2023-02-14 06:00:00+00:00,13157.0000000000)
  (2023-02-14 16:00:00+00:00,13253.5000000000)
  (2023-02-15 02:00:00+00:00,10642.2500000000)
  (2023-02-15 12:00:00+00:00,13271.5000000000)
  (2023-02-15 22:00:00+00:00,11481.0000000000)
  (2023-02-16 08:00:00+00:00,14448.2500000000)
  (2023-02-16 18:00:00+00:00,14829.2500000000)
  (2023-02-17 04:00:00+00:00,11836.5000000000)
  (2023-02-17 14:00:00+00:00,14721.5000000000)
  (2023-02-18 00:00:00+00:00,10818.2500000000)
  (2023-02-18 10:00:00+00:00,13059.2500000000)
  (2023-02-18 20:00:00+00:00,11462.2500000000)
  (2023-02-19 06:00:00+00:00,9875.7500000000)
  (2023-02-19 16:00:00+00:00,12017.0000000000)
  (2023-02-20 02:00:00+00:00,11030.5000000000)
  (2023-02-20 12:00:00+00:00,14783.2500000000)
  (2023-02-20 22:00:00+00:00,12095.5000000000)
  (2023-02-21 08:00:00+00:00,14975.5000000000)
  (2023-02-21 18:00:00+00:00,13823.2500000000)
  (2023-02-22 04:00:00+00:00,10710.7500000000)
  (2023-02-22 14:00:00+00:00,13308.0000000000)
  (2023-02-23 00:00:00+00:00,10974.5000000000)
  (2023-02-23 10:00:00+00:00,13637.7500000000)
  (2023-02-23 20:00:00+00:00,12494.7500000000)
  (2023-02-24 06:00:00+00:00,14075.0000000000)
  (2023-02-24 16:00:00+00:00,14520.0000000000)
  (2023-02-25 02:00:00+00:00,11221.2500000000)
  (2023-02-25 12:00:00+00:00,13322.7500000000)
  (2023-02-25 22:00:00+00:00,11229.0000000000)
  (2023-02-26 08:00:00+00:00,11466.0000000000)
  (2023-02-26 18:00:00+00:00,12088.7500000000)
  (2023-02-27 04:00:00+00:00,10714.7500000000)
  (2023-02-27 14:00:00+00:00,12833.2500000000)
  (2023-02-28 00:00:00+00:00,9937.5000000000)
  (2023-02-28 10:00:00+00:00,13277.2500000000)
  (2023-02-28 20:00:00+00:00,11987.2500000000)
  (2023-03-01 06:00:00+00:00,13363.2500000000)
  (2023-03-01 16:00:00+00:00,13129.2500000000)
  (2023-03-02 02:00:00+00:00,10614.2500000000)
  (2023-03-02 12:00:00+00:00,13019.2500000000)
  (2023-03-02 22:00:00+00:00,11189.7500000000)
  (2023-03-03 08:00:00+00:00,13847.7500000000)
  (2023-03-03 18:00:00+00:00,14784.2500000000)
  (2023-03-04 04:00:00+00:00,10508.7500000000)
  (2023-03-04 14:00:00+00:00,12056.0000000000)
  (2023-03-05 00:00:00+00:00,9608.2500000000)
  (2023-03-05 10:00:00+00:00,12791.0000000000)
  (2023-03-05 20:00:00+00:00,10715.7500000000)
  (2023-03-06 06:00:00+00:00,13109.5000000000)
  (2023-03-06 16:00:00+00:00,13581.5000000000)
  (2023-03-07 02:00:00+00:00,11335.2500000000)
  (2023-03-07 12:00:00+00:00,14005.0000000000)
  (2023-03-07 22:00:00+00:00,12887.5000000000)
  (2023-03-08 08:00:00+00:00,13302.2500000000)
  (2023-03-08 18:00:00+00:00,12978.5000000000)
  (2023-03-09 04:00:00+00:00,11224.0000000000)
  (2023-03-09 14:00:00+00:00,13089.0000000000)
  (2023-03-10 00:00:00+00:00,9817.5000000000)
  (2023-03-10 10:00:00+00:00,14268.2500000000)
  (2023-03-10 20:00:00+00:00,12310.7500000000)
  (2023-03-11 06:00:00+00:00,11538.7500000000)
  (2023-03-11 16:00:00+00:00,11863.2500000000)
  (2023-03-12 02:00:00+00:00,10373.0000000000)
  (2023-03-12 12:00:00+00:00,10671.7500000000)
  (2023-03-12 22:00:00+00:00,10564.7500000000)
  (2023-03-13 08:00:00+00:00,15320.0000000000)
  (2023-03-13 18:00:00+00:00,14326.5000000000)
  (2023-03-14 04:00:00+00:00,12064.2500000000)
  (2023-03-14 14:00:00+00:00,13974.2500000000)
  (2023-03-15 00:00:00+00:00,11661.0000000000)
  (2023-03-15 10:00:00+00:00,14077.7500000000)
  (2023-03-15 20:00:00+00:00,13069.5000000000)
  (2023-03-16 06:00:00+00:00,13376.0000000000)
  (2023-03-16 16:00:00+00:00,13278.0000000000)
  (2023-03-17 02:00:00+00:00,11673.0000000000)
  (2023-03-17 12:00:00+00:00,12452.2500000000)
  (2023-03-17 22:00:00+00:00,11798.7500000000)
  (2023-03-18 08:00:00+00:00,11508.5000000000)
  (2023-03-18 18:00:00+00:00,11547.2500000000)
  (2023-03-19 04:00:00+00:00,8522.2500000000)
  (2023-03-19 14:00:00+00:00,9943.5000000000)
  (2023-03-20 00:00:00+00:00,9031.0000000000)
  (2023-03-20 10:00:00+00:00,14229.5000000000)
  (2023-03-20 20:00:00+00:00,12293.0000000000)
  (2023-03-21 06:00:00+00:00,12326.2500000000)
  (2023-03-21 16:00:00+00:00,13099.2500000000)
  (2023-03-22 02:00:00+00:00,9922.5000000000)
  (2023-03-22 12:00:00+00:00,13317.5000000000)
  (2023-03-22 22:00:00+00:00,11174.2500000000)
  (2023-03-23 08:00:00+00:00,13522.2500000000)
  (2023-03-23 18:00:00+00:00,13489.7500000000)
  (2023-03-24 04:00:00+00:00,11861.5000000000)
  (2023-03-24 14:00:00+00:00,13463.0000000000)
  (2023-03-25 00:00:00+00:00,9961.0000000000)
  (2023-03-25 10:00:00+00:00,12349.5000000000)
  (2023-03-25 20:00:00+00:00,11261.5000000000)
  (2023-03-26 06:00:00+00:00,10106.5000000000)
  (2023-03-26 16:00:00+00:00,11693.2500000000)
  (2023-03-27 02:00:00+00:00,10719.0000000000)
  (2023-03-27 12:00:00+00:00,13007.5000000000)
  (2023-03-27 22:00:00+00:00,11840.2500000000)
  (2023-03-28 08:00:00+00:00,13779.7500000000)
  (2023-03-28 18:00:00+00:00,12745.7500000000)
  (2023-03-29 04:00:00+00:00,12833.5000000000)
  (2023-03-29 14:00:00+00:00,12577.2500000000)
  (2023-03-30 00:00:00+00:00,8875.7500000000)
  (2023-03-30 10:00:00+00:00,13988.7500000000)
  (2023-03-30 20:00:00+00:00,12141.0000000000)
  (2023-03-31 06:00:00+00:00,12940.5000000000)
  (2023-03-31 16:00:00+00:00,12430.0000000000)
  (2023-04-01 02:00:00+00:00,9724.7500000000)
  (2023-04-01 12:00:00+00:00,12050.2500000000)
  (2023-04-01 22:00:00+00:00,10665.0000000000)
  (2023-04-02 08:00:00+00:00,11516.2500000000)
  (2023-04-02 18:00:00+00:00,11954.0000000000)
  (2023-04-03 04:00:00+00:00,12869.0000000000)
  (2023-04-03 14:00:00+00:00,12532.0000000000)
  (2023-04-04 00:00:00+00:00,10315.0000000000)
  (2023-04-04 10:00:00+00:00,12925.0000000000)
  (2023-04-04 20:00:00+00:00,11510.5000000000)
  (2023-04-05 06:00:00+00:00,12962.7500000000)
  (2023-04-05 16:00:00+00:00,12645.0000000000)
  (2023-04-06 02:00:00+00:00,10598.0000000000)
  (2023-04-06 12:00:00+00:00,11906.7500000000)
  (2023-04-06 22:00:00+00:00,10037.0000000000)
  (2023-04-07 08:00:00+00:00,10924.5000000000)
  (2023-04-07 18:00:00+00:00,11690.2500000000)
  (2023-04-08 04:00:00+00:00,9465.7500000000)
  (2023-04-08 14:00:00+00:00,10164.2500000000)
  (2023-04-09 00:00:00+00:00,8468.7500000000)
  (2023-04-09 10:00:00+00:00,9201.0000000000)
  (2023-04-09 20:00:00+00:00,8655.7500000000)
  (2023-04-10 06:00:00+00:00,9060.0000000000)
  (2023-04-10 16:00:00+00:00,10620.5000000000)
  (2023-04-11 02:00:00+00:00,10188.0000000000)
  (2023-04-11 12:00:00+00:00,12168.7500000000)
  (2023-04-11 22:00:00+00:00,9416.7500000000)
  (2023-04-12 08:00:00+00:00,12672.5000000000)
  (2023-04-12 18:00:00+00:00,13437.5000000000)
  (2023-04-13 04:00:00+00:00,11826.2500000000)
  (2023-04-13 14:00:00+00:00,11755.7500000000)
  (2023-04-14 00:00:00+00:00,9400.5000000000)
  (2023-04-14 10:00:00+00:00,13208.0000000000)
  (2023-04-14 20:00:00+00:00,11571.0000000000)
  (2023-04-15 06:00:00+00:00,10864.0000000000)
  (2023-04-15 16:00:00+00:00,11962.7500000000)
  (2023-04-16 02:00:00+00:00,8838.7500000000)
  (2023-04-16 12:00:00+00:00,10079.7500000000)
  (2023-04-16 22:00:00+00:00,9573.2500000000)
  (2023-04-17 08:00:00+00:00,12906.5000000000)
  (2023-04-17 18:00:00+00:00,12760.5000000000)
  (2023-04-18 04:00:00+00:00,12499.2500000000)
  (2023-04-18 14:00:00+00:00,13069.5000000000)
  (2023-04-19 00:00:00+00:00,10181.7500000000)
  (2023-04-19 10:00:00+00:00,12755.2500000000)
  (2023-04-19 20:00:00+00:00,11944.2500000000)
  (2023-04-20 06:00:00+00:00,13468.5000000000)
  (2023-04-20 16:00:00+00:00,13265.5000000000)
  (2023-04-21 02:00:00+00:00,10704.0000000000)
  (2023-04-21 12:00:00+00:00,11601.5000000000)
  (2023-04-21 22:00:00+00:00,9902.2500000000)
  (2023-04-22 08:00:00+00:00,10721.0000000000)
  (2023-04-22 18:00:00+00:00,11197.7500000000)
  (2023-04-23 04:00:00+00:00,8551.5000000000)
  (2023-04-23 14:00:00+00:00,9121.0000000000)
  (2023-04-24 00:00:00+00:00,8925.7500000000)
  (2023-04-24 10:00:00+00:00,12452.7500000000)
  (2023-04-24 20:00:00+00:00,11534.7500000000)
  (2023-04-25 06:00:00+00:00,13047.7500000000)
  (2023-04-25 16:00:00+00:00,12969.5000000000)
  (2023-04-26 02:00:00+00:00,10688.2500000000)
  (2023-04-26 12:00:00+00:00,12469.7500000000)
  (2023-04-26 22:00:00+00:00,10695.2500000000)
  (2023-04-27 08:00:00+00:00,13573.0000000000)
  (2023-04-27 18:00:00+00:00,12793.0000000000)
  (2023-04-28 04:00:00+00:00,11702.5000000000)
  (2023-04-28 14:00:00+00:00,12269.0000000000)
  (2023-04-29 00:00:00+00:00,9126.2500000000)
  (2023-04-29 10:00:00+00:00,11221.2500000000)
  (2023-04-29 20:00:00+00:00,11113.7500000000)
  (2023-04-30 06:00:00+00:00,9599.7500000000)
  (2023-04-30 16:00:00+00:00,9267.7500000000)
  (2023-05-01 02:00:00+00:00,7790.7500000000)
  (2023-05-01 12:00:00+00:00,9029.0000000000)
  (2023-05-01 22:00:00+00:00,8192.2500000000)
  (2023-05-02 08:00:00+00:00,13138.7500000000)
  (2023-05-02 18:00:00+00:00,11848.7500000000)
  (2023-05-03 04:00:00+00:00,11335.5000000000)
  (2023-05-03 14:00:00+00:00,11159.2500000000)
  (2023-05-04 00:00:00+00:00,8600.5000000000)
  (2023-05-04 10:00:00+00:00,11626.7500000000)
  (2023-05-04 20:00:00+00:00,11829.5000000000)
  (2023-05-05 06:00:00+00:00,13723.7500000000)
  (2023-05-05 16:00:00+00:00,13390.5000000000)
  (2023-05-06 02:00:00+00:00,9642.7500000000)
  (2023-05-06 12:00:00+00:00,11140.2500000000)
  (2023-05-06 22:00:00+00:00,8903.2500000000)
  (2023-05-07 08:00:00+00:00,10319.2500000000)
  (2023-05-07 18:00:00+00:00,11242.7500000000)
  (2023-05-08 04:00:00+00:00,10957.5000000000)
  (2023-05-08 14:00:00+00:00,11861.0000000000)
  (2023-05-09 00:00:00+00:00,10317.0000000000)
  (2023-05-09 10:00:00+00:00,12199.2500000000)
  (2023-05-09 20:00:00+00:00,12049.7500000000)
  (2023-05-10 06:00:00+00:00,13274.2500000000)
  (2023-05-10 16:00:00+00:00,12600.5000000000)
  (2023-05-11 02:00:00+00:00,10107.5000000000)
  (2023-05-11 12:00:00+00:00,12687.2500000000)
  (2023-05-11 22:00:00+00:00,9981.0000000000)
  (2023-05-12 08:00:00+00:00,12479.5000000000)
  (2023-05-12 18:00:00+00:00,12128.0000000000)
  (2023-05-13 04:00:00+00:00,8708.2500000000)
  (2023-05-13 14:00:00+00:00,9973.7500000000)
  (2023-05-14 00:00:00+00:00,7851.7500000000)
  (2023-05-14 10:00:00+00:00,9656.5000000000)
  (2023-05-14 20:00:00+00:00,9495.5000000000)
  (2023-05-15 06:00:00+00:00,11971.0000000000)
  (2023-05-15 16:00:00+00:00,12338.0000000000)
  (2023-05-16 02:00:00+00:00,9485.0000000000)
  (2023-05-16 12:00:00+00:00,12344.0000000000)
  (2023-05-16 22:00:00+00:00,10417.7500000000)
  (2023-05-17 08:00:00+00:00,13698.7500000000)
  (2023-05-17 18:00:00+00:00,10949.2500000000)
  (2023-05-18 04:00:00+00:00,7680.5000000000)
  (2023-05-18 14:00:00+00:00,8344.0000000000)
  (2023-05-19 00:00:00+00:00,7377.5000000000)
  (2023-05-19 10:00:00+00:00,10515.0000000000)
  (2023-05-19 20:00:00+00:00,9506.5000000000)
  (2023-05-20 06:00:00+00:00,9431.5000000000)
  (2023-05-20 16:00:00+00:00,10488.2500000000)
  (2023-05-21 02:00:00+00:00,7606.7500000000)
  (2023-05-21 12:00:00+00:00,9066.7500000000)
  (2023-05-21 22:00:00+00:00,9426.7500000000)
  (2023-05-22 08:00:00+00:00,12167.2500000000)
  (2023-05-22 18:00:00+00:00,12168.5000000000)
  (2023-05-23 04:00:00+00:00,10710.5000000000)
  (2023-05-23 14:00:00+00:00,12272.7500000000)
  (2023-05-24 00:00:00+00:00,8984.2500000000)
  (2023-05-24 10:00:00+00:00,11893.0000000000)
  (2023-05-24 20:00:00+00:00,10017.2500000000)
  (2023-05-25 06:00:00+00:00,11791.0000000000)
  (2023-05-25 16:00:00+00:00,12275.0000000000)
  (2023-05-26 02:00:00+00:00,9271.0000000000)
  (2023-05-26 12:00:00+00:00,11811.7500000000)
  (2023-05-26 22:00:00+00:00,8461.5000000000)
  (2023-05-27 08:00:00+00:00,9723.7500000000)
  (2023-05-27 18:00:00+00:00,8671.2500000000)
  (2023-05-28 04:00:00+00:00,6778.2500000000)
  (2023-05-28 14:00:00+00:00,7915.0000000000)
  (2023-05-29 00:00:00+00:00,7054.0000000000)
  (2023-05-29 10:00:00+00:00,9047.0000000000)
  (2023-05-29 20:00:00+00:00,9060.0000000000)
  (2023-05-30 06:00:00+00:00,11452.5000000000)
  (2023-05-30 16:00:00+00:00,11155.7500000000)
  (2023-05-31 02:00:00+00:00,8350.2500000000)
  (2023-05-31 12:00:00+00:00,11888.0000000000)
  (2023-05-31 22:00:00+00:00,9581.2500000000)
  (2023-06-01 08:00:00+00:00,11852.2500000000)
  (2023-06-01 18:00:00+00:00,11524.0000000000)
  (2023-06-02 04:00:00+00:00,10693.0000000000)
  (2023-06-02 14:00:00+00:00,11396.5000000000)
  (2023-06-03 00:00:00+00:00,7988.5000000000)
  (2023-06-03 10:00:00+00:00,10013.2500000000)
  (2023-06-03 20:00:00+00:00,8723.0000000000)
  (2023-06-04 06:00:00+00:00,8530.2500000000)
  (2023-06-04 16:00:00+00:00,9259.5000000000)
  (2023-06-05 02:00:00+00:00,7831.0000000000)
  (2023-06-05 12:00:00+00:00,11189.2500000000)
  (2023-06-05 22:00:00+00:00,8741.5000000000)
  (2023-06-06 08:00:00+00:00,11899.5000000000)
  (2023-06-06 18:00:00+00:00,11088.2500000000)
  (2023-06-07 04:00:00+00:00,10182.0000000000)
  (2023-06-07 14:00:00+00:00,11212.7500000000)
  (2023-06-08 00:00:00+00:00,8209.0000000000)
  (2023-06-08 10:00:00+00:00,11942.2500000000)
  (2023-06-08 20:00:00+00:00,10331.2500000000)
  (2023-06-09 06:00:00+00:00,12197.7500000000)
  (2023-06-09 16:00:00+00:00,12620.5000000000)
  (2023-06-10 02:00:00+00:00,8848.2500000000)
  (2023-06-10 12:00:00+00:00,10856.2500000000)
  (2023-06-10 22:00:00+00:00,9185.2500000000)
  (2023-06-11 08:00:00+00:00,10061.0000000000)
  (2023-06-11 18:00:00+00:00,10919.2500000000)
  (2023-06-12 04:00:00+00:00,10557.5000000000)
  (2023-06-12 14:00:00+00:00,11775.2500000000)
  (2023-06-13 00:00:00+00:00,9362.5000000000)
  (2023-06-13 10:00:00+00:00,12533.2500000000)
  (2023-06-13 20:00:00+00:00,10872.0000000000)
  (2023-06-14 06:00:00+00:00,11917.5000000000)
  (2023-06-14 16:00:00+00:00,11752.7500000000)
  (2023-06-15 02:00:00+00:00,8733.0000000000)
  (2023-06-15 12:00:00+00:00,11386.0000000000)
  (2023-06-15 22:00:00+00:00,9008.2500000000)
  (2023-06-16 08:00:00+00:00,12087.2500000000)
  (2023-06-16 18:00:00+00:00,10436.5000000000)
  (2023-06-17 04:00:00+00:00,8075.2500000000)
  (2023-06-17 14:00:00+00:00,9191.5000000000)
  (2023-06-18 00:00:00+00:00,7594.5000000000)
  (2023-06-18 10:00:00+00:00,9387.5000000000)
  (2023-06-18 20:00:00+00:00,9257.2500000000)
  (2023-06-19 06:00:00+00:00,11591.2500000000)
  (2023-06-19 16:00:00+00:00,12010.2500000000)
  (2023-06-20 02:00:00+00:00,9114.0000000000)
  (2023-06-20 12:00:00+00:00,12067.5000000000)
  (2023-06-20 22:00:00+00:00,10043.0000000000)
  (2023-06-21 08:00:00+00:00,12533.0000000000)
  (2023-06-21 18:00:00+00:00,12396.0000000000)
  (2023-06-22 04:00:00+00:00,10988.5000000000)
  (2023-06-22 14:00:00+00:00,12144.2500000000)
  (2023-06-23 00:00:00+00:00,9031.0000000000)
  (2023-06-23 10:00:00+00:00,12806.7500000000)
  (2023-06-23 20:00:00+00:00,10241.2500000000)
  (2023-06-24 06:00:00+00:00,9479.2500000000)
  (2023-06-24 16:00:00+00:00,10056.2500000000)
  (2023-06-25 02:00:00+00:00,7255.0000000000)
  (2023-06-25 12:00:00+00:00,8869.2500000000)
  (2023-06-25 22:00:00+00:00,8458.7500000000)
  (2023-06-26 08:00:00+00:00,12313.7500000000)
  (2023-06-26 18:00:00+00:00,12484.2500000000)
  (2023-06-27 04:00:00+00:00,11472.2500000000)
  (2023-06-27 14:00:00+00:00,11789.2500000000)
  (2023-06-28 00:00:00+00:00,10063.7500000000)
  (2023-06-28 10:00:00+00:00,11950.2500000000)
  (2023-06-28 20:00:00+00:00,10353.0000000000)
  (2023-06-29 06:00:00+00:00,12065.2500000000)
  (2023-06-29 16:00:00+00:00,12028.2500000000)
  (2023-06-30 02:00:00+00:00,9377.0000000000)
  (2023-06-30 12:00:00+00:00,11923.0000000000)
  (2023-06-30 22:00:00+00:00,8599.7500000000)
  (2023-07-01 08:00:00+00:00,10467.5000000000)
  (2023-07-01 18:00:00+00:00,10098.2500000000)
  (2023-07-02 04:00:00+00:00,8541.2500000000)
  (2023-07-02 14:00:00+00:00,9882.7500000000)
  (2023-07-03 00:00:00+00:00,9206.2500000000)
  (2023-07-03 10:00:00+00:00,13184.7500000000)
  (2023-07-03 20:00:00+00:00,10995.0000000000)
  (2023-07-04 06:00:00+00:00,12790.7500000000)
  (2023-07-04 16:00:00+00:00,12855.2500000000)
  (2023-07-05 02:00:00+00:00,8700.5000000000)
  (2023-07-05 12:00:00+00:00,12298.7500000000)
  (2023-07-05 22:00:00+00:00,10662.5000000000)
  (2023-07-06 08:00:00+00:00,11897.5000000000)
  (2023-07-06 18:00:00+00:00,11348.2500000000)
  (2023-07-07 04:00:00+00:00,10199.0000000000)
  (2023-07-07 14:00:00+00:00,10932.5000000000)
  (2023-07-08 00:00:00+00:00,7847.5000000000)
  (2023-07-08 10:00:00+00:00,10209.5000000000)
  (2023-07-08 20:00:00+00:00,9226.5000000000)
  (2023-07-09 06:00:00+00:00,8749.5000000000)
  (2023-07-09 16:00:00+00:00,9887.5000000000)
  (2023-07-10 02:00:00+00:00,8617.2500000000)
  (2023-07-10 12:00:00+00:00,12231.7500000000)
  (2023-07-10 22:00:00+00:00,9445.7500000000)
  (2023-07-11 08:00:00+00:00,11964.5000000000)
  (2023-07-11 18:00:00+00:00,12068.0000000000)
  (2023-07-12 04:00:00+00:00,11953.0000000000)
  (2023-07-12 14:00:00+00:00,12803.5000000000)
  (2023-07-13 00:00:00+00:00,8723.7500000000)
  (2023-07-13 10:00:00+00:00,12459.7500000000)
  (2023-07-13 20:00:00+00:00,10716.0000000000)
  (2023-07-14 06:00:00+00:00,11465.5000000000)
  (2023-07-14 16:00:00+00:00,10959.7500000000)
  (2023-07-15 02:00:00+00:00,8865.0000000000)
  (2023-07-15 12:00:00+00:00,10737.5000000000)
  (2023-07-15 22:00:00+00:00,8578.2500000000)
  (2023-07-16 08:00:00+00:00,9440.7500000000)
  (2023-07-16 18:00:00+00:00,9403.2500000000)
  (2023-07-17 04:00:00+00:00,9841.5000000000)
  (2023-07-17 14:00:00+00:00,11834.5000000000)
  (2023-07-18 00:00:00+00:00,9559.2500000000)
  (2023-07-18 10:00:00+00:00,12514.2500000000)
  (2023-07-18 20:00:00+00:00,10223.5000000000)
  (2023-07-19 06:00:00+00:00,11161.7500000000)
  (2023-07-19 16:00:00+00:00,11358.0000000000)
  (2023-07-20 02:00:00+00:00,8604.7500000000)
  (2023-07-20 12:00:00+00:00,11937.7500000000)
  (2023-07-20 22:00:00+00:00,9389.0000000000)
  (2023-07-21 08:00:00+00:00,11544.7500000000)
  (2023-07-21 18:00:00+00:00,9965.7500000000)
  (2023-07-22 04:00:00+00:00,7953.0000000000)
  (2023-07-22 14:00:00+00:00,9801.2500000000)
  (2023-07-23 00:00:00+00:00,7744.0000000000)
  (2023-07-23 10:00:00+00:00,9802.7500000000)
  (2023-07-23 20:00:00+00:00,9447.0000000000)
  (2023-07-24 06:00:00+00:00,11526.2500000000)
  (2023-07-24 16:00:00+00:00,10923.2500000000)
  (2023-07-25 02:00:00+00:00,9105.7500000000)
  (2023-07-25 12:00:00+00:00,11734.2500000000)
  (2023-07-25 22:00:00+00:00,9047.0000000000)
  (2023-07-26 08:00:00+00:00,11983.2500000000)
  (2023-07-26 18:00:00+00:00,10955.5000000000)
  (2023-07-27 04:00:00+00:00,10502.0000000000)
  (2023-07-27 14:00:00+00:00,11006.2500000000)
  (2023-07-28 00:00:00+00:00,8019.2500000000)
  (2023-07-28 10:00:00+00:00,11518.2500000000)
  (2023-07-28 20:00:00+00:00,9197.0000000000)
  (2023-07-29 06:00:00+00:00,8934.2500000000)
  (2023-07-29 16:00:00+00:00,9449.5000000000)
  (2023-07-30 02:00:00+00:00,7159.2500000000)
  (2023-07-30 12:00:00+00:00,8865.5000000000)
  (2023-07-30 22:00:00+00:00,8780.7500000000)
  (2023-07-31 08:00:00+00:00,12014.2500000000)
  (2023-07-31 18:00:00+00:00,10712.5000000000)
  (2023-08-01 04:00:00+00:00,10452.5000000000)
  (2023-08-01 14:00:00+00:00,11029.2500000000)
  (2023-08-02 00:00:00+00:00,9347.7500000000)
  (2023-08-02 10:00:00+00:00,12014.0000000000)
  (2023-08-02 20:00:00+00:00,10563.2500000000)
  (2023-08-03 06:00:00+00:00,12009.0000000000)
  (2023-08-03 16:00:00+00:00,11825.0000000000)
  (2023-08-04 02:00:00+00:00,8369.2500000000)
  (2023-08-04 12:00:00+00:00,11717.7500000000)
  (2023-08-04 22:00:00+00:00,8130.0000000000)
  (2023-08-05 08:00:00+00:00,9467.5000000000)
  (2023-08-05 18:00:00+00:00,8832.7500000000)
  (2023-08-06 04:00:00+00:00,7066.2500000000)
  (2023-08-06 14:00:00+00:00,9889.2500000000)
  (2023-08-07 00:00:00+00:00,9106.0000000000)
  (2023-08-07 10:00:00+00:00,12368.5000000000)
  (2023-08-07 20:00:00+00:00,10306.5000000000)
  (2023-08-08 06:00:00+00:00,11853.5000000000)
  (2023-08-08 16:00:00+00:00,12028.7500000000)
  (2023-08-09 02:00:00+00:00,9526.7500000000)
  (2023-08-09 12:00:00+00:00,11879.2500000000)
  (2023-08-09 22:00:00+00:00,9710.0000000000)
  (2023-08-10 08:00:00+00:00,12239.0000000000)
  (2023-08-10 18:00:00+00:00,10949.2500000000)
  (2023-08-11 04:00:00+00:00,9947.2500000000)
  (2023-08-11 14:00:00+00:00,10757.7500000000)
  (2023-08-12 00:00:00+00:00,8289.2500000000)
  (2023-08-12 10:00:00+00:00,10543.2500000000)
  (2023-08-12 20:00:00+00:00,9651.2500000000)
  (2023-08-13 06:00:00+00:00,8771.7500000000)
  (2023-08-13 16:00:00+00:00,9083.7500000000)
  (2023-08-14 02:00:00+00:00,8225.2500000000)
  (2023-08-14 12:00:00+00:00,11654.0000000000)
  (2023-08-14 22:00:00+00:00,9375.0000000000)
  (2023-08-15 08:00:00+00:00,12845.0000000000)
  (2023-08-15 18:00:00+00:00,11145.2500000000)
  (2023-08-16 04:00:00+00:00,10279.2500000000)
  (2023-08-16 14:00:00+00:00,11319.2500000000)
  (2023-08-17 00:00:00+00:00,8327.2500000000)
  (2023-08-17 10:00:00+00:00,12410.5000000000)
  (2023-08-17 20:00:00+00:00,10506.0000000000)
  (2023-08-18 06:00:00+00:00,12055.2500000000)
  (2023-08-18 16:00:00+00:00,11875.5000000000)
  (2023-08-19 02:00:00+00:00,7978.7500000000)
  (2023-08-19 12:00:00+00:00,10268.7500000000)
  (2023-08-19 22:00:00+00:00,8560.7500000000)
  (2023-08-20 08:00:00+00:00,9634.0000000000)
  (2023-08-20 18:00:00+00:00,9927.5000000000)
  (2023-08-21 04:00:00+00:00,10299.7500000000)
  (2023-08-21 14:00:00+00:00,12207.7500000000)
  (2023-08-22 00:00:00+00:00,8585.0000000000)
  (2023-08-22 10:00:00+00:00,12390.2500000000)
  (2023-08-22 20:00:00+00:00,10024.2500000000)
  (2023-08-23 06:00:00+00:00,11681.0000000000)
  (2023-08-23 16:00:00+00:00,11912.0000000000)
  (2023-08-24 02:00:00+00:00,8941.2500000000)
  (2023-08-24 12:00:00+00:00,11965.0000000000)
  (2023-08-24 22:00:00+00:00,9317.0000000000)
  (2023-08-25 08:00:00+00:00,13084.7500000000)
  (2023-08-25 18:00:00+00:00,10903.5000000000)
  (2023-08-26 04:00:00+00:00,8857.0000000000)
  (2023-08-26 14:00:00+00:00,10666.7500000000)
  (2023-08-27 00:00:00+00:00,6969.5000000000)
  (2023-08-27 10:00:00+00:00,9808.7500000000)
  (2023-08-27 20:00:00+00:00,9008.2500000000)
  (2023-08-28 06:00:00+00:00,12009.0000000000)
  (2023-08-28 16:00:00+00:00,11724.5000000000)
  (2023-08-29 02:00:00+00:00,8601.5000000000)
  (2023-08-29 12:00:00+00:00,11679.2500000000)
  (2023-08-29 22:00:00+00:00,8768.7500000000)
  (2023-08-30 08:00:00+00:00,11892.2500000000)
  (2023-08-30 18:00:00+00:00,11553.2500000000)
  (2023-08-31 04:00:00+00:00,10759.7500000000)
  (2023-08-31 14:00:00+00:00,12132.7500000000)
  (2023-09-01 00:00:00+00:00,9664.2500000000)
  (2023-09-01 10:00:00+00:00,12497.0000000000)
  (2023-09-01 20:00:00+00:00,10099.5000000000)
  (2023-09-02 06:00:00+00:00,9704.2500000000)
  (2023-09-02 16:00:00+00:00,10355.5000000000)
  (2023-09-03 02:00:00+00:00,7588.7500000000)
  (2023-09-03 12:00:00+00:00,9418.0000000000)
  (2023-09-03 22:00:00+00:00,8837.7500000000)
  (2023-09-04 08:00:00+00:00,12124.5000000000)
  (2023-09-04 18:00:00+00:00,11887.2500000000)
  (2023-09-05 04:00:00+00:00,11118.7500000000)
  (2023-09-05 14:00:00+00:00,11445.0000000000)
  (2023-09-06 00:00:00+00:00,8384.2500000000)
  (2023-09-06 10:00:00+00:00,12021.7500000000)
  (2023-09-06 20:00:00+00:00,10802.5000000000)
  (2023-09-07 06:00:00+00:00,12508.2500000000)
  (2023-09-07 16:00:00+00:00,12100.0000000000)
  (2023-09-08 02:00:00+00:00,8856.5000000000)
  (2023-09-08 12:00:00+00:00,11438.5000000000)
  (2023-09-08 22:00:00+00:00,9116.2500000000)
  (2023-09-09 08:00:00+00:00,10459.0000000000)
  (2023-09-09 18:00:00+00:00,10240.7500000000)
  (2023-09-10 04:00:00+00:00,7618.2500000000)
  (2023-09-10 14:00:00+00:00,8761.0000000000)
  (2023-09-11 00:00:00+00:00,8131.7500000000)
  (2023-09-11 10:00:00+00:00,11857.7500000000)
  (2023-09-11 20:00:00+00:00,9848.5000000000)
  (2023-09-12 06:00:00+00:00,12473.5000000000)
  (2023-09-12 16:00:00+00:00,12252.2500000000)
  (2023-09-13 02:00:00+00:00,9283.5000000000)
  (2023-09-13 12:00:00+00:00,12361.0000000000)
  (2023-09-13 22:00:00+00:00,10021.2500000000)
  (2023-09-14 08:00:00+00:00,12315.7500000000)
  (2023-09-14 18:00:00+00:00,11704.5000000000)
  (2023-09-15 04:00:00+00:00,10870.7500000000)
  (2023-09-15 14:00:00+00:00,11233.2500000000)
  (2023-09-16 00:00:00+00:00,8605.2500000000)
  (2023-09-16 10:00:00+00:00,10321.2500000000)
  (2023-09-16 20:00:00+00:00,9025.0000000000)
  (2023-09-17 06:00:00+00:00,8585.5000000000)
  (2023-09-17 16:00:00+00:00,10557.0000000000)
  (2023-09-18 02:00:00+00:00,9494.0000000000)
  (2023-09-18 12:00:00+00:00,13355.2500000000)
  (2023-09-18 22:00:00+00:00,9815.5000000000)
  (2023-09-19 08:00:00+00:00,12684.5000000000)
  (2023-09-19 18:00:00+00:00,12826.7500000000)
  (2023-09-20 04:00:00+00:00,10968.2500000000)
  (2023-09-20 14:00:00+00:00,11370.0000000000)
  (2023-09-21 00:00:00+00:00,9346.5000000000)
  (2023-09-21 10:00:00+00:00,12447.2500000000)
  (2023-09-21 20:00:00+00:00,11870.0000000000)
  (2023-09-22 06:00:00+00:00,12824.0000000000)
  (2023-09-22 16:00:00+00:00,11907.7500000000)
  (2023-09-23 02:00:00+00:00,7995.5000000000)
  (2023-09-23 12:00:00+00:00,10176.7500000000)
  (2023-09-23 22:00:00+00:00,8974.2500000000)
  (2023-09-24 08:00:00+00:00,9462.5000000000)
  (2023-09-24 18:00:00+00:00,9836.0000000000)
  (2023-09-25 04:00:00+00:00,11513.7500000000)
  (2023-09-25 14:00:00+00:00,11670.2500000000)
  (2023-09-26 00:00:00+00:00,8783.0000000000)
  (2023-09-26 10:00:00+00:00,11955.5000000000)
  (2023-09-26 20:00:00+00:00,10561.7500000000)
  (2023-09-27 06:00:00+00:00,12395.0000000000)
  (2023-09-27 16:00:00+00:00,12311.0000000000)
  (2023-09-28 02:00:00+00:00,9221.7500000000)
  (2023-09-28 12:00:00+00:00,11959.5000000000)
  (2023-09-28 22:00:00+00:00,9377.0000000000)
  (2023-09-29 08:00:00+00:00,12556.2500000000)
  (2023-09-29 18:00:00+00:00,12386.0000000000)
  (2023-09-30 04:00:00+00:00,10169.2500000000)
  (2023-09-30 14:00:00+00:00,10581.2500000000)
  (2023-10-01 00:00:00+00:00,7942.5000000000)
  (2023-10-01 10:00:00+00:00,9583.7500000000)
  (2023-10-01 20:00:00+00:00,9267.0000000000)
  (2023-10-02 06:00:00+00:00,11171.7500000000)
  (2023-10-02 16:00:00+00:00,11739.0000000000)
  (2023-10-03 02:00:00+00:00,7782.5000000000)
  (2023-10-03 12:00:00+00:00,10716.2500000000)
  (2023-10-03 22:00:00+00:00,8917.7500000000)
  (2023-10-04 08:00:00+00:00,12698.5000000000)
  (2023-10-04 18:00:00+00:00,13126.2500000000)
  (2023-10-05 04:00:00+00:00,11998.5000000000)
  (2023-10-05 14:00:00+00:00,13181.7500000000)
  (2023-10-06 00:00:00+00:00,8760.5000000000)
  (2023-10-06 10:00:00+00:00,13838.7500000000)
  (2023-10-06 20:00:00+00:00,10178.5000000000)
  (2023-10-07 06:00:00+00:00,11358.7500000000)
  (2023-10-07 16:00:00+00:00,11351.0000000000)
  (2023-10-08 02:00:00+00:00,7609.0000000000)
  (2023-10-08 12:00:00+00:00,9490.2500000000)
  (2023-10-08 22:00:00+00:00,8116.2500000000)
  (2023-10-09 08:00:00+00:00,12709.0000000000)
  (2023-10-09 18:00:00+00:00,11976.2500000000)
  (2023-10-10 04:00:00+00:00,11375.2500000000)
  (2023-10-10 14:00:00+00:00,12330.0000000000)
  (2023-10-11 00:00:00+00:00,9096.7500000000)
  (2023-10-11 10:00:00+00:00,13340.5000000000)
  (2023-10-11 20:00:00+00:00,11684.0000000000)
  (2023-10-12 06:00:00+00:00,13667.5000000000)
  (2023-10-12 16:00:00+00:00,13053.7500000000)
  (2023-10-13 02:00:00+00:00,8848.2500000000)
  (2023-10-13 12:00:00+00:00,12493.0000000000)
  (2023-10-13 22:00:00+00:00,9241.0000000000)
  (2023-10-14 08:00:00+00:00,11434.0000000000)
  (2023-10-14 18:00:00+00:00,10799.7500000000)
  (2023-10-15 04:00:00+00:00,8713.7500000000)
  (2023-10-15 14:00:00+00:00,10339.2500000000)
  (2023-10-16 00:00:00+00:00,9678.5000000000)
  (2023-10-16 10:00:00+00:00,13359.7500000000)
  (2023-10-16 20:00:00+00:00,10723.2500000000)
  (2023-10-17 06:00:00+00:00,13091.2500000000)
  (2023-10-17 16:00:00+00:00,13643.5000000000)
  (2023-10-18 02:00:00+00:00,10052.7500000000)
  (2023-10-18 12:00:00+00:00,11878.7500000000)
  (2023-10-18 22:00:00+00:00,11251.5000000000)
  (2023-10-19 08:00:00+00:00,14594.7500000000)
  (2023-10-19 18:00:00+00:00,13221.2500000000)
  (2023-10-20 04:00:00+00:00,12188.7500000000)
  (2023-10-20 14:00:00+00:00,13352.0000000000)
  (2023-10-21 00:00:00+00:00,9444.5000000000)
  (2023-10-21 10:00:00+00:00,11001.7500000000)
  (2023-10-21 20:00:00+00:00,9247.7500000000)
  (2023-10-22 06:00:00+00:00,9481.2500000000)
  (2023-10-22 16:00:00+00:00,11809.0000000000)
  (2023-10-23 02:00:00+00:00,9651.7500000000)
  (2023-10-23 12:00:00+00:00,12450.0000000000)
  (2023-10-23 22:00:00+00:00,9707.2500000000)
  (2023-10-24 08:00:00+00:00,14257.0000000000)
  (2023-10-24 18:00:00+00:00,13527.0000000000)
  (2023-10-25 04:00:00+00:00,12567.7500000000)
  (2023-10-25 14:00:00+00:00,11907.7500000000)
  (2023-10-26 00:00:00+00:00,9745.2500000000)
  (2023-10-26 10:00:00+00:00,13946.5000000000)
  (2023-10-26 20:00:00+00:00,11415.7500000000)
  (2023-10-27 06:00:00+00:00,13503.5000000000)
  (2023-10-27 16:00:00+00:00,13561.5000000000)
  (2023-10-28 02:00:00+00:00,8681.2500000000)
  (2023-10-28 12:00:00+00:00,10443.5000000000)
  (2023-10-28 22:00:00+00:00,8917.5000000000)
  (2023-10-29 08:00:00+00:00,11955.2500000000)
  (2023-10-29 18:00:00+00:00,12173.2500000000)
  (2023-10-30 04:00:00+00:00,10160.2500000000)
  (2023-10-30 14:00:00+00:00,11589.5000000000)
  (2023-10-31 00:00:00+00:00,8045.2500000000)
  (2023-10-31 10:00:00+00:00,12087.2500000000)
  (2023-10-31 20:00:00+00:00,11789.5000000000)
  (2023-11-01 06:00:00+00:00,11554.5000000000)
  (2023-11-01 16:00:00+00:00,13931.0000000000)
  (2023-11-02 02:00:00+00:00,9915.5000000000)
  (2023-11-02 12:00:00+00:00,14084.2500000000)
  (2023-11-02 22:00:00+00:00,10632.2500000000)
  (2023-11-03 08:00:00+00:00,13031.5000000000)
  (2023-11-03 18:00:00+00:00,12249.5000000000)
  (2023-11-04 04:00:00+00:00,9906.2500000000)
  (2023-11-04 14:00:00+00:00,11344.2500000000)
  (2023-11-05 00:00:00+00:00,9726.0000000000)
  (2023-11-05 10:00:00+00:00,11432.7500000000)
  (2023-11-05 20:00:00+00:00,10571.7500000000)
  (2023-11-06 06:00:00+00:00,13362.5000000000)
  (2023-11-06 16:00:00+00:00,14818.0000000000)
  (2023-11-07 02:00:00+00:00,10525.7500000000)
  (2023-11-07 12:00:00+00:00,14260.7500000000)
  (2023-11-07 22:00:00+00:00,10978.7500000000)
  (2023-11-08 08:00:00+00:00,14795.7500000000)
  (2023-11-08 18:00:00+00:00,14656.2500000000)
  (2023-11-09 04:00:00+00:00,12063.5000000000)
  (2023-11-09 14:00:00+00:00,14578.0000000000)
  (2023-11-10 00:00:00+00:00,10063.2500000000)
  (2023-11-10 10:00:00+00:00,14611.5000000000)
  (2023-11-10 20:00:00+00:00,11923.0000000000)
  (2023-11-11 06:00:00+00:00,9886.7500000000)
  (2023-11-11 16:00:00+00:00,12026.5000000000)
  (2023-11-12 02:00:00+00:00,8132.2500000000)
  (2023-11-12 12:00:00+00:00,10933.2500000000)
  (2023-11-12 22:00:00+00:00,9686.2500000000)
  (2023-11-13 08:00:00+00:00,13199.0000000000)
  (2023-11-13 18:00:00+00:00,14314.2500000000)
  (2023-11-14 04:00:00+00:00,11387.0000000000)
  (2023-11-14 14:00:00+00:00,13895.7500000000)
  (2023-11-15 00:00:00+00:00,10309.2500000000)
  (2023-11-15 10:00:00+00:00,14284.0000000000)
  (2023-11-15 20:00:00+00:00,12630.7500000000)
  (2023-11-16 06:00:00+00:00,15045.5000000000)
  (2023-11-16 16:00:00+00:00,14153.7500000000)
  (2023-11-17 02:00:00+00:00,11073.5000000000)
  (2023-11-17 12:00:00+00:00,14369.2500000000)
  (2023-11-17 22:00:00+00:00,11334.0000000000)
  (2023-11-18 08:00:00+00:00,11384.7500000000)
  (2023-11-18 18:00:00+00:00,11781.5000000000)
  (2023-11-19 04:00:00+00:00,10328.5000000000)
  (2023-11-19 14:00:00+00:00,11186.5000000000)
  (2023-11-20 00:00:00+00:00,8939.2500000000)
  (2023-11-20 10:00:00+00:00,14024.2500000000)
  (2023-11-20 20:00:00+00:00,12113.2500000000)
  (2023-11-21 06:00:00+00:00,14311.5000000000)
  (2023-11-21 16:00:00+00:00,15187.0000000000)
  (2023-11-22 02:00:00+00:00,10189.5000000000)
  (2023-11-22 12:00:00+00:00,14220.5000000000)
  (2023-11-22 22:00:00+00:00,11425.7500000000)
  (2023-11-23 08:00:00+00:00,14046.5000000000)
  (2023-11-23 18:00:00+00:00,14137.7500000000)
  (2023-11-24 04:00:00+00:00,12287.2500000000)
  (2023-11-24 14:00:00+00:00,14000.7500000000)
  (2023-11-25 00:00:00+00:00,11461.0000000000)
  (2023-11-25 10:00:00+00:00,12896.0000000000)
  (2023-11-25 20:00:00+00:00,11912.0000000000)
  (2023-11-26 06:00:00+00:00,10693.0000000000)
  (2023-11-26 16:00:00+00:00,12966.7500000000)
  (2023-11-27 02:00:00+00:00,9836.0000000000)
  (2023-11-27 12:00:00+00:00,14056.2500000000)
  (2023-11-27 22:00:00+00:00,12704.0000000000)
  (2023-11-28 08:00:00+00:00,15113.2500000000)
  (2023-11-28 18:00:00+00:00,14615.2500000000)
  (2023-11-29 04:00:00+00:00,12456.7500000000)
  (2023-11-29 14:00:00+00:00,15016.5000000000)
  (2023-11-30 00:00:00+00:00,11165.2500000000)
  (2023-11-30 10:00:00+00:00,14611.5000000000)
  (2023-11-30 20:00:00+00:00,13290.0000000000)
  (2023-12-01 06:00:00+00:00,13989.5000000000)
  (2023-12-01 16:00:00+00:00,14597.2500000000)
  (2023-12-02 02:00:00+00:00,10040.5000000000)
  (2023-12-02 12:00:00+00:00,12654.2500000000)
  (2023-12-02 22:00:00+00:00,10981.0000000000)
  (2023-12-03 08:00:00+00:00,11923.7500000000)
  (2023-12-03 18:00:00+00:00,12420.5000000000)
  (2023-12-04 04:00:00+00:00,11640.7500000000)
  (2023-12-04 14:00:00+00:00,15141.7500000000)
  (2023-12-05 00:00:00+00:00,12003.2500000000)
  (2023-12-05 10:00:00+00:00,16305.0000000000)
  (2023-12-05 20:00:00+00:00,14420.0000000000)
  (2023-12-06 06:00:00+00:00,15223.0000000000)
  (2023-12-06 16:00:00+00:00,15370.2500000000)
  (2023-12-07 02:00:00+00:00,10762.2500000000)
  (2023-12-07 12:00:00+00:00,14424.2500000000)
  (2023-12-07 22:00:00+00:00,11674.5000000000)
  (2023-12-08 08:00:00+00:00,15583.2500000000)
  (2023-12-08 18:00:00+00:00,14935.2500000000)
  (2023-12-09 04:00:00+00:00,9558.2500000000)
  (2023-12-09 14:00:00+00:00,12485.0000000000)
  (2023-12-10 00:00:00+00:00,10730.2500000000)
  (2023-12-10 10:00:00+00:00,13188.7500000000)
  (2023-12-10 20:00:00+00:00,11196.7500000000)
  (2023-12-11 06:00:00+00:00,13884.2500000000)
  (2023-12-11 16:00:00+00:00,14949.7500000000)
  (2023-12-12 02:00:00+00:00,10709.2500000000)
  (2023-12-12 12:00:00+00:00,13284.2500000000)
  (2023-12-12 22:00:00+00:00,11434.7500000000)
  (2023-12-13 08:00:00+00:00,15321.2500000000)
  (2023-12-13 18:00:00+00:00,15453.7500000000)
  (2023-12-14 04:00:00+00:00,12218.2500000000)
  (2023-12-14 14:00:00+00:00,13879.2500000000)
  (2023-12-15 00:00:00+00:00,10592.2500000000)
  (2023-12-15 10:00:00+00:00,15310.0000000000)
  (2023-12-15 20:00:00+00:00,13030.0000000000)
  (2023-12-16 06:00:00+00:00,11721.0000000000)
  (2023-12-16 16:00:00+00:00,13811.5000000000)
  (2023-12-17 02:00:00+00:00,9846.2500000000)
  (2023-12-17 12:00:00+00:00,12530.7500000000)
  (2023-12-17 22:00:00+00:00,11565.5000000000)
  (2023-12-18 08:00:00+00:00,14831.7500000000)
  (2023-12-18 18:00:00+00:00,14976.5000000000)
  (2023-12-19 04:00:00+00:00,12085.5000000000)
  (2023-12-19 14:00:00+00:00,15102.2500000000)
  (2023-12-20 00:00:00+00:00,9789.7500000000)
  (2023-12-20 10:00:00+00:00,14623.0000000000)
  (2023-12-20 20:00:00+00:00,12992.2500000000)
  (2023-12-21 06:00:00+00:00,13476.0000000000)
  (2023-12-21 16:00:00+00:00,14658.5000000000)
  (2023-12-22 02:00:00+00:00,10310.2500000000)
  (2023-12-22 12:00:00+00:00,14142.0000000000)
  (2023-12-22 22:00:00+00:00,11829.2500000000)
  (2023-12-23 08:00:00+00:00,12711.0000000000)
  (2023-12-23 18:00:00+00:00,12818.2500000000)
  (2023-12-24 04:00:00+00:00,8192.0000000000)
  (2023-12-24 14:00:00+00:00,11206.2500000000)
  (2023-12-25 00:00:00+00:00,8872.2500000000)
  (2023-12-25 10:00:00+00:00,12142.7500000000)
  (2023-12-25 20:00:00+00:00,10263.5000000000)
  (2023-12-26 06:00:00+00:00,9182.0000000000)
  (2023-12-26 16:00:00+00:00,11948.0000000000)
  (2023-12-27 02:00:00+00:00,9072.5000000000)
  (2023-12-27 12:00:00+00:00,10801.5000000000)
  (2023-12-27 22:00:00+00:00,10975.7500000000)
  (2023-12-28 08:00:00+00:00,12639.2500000000)
  (2023-12-28 18:00:00+00:00,13179.2500000000)
  (2023-12-29 04:00:00+00:00,9805.7500000000)
  (2023-12-29 14:00:00+00:00,12801.7500000000)
  (2023-12-30 00:00:00+00:00,9652.7500000000)
  (2023-12-30 10:00:00+00:00,12090.2500000000)
  (2023-12-30 20:00:00+00:00,10385.2500000000)
  (2023-12-31 06:00:00+00:00,8519.0000000000)
  (2023-12-31 16:00:00+00:00,11984.2500000000)
  (2024-01-01 02:00:00+00:00,8130.7500000000)
  (2024-01-01 12:00:00+00:00,9332.2500000000)
  (2024-01-01 22:00:00+00:00,8878.2500000000)
  (2024-01-02 08:00:00+00:00,12202.2500000000)
  (2024-01-02 18:00:00+00:00,13914.7500000000)
  (2024-01-03 04:00:00+00:00,11075.5000000000)
  (2024-01-03 14:00:00+00:00,13804.0000000000)
  (2024-01-04 00:00:00+00:00,10475.7500000000)
  (2024-01-04 10:00:00+00:00,15027.5000000000)
  (2024-01-04 20:00:00+00:00,12201.5000000000)
  (2024-01-05 06:00:00+00:00,12754.0000000000)
  (2024-01-05 16:00:00+00:00,14953.2500000000)
  (2024-01-06 02:00:00+00:00,10575.2500000000)
  (2024-01-06 12:00:00+00:00,13559.0000000000)
  (2024-01-06 22:00:00+00:00,10960.7500000000)
  (2024-01-07 08:00:00+00:00,11864.7500000000)
  (2024-01-07 18:00:00+00:00,13587.7500000000)
  (2024-01-08 04:00:00+00:00,12054.5000000000)
  (2024-01-08 14:00:00+00:00,13714.0000000000)
  (2024-01-09 00:00:00+00:00,10910.2500000000)
  (2024-01-09 10:00:00+00:00,14478.5000000000)
  (2024-01-09 20:00:00+00:00,12886.7500000000)
  (2024-01-10 06:00:00+00:00,14084.0000000000)
  (2024-01-10 16:00:00+00:00,14923.2500000000)
  (2024-01-11 02:00:00+00:00,11926.7500000000)
  (2024-01-11 12:00:00+00:00,15266.7500000000)
  (2024-01-11 22:00:00+00:00,11551.0000000000)
  (2024-01-12 08:00:00+00:00,14904.5000000000)
  (2024-01-12 18:00:00+00:00,14471.2500000000)
  (2024-01-13 04:00:00+00:00,10730.7500000000)
  (2024-01-13 14:00:00+00:00,13909.2500000000)
  (2024-01-14 00:00:00+00:00,10050.7500000000)
  (2024-01-14 10:00:00+00:00,13660.2500000000)
  (2024-01-14 20:00:00+00:00,12888.7500000000)
  (2024-01-15 06:00:00+00:00,15209.2500000000)
  (2024-01-15 16:00:00+00:00,15992.7500000000)
  (2024-01-16 02:00:00+00:00,11031.2500000000)
  (2024-01-16 12:00:00+00:00,14328.2500000000)
  (2024-01-16 22:00:00+00:00,12549.7500000000)
  (2024-01-17 08:00:00+00:00,16450.5000000000)
  (2024-01-17 18:00:00+00:00,14844.5000000000)
  (2024-01-18 04:00:00+00:00,12217.7500000000)
  (2024-01-18 14:00:00+00:00,14778.0000000000)
  (2024-01-19 00:00:00+00:00,12517.7500000000)
  (2024-01-19 10:00:00+00:00,15933.0000000000)
  (2024-01-19 20:00:00+00:00,13457.5000000000)
  (2024-01-20 06:00:00+00:00,12623.2500000000)
  (2024-01-20 16:00:00+00:00,13263.2500000000)
  (2024-01-21 02:00:00+00:00,10823.2500000000)
  (2024-01-21 12:00:00+00:00,13841.2500000000)
  (2024-01-21 22:00:00+00:00,11660.7500000000)
  (2024-01-22 08:00:00+00:00,15712.0000000000)
  (2024-01-22 18:00:00+00:00,15350.7500000000)
  (2024-01-23 04:00:00+00:00,12203.5000000000)
  (2024-01-23 14:00:00+00:00,15201.7500000000)
  (2024-01-24 00:00:00+00:00,12161.7500000000)
  (2024-01-24 10:00:00+00:00,14461.7500000000)
  (2024-01-24 20:00:00+00:00,12880.5000000000)
  (2024-01-25 06:00:00+00:00,14574.0000000000)
  (2024-01-25 16:00:00+00:00,14698.7500000000)
  (2024-01-26 02:00:00+00:00,10007.7500000000)
  (2024-01-26 12:00:00+00:00,14865.7500000000)
  (2024-01-26 22:00:00+00:00,11352.5000000000)
  (2024-01-27 08:00:00+00:00,13128.2500000000)
  (2024-01-27 18:00:00+00:00,12975.7500000000)
  (2024-01-28 04:00:00+00:00,9342.0000000000)
  (2024-01-28 14:00:00+00:00,10439.5000000000)
  (2024-01-29 00:00:00+00:00,11222.7500000000)
  (2024-01-29 10:00:00+00:00,14189.5000000000)
  (2024-01-29 20:00:00+00:00,12893.5000000000)
  (2024-01-30 06:00:00+00:00,13463.2500000000)
  (2024-01-30 16:00:00+00:00,14721.7500000000)
  (2024-01-31 02:00:00+00:00,11720.7500000000)
  (2024-01-31 12:00:00+00:00,15247.5000000000)
  (2024-01-31 22:00:00+00:00,12869.7500000000)
  (2024-02-01 08:00:00+00:00,15282.5000000000)
  (2024-02-01 18:00:00+00:00,15823.5000000000)
  (2024-02-02 04:00:00+00:00,12573.0000000000)
  (2024-02-02 14:00:00+00:00,14807.2500000000)
  (2024-02-03 00:00:00+00:00,11228.7500000000)
  (2024-02-03 10:00:00+00:00,14261.7500000000)
  (2024-02-03 20:00:00+00:00,12295.2500000000)
  (2024-02-04 06:00:00+00:00,10989.0000000000)
  (2024-02-04 16:00:00+00:00,13559.2500000000)
  (2024-02-05 02:00:00+00:00,10622.2500000000)
  (2024-02-05 12:00:00+00:00,15233.0000000000)
  (2024-02-05 22:00:00+00:00,12645.0000000000)
  (2024-02-06 08:00:00+00:00,14705.2500000000)
  (2024-02-06 18:00:00+00:00,14264.0000000000)
  (2024-02-07 04:00:00+00:00,12503.0000000000)
  (2024-02-07 14:00:00+00:00,14121.0000000000)
  (2024-02-08 00:00:00+00:00,10473.2500000000)
  (2024-02-08 10:00:00+00:00,14752.5000000000)
  (2024-02-08 20:00:00+00:00,12274.0000000000)
  (2024-02-09 06:00:00+00:00,14868.2500000000)
  (2024-02-09 16:00:00+00:00,14646.2500000000)
  (2024-02-10 02:00:00+00:00,10204.7500000000)
  (2024-02-10 12:00:00+00:00,12176.2500000000)
  (2024-02-10 22:00:00+00:00,10590.0000000000)
  (2024-02-11 08:00:00+00:00,11717.5000000000)
  (2024-02-11 18:00:00+00:00,12435.7500000000)
  (2024-02-12 04:00:00+00:00,10971.7500000000)
  (2024-02-12 14:00:00+00:00,13663.5000000000)
  (2024-02-13 00:00:00+00:00,10741.7500000000)
  (2024-02-13 10:00:00+00:00,14395.7500000000)
  (2024-02-13 20:00:00+00:00,12830.7500000000)
  (2024-02-14 06:00:00+00:00,13350.5000000000)
  (2024-02-14 16:00:00+00:00,13584.5000000000)
  (2024-02-15 02:00:00+00:00,9802.2500000000)
  (2024-02-15 12:00:00+00:00,13049.0000000000)
  (2024-02-15 22:00:00+00:00,11226.5000000000)
  (2024-02-16 08:00:00+00:00,14562.7500000000)
  (2024-02-16 18:00:00+00:00,14458.5000000000)
  (2024-02-17 04:00:00+00:00,10894.2500000000)
  (2024-02-17 14:00:00+00:00,11725.7500000000)
  (2024-02-18 00:00:00+00:00,8381.7500000000)
  (2024-02-18 10:00:00+00:00,11941.5000000000)
  (2024-02-18 20:00:00+00:00,11399.2500000000)
  (2024-02-19 06:00:00+00:00,13685.7500000000)
  (2024-02-19 16:00:00+00:00,14190.5000000000)
  (2024-02-20 02:00:00+00:00,10184.7500000000)
  (2024-02-20 12:00:00+00:00,14141.5000000000)
  (2024-02-20 22:00:00+00:00,11508.0000000000)
  (2024-02-21 08:00:00+00:00,14841.7500000000)
  (2024-02-21 18:00:00+00:00,13918.0000000000)
  (2024-02-22 04:00:00+00:00,11808.2500000000)
  (2024-02-22 14:00:00+00:00,13160.5000000000)
  (2024-02-23 00:00:00+00:00,11456.7500000000)
  (2024-02-23 10:00:00+00:00,14297.0000000000)
  (2024-02-23 20:00:00+00:00,12410.0000000000)
  (2024-02-24 06:00:00+00:00,9852.5000000000)
  (2024-02-24 16:00:00+00:00,11703.7500000000)
  (2024-02-25 02:00:00+00:00,9598.0000000000)
  (2024-02-25 12:00:00+00:00,10992.0000000000)
  (2024-02-25 22:00:00+00:00,9857.0000000000)
  (2024-02-26 08:00:00+00:00,13446.0000000000)
  (2024-02-26 18:00:00+00:00,13950.5000000000)
  (2024-02-27 04:00:00+00:00,11323.7500000000)
  (2024-02-27 14:00:00+00:00,13392.2500000000)
  (2024-02-28 00:00:00+00:00,9670.7500000000)
  (2024-02-28 10:00:00+00:00,13852.2500000000)
  (2024-02-28 20:00:00+00:00,12386.2500000000)
  (2024-02-29 06:00:00+00:00,14617.2500000000)
  (2024-02-29 16:00:00+00:00,14160.0000000000)
  (2024-03-01 02:00:00+00:00,10694.5000000000)
  (2024-03-01 12:00:00+00:00,12983.2500000000)
  (2024-03-01 22:00:00+00:00,9827.2500000000)
  (2024-03-02 08:00:00+00:00,11784.0000000000)
  (2024-03-02 18:00:00+00:00,12119.0000000000)
  (2024-03-03 04:00:00+00:00,8852.7500000000)
  (2024-03-03 14:00:00+00:00,10525.0000000000)
  (2024-03-04 00:00:00+00:00,9470.2500000000)
  (2024-03-04 10:00:00+00:00,14246.7500000000)
  (2024-03-04 20:00:00+00:00,13471.7500000000)
  (2024-03-05 06:00:00+00:00,14172.2500000000)
  (2024-03-05 16:00:00+00:00,14592.7500000000)
  (2024-03-06 02:00:00+00:00,10188.7500000000)
  (2024-03-06 12:00:00+00:00,12889.7500000000)
  (2024-03-06 22:00:00+00:00,11551.7500000000)
  (2024-03-07 08:00:00+00:00,13950.0000000000)
  (2024-03-07 18:00:00+00:00,14014.5000000000)
  (2024-03-08 04:00:00+00:00,10613.2500000000)
  (2024-03-08 14:00:00+00:00,11532.7500000000)
  (2024-03-09 00:00:00+00:00,10169.0000000000)
  (2024-03-09 10:00:00+00:00,12918.7500000000)
  (2024-03-09 20:00:00+00:00,11848.5000000000)
  (2024-03-10 06:00:00+00:00,10773.0000000000)
  (2024-03-10 16:00:00+00:00,12540.7500000000)
  (2024-03-11 02:00:00+00:00,10584.7500000000)
  (2024-03-11 12:00:00+00:00,13918.5000000000)
  (2024-03-11 22:00:00+00:00,10574.2500000000)
  (2024-03-12 08:00:00+00:00,13276.0000000000)
  (2024-03-12 18:00:00+00:00,13221.7500000000)
  (2024-03-13 04:00:00+00:00,10567.0000000000)
  (2024-03-13 14:00:00+00:00,12417.5000000000)
  (2024-03-14 00:00:00+00:00,9925.0000000000)
  (2024-03-14 10:00:00+00:00,13990.7500000000)
  (2024-03-14 20:00:00+00:00,13227.2500000000)
  (2024-03-15 06:00:00+00:00,13578.7500000000)
  (2024-03-15 16:00:00+00:00,12905.7500000000)
  (2024-03-16 02:00:00+00:00,9794.2500000000)
  (2024-03-16 12:00:00+00:00,13013.2500000000)
  (2024-03-16 22:00:00+00:00,11224.2500000000)
  (2024-03-17 08:00:00+00:00,10804.7500000000)
  (2024-03-17 18:00:00+00:00,11462.7500000000)
  (2024-03-18 04:00:00+00:00,10993.2500000000)
  (2024-03-18 14:00:00+00:00,12546.7500000000)
  (2024-03-19 00:00:00+00:00,10271.0000000000)
  (2024-03-19 10:00:00+00:00,13275.5000000000)
  (2024-03-19 20:00:00+00:00,12119.2500000000)
  (2024-03-20 06:00:00+00:00,12768.7500000000)
  (2024-03-20 16:00:00+00:00,12666.2500000000)
  (2024-03-21 02:00:00+00:00,9460.2500000000)
  (2024-03-21 12:00:00+00:00,12993.7500000000)
  (2024-03-21 22:00:00+00:00,10869.0000000000)
  (2024-03-22 08:00:00+00:00,13967.5000000000)
  (2024-03-22 18:00:00+00:00,14223.2500000000)
  (2024-03-23 04:00:00+00:00,9120.5000000000)
  (2024-03-23 14:00:00+00:00,11528.7500000000)
  (2024-03-24 00:00:00+00:00,10169.7500000000)
  (2024-03-24 10:00:00+00:00,12003.2500000000)
  (2024-03-24 20:00:00+00:00,10478.2500000000)
  (2024-03-25 06:00:00+00:00,12599.5000000000)
  (2024-03-25 16:00:00+00:00,12731.5000000000)
  (2024-03-26 02:00:00+00:00,9932.2500000000)
  (2024-03-26 12:00:00+00:00,12991.2500000000)
  (2024-03-26 22:00:00+00:00,11655.7500000000)
  (2024-03-27 08:00:00+00:00,13395.5000000000)
  (2024-03-27 18:00:00+00:00,13137.0000000000)
  (2024-03-28 04:00:00+00:00,11365.2500000000)
  (2024-03-28 14:00:00+00:00,13581.0000000000)
  (2024-03-29 00:00:00+00:00,10019.5000000000)
  (2024-03-29 10:00:00+00:00,12324.0000000000)
  (2024-03-29 20:00:00+00:00,9564.2500000000)
  (2024-03-30 06:00:00+00:00,9803.7500000000)
  (2024-03-30 16:00:00+00:00,11213.5000000000)
  (2024-03-31 02:00:00+00:00,7275.5000000000)
  (2024-03-31 12:00:00+00:00,8652.2500000000)
  (2024-03-31 22:00:00+00:00,7560.2500000000)
  (2024-04-01 08:00:00+00:00,9903.0000000000)
  (2024-04-01 18:00:00+00:00,11461.2500000000)
  (2024-04-02 04:00:00+00:00,11977.0000000000)
  (2024-04-02 14:00:00+00:00,13202.2500000000)
  (2024-04-03 00:00:00+00:00,9574.5000000000)
  (2024-04-03 10:00:00+00:00,11945.5000000000)
  (2024-04-03 20:00:00+00:00,11214.2500000000)
  (2024-04-04 06:00:00+00:00,12693.0000000000)
  (2024-04-04 16:00:00+00:00,12828.5000000000)
  (2024-04-05 02:00:00+00:00,10385.7500000000)
  (2024-04-05 12:00:00+00:00,12099.2500000000)
  (2024-04-05 22:00:00+00:00,10537.0000000000)
  (2024-04-06 08:00:00+00:00,11779.0000000000)
  (2024-04-06 18:00:00+00:00,11257.5000000000)
  (2024-04-07 04:00:00+00:00,9432.7500000000)
  (2024-04-07 14:00:00+00:00,8899.2500000000)
  (2024-04-08 00:00:00+00:00,8230.5000000000)
  (2024-04-08 10:00:00+00:00,12016.2500000000)
  (2024-04-08 20:00:00+00:00,10241.5000000000)
  (2024-04-09 06:00:00+00:00,12814.7500000000)
  (2024-04-09 16:00:00+00:00,13717.0000000000)
  (2024-04-10 02:00:00+00:00,10348.0000000000)
  (2024-04-10 12:00:00+00:00,13183.7500000000)
  (2024-04-10 22:00:00+00:00,9690.2500000000)
  (2024-04-11 08:00:00+00:00,14140.0000000000)
  (2024-04-11 18:00:00+00:00,13114.2500000000)
  (2024-04-12 04:00:00+00:00,12008.7500000000)
  (2024-04-12 14:00:00+00:00,12797.5000000000)
  (2024-04-13 00:00:00+00:00,9155.0000000000)
  (2024-04-13 10:00:00+00:00,10822.2500000000)
  (2024-04-13 20:00:00+00:00,9729.2500000000)
  (2024-04-14 06:00:00+00:00,11005.0000000000)
  (2024-04-14 16:00:00+00:00,11741.5000000000)
  (2024-04-15 02:00:00+00:00,8897.5000000000)
  (2024-04-15 12:00:00+00:00,11599.2500000000)
  (2024-04-15 22:00:00+00:00,11079.5000000000)
  (2024-04-16 08:00:00+00:00,12971.5000000000)
  (2024-04-16 18:00:00+00:00,12346.0000000000)
  (2024-04-17 04:00:00+00:00,11079.7500000000)
  (2024-04-17 14:00:00+00:00,11718.7500000000)
  (2024-04-18 00:00:00+00:00,9103.5000000000)
  (2024-04-18 10:00:00+00:00,12638.5000000000)
  (2024-04-18 20:00:00+00:00,10856.7500000000)
  (2024-04-19 06:00:00+00:00,12976.2500000000)
  (2024-04-19 16:00:00+00:00,13015.5000000000)
  (2024-04-20 02:00:00+00:00,10146.0000000000)
  (2024-04-20 12:00:00+00:00,10081.7500000000)
  (2024-04-20 22:00:00+00:00,8653.5000000000)
  (2024-04-21 08:00:00+00:00,11150.7500000000)
  (2024-04-21 18:00:00+00:00,10618.2500000000)
  (2024-04-22 04:00:00+00:00,11359.5000000000)
  (2024-04-22 14:00:00+00:00,11961.5000000000)
  (2024-04-23 00:00:00+00:00,9412.2500000000)
  (2024-04-23 10:00:00+00:00,12458.2500000000)
  (2024-04-23 20:00:00+00:00,10689.7500000000)
  (2024-04-24 06:00:00+00:00,13042.2500000000)
  (2024-04-24 16:00:00+00:00,12628.5000000000)
  (2024-04-25 02:00:00+00:00,10029.5000000000)
  (2024-04-25 12:00:00+00:00,12416.5000000000)
  (2024-04-25 22:00:00+00:00,9858.5000000000)
  (2024-04-26 08:00:00+00:00,12628.7500000000)
  (2024-04-26 18:00:00+00:00,11476.2500000000)
  (2024-04-27 04:00:00+00:00,8858.0000000000)
  (2024-04-27 14:00:00+00:00,9968.0000000000)
  (2024-04-28 00:00:00+00:00,9105.5000000000)
  (2024-04-28 10:00:00+00:00,10205.5000000000)
  (2024-04-28 20:00:00+00:00,10797.2500000000)
  (2024-04-29 06:00:00+00:00,12353.7500000000)
  (2024-04-29 16:00:00+00:00,11607.7500000000)
  (2024-04-30 02:00:00+00:00,9717.2500000000)
  (2024-04-30 12:00:00+00:00,11256.2500000000)
  (2024-04-30 22:00:00+00:00,9592.2500000000)
  (2024-05-01 08:00:00+00:00,9947.0000000000)
  (2024-05-01 18:00:00+00:00,10803.7500000000)
  (2024-05-02 04:00:00+00:00,11325.5000000000)
  (2024-05-02 14:00:00+00:00,12165.7500000000)
  (2024-05-03 00:00:00+00:00,9207.2500000000)
  (2024-05-03 10:00:00+00:00,12788.7500000000)
  (2024-05-03 20:00:00+00:00,9789.2500000000)
  (2024-05-04 06:00:00+00:00,9730.0000000000)
  (2024-05-04 16:00:00+00:00,10073.7500000000)
  (2024-05-05 02:00:00+00:00,7431.7500000000)
  (2024-05-05 12:00:00+00:00,9795.0000000000)
  (2024-05-05 22:00:00+00:00,8814.5000000000)
  (2024-05-06 08:00:00+00:00,12471.0000000000)
  (2024-05-06 18:00:00+00:00,11474.2500000000)
  (2024-05-07 04:00:00+00:00,11244.2500000000)
  (2024-05-07 14:00:00+00:00,11784.0000000000)
  (2024-05-08 00:00:00+00:00,8544.5000000000)
  (2024-05-08 10:00:00+00:00,12001.0000000000)
  (2024-05-08 20:00:00+00:00,9830.7500000000)
  (2024-05-09 06:00:00+00:00,9504.7500000000)
  (2024-05-09 16:00:00+00:00,9730.5000000000)
  (2024-05-10 02:00:00+00:00,8291.0000000000)
  (2024-05-10 12:00:00+00:00,9658.0000000000)
  (2024-05-10 22:00:00+00:00,8050.0000000000)
  (2024-05-11 08:00:00+00:00,9899.2500000000)
  (2024-05-11 18:00:00+00:00,9371.2500000000)
  (2024-05-12 04:00:00+00:00,7489.7500000000)
  (2024-05-12 14:00:00+00:00,8292.5000000000)
  (2024-05-13 00:00:00+00:00,8482.2500000000)
  (2024-05-13 10:00:00+00:00,11962.2500000000)
  (2024-05-13 20:00:00+00:00,11197.7500000000)
  (2024-05-14 06:00:00+00:00,12217.0000000000)
  (2024-05-14 16:00:00+00:00,12605.7500000000)
  (2024-05-15 02:00:00+00:00,9458.0000000000)
  (2024-05-15 12:00:00+00:00,11582.7500000000)
  (2024-05-15 22:00:00+00:00,10157.5000000000)
  (2024-05-16 08:00:00+00:00,12475.5000000000)
  (2024-05-16 18:00:00+00:00,12872.0000000000)
  (2024-05-17 04:00:00+00:00,11902.0000000000)
  (2024-05-17 14:00:00+00:00,11816.5000000000)
  (2024-05-18 00:00:00+00:00,8760.7500000000)
  (2024-05-18 10:00:00+00:00,10862.2500000000)
  (2024-05-18 20:00:00+00:00,9601.2500000000)
  (2024-05-19 06:00:00+00:00,8583.5000000000)
  (2024-05-19 16:00:00+00:00,9139.7500000000)
  (2024-05-20 02:00:00+00:00,7176.2500000000)
  (2024-05-20 12:00:00+00:00,9259.2500000000)
  (2024-05-20 22:00:00+00:00,9320.2500000000)
  (2024-05-21 08:00:00+00:00,12618.7500000000)
  (2024-05-21 18:00:00+00:00,12314.5000000000)
  (2024-05-22 04:00:00+00:00,11870.2500000000)
  (2024-05-22 14:00:00+00:00,13341.7500000000)
  (2024-05-23 00:00:00+00:00,8754.5000000000)
  (2024-05-23 10:00:00+00:00,12047.2500000000)
  (2024-05-23 20:00:00+00:00,10245.5000000000)
  (2024-05-24 06:00:00+00:00,12465.0000000000)
  (2024-05-24 16:00:00+00:00,12069.7500000000)
  (2024-05-25 02:00:00+00:00,8127.2500000000)
  (2024-05-25 12:00:00+00:00,9829.5000000000)
  (2024-05-25 22:00:00+00:00,8100.5000000000)
  (2024-05-26 08:00:00+00:00,9781.5000000000)
  (2024-05-26 18:00:00+00:00,10170.0000000000)
  (2024-05-27 04:00:00+00:00,10945.7500000000)
  (2024-05-27 14:00:00+00:00,11397.5000000000)
  (2024-05-28 00:00:00+00:00,8495.0000000000)
  (2024-05-28 10:00:00+00:00,12694.7500000000)
  (2024-05-28 20:00:00+00:00,9960.2500000000)
  (2024-05-29 06:00:00+00:00,12516.7500000000)
  (2024-05-29 16:00:00+00:00,13034.5000000000)
  (2024-05-30 02:00:00+00:00,8517.2500000000)
  (2024-05-30 12:00:00+00:00,11577.7500000000)
  (2024-05-30 22:00:00+00:00,9679.7500000000)
  (2024-05-31 08:00:00+00:00,12860.5000000000)
  (2024-05-31 18:00:00+00:00,11416.7500000000)
  (2024-06-01 04:00:00+00:00,8555.2500000000)
  (2024-06-01 14:00:00+00:00,10379.7500000000)
  (2024-06-02 00:00:00+00:00,7589.0000000000)
  (2024-06-02 10:00:00+00:00,10316.7500000000)
  (2024-06-02 20:00:00+00:00,9752.5000000000)
  (2024-06-03 06:00:00+00:00,12531.7500000000)
  (2024-06-03 16:00:00+00:00,12559.2500000000)
  (2024-06-04 02:00:00+00:00,9102.2500000000)
  (2024-06-04 12:00:00+00:00,11189.0000000000)
  (2024-06-04 22:00:00+00:00,8986.2500000000)
  (2024-06-05 08:00:00+00:00,13126.0000000000)
  (2024-06-05 18:00:00+00:00,11643.5000000000)
  (2024-06-06 04:00:00+00:00,10932.7500000000)
  (2024-06-06 14:00:00+00:00,11397.2500000000)
  (2024-06-07 00:00:00+00:00,8771.5000000000)
  (2024-06-07 10:00:00+00:00,11697.7500000000)
  (2024-06-07 20:00:00+00:00,10147.5000000000)
  (2024-06-08 06:00:00+00:00,9751.2500000000)
  (2024-06-08 16:00:00+00:00,10215.7500000000)
  (2024-06-09 02:00:00+00:00,8437.5000000000)
  (2024-06-09 12:00:00+00:00,9176.2500000000)
  (2024-06-09 22:00:00+00:00,9261.0000000000)
  (2024-06-10 08:00:00+00:00,12544.2500000000)
  (2024-06-10 18:00:00+00:00,11398.5000000000)
  (2024-06-11 04:00:00+00:00,11156.0000000000)
  (2024-06-11 14:00:00+00:00,12740.0000000000)
  (2024-06-12 00:00:00+00:00,9891.2500000000)
  (2024-06-12 10:00:00+00:00,13414.7500000000)
  (2024-06-12 20:00:00+00:00,11405.7500000000)
  (2024-06-13 06:00:00+00:00,12813.5000000000)
  (2024-06-13 16:00:00+00:00,12555.5000000000)
  (2024-06-14 02:00:00+00:00,8846.7500000000)
  (2024-06-14 12:00:00+00:00,12528.5000000000)
  (2024-06-14 22:00:00+00:00,9771.7500000000)
  (2024-06-15 08:00:00+00:00,11200.5000000000)
  (2024-06-15 18:00:00+00:00,9548.5000000000)
  (2024-06-16 04:00:00+00:00,7874.0000000000)
  (2024-06-16 14:00:00+00:00,9073.2500000000)
  (2024-06-17 00:00:00+00:00,8535.0000000000)
  (2024-06-17 10:00:00+00:00,12205.2500000000)
  (2024-06-17 20:00:00+00:00,10763.0000000000)
  (2024-06-18 06:00:00+00:00,12643.5000000000)
  (2024-06-18 16:00:00+00:00,12670.5000000000)
  (2024-06-19 02:00:00+00:00,9441.5000000000)
  (2024-06-19 12:00:00+00:00,12609.2500000000)
  (2024-06-19 22:00:00+00:00,10370.0000000000)
  (2024-06-20 08:00:00+00:00,12428.5000000000)
  (2024-06-20 18:00:00+00:00,11278.0000000000)
  (2024-06-21 04:00:00+00:00,10458.2500000000)
  (2024-06-21 14:00:00+00:00,12170.7500000000)
  (2024-06-22 00:00:00+00:00,8788.2500000000)
  (2024-06-22 10:00:00+00:00,11960.0000000000)
  (2024-06-22 20:00:00+00:00,10264.2500000000)
  (2024-06-23 06:00:00+00:00,9743.7500000000)
  (2024-06-23 16:00:00+00:00,10291.2500000000)
  (2024-06-24 02:00:00+00:00,8292.5000000000)
  (2024-06-24 12:00:00+00:00,11265.5000000000)
  (2024-06-24 22:00:00+00:00,9031.0000000000)
  (2024-06-25 08:00:00+00:00,11953.7500000000)
  (2024-06-25 18:00:00+00:00,11481.7500000000)
  (2024-06-26 04:00:00+00:00,10749.5000000000)
  (2024-06-26 14:00:00+00:00,11717.2500000000)
  (2024-06-27 00:00:00+00:00,9369.7500000000)
  (2024-06-27 10:00:00+00:00,12715.7500000000)
  (2024-06-27 20:00:00+00:00,11132.5000000000)
  (2024-06-28 06:00:00+00:00,12543.0000000000)
  (2024-06-28 16:00:00+00:00,12586.7500000000)
  (2024-06-29 02:00:00+00:00,8046.7500000000)
  (2024-06-29 12:00:00+00:00,9693.7500000000)
  (2024-06-29 22:00:00+00:00,9289.7500000000)
  (2024-06-30 08:00:00+00:00,9799.2500000000)
  (2024-06-30 18:00:00+00:00,11612.2500000000)
  (2024-07-01 04:00:00+00:00,10682.0000000000)
  (2024-07-01 14:00:00+00:00,11556.5000000000)
  (2024-07-02 00:00:00+00:00,9242.7500000000)
  (2024-07-02 10:00:00+00:00,13106.0000000000)
  (2024-07-02 20:00:00+00:00,11279.2500000000)
  (2024-07-03 06:00:00+00:00,13200.0000000000)
  (2024-07-03 16:00:00+00:00,12416.7500000000)
  (2024-07-04 02:00:00+00:00,9185.0000000000)
  (2024-07-04 12:00:00+00:00,12243.0000000000)
  (2024-07-04 22:00:00+00:00,9978.7500000000)
  (2024-07-05 08:00:00+00:00,13167.0000000000)
  (2024-07-05 18:00:00+00:00,10653.0000000000)
  (2024-07-06 04:00:00+00:00,8564.0000000000)
  (2024-07-06 14:00:00+00:00,10438.2500000000)
  (2024-07-07 00:00:00+00:00,8322.5000000000)
  (2024-07-07 10:00:00+00:00,9697.0000000000)
  (2024-07-07 20:00:00+00:00,9558.7500000000)
  (2024-07-08 06:00:00+00:00,12152.2500000000)
  (2024-07-08 16:00:00+00:00,11919.7500000000)
  (2024-07-09 02:00:00+00:00,8963.2500000000)
  (2024-07-09 12:00:00+00:00,12255.2500000000)
  (2024-07-09 22:00:00+00:00,10863.7500000000)
  (2024-07-10 08:00:00+00:00,13760.7500000000)
  (2024-07-10 18:00:00+00:00,11957.5000000000)
  (2024-07-11 04:00:00+00:00,11012.2500000000)
  (2024-07-11 14:00:00+00:00,11976.5000000000)
  (2024-07-12 00:00:00+00:00,9131.7500000000)
  (2024-07-12 10:00:00+00:00,13156.2500000000)
  (2024-07-12 20:00:00+00:00,11353.2500000000)
  (2024-07-13 06:00:00+00:00,10120.2500000000)
  (2024-07-13 16:00:00+00:00,10134.0000000000)
  (2024-07-14 02:00:00+00:00,8058.0000000000)
  (2024-07-14 12:00:00+00:00,8868.0000000000)
  (2024-07-14 22:00:00+00:00,9281.2500000000)
  (2024-07-15 08:00:00+00:00,11682.7500000000)
  (2024-07-15 18:00:00+00:00,11576.7500000000)
  (2024-07-16 04:00:00+00:00,11250.5000000000)
  (2024-07-16 14:00:00+00:00,11559.0000000000)
  (2024-07-17 00:00:00+00:00,8804.0000000000)
  (2024-07-17 10:00:00+00:00,12346.5000000000)
  (2024-07-17 20:00:00+00:00,11657.5000000000)
  (2024-07-18 06:00:00+00:00,12850.2500000000)
  (2024-07-18 16:00:00+00:00,12296.5000000000)
  (2024-07-19 02:00:00+00:00,9108.2500000000)
  (2024-07-19 12:00:00+00:00,11579.7500000000)
  (2024-07-19 22:00:00+00:00,9132.5000000000)
  (2024-07-20 08:00:00+00:00,10264.5000000000)
  (2024-07-20 18:00:00+00:00,9760.7500000000)
  (2024-07-21 04:00:00+00:00,7838.7500000000)
  (2024-07-21 14:00:00+00:00,9566.7500000000)
  (2024-07-22 00:00:00+00:00,8641.2500000000)
  (2024-07-22 10:00:00+00:00,12769.5000000000)
  (2024-07-22 20:00:00+00:00,11098.0000000000)
  (2024-07-23 06:00:00+00:00,11907.2500000000)
  (2024-07-23 16:00:00+00:00,12244.5000000000)
  (2024-07-24 02:00:00+00:00,9064.5000000000)
  (2024-07-24 12:00:00+00:00,11819.2500000000)
  (2024-07-24 22:00:00+00:00,9635.7500000000)
  (2024-07-25 08:00:00+00:00,11265.2500000000)
  (2024-07-25 18:00:00+00:00,10767.5000000000)
  (2024-07-26 04:00:00+00:00,10246.7500000000)
  (2024-07-26 14:00:00+00:00,11016.2500000000)
  (2024-07-27 00:00:00+00:00,8088.0000000000)
  (2024-07-27 10:00:00+00:00,10392.2500000000)
  (2024-07-27 20:00:00+00:00,9048.2500000000)
  (2024-07-28 06:00:00+00:00,8916.5000000000)
  (2024-07-28 16:00:00+00:00,11232.7500000000)
  (2024-07-29 02:00:00+00:00,9698.7500000000)
  (2024-07-29 12:00:00+00:00,10861.2500000000)
  (2024-07-29 22:00:00+00:00,8973.5000000000)
  (2024-07-30 08:00:00+00:00,11292.2500000000)
  (2024-07-30 18:00:00+00:00,11150.7500000000)
  (2024-07-31 04:00:00+00:00,10529.0000000000)
  (2024-07-31 14:00:00+00:00,11386.0000000000)
  (2024-08-01 00:00:00+00:00,8869.2500000000)
  (2024-08-01 10:00:00+00:00,12284.2500000000)
  (2024-08-01 20:00:00+00:00,10648.2500000000)
  (2024-08-02 06:00:00+00:00,12179.0000000000)
  (2024-08-02 16:00:00+00:00,11420.7500000000)
  (2024-08-03 02:00:00+00:00,8243.2500000000)
  (2024-08-03 12:00:00+00:00,9473.2500000000)
  (2024-08-03 22:00:00+00:00,8082.7500000000)
  (2024-08-04 08:00:00+00:00,9432.0000000000)
  (2024-08-04 18:00:00+00:00,9948.2500000000)
  (2024-08-05 04:00:00+00:00,11546.5000000000)
  (2024-08-05 14:00:00+00:00,11519.5000000000)
  (2024-08-06 00:00:00+00:00,8633.2500000000)
  (2024-08-06 10:00:00+00:00,11806.0000000000)
  (2024-08-06 20:00:00+00:00,10605.2500000000)
  (2024-08-07 06:00:00+00:00,12037.7500000000)
  (2024-08-07 16:00:00+00:00,12612.5000000000)
  (2024-08-08 02:00:00+00:00,9515.5000000000)
  (2024-08-08 12:00:00+00:00,12271.7500000000)
  (2024-08-08 22:00:00+00:00,9561.7500000000)
  (2024-08-09 08:00:00+00:00,13221.2500000000)
  (2024-08-09 18:00:00+00:00,12418.0000000000)
  (2024-08-10 04:00:00+00:00,9558.5000000000)
  (2024-08-10 14:00:00+00:00,9837.0000000000)
  (2024-08-11 00:00:00+00:00,8927.0000000000)
  (2024-08-11 10:00:00+00:00,10095.7500000000)
  (2024-08-11 20:00:00+00:00,9663.5000000000)
  (2024-08-12 06:00:00+00:00,11978.7500000000)
  (2024-08-12 16:00:00+00:00,12050.0000000000)
  (2024-08-13 02:00:00+00:00,9581.7500000000)
  (2024-08-13 12:00:00+00:00,12413.7500000000)
  (2024-08-13 22:00:00+00:00,11333.2500000000)
  (2024-08-14 08:00:00+00:00,12473.2500000000)
  (2024-08-14 18:00:00+00:00,12564.2500000000)
  (2024-08-15 04:00:00+00:00,11232.5000000000)
  (2024-08-15 14:00:00+00:00,11716.5000000000)
  (2024-08-16 00:00:00+00:00,9396.5000000000)
  (2024-08-16 10:00:00+00:00,12867.0000000000)
  (2024-08-16 20:00:00+00:00,10286.5000000000)
  (2024-08-17 06:00:00+00:00,10080.0000000000)
  (2024-08-17 16:00:00+00:00,10552.2500000000)
  (2024-08-18 02:00:00+00:00,7779.7500000000)
  (2024-08-18 12:00:00+00:00,9221.2500000000)
  (2024-08-18 22:00:00+00:00,9028.5000000000)
  (2024-08-19 08:00:00+00:00,12384.0000000000)
  (2024-08-19 18:00:00+00:00,11210.5000000000)
  (2024-08-20 04:00:00+00:00,10541.0000000000)
  (2024-08-20 14:00:00+00:00,11823.5000000000)
  (2024-08-21 00:00:00+00:00,9332.2500000000)
  (2024-08-21 10:00:00+00:00,12946.7500000000)
  (2024-08-21 20:00:00+00:00,11457.2500000000)
  (2024-08-22 06:00:00+00:00,12319.5000000000)
  (2024-08-22 16:00:00+00:00,11936.2500000000)
  (2024-08-23 02:00:00+00:00,9647.2500000000)
  (2024-08-23 12:00:00+00:00,12065.7500000000)
  (2024-08-23 22:00:00+00:00,9598.2500000000)
  (2024-08-24 08:00:00+00:00,10803.7500000000)
  (2024-08-24 18:00:00+00:00,11295.7500000000)
  (2024-08-25 04:00:00+00:00,8152.7500000000)
  (2024-08-25 14:00:00+00:00,9308.2500000000)
  (2024-08-26 00:00:00+00:00,7780.2500000000)
  (2024-08-26 10:00:00+00:00,12352.7500000000)
  (2024-08-26 20:00:00+00:00,10095.2500000000)
  (2024-08-27 06:00:00+00:00,11927.7500000000)
  (2024-08-27 16:00:00+00:00,12082.2500000000)
  (2024-08-28 02:00:00+00:00,9907.5000000000)
  (2024-08-28 12:00:00+00:00,11979.2500000000)
  (2024-08-28 22:00:00+00:00,11180.7500000000)
  (2024-08-29 08:00:00+00:00,13247.7500000000)
  (2024-08-29 18:00:00+00:00,12186.0000000000)
  (2024-08-30 04:00:00+00:00,11411.2500000000)
  (2024-08-30 14:00:00+00:00,12235.2500000000)
  (2024-08-31 00:00:00+00:00,8717.0000000000)
  (2024-08-31 10:00:00+00:00,9713.0000000000)
  (2024-08-31 20:00:00+00:00,9103.2500000000)
  (2024-09-01 06:00:00+00:00,8762.5000000000)
  (2024-09-01 16:00:00+00:00,10196.7500000000)
  (2024-09-02 02:00:00+00:00,8827.7500000000)
  (2024-09-02 12:00:00+00:00,12024.7500000000)
  (2024-09-02 22:00:00+00:00,10640.2500000000)
  (2024-09-03 08:00:00+00:00,13547.7500000000)
  (2024-09-03 18:00:00+00:00,13114.0000000000)
  (2024-09-04 04:00:00+00:00,12161.0000000000)
  (2024-09-04 14:00:00+00:00,13711.0000000000)
  (2024-09-05 00:00:00+00:00,10670.2500000000)
  (2024-09-05 10:00:00+00:00,12946.5000000000)
  (2024-09-05 20:00:00+00:00,10977.5000000000)
  (2024-09-06 06:00:00+00:00,13751.0000000000)
  (2024-09-06 16:00:00+00:00,13545.2500000000)
  (2024-09-07 02:00:00+00:00,9353.5000000000)
  (2024-09-07 12:00:00+00:00,10291.7500000000)
  (2024-09-07 22:00:00+00:00,9576.2500000000)
  (2024-09-08 08:00:00+00:00,10186.7500000000)
  (2024-09-08 18:00:00+00:00,11693.2500000000)
  (2024-09-09 04:00:00+00:00,11732.0000000000)
  (2024-09-09 14:00:00+00:00,12418.5000000000)
  (2024-09-10 00:00:00+00:00,9930.7500000000)
  (2024-09-10 10:00:00+00:00,12425.7500000000)
  (2024-09-10 20:00:00+00:00,11529.7500000000)
  (2024-09-11 06:00:00+00:00,13802.5000000000)
  (2024-09-11 16:00:00+00:00,12425.2500000000)
  (2024-09-12 02:00:00+00:00,9673.7500000000)
  (2024-09-12 12:00:00+00:00,11950.5000000000)
  (2024-09-12 22:00:00+00:00,9187.0000000000)
  (2024-09-13 08:00:00+00:00,13166.0000000000)
  (2024-09-13 18:00:00+00:00,13378.5000000000)
  (2024-09-14 04:00:00+00:00,9399.7500000000)
  (2024-09-14 14:00:00+00:00,9875.0000000000)
  (2024-09-15 00:00:00+00:00,7928.2500000000)
  (2024-09-15 10:00:00+00:00,10001.2500000000)
  (2024-09-15 20:00:00+00:00,10172.7500000000)
  (2024-09-16 06:00:00+00:00,13833.2500000000)
  (2024-09-16 16:00:00+00:00,13550.0000000000)
  (2024-09-17 02:00:00+00:00,9639.7500000000)
  (2024-09-17 12:00:00+00:00,11895.2500000000)
  (2024-09-17 22:00:00+00:00,9224.2500000000)
  (2024-09-18 08:00:00+00:00,12705.2500000000)
  (2024-09-18 18:00:00+00:00,12419.7500000000)
  (2024-09-19 04:00:00+00:00,11643.7500000000)
  (2024-09-19 14:00:00+00:00,11949.2500000000)
  (2024-09-20 00:00:00+00:00,9329.5000000000)
  (2024-09-20 10:00:00+00:00,11257.0000000000)
  (2024-09-20 20:00:00+00:00,10406.5000000000)
  (2024-09-21 06:00:00+00:00,10319.5000000000)
  (2024-09-21 16:00:00+00:00,10332.5000000000)
  (2024-09-22 02:00:00+00:00,8112.7500000000)
  (2024-09-22 12:00:00+00:00,8654.7500000000)
  (2024-09-22 22:00:00+00:00,8777.7500000000)
  (2024-09-23 08:00:00+00:00,12641.5000000000)
  (2024-09-23 18:00:00+00:00,11919.0000000000)
  (2024-09-24 04:00:00+00:00,11349.7500000000)
  (2024-09-24 14:00:00+00:00,12194.2500000000)
  (2024-09-25 00:00:00+00:00,10127.0000000000)
  (2024-09-25 10:00:00+00:00,13888.5000000000)
  (2024-09-25 20:00:00+00:00,11402.0000000000)
  (2024-09-26 06:00:00+00:00,14724.5000000000)
  (2024-09-26 16:00:00+00:00,14463.0000000000)
  (2024-09-27 02:00:00+00:00,10703.2500000000)
  (2024-09-27 12:00:00+00:00,12542.5000000000)
  (2024-09-27 22:00:00+00:00,10685.5000000000)
  (2024-09-28 08:00:00+00:00,12409.2500000000)
  (2024-09-28 18:00:00+00:00,11955.0000000000)
  (2024-09-29 04:00:00+00:00,9453.0000000000)
  (2024-09-29 14:00:00+00:00,9142.5000000000)
  (2024-09-30 00:00:00+00:00,8026.2500000000)
  (2024-09-30 10:00:00+00:00,12324.0000000000)
  (2024-09-30 20:00:00+00:00,11694.5000000000)
  (2024-10-01 06:00:00+00:00,14358.5000000000)
  (2024-10-01 16:00:00+00:00,14206.0000000000)
  (2024-10-02 02:00:00+00:00,10224.5000000000)
  (2024-10-02 12:00:00+00:00,14056.0000000000)
  (2024-10-02 22:00:00+00:00,10588.2500000000)
  (2024-10-03 08:00:00+00:00,12107.5000000000)
  (2024-10-03 18:00:00+00:00,12061.5000000000)
  (2024-10-04 04:00:00+00:00,11020.7500000000)
  (2024-10-04 14:00:00+00:00,11485.2500000000)
  (2024-10-05 00:00:00+00:00,8720.5000000000)
  (2024-10-05 10:00:00+00:00,10745.7500000000)
  (2024-10-05 20:00:00+00:00,9564.0000000000)
  (2024-10-06 06:00:00+00:00,9079.0000000000)
  (2024-10-06 16:00:00+00:00,10918.7500000000)
  (2024-10-07 02:00:00+00:00,9418.0000000000)
  (2024-10-07 12:00:00+00:00,13649.7500000000)
  (2024-10-07 22:00:00+00:00,9427.5000000000)
  (2024-10-08 08:00:00+00:00,13510.7500000000)
  (2024-10-08 18:00:00+00:00,12574.0000000000)
  (2024-10-09 04:00:00+00:00,11567.5000000000)
  (2024-10-09 14:00:00+00:00,13321.0000000000)
  (2024-10-10 00:00:00+00:00,9216.2500000000)
  (2024-10-10 10:00:00+00:00,13862.5000000000)
  (2024-10-10 20:00:00+00:00,12235.5000000000)
  (2024-10-11 06:00:00+00:00,14584.5000000000)
  (2024-10-11 16:00:00+00:00,13927.2500000000)
  (2024-10-12 02:00:00+00:00,8998.0000000000)
  (2024-10-12 12:00:00+00:00,10482.2500000000)
  (2024-10-12 22:00:00+00:00,10227.0000000000)
  (2024-10-13 08:00:00+00:00,11018.0000000000)
  (2024-10-13 18:00:00+00:00,11404.7500000000)
  (2024-10-14 04:00:00+00:00,12956.5000000000)
  (2024-10-14 14:00:00+00:00,12929.2500000000)
  (2024-10-15 00:00:00+00:00,10206.5000000000)
  (2024-10-15 10:00:00+00:00,12852.7500000000)
  (2024-10-15 20:00:00+00:00,10564.5000000000)
  (2024-10-16 06:00:00+00:00,14158.5000000000)
  (2024-10-16 16:00:00+00:00,14548.2500000000)
  (2024-10-17 02:00:00+00:00,11229.7500000000)
  (2024-10-17 12:00:00+00:00,12372.5000000000)
  (2024-10-17 22:00:00+00:00,10643.2500000000)
  (2024-10-18 08:00:00+00:00,13411.0000000000)
  (2024-10-18 18:00:00+00:00,12285.5000000000)
  (2024-10-19 04:00:00+00:00,10439.5000000000)
  (2024-10-19 14:00:00+00:00,11553.5000000000)
  (2024-10-20 00:00:00+00:00,8562.7500000000)
  (2024-10-20 10:00:00+00:00,10407.0000000000)
  (2024-10-20 20:00:00+00:00,11089.2500000000)
  (2024-10-21 06:00:00+00:00,14737.7500000000)
  (2024-10-21 16:00:00+00:00,13118.5000000000)
  (2024-10-22 02:00:00+00:00,9661.5000000000)
  (2024-10-22 12:00:00+00:00,13650.2500000000)
  (2024-10-22 22:00:00+00:00,11517.5000000000)
  (2024-10-23 08:00:00+00:00,13487.0000000000)
  (2024-10-23 18:00:00+00:00,12803.0000000000)
  (2024-10-24 04:00:00+00:00,11688.2500000000)
  (2024-10-24 14:00:00+00:00,12363.5000000000)
  (2024-10-25 00:00:00+00:00,9989.5000000000)
  (2024-10-25 10:00:00+00:00,12927.7500000000)
  (2024-10-25 20:00:00+00:00,10526.0000000000)
  (2024-10-26 06:00:00+00:00,10061.0000000000)
  (2024-10-26 16:00:00+00:00,11047.2500000000)
  (2024-10-27 02:00:00+00:00,8457.5000000000)
  (2024-10-27 12:00:00+00:00,10884.5000000000)
  (2024-10-27 22:00:00+00:00,9583.5000000000)
  (2024-10-28 08:00:00+00:00,13040.5000000000)
  (2024-10-28 18:00:00+00:00,13122.2500000000)
  (2024-10-29 04:00:00+00:00,10348.5000000000)
  (2024-10-29 14:00:00+00:00,12906.7500000000)
  (2024-10-30 00:00:00+00:00,10227.0000000000)
  (2024-10-30 10:00:00+00:00,14772.0000000000)
  (2024-10-30 20:00:00+00:00,12458.0000000000)
  (2024-10-31 06:00:00+00:00,10504.0000000000)
  (2024-10-31 16:00:00+00:00,12657.5000000000)
  (2024-11-01 02:00:00+00:00,9472.2500000000)
  (2024-11-01 12:00:00+00:00,13376.7500000000)
  (2024-11-01 22:00:00+00:00,11476.2500000000)
  (2024-11-02 08:00:00+00:00,11153.7500000000)
  (2024-11-02 18:00:00+00:00,11102.7500000000)
  (2024-11-03 04:00:00+00:00,9476.5000000000)
  (2024-11-03 14:00:00+00:00,11540.0000000000)
  (2024-11-04 00:00:00+00:00,9204.7500000000)
  (2024-11-04 10:00:00+00:00,13192.2500000000)
  (2024-11-04 20:00:00+00:00,11342.7500000000)
  (2024-11-05 06:00:00+00:00,12672.5000000000)
  (2024-11-05 16:00:00+00:00,13520.7500000000)
  (2024-11-06 02:00:00+00:00,9636.2500000000)
  (2024-11-06 12:00:00+00:00,13678.2500000000)
  (2024-11-06 22:00:00+00:00,10625.7500000000)
  (2024-11-07 08:00:00+00:00,13887.2500000000)
  (2024-11-07 18:00:00+00:00,13497.7500000000)
  (2024-11-08 04:00:00+00:00,11012.2500000000)
  (2024-11-08 14:00:00+00:00,13868.2500000000)
  (2024-11-09 00:00:00+00:00,10038.2500000000)
  (2024-11-09 10:00:00+00:00,11969.7500000000)
  (2024-11-09 20:00:00+00:00,10515.2500000000)
  (2024-11-10 06:00:00+00:00,9473.0000000000)
  (2024-11-10 16:00:00+00:00,11897.2500000000)
  (2024-11-11 02:00:00+00:00,9187.0000000000)
  (2024-11-11 12:00:00+00:00,14109.5000000000)
  (2024-11-11 22:00:00+00:00,10854.5000000000)
  (2024-11-12 08:00:00+00:00,13699.5000000000)
  (2024-11-12 18:00:00+00:00,13810.5000000000)
  (2024-11-13 04:00:00+00:00,11112.0000000000)
  (2024-11-13 14:00:00+00:00,13743.7500000000)
  (2024-11-14 00:00:00+00:00,11657.2500000000)
  (2024-11-14 10:00:00+00:00,15147.0000000000)
  (2024-11-14 20:00:00+00:00,12296.2500000000)
  (2024-11-15 06:00:00+00:00,14275.5000000000)
  (2024-11-15 16:00:00+00:00,14593.2500000000)
  (2024-11-16 02:00:00+00:00,10968.0000000000)
  (2024-11-16 12:00:00+00:00,13635.7500000000)
  (2024-11-16 22:00:00+00:00,11204.7500000000)
  (2024-11-17 08:00:00+00:00,12484.2500000000)
  (2024-11-17 18:00:00+00:00,12639.2500000000)
  (2024-11-18 04:00:00+00:00,12346.7500000000)
  (2024-11-18 14:00:00+00:00,15537.2500000000)
  (2024-11-19 00:00:00+00:00,12261.7500000000)
  (2024-11-19 10:00:00+00:00,14568.0000000000)
  (2024-11-19 20:00:00+00:00,13142.2500000000)
  (2024-11-20 06:00:00+00:00,13827.2500000000)
  (2024-11-20 16:00:00+00:00,14116.7500000000)
  (2024-11-21 02:00:00+00:00,11823.7500000000)
  (2024-11-21 12:00:00+00:00,14971.5000000000)
  (2024-11-21 22:00:00+00:00,11544.2500000000)
  (2024-11-22 08:00:00+00:00,14215.7500000000)
  (2024-11-22 18:00:00+00:00,15700.2500000000)
  (2024-11-23 04:00:00+00:00,11828.7500000000)
  (2024-11-23 14:00:00+00:00,12937.2500000000)
  (2024-11-24 00:00:00+00:00,11219.0000000000)
  (2024-11-24 10:00:00+00:00,13490.2500000000)
  (2024-11-24 20:00:00+00:00,12426.7500000000)
  (2024-11-25 06:00:00+00:00,13827.2500000000)
  (2024-11-25 16:00:00+00:00,15875.7500000000)
  (2024-11-26 02:00:00+00:00,11194.7500000000)
  (2024-11-26 12:00:00+00:00,14666.2500000000)
  (2024-11-26 22:00:00+00:00,12056.5000000000)
  (2024-11-27 08:00:00+00:00,14335.7500000000)
  (2024-11-27 18:00:00+00:00,15008.0000000000)
  (2024-11-28 04:00:00+00:00,11755.2500000000)
  (2024-11-28 14:00:00+00:00,16055.2500000000)
  (2024-11-29 00:00:00+00:00,11123.0000000000)
  (2024-11-29 10:00:00+00:00,14017.7500000000)
  (2024-11-29 20:00:00+00:00,12348.7500000000)
  (2024-11-30 06:00:00+00:00,12067.0000000000)
  (2024-11-30 16:00:00+00:00,13801.5000000000)
  (2024-12-01 02:00:00+00:00,10169.0000000000)
  (2024-12-01 12:00:00+00:00,11644.2500000000)
  (2024-12-01 22:00:00+00:00,12292.5000000000)
  (2024-12-02 08:00:00+00:00,14870.7500000000)
  (2024-12-02 18:00:00+00:00,15162.5000000000)
  (2024-12-03 04:00:00+00:00,11927.2500000000)
  (2024-12-03 14:00:00+00:00,14049.0000000000)
  (2024-12-04 00:00:00+00:00,11514.5000000000)
  (2024-12-04 10:00:00+00:00,15439.5000000000)
  (2024-12-04 20:00:00+00:00,13067.7500000000)
  (2024-12-05 06:00:00+00:00,13331.2500000000)
  (2024-12-05 16:00:00+00:00,16357.0000000000)
  (2024-12-06 02:00:00+00:00,12428.2500000000)
  (2024-12-06 12:00:00+00:00,14660.5000000000)
  (2024-12-06 22:00:00+00:00,10473.0000000000)
  (2024-12-07 08:00:00+00:00,12132.5000000000)
  (2024-12-07 18:00:00+00:00,13761.2500000000)
  (2024-12-08 04:00:00+00:00,10184.0000000000)
  (2024-12-08 14:00:00+00:00,12364.5000000000)
  (2024-12-09 00:00:00+00:00,10966.0000000000)
  (2024-12-09 10:00:00+00:00,16187.5000000000)
  (2024-12-09 20:00:00+00:00,14703.7500000000)
  (2024-12-10 06:00:00+00:00,15059.5000000000)
  (2024-12-10 16:00:00+00:00,14229.2500000000)
  (2024-12-11 02:00:00+00:00,10325.0000000000)
  (2024-12-11 12:00:00+00:00,14122.7500000000)
  (2024-12-11 22:00:00+00:00,11272.5000000000)
  (2024-12-12 08:00:00+00:00,13828.0000000000)
  (2024-12-12 18:00:00+00:00,13676.5000000000)
  (2024-12-13 04:00:00+00:00,11440.0000000000)
  (2024-12-13 14:00:00+00:00,13801.5000000000)
  (2024-12-14 00:00:00+00:00,10950.5000000000)
  (2024-12-14 10:00:00+00:00,14733.7500000000)
  (2024-12-14 20:00:00+00:00,13381.2500000000)
  (2024-12-15 06:00:00+00:00,11819.7500000000)
  (2024-12-15 16:00:00+00:00,13075.2500000000)
  (2024-12-16 02:00:00+00:00,10262.5000000000)
  (2024-12-16 12:00:00+00:00,14578.2500000000)
  (2024-12-16 22:00:00+00:00,11446.2500000000)
  (2024-12-17 08:00:00+00:00,15377.5000000000)
  (2024-12-17 18:00:00+00:00,14535.5000000000)
  (2024-12-18 04:00:00+00:00,11544.0000000000)
  (2024-12-18 14:00:00+00:00,15147.2500000000)
  (2024-12-19 00:00:00+00:00,11341.5000000000)
  (2024-12-19 10:00:00+00:00,15527.7500000000)
  (2024-12-19 20:00:00+00:00,12250.5000000000)
  (2024-12-20 06:00:00+00:00,14766.5000000000)
  (2024-12-20 16:00:00+00:00,14806.2500000000)
  (2024-12-21 02:00:00+00:00,10847.2500000000)
  (2024-12-21 12:00:00+00:00,13451.5000000000)
  (2024-12-21 22:00:00+00:00,11720.7500000000)
  (2024-12-22 08:00:00+00:00,11594.5000000000)
  (2024-12-22 18:00:00+00:00,12656.0000000000)
  (2024-12-23 04:00:00+00:00,10745.5000000000)
  (2024-12-23 14:00:00+00:00,12760.0000000000)
  (2024-12-24 00:00:00+00:00,8891.5000000000)
  (2024-12-24 10:00:00+00:00,11481.2500000000)
  (2024-12-24 20:00:00+00:00,9990.7500000000)
  (2024-12-25 06:00:00+00:00,9703.2500000000)
  (2024-12-25 16:00:00+00:00,10808.0000000000)
  (2024-12-26 02:00:00+00:00,8317.5000000000)
  (2024-12-26 12:00:00+00:00,10402.7500000000)
  (2024-12-26 22:00:00+00:00,8989.0000000000)
  (2024-12-27 08:00:00+00:00,10746.0000000000)
  (2024-12-27 18:00:00+00:00,11570.2500000000)
  (2024-12-28 04:00:00+00:00,8879.5000000000)
  (2024-12-28 14:00:00+00:00,10819.7500000000)
  (2024-12-29 00:00:00+00:00,9162.5000000000)
  (2024-12-29 10:00:00+00:00,12561.5000000000)
  (2024-12-29 20:00:00+00:00,11811.0000000000)
  (2024-12-30 06:00:00+00:00,11890.7500000000)
  (2024-12-30 16:00:00+00:00,13345.5000000000)
  (2024-12-31 02:00:00+00:00,9699.0000000000)
  (2024-12-31 12:00:00+00:00,12130.0000000000)
  (2024-12-31 22:00:00+00:00,11462.2500000000)
};
    \end{axis}
    \end{scope}
\end{tikzpicture}%
        \caption{Grid load from 2015 to 2024.}
        \label{fig:grid_load_full}
    \end{subfigure}%
    \hfill
    \begin{subfigure}[a]{0.49\linewidth}
        \centering
        \input{figures/load_annual}%
        \caption{Average annual grid load.}
        \label{fig:grid_load_year}
    \end{subfigure}
    \begin{subfigure}[b]{0.49\linewidth}
        \centering
        \begin{tikzpicture}
    \begin{groupplot}[
      group style={
        group size=2 by 2,
        horizontal sep=0.05cm,
        vertical sep=0.05cm,
      },
      small axis,
      width=0.6\linewidth,
      height=0.5\imageheight,
      grid=both,
      ymin=8000, ymax=16000,
      scaled y ticks=false,
      xtick={0,1,2,3,4,5,6},
      xticklabels={ $\mathrm{Mon.}$, $\mathrm{Tue.}$, $\mathrm{Wed.}$, $\mathrm{Thu.}$, $\mathrm{Fri.}$, $\mathrm{Sat.}$, $\mathrm{Sun.}$},
      error bars/y dir=both,
      error bars/y explicit,
    ]
    
    \nextgroupplot[
      xticklabels={},
      ytick={10000,14000},
    ]
    \addplot[name path=lowS, draw=none, color=tabblue] coordinates {
      (0,11700.5509215126-1885.4994100763)
      (1,12246.8739741162-1636.6140234017)
      (2,12215.8201104101-1598.5330199857)
      (3,12163.7301398157-1760.7676988706)
      (4,11966.0550820707-1764.6920651783)
      (5,10656.3617157490-1288.8491039263)
      (6,10022.0004752852-1326.6263477263)
    };
    \addplot[name path=highS, draw=none, color=tabblue] coordinates {
      (0,11700.5509215126+1885.4994100763)
      (1,12246.8739741162+1636.6140234017)
      (2,12215.8201104101+1598.5330199857)
      (3,12163.7301398157+1760.7676988706)
      (4,11966.0550820707+1764.6920651783)
      (5,10656.3617157490+1288.8491039263)
      (6,10022.0004752852+1326.6263477263)
    };
    \addplot[draw=none, color=tabblue, fill opacity=0.2] fill between[of=lowS and highS];
    \addplot[thick, color=tabblue] coordinates {
      (0,11700.5509215126)
      (1,12246.8739741162)
      (2,12215.8201104101)
      (3,12163.7301398157)
      (4,11966.0550820707)
      (5,10656.3617157490)
      (6,10022.0004752852)
    };
    \node[anchor=north east, font=\scriptsize] at (axis description cs:1,1) {Spring};
    
    \nextgroupplot[
      yticklabels={},
      xticklabels={},
    ]
    \addplot[name path=lowSu, draw=none, color=tabblue] coordinates {
      (0,11289.4928097345-1758.9508447551)
      (1,11576.0158137317-1557.4125577151)
      (2,11626.8676563487-1537.7452479029)
      (3,11620.4147727273-1521.0103018355)
      (4,11463.5467260013-1522.4054897042)
      (5,9998.3093019583-1121.5123436609)
      (6,9337.1979833547-1062.9314699704)
    };
    \addplot[name path=highSu, draw=none, color=tabblue] coordinates {
      (0,11289.4928097345+1758.9508447551)
      (1,11576.0158137317+1557.4125577151)
      (2,11626.8676563487+1537.7452479029)
      (3,11620.4147727273+1521.0103018355)
      (4,11463.5467260013+1522.4054897042)
      (5,9998.3093019583+1121.5123436609)
      (6,9337.1979833547+1062.9314699704)
    };
    \addplot[draw=none, color=tabblue, fill opacity=0.2] fill between[of=lowSu and highSu];
    \addplot[thick, color=tabblue] coordinates {
      (0,11289.4928097345)
      (1,11576.0158137317)
      (2,11626.8676563487)
      (3,11620.4147727273)
      (4,11463.5467260013)
      (5,9998.3093019583)
      (6,9337.1979833547)
    };
    \node[anchor=north east, font=\scriptsize] at (axis description cs:1,1) {Summer};
    
    \nextgroupplot[
      xlabel={},
      ytick={10000,14000},
      xticklabel style={rotate=90}
    ]
    \addplot[name path=lowA, draw=none, color=tabblue] coordinates {
      (0,12213.0917467949-1911.4546982594)
      (1,12479.5726427197-1778.6701457275)
      (2,12470.7385379929-1679.9727490495)
      (3,12576.6769563823-1739.2532337336)
      (4,12397.4703335471-1681.9965967307)
      (5,10886.2099230523-1299.0548962355)
      (6,10242.5913523018-1317.6637702947)
    };
    \addplot[name path=highA, draw=none, color=tabblue] coordinates {
      (0,12213.0917467949+1911.4546982594)
      (1,12479.5726427197+1778.6701457275)
      (2,12470.7385379929+1679.9727490495)
      (3,12576.6769563823+1739.2532337336)
      (4,12397.4703335471+1681.9965967307)
      (5,10886.2099230523+1299.0548962355)
      (6,10242.5913523018+1317.6637702947)
    };
    \addplot[draw=none, color=tabblue, fill opacity=0.2] fill between[of=lowA and highA];
    \addplot[thick, color=tabblue] coordinates {
      (0,12213.0917467949)
      (1,12479.5726427197)
      (2,12470.7385379929)
      (3,12576.6769563823)
      (4,12397.4703335471)
      (5,10886.2099230523)
      (6,10242.5913523018)
    };
    \node[anchor=north east, font=\scriptsize] at (axis description cs:1,1) {Autumn};
    
    \nextgroupplot[
      xlabel={},
      yticklabels={},
      xticklabel style={rotate=90}
    ]
    \addplot[name path=lowW, draw=none, color=tabblue] coordinates {
      (0,13069.0817726832-2052.8061840320)
      (1,13252.2588824289-1895.5014777154)
      (2,13377.6900000000-1912.8380539171)
      (3,13289.2585649644-1875.3554274632)
      (4,13158.9607558140-1864.8635028232)
      (5,11770.2330073742-1378.7719075799)
      (6,11237.1205929487-1422.5128227133)
    };
    \addplot[name path=highW, draw=none, color=tabblue] coordinates {
      (0,13069.0817726832+2052.8061840320)
      (1,13252.2588824289+1895.5014777154)
      (2,13377.6900000000+1912.8380539171)
      (3,13289.2585649644+1875.3554274632)
      (4,13158.9607558140+1864.8635028232)
      (5,11770.2330073742+1378.7719075799)
      (6,11237.1205929487+1422.5128227133)
    };
    \addplot[draw=none, color=tabblue, fill opacity=0.2] fill between[of=lowW and highW];
    \addplot[thick, color=tabblue] coordinates {
      (0,13069.0817726832)
      (1,13252.2588824289)
      (2,13377.6900000000)
      (3,13289.2585649644)
      (4,13158.9607558140)
      (5,11770.2330073742)
      (6,11237.1205929487)
    };
    \node[anchor=north east, font=\scriptsize] at (axis description cs:1,1) {Winter};
    
    \end{groupplot}
    \node[anchor=north, font=\scriptsize, yshift=-20pt] at ($(group c1r2.south east)!0.5!(group c2r2.south west)$) {Day of the Week};

    \node[anchor=south, rotate=90, font=\scriptsize, yshift=22pt] at ($(group c1r1.north west)!0.5!(group c1r2.south west)$){Grid Load in MWh};
\end{tikzpicture}%
       \caption{Average weekly grid load by season.}
        \label{fig:average_load_by_day_and_season}
    \end{subfigure}%
    \hfill
    \begin{subfigure}[b]{0.49\linewidth}
        \centering
        \begin{tikzpicture}
    \begin{groupplot}[
      group style={
        group size=2 by 2,
        horizontal sep=0.05cm,
        vertical sep=0.05cm,
      },
      small axis,
      width=0.6\linewidth,
      height=0.5\imageheight,
      grid=both,
      ymin=8000, ymax=16000,
      scaled y ticks=false,
      xtick={0,1,2,3,4,5,6,7,8,9,10,11,12,13,14,15,16,17,18,19,20,21,22,23},
      xticklabels={$0$,,,,,,$6$,,,,,,$12$,,,,,,$18$,,,,,},
      error bars/y dir=both,
      error bars/y explicit,
    ]

    % Spring
    \nextgroupplot[
      xticklabels={},
      ytick={10000,14000},
    ]
    
    \addplot[name path=lowS, draw=none, color=tabblue] coordinates {
      (0,9548.5206521739-1052.0633231110)
      (1,9553.4008152174-1044.2378937352)
      (2,9777.6644021739-1071.4250777200)
      (3,10183.6000000000-1224.9081208132)
      (4,10900.0355978261-1578.2725470860)
      (5,11683.2923913043-1789.6239718232)
      (6,12316.4980978261-1724.6817719301)
      (7,12617.8788043478-1585.7293908114)
      (8,12779.9073369565-1489.4995144330)
      (9,12979.1092391304-1398.0139238701)
      (10,12871.5266304348-1482.7663733823)
      (11,12547.0820652174-1582.7021010346)
      (12,12278.1043478261-1609.6607465032)
      (13,12105.0010869565-1595.0132239923)
      (14,11983.6578804348-1541.6303660520)
      (15,12094.1937500000-1443.7071633114)
      (16,12345.0896739130-1367.9753945879)
      (17,12600.2190217391-1468.2983828369)
      (18,12496.2274456522-1505.9102661966)
      (19,11996.8252717391-1309.1048787196)
      (20,11368.5793478261-1164.7584195705)
      (21,10745.9127717391-1163.3263529853)
      (22,10151.5828804348-1111.2609786740)
      (23,9740.7741847826-1063.1001966805)
    };
    
    \addplot[name path=highS, draw=none, color=tabblue] coordinates {
      (0,9548.5206521739+1052.0633231110)
      (1,9553.4008152174+1044.2378937352)
      (2,9777.6644021739+1071.4250777200)
      (3,10183.6000000000+1224.9081208132)
      (4,10900.0355978261+1578.2725470860)
      (5,11683.2923913043+1789.6239718232)
      (6,12316.4980978261+1724.6817719301)
      (7,12617.8788043478+1585.7293908114)
      (8,12779.9073369565+1489.4995144330)
      (9,12979.1092391304+1398.0139238701)
      (10,12871.5266304348+1482.7663733823)
      (11,12547.0820652174+1582.7021010346)
      (12,12278.1043478261+1609.6607465032)
      (13,12105.0010869565+1595.0132239923)
      (14,11983.6578804348+1541.6303660520)
      (15,12094.1937500000+1443.7071633114)
      (16,12345.0896739130+1367.9753945879)
      (17,12600.2190217391+1468.2983828369)
      (18,12496.2274456522+1505.9102661966)
      (19,11996.8252717391+1309.1048787196)
      (20,11368.5793478261+1164.7584195705)
      (21,10745.9127717391+1163.3263529853)
      (22,10151.5828804348+1111.2609786740)
      (23,9740.7741847826+1063.1001966805)
    };
    
    \addplot[draw=none, color=tabblue, fill opacity=0.2]
      fill between[of=lowS and highS];
    
    \addplot[thick, color=tabblue] coordinates {
      (0,9548.5206521739)
      (1,9553.4008152174)
      (2,9777.6644021739)
      (3,10183.6000000000)
      (4,10900.0355978261)
      (5,11683.2923913043)
      (6,12316.4980978261)
      (7,12617.8788043478)
      (8,12779.9073369565)
      (9,12979.1092391304)
      (10,12871.5266304348)
      (11,12547.0820652174)
      (12,12278.1043478261)
      (13,12105.0010869565)
      (14,11983.6578804348)
      (15,12094.1937500000)
      (16,12345.0896739130)
      (17,12600.2190217391)
      (18,12496.2274456522)
      (19,11996.8252717391)
      (20,11368.5793478261)
      (21,10745.9127717391)
      (22,10151.5828804348)
      (23,9740.7741847826)
    };
    \node[anchor=north east, font=\scriptsize]
      at (axis description cs:1,1) {Spring};

    % Sommer
    \nextgroupplot[
      yticklabels={},
      xticklabels={},
    ]
    
    \addplot[name path=lowSu, draw=none, color=tabblue] coordinates {
      (0,8835.4896739130-655.1332286087)
      (1,8824.0040760870-679.9500825164)
      (2,8984.6619565217-738.7975322472)
      (3,9381.5823369565-941.2492874422)
      (4,10316.5350543478-1354.8553717678)
      (5,11200.8127717391-1540.3034716420)
      (6,11802.1551630435-1463.4321638596)
      (7,12095.2010869565-1330.5508394533)
      (8,12354.3763586957-1280.5983514685)
      (9,12626.1171195652-1220.0904306969)
      (10,12470.0415760870-1266.5833505186)
      (11,12182.1619565217-1346.5976721583)
      (12,11973.1989130435-1364.8857900117)
      (13,11838.8244565217-1335.3318829090)
      (14,11727.9282608696-1274.3535855297)
      (15,11845.1505434783-1198.8129717517)
      (16,11934.3201086957-1105.1956542367)
      (17,11844.4923913043-1081.7474522293)
      (18,11454.7540760870-1013.0914307500)
      (19,11079.0622282609-910.2851016399)
      (20,10596.2472826087-782.5235154443)
      (21,9932.7059782609-712.5969294106)
      (22,9405.2964673913-687.8860088178)
      (23,9045.5429347826-661.9085698955)
    };
    
    \addplot[name path=highSu, draw=none, color=tabblue] coordinates {
      (0,8835.4896739130+655.1332286087)
      (1,8824.0040760870+679.9500825164)
      (2,8984.6619565217+738.7975322472)
      (3,9381.5823369565+941.2492874422)
      (4,10316.5350543478+1354.8553717678)
      (5,11200.8127717391+1540.3034716420)
      (6,11802.1551630435+1463.4321638596)
      (7,12095.2010869565+1330.5508394533)
      (8,12354.3763586957+1280.5983514685)
      (9,12626.1171195652+1220.0904306969)
      (10,12470.0415760870+1266.5833505186)
      (11,12182.1619565217+1346.5976721583)
      (12,11973.1989130435+1364.8857900117)
      (13,11838.8244565217+1335.3318829090)
      (14,11727.9282608696+1274.3535855297)
      (15,11845.1505434783+1198.8129717517)
      (16,11934.3201086957+1105.1956542367)
      (17,11844.4923913043+1081.7474522293)
      (18,11454.7540760870+1013.0914307500)
      (19,11079.0622282609+910.2851016399)
      (20,10596.2472826087+782.5235154443)
      (21,9932.7059782609+712.5969294106)
      (22,9405.2964673913+687.8860088178)
      (23,9045.5429347826+661.9085698955)
    };
    
    \addplot[draw=none, color=tabblue, fill opacity=0.2]
      fill between[of=lowSu and highSu];
    
    \addplot[thick, color=tabblue] coordinates {
      (0,8835.4896739130)
      (1,8824.0040760870)
      (2,8984.6619565217)
      (3,9381.5823369565)
      (4,10316.5350543478)
      (5,11200.8127717391)
      (6,11802.1551630435)
      (7,12095.2010869565)
      (8,12354.3763586957)
      (9,12626.1171195652)
      (10,12470.0415760870)
      (11,12182.1619565217)
      (12,11973.1989130435)
      (13,11838.8244565217)
      (14,11727.9282608696)
      (15,11845.1505434783)
      (16,11934.3201086957)
      (17,11844.4923913043)
      (18,11454.7540760870)
      (19,11079.0622282609)
      (20,10596.2472826087)
      (21,9932.7059782609)
      (22,9405.2964673913)
      (23,9045.5429347826)
    };
    
    \node[anchor=north east, font=\scriptsize]
      at (axis description cs:1,1) {Summer};

    % Autuum
    \nextgroupplot[
    ytick={10000,14000},
    ]
    
    \addplot[name path=lowA, draw=none, color=tabblue] coordinates {
      (0,9607.8431318681-891.1805935368)
      (1,9587.7961538462-879.2663776291)
      (2,9790.5486263736-908.7367459827)
      (3,10296.1975274725-1031.8871086573)
      (4,11180.9532967033-1423.6521678165)
      (5,12048.8296703297-1702.1240231811)
      (6,12679.3225274725-1680.5332578158)
      (7,12961.4815934066-1531.8591819167)
      (8,13113.8456043956-1419.0316420779)
      (9,13284.8714285714-1341.8512605244)
      (10,13223.3810439560-1401.5850036098)
      (11,12934.6414835165-1499.9107381661)
      (12,12679.5766483516-1520.2800714219)
      (13,12538.1192307692-1496.8410554972)
      (14,12476.3478021978-1445.7344215015)
      (15,12759.0626373626-1408.5261786473)
      (16,13236.7675824176-1411.2163892048)
      (17,13373.0791208791-1319.6374561211)
      (18,12982.8543956044-1298.8706045303)
      (19,12226.5351648352-1224.1208934459)
      (20,11489.3793956044-1107.0333260622)
      (21,10860.0313186813-1075.5089340005)
      (22,10274.2076923077-1001.4607997791)
      (23,9861.0225274725-932.4648172213)
    };
    
    \addplot[name path=highA, draw=none, color=tabblue] coordinates {
      (0,9607.8431318681+891.1805935368)
      (1,9587.7961538462+879.2663776291)
      (2,9790.5486263736+908.7367459827)
      (3,10296.1975274725+1031.8871086573)
      (4,11180.9532967033+1423.6521678165)
      (5,12048.8296703297+1702.1240231811)
      (6,12679.3225274725+1680.5332578158)
      (7,12961.4815934066+1531.8591819167)
      (8,13113.8456043956+1419.0316420779)
      (9,13284.8714285714+1341.8512605244)
      (10,13223.3810439560+1401.5850036098)
      (11,12934.6414835165+1499.9107381661)
      (12,12679.5766483516+1520.2800714219)
      (13,12538.1192307692+1496.8410554972)
      (14,12476.3478021978+1445.7344215015)
      (15,12759.0626373626+1408.5261786473)
      (16,13236.7675824176+1411.2163892048)
      (17,13373.0791208791+1319.6374561211)
      (18,12982.8543956044+1298.8706045303)
      (19,12226.5351648352+1224.1208934459)
      (20,11489.3793956044+1107.0333260622)
      (21,10860.0313186813+1075.5089340005)
      (22,10274.2076923077+1001.4607997791)
      (23,9861.0225274725+932.4648172213)
    };
    
    \addplot[draw=none, color=tabblue, fill opacity=0.2]
      fill between[of=lowA and highA];
    
    \addplot[thick, color=tabblue] coordinates {
      (0,9607.8431318681)
      (1,9587.7961538462)
      (2,9790.5486263736)
      (3,10296.1975274725)
      (4,11180.9532967033)
      (5,12048.8296703297)
      (6,12679.3225274725)
      (7,12961.4815934066)
      (8,13113.8456043956)
      (9,13284.8714285714)
      (10,13223.3810439560)
      (11,12934.6414835165)
      (12,12679.5766483516)
      (13,12538.1192307692)
      (14,12476.3478021978)
      (15,12759.0626373626)
      (16,13236.7675824176)
      (17,13373.0791208791)
      (18,12982.8543956044)
      (19,12226.5351648352)
      (20,11489.3793956044)
      (21,10860.0313186813)
      (22,10274.2076923077)
      (23,9861.0225274725)
    };
    
    \node[anchor=north east, font=\scriptsize]
      at (axis description cs:1,1) {Autumn};
    
    \nextgroupplot[
      % xticklabels={},
      yticklabels={},
    ]
    
    \addplot[name path=lowW, draw=none, color=tabblue] coordinates {
      (0,10605.6785714286-1081.2084010263)
      (1,10501.4828349945-1087.2089957690)
      (2,10583.9695459579-1109.7656760056)
      (3,10824.6245847176-1169.1976703419)
      (4,11251.3862126246-1375.9537516624)
      (5,12064.1115725360-1866.3570837161)
      (6,13013.1594684385-2103.9015862775)
      (7,13535.0877630122-1945.5641645436)
      (8,13791.0440199336-1707.7848943464)
      (9,13987.8823366556-1546.5801477253)
      (10,14216.2635658915-1406.1068315971)
      (11,13985.6824473976-1469.0521107904)
      (12,13664.4188815061-1558.3386139924)
      (13,13474.1769102990-1585.0416561552)
      (14,13409.4975083056-1551.6194103412)
      (15,13633.0902547065-1459.7333422422)
      (16,14177.9792358804-1410.7420004654)
      (17,14345.8444075305-1349.9958604331)
      (18,14071.6240310078-1376.7381693033)
      (19,13317.3853820598-1302.4547479966)
      (20,12578.6270764120-1188.4714245742)
      (21,12137.0794573643-1090.9295652631)
      (22,11471.0052602436-1051.9844861521)
      (23,10933.8441306755-1067.4221851712)
    };
    
    \addplot[name path=highW, draw=none, color=tabblue] coordinates {
      (0,10605.6785714286+1081.2084010263)
      (1,10501.4828349945+1087.2089957690)
      (2,10583.9695459579+1109.7656760056)
      (3,10824.6245847176+1169.1976703419)
      (4,11251.3862126246+1375.9537516624)
      (5,12064.1115725360+1866.3570837161)
      (6,13013.1594684385+2103.9015862775)
      (7,13535.0877630122+1945.5641645436)
      (8,13791.0440199336+1707.7848943464)
      (9,13987.8823366556+1546.5801477253)
      (10,14216.2635658915+1406.1068315971)
      (11,13985.6824473976+1469.0521107904)
      (12,13664.4188815061+1558.3386139924)
      (13,13474.1769102990+1585.0416561552)
      (14,13409.4975083056+1551.6194103412)
      (15,13633.0902547065+1459.7333422422)
      (16,14177.9792358804+1410.7420004654)
      (17,14345.8444075305+1349.9958604331)
      (18,14071.6240310078+1376.7381693033)
      (19,13317.3853820598+1302.4547479966)
      (20,12578.6270764120+1188.4714245742)
      (21,12137.0794573643+1090.9295652631)
      (22,11471.0052602436+1051.9844861521)
      (23,10933.8441306755+1067.4221851712)
    };
    
    \addplot[draw=none, color=tabblue, fill opacity=0.2]
      fill between[of=lowW and highW];
    
    \addplot[thick, color=tabblue] coordinates {
      (0,10605.6785714286)
      (1,10501.4828349945)
      (2,10583.9695459579)
      (3,10824.6245847176)
      (4,11251.3862126246)
      (5,12064.1115725360)
      (6,13013.1594684385)
      (7,13535.0877630122)
      (8,13791.0440199336)
      (9,13987.8823366556)
      (10,14216.2635658915)
      (11,13985.6824473976)
      (12,13664.4188815061)
      (13,13474.1769102990)
      (14,13409.4975083056)
      (15,13633.0902547065)
      (16,14177.9792358804)
      (17,14345.8444075305)
      (18,14071.6240310078)
      (19,13317.3853820598)
      (20,12578.6270764120)
      (21,12137.0794573643)
      (22,11471.0052602436)
      (23,10933.8441306755)
    };
    
    \node[anchor=north east, font=\scriptsize]
      at (axis description cs:1,1) {Winter};
    
    \end{groupplot}
    \node[anchor=north, font=\scriptsize, yshift=-20pt] at ($(group c1r2.south east)!0.5!(group c2r2.south west)$) {Hour of the Day};

    \node[anchor=south, rotate=90, font=\scriptsize, yshift=22pt] at ($(group c1r1.north west)!0.5!(group c1r2.south west)$){Grid Load in MWh};

\end{tikzpicture}%
        \caption{Average daily grid load by season.}
        \label{fig:average_load_by_hour_and_season}
    \end{subfigure}
    \caption{Temporal patterns of electricity demand across multiple time scales.}
    \label{fig:grid_load_overview}
\end{figure}

\paragraph*{Temporal Features}
Temporal features are included to capture recurring demand patterns driven by human behavior and institutional schedules. Electricity consumption exhibits pronounced daily, weekly, and seasonal cycles that are not fully explained by physical variables alone. Accordingly, the dataset contains features representing the hour of the day, day of the week, week of the year and month.

Several temporal features (e.g., hour of day, day of week) are inherently cyclical. To preserve their periodic structure and avoid artificial discontinuities at time boundaries, these variables are encoded using sine and cosine transformations of the corresponding period. This trigonometric encoding efficiently captures temporal proximity and enables the model to learn smooth recurring patterns~\cite{rodrigo_2025_skforecast, pedregosa_2011_scikit}. The resulting encoding is illustrated in Figure~\ref{fig:cyclic_encoding}.
\begin{figure}[ht]
    \centering
    
\begin{tikzpicture}
    \begin{axis}[
        small axis,
        width=0.5\linewidth,
        xlabel={sin(hour)},
        ylabel={cos(hour)},
        grid=both,
        axis equal image,
        xmin=-1.2, xmax=1.2,
        ymin=-1.2, ymax=1.2,
        colormap/viridis,
        point meta min=0,
        point meta max=23,
        colorbar,
        colorbar style={font=\scriptsize, ylabel={Hour of day}, yticklabel style={/pgf/number format/assume math mode=false}},
    ]

    \addplot[
    only marks,
    scatter,
    scatter src=explicit,
    mark=*,
    mark size=2.8pt
    ] coordinates {
    ( {sin(deg(2*pi*0/24))},  {cos(deg(2*pi*0/24))} )  [0]
    ( {sin(deg(2*pi*1/24))},  {cos(deg(2*pi*1/24))} )  [1]
    ( {sin(deg(2*pi*2/24))},  {cos(deg(2*pi*2/24))} )  [2]
    ( {sin(deg(2*pi*3/24))},  {cos(deg(2*pi*3/24))} )  [3]
    ( {sin(deg(2*pi*4/24))},  {cos(deg(2*pi*4/24))} )  [4]
    ( {sin(deg(2*pi*5/24))},  {cos(deg(2*pi*5/24))} )  [5]
    ( {sin(deg(2*pi*6/24))},  {cos(deg(2*pi*6/24))} )  [6]
    ( {sin(deg(2*pi*7/24))},  {cos(deg(2*pi*7/24))} )  [7]
    ( {sin(deg(2*pi*8/24))},  {cos(deg(2*pi*8/24))} )  [8]
    ( {sin(deg(2*pi*9/24))},  {cos(deg(2*pi*9/24))} )  [9]
    ( {sin(deg(2*pi*10/24))}, {cos(deg(2*pi*10/24))} ) [10]
    ( {sin(deg(2*pi*11/24))}, {cos(deg(2*pi*11/24))} ) [11]
    ( {sin(deg(2*pi*12/24))}, {cos(deg(2*pi*12/24))} ) [12]
    ( {sin(deg(2*pi*13/24))}, {cos(deg(2*pi*13/24))} ) [13]
    ( {sin(deg(2*pi*14/24))}, {cos(deg(2*pi*14/24))} ) [14]
    ( {sin(deg(2*pi*15/24))}, {cos(deg(2*pi*15/24))} ) [15]
    ( {sin(deg(2*pi*16/24))}, {cos(deg(2*pi*16/24))} ) [16]
    ( {sin(deg(2*pi*17/24))}, {cos(deg(2*pi*17/24))} ) [17]
    ( {sin(deg(2*pi*18/24))}, {cos(deg(2*pi*18/24))} ) [18]
    ( {sin(deg(2*pi*19/24))}, {cos(deg(2*pi*19/24))} ) [19]
    ( {sin(deg(2*pi*20/24))}, {cos(deg(2*pi*20/24))} ) [20]
    ( {sin(deg(2*pi*21/24))}, {cos(deg(2*pi*21/24))} ) [21]
    ( {sin(deg(2*pi*22/24))}, {cos(deg(2*pi*22/24))} ) [22]
    ( {sin(deg(2*pi*23/24))}, {cos(deg(2*pi*23/24))} ) [23]
    };

\end{axis}
\end{tikzpicture}
    \caption{Cyclic Encoding of the hour feature.}
    \label{fig:cyclic_encoding}
\end{figure}

Calendar effects are represented through public and school holidays as well as bridge days. Since public holidays vary across German federal states, a single feature is constructed by aggregating the number of states observing a holiday on each date. School holidays, which differ in timing and duration across regions, are included to account for systematic demand changes associated with reduced industrial activity and altered household behavior~\cite{johannsen_2025_schulferien}.

In addition, statistical rolling features such as the mean and standard deviation over the preceding seven days are included, allowing the model to access summary statistics of recent demand behavior without learning them implicitly from raw lag values. The dataset further contains the grid load from the equivalent weekday one year prior, where the equivalent date refers not to a fixed 365-day offset but to the nearest matching weekday, ensuring comparability in weekly demand structure across years. 

\paragraph*{Weather Data}
Weather variables are included to capture the impact of meteorological conditions on electricity demand and renewable generation. The data is obtained via the Meteostat Python library~\cite{lamprecht_2025_meteostat} and compiled from multiple stations distributed across the 50Hertz control area to ensure adequate spatial coverage. The selected stations represent the north-eastern, north-western, central, south-eastern, and south-western parts of the region. For each timestamp weather observations from the nearest station are assigned, allowing regional variation in temperature, wind and solar conditions to be reflected in the dataset. While the data consists of realized weather observations, it is treated as a proxy for weather forecasts, as actual forecast data was not available.

\paragraph*{Market Data}
Market price data is obtained from the Energy-Charts API~\cite{burger_2025_energy-charts} and covers the DE/LU price zone. Until October 2018, Germany, Austria, and Luxembourg formed a joint bidding zone (DE/AT/LU), after which Austria became a separate market area. Consequently, two price series are available in the dataset. For consistency over the full observation period and since the focus is on electricity demand in Germany, these series are combined into a single day-ahead market price variable.

The complete dataset can be found in the Appendix~\ref{app:dataset}.